\documentclass[12pt,oneside,openany]{book}
\usepackage[T1]{fontenc}
\usepackage[utf8]{inputenc}
\usepackage{amsmath,amssymb,amsthm,mathtools,mathrsfs}
\DeclareFontFamily{U}{rsfs}{\skewchar\font127 }
\DeclareFontShape{U}{rsfs}{m}{n}{<-6.5>rsfs5<6.5-7.5>rsfs7<7.5->rsfs10}{}
\usepackage{booktabs,array}
\usepackage[numbers,sort&compress]{natbib}
\usepackage{imakeidx}%
\usepackage{environ}
\PassOptionsToPackage{hypertexnames=false}{hyperref}

\makeatletter

\usepackage{lmodern}
\usepackage[sc]{mathpazo}
\usepackage{bm}
\linespread{1.1}
\usepackage{microtype}
\microtypesetup{stretch=15,shrink=15}

\usepackage{geometry}
\geometry{letterpaper,left=1in,right=1in,top=0.75in,bottom=0.75in,includeheadfoot,
  headheight=15pt,headsep=16pt,footskip=26pt}
\setlength{\parindent}{1.5em}
\setlength{\parskip}{0pt}
\setlength{\emergencystretch}{3em}
\tolerance=1000
\clubpenalty=3000 \widowpenalty=3000 \displaywidowpenalty=1500
\raggedbottom
\allowdisplaybreaks[2]

\usepackage{enumitem}
\setlist{topsep=4pt plus 2pt minus 1pt,itemsep=2pt plus 1pt,parsep=0pt}
\setlist[itemize]{leftmargin=1.5em}
\setlist[enumerate]{leftmargin=2em}

\usepackage[newparttoc]{titlesec}
\usepackage{titletoc,fancyhdr}

\newcommand{\trilogy@spacedsc}[1]{{\scshape\textls[90]{\MakeLowercase{#1}}}}
\newcommand{\trilogy@headpar}{\rightskip=0pt plus 0.5\hsize\relax
  \parfillskip=0pt plus 0.5\hsize\relax\hyphenpenalty=10000\relax\exhyphenpenalty=10000\relax}

\titleformat{\part}[display]
  {\normalfont\centering}
  {\large\trilogy@spacedsc{\partname\ \thepart}}
  {22pt}
  {\fontsize{22}{28}\selectfont}
\assignpagestyle{\part}{empty}

\titleformat{\chapter}[display]
  {\normalfont\trilogy@headpar}
  {\trilogy@spacedsc{\chaptertitlename\ \thechapter}}
  {12pt}
  {\fontsize{20}{25}\selectfont}
\titleformat{name=\chapter,numberless}[display]
  {\normalfont\trilogy@headpar}
  {}
  {0pt}
  {\fontsize{20}{25}\selectfont\trilogy@starmark}
\titlespacing*{\chapter}{0pt}{12pt}{30pt}
\newcommand{\trilogy@starmark}[1]{#1\markboth{#1}{#1}}

\titleformat{\section}
  {\normalfont\fontsize{13.5}{17}\selectfont\bfseries\trilogy@headpar}{\thesection}{0.75em}{}
\titlespacing*{\section}{0pt}{22pt plus 4pt minus 2pt}{9pt plus 2pt}
\titleformat{\subsection}
  {\normalfont\normalsize\bfseries\trilogy@headpar}{\thesubsection}{0.75em}{}
\titlespacing*{\subsection}{0pt}{15pt plus 3pt minus 2pt}{6pt plus 1pt}
\titleformat{\paragraph}[runin]
  {\normalfont\normalsize\bfseries}{}{0pt}{}
\titlespacing*{\paragraph}{0pt}{9pt plus 2pt minus 1pt}{0.6em}

\setcounter{tocdepth}{1}
\setcounter{secnumdepth}{1}

\newcommand{\shorthead}[2]{\expandafter\def\csname trilogy@sh@\detokenize{#1}\endcsname{#2}}
\newcommand{\trilogy@lookup}[1]{%
  \ifcsname trilogy@sh@\detokenize{#1}\endcsname
    \expandafter\let\expandafter\trilogy@headtext\csname trilogy@sh@\detokenize{#1}\endcsname
  \else
    \def\trilogy@headtext{#1}%
  \fi}
\newcommand{\trilogy@headfont}{\normalfont\small\itshape}
\DeclareRobustCommand{\trilogy@hd}[2]{%
  \ifx\relax#1\relax\else{\upshape#1}\hspace{0.8em}\fi#2}
\newsavebox{\trilogy@chapbox}
\newsavebox{\trilogy@secbox}
\newsavebox{\trilogy@cachedchapbox}
\newsavebox{\trilogy@cachedsecbox}
\newif\iftrilogy@headtoolong
\newif\iftrilogy@cachedheadtoolong
\let\trilogy@cachedheadkey\relax
\newcommand{\trilogy@computeheads}{%
  \trilogy@headtoolongfalse
  \sbox\trilogy@chapbox{\trilogy@headfont\nouppercase{\leftmark}}%
  \sbox\trilogy@secbox{}%
  \protected@edef\trilogy@l{\leftmark}%
  \protected@edef\trilogy@r{\rightmark}%
  \ifx\trilogy@r\@empty\else\ifx\trilogy@r\trilogy@l\else
    \sbox\trilogy@secbox{\trilogy@headfont\nouppercase{\rightmark}}%
  \fi\fi
  \dimen@=\wd\trilogy@chapbox\advance\dimen@ 4em\relax
  \ifdim\dimen@>\headwidth
    \sbox\trilogy@chapbox{\trilogy@headfont\@chapapp\ \thechapter}%
    \trilogy@headtoolongtrue
  \fi
  \dimen@=\wd\trilogy@chapbox\advance\dimen@\wd\trilogy@secbox\advance\dimen@ 8em\relax
  \ifdim\dimen@>\headwidth
    \sbox\trilogy@secbox{\trilogy@headfont
      \expandafter\let\csname trilogy@hd \endcsname\trilogy@hdnum\nouppercase{\rightmark}}%
  \fi}
\newcommand{\trilogy@heads}[1]{%
  \protected@edef\trilogy@headkey{%
    {\leftmark}{\rightmark}{\the\headwidth}%
    {\fontname\font}{\the\dimexpr1em\relax}{\thechapter}{\@chapapp}}%
  \ifx\trilogy@headkey\trilogy@cachedheadkey\else
    \trilogy@computeheads
    \global\setbox\trilogy@cachedchapbox=\copy\trilogy@chapbox
    \global\setbox\trilogy@cachedsecbox=\copy\trilogy@secbox
    \global\let\trilogy@cachedheadkey\trilogy@headkey
    \iftrilogy@headtoolong
      \global\trilogy@cachedheadtoolongtrue
    \else
      \global\trilogy@cachedheadtoolongfalse
    \fi
  \fi
  \setbox\trilogy@chapbox=\copy\trilogy@cachedchapbox
  \setbox\trilogy@secbox=\copy\trilogy@cachedsecbox
  \iftrilogy@cachedheadtoolong
    \ifx\relax#1\relax\else\PackageWarningNoLine{trilogy-style}{Running head too long on
      page \thepage; give the chapter a \string\shorthead}\fi
  \fi}
\newcommand{\trilogy@hdnum}[2]{\ifx\relax#1\relax\else\S\,{\upshape#1}\fi}

\pagestyle{fancy}
\fancyhf{}
\fancyhead[L]{\trilogy@heads{warn}\usebox\trilogy@chapbox}
\fancyhead[R]{\trilogy@heads{}%
  \ifdim\wd\trilogy@secbox>\z@\usebox\trilogy@secbox\hspace{2em}\fi
  \normalfont\small\thepage}
\renewcommand{\headrulewidth}{0pt}
\renewcommand{\footrulewidth}{0pt}
\fancypagestyle{plain}{\fancyhf{}\fancyfoot[C]{\normalfont\small\thepage}%
  \renewcommand{\headrulewidth}{0pt}\renewcommand{\footrulewidth}{0pt}}
\renewcommand{\chaptermark}[1]{\trilogy@lookup{#1}%
  \markboth{\trilogy@hd{\ifnum\c@secnumdepth>\m@ne\if@mainmatter\thechapter\fi\fi}{\trilogy@headtext}}{}}
\renewcommand{\sectionmark}[1]{\trilogy@lookup{#1}%
  \markright{\trilogy@hd{\ifnum\c@secnumdepth>\z@\thesection\fi}{\trilogy@headtext}}}

\contentsmargin{2.2em}
\titlecontents{part}[0pt]
  {\addvspace{16pt plus 2pt}\normalfont}
  {\trilogy@spacedsc{\partname\ \thecontentslabel}\hspace{1em}}
  {}
  {\hfill\contentspage}
  [\addvspace{4pt}]
\titlecontents{chapter}[2.2em]
  {\addvspace{7pt plus 1pt}\normalfont}
  {\contentslabel{2.2em}}
  {\hspace*{-2.2em}}
  {\hfill\contentspage}
\titlecontents{section}[4.8em]
  {\normalfont\small}
  {\contentslabel{2.6em}}
  {\hspace*{-2.6em}}
  {\titlerule*[0.9em]{.}\contentspage}

\indexsetup{othercode={\small\raggedcolumns\maxbalancingoverflow=\z@}}
\@ifpackageloaded{natbib}{%
  \renewcommand{\bibsection}{\chapter*{References}%
    \addcontentsline{toc}{chapter}{References}}%
}{}

\newlength{\trilogyenvskip}
\setlength{\trilogyenvskip}{10pt plus 3pt minus 2pt}
\newtheoremstyle{plain}{\trilogyenvskip}{\trilogyenvskip}
  {\itshape}{}{\bfseries}{.}{0.5em}{}
\newtheoremstyle{definition}{\trilogyenvskip}{\trilogyenvskip}
  {\normalfont}{}{\bfseries}{.}{0.5em}{}
\newtheoremstyle{remark}{\trilogyenvskip}{\trilogyenvskip}
  {\normalfont}{}{\itshape}{.}{0.5em}{}

\newcommand{\trilogynote}[2]{\par\addvspace{5pt plus 1pt}\noindent{\itshape #1}\enspace #2\par}

\newcommand{\usedin}[1]{\trilogynote{Used in:}{#1}}

\newcommand{\walkthroughname}{Technical walkthrough}
\newcommand{\walkthroughend}{\ensuremath{\square}}
\newtheoremstyle{trilogysetoff}{\trilogyenvskip}{\trilogyenvskip}
  {\normalfont\leftskip=1.5em\rightskip=1.5em\relax}{}{\bfseries}{.}{0.5em}{}

\theoremstyle{definition}
\newtheorem{trilogywalk}{\walkthroughname}[chapter]
\newenvironment{walkthrough}[1][]
  {\ifx\relax#1\relax\begin{trilogywalk}\else\begin{trilogywalk}[#1]\fi
   \pushQED{\qed}\renewcommand{\qedsymbol}{\walkthroughend}}
  {\popQED\end{trilogywalk}}
\theoremstyle{plain}

\newenvironment{trilogyaside}[1]
  {\par\addvspace{11pt plus 3pt minus 2pt}%
   \list{}{\leftmargin=1.5em\rightmargin=1.5em\topsep=0pt\partopsep=0pt\parsep=0pt
     \listparindent=1.5em\itemsep=0pt}%
   \item\relax{\trilogy@spacedsc{#1}}\par\nopagebreak\vspace{3pt}\noindent\ignorespaces}
  {\endlist\par\addvspace{11pt plus 3pt minus 2pt}}
\newenvironment{sayit}{\begin{trilogyaside}{How to say it}}{\end{trilogyaside}}
\newenvironment{fiveminute}{\begin{trilogyaside}{Five-minute explanation}}{\end{trilogyaside}}
\newcommand{\tophysicist}{\par\addvspace{2pt}\noindent{\itshape To a physicist.}\enspace\ignorespaces}
\newcommand{\tomathematician}{\par\addvspace{4pt}\noindent{\itshape To a mathematician.}\enspace\ignorespaces}

\newcommand{\trilogyneedspace}[1]{\par\penalty-100\begingroup
  \setlength{\dimen@}{#1}\dimen@ii\pagegoal\advance\dimen@ii-\pagetotal
  \ifdim\dimen@>\dimen@ii\ifdim\dimen@ii>\z@\vfil\fi\break\fi\endgroup}

\usepackage[font=small,labelsep=period,skip=8pt]{caption}

\usepackage{etoolbox}
\patchcmd{\@footnotetext}{\footnotesize}{\small}{}{}

\newcommand{\trilogy@overhang}{0pt}
\newcommand{\widetablesetup}{%
  \setlength{\LTleft}{-\trilogy@overhang plus 1fill}%
  \setlength{\LTright}{-\trilogy@overhang plus 1fill}%
  \setlength{\linewidth}{\dimexpr\textwidth+2\dimexpr\trilogy@overhang\relax\relax}}

\usepackage[hidelinks,bookmarksnumbered=true]{hyperref}
\usepackage{bookmark}
\hypersetup{pdfauthor={},pdfpagemode=UseNone,pdfdisplaydoctitle=true,
  pdfstartview={FitH \hypercalcbp{\paperheight}}}

\newcommand{\trilogyseries}{Notes on Spin Glasses}
\newcommand{\trilogytitlepage}[3]{%
  \hypersetup{pageanchor=false}%
  \begin{titlepage}
  \centering\normalfont
  \vspace*{1.7in}
  {\large\trilogy@spacedsc{\trilogyseries}\par}
  \vspace{12pt}
  {\normalsize\itshape Volume #1\par}
  \vspace{1.05in}
  {\fontsize{28}{35}\selectfont #2\par}
  \vspace{24pt}
  {\Large\itshape #3\par}
  \vfill
  \end{titlepage}%
  \hypersetup{pageanchor=true}}

\makeatother
%
% Short running heads for long titles: \shorthead{exact title}{short head}.
\shorthead{Self-overlap moments on quantile intervals}{Self-overlap moments}
\shorthead{Gaussian tools for arbitrary profiles}{Gaussian tools}
\shorthead{Localization of the lowest quantiles}{The lowest quantiles}
\shorthead{The cavity interpolation and its tree laws}{The cavity interpolation}
\shorthead{The derivative at the independent endpoint}{The endpoint derivative}

\usepackage{xurl}%
\usepackage[nameinlink,capitalize,noabbrev]{cleveref}
\usepackage{aliascnt}
\hypersetup{pdftitle={An N\textasciicircum(-7/10) Bound for the Finite-Size Bias of the Sherrington-Kirkpatrick Free Energy},pdfauthor={P. M. Aronow and Patrick Lopatto}}
\renewcommand{\chaptername}{Section}
\numberwithin{equation}{chapter}

\theoremstyle{plain}
\newtheorem{theorem}{Theorem}[chapter]
\newaliascnt{proposition}{theorem}
\newtheorem{proposition}[proposition]{Proposition}
\aliascntresetthe{proposition}
\newaliascnt{lemma}{theorem}
\newtheorem{lemma}[lemma]{Lemma}
\aliascntresetthe{lemma}
\newaliascnt{corollary}{theorem}
\newtheorem{corollary}[corollary]{Corollary}
\aliascntresetthe{corollary}
\theoremstyle{definition}
\newaliascnt{definition}{theorem}
\newtheorem{definition}[definition]{Definition}
\aliascntresetthe{definition}
\newaliascnt{example}{theorem}
\newtheorem{example}[example]{Example}
\aliascntresetthe{example}
\theoremstyle{remark}
\newaliascnt{remark}{theorem}
\newtheorem{remark}[remark]{Remark}
\aliascntresetthe{remark}
\makeatletter
\@addtoreset{proposition}{chapter}\@addtoreset{lemma}{chapter}\@addtoreset{corollary}{chapter}
\@addtoreset{definition}{chapter}\@addtoreset{example}{chapter}\@addtoreset{remark}{chapter}
\makeatother
\crefname{theorem}{Theorem}{Theorems}
\crefname{proposition}{Proposition}{Propositions}
\crefname{lemma}{Lemma}{Lemmas}
\crefname{corollary}{Corollary}{Corollaries}
\crefname{definition}{Definition}{Definitions}
\crefname{example}{Example}{Examples}
\crefname{remark}{Remark}{Remarks}
\crefname{chapter}{Section}{Sections}
\crefname{part}{Part}{Parts}
\crefname{section}{Section}{Sections}
\crefname{subsection}{Section}{Sections}
\crefname{table}{Table}{Tables}
\crefname{equation}{}{}
\crefformat{equation}{(#2#1#3)}
\Crefformat{equation}{(#2#1#3)}
\crefrangeformat{equation}{(#3#1#4)--(#5#2#6)}
\Crefrangeformat{equation}{(#3#1#4)--(#5#2#6)}
\crefmultiformat{equation}{(#2#1#3)}{ and~(#2#1#3)}{, (#2#1#3)}{, and~(#2#1#3)}
\Crefmultiformat{equation}{(#2#1#3)}{ and~(#2#1#3)}{, (#2#1#3)}{, and~(#2#1#3)}
\newcommand{\creflastconjunction}{, and~}
\newcommand{\creflastgroupconjunction}{, and~}

\makeatletter
\newcommand{\npp@noteblock}[2]{%
  \ifx\relax#1\relax\else{\centering\trilogy@spacedsc{#1}\par}\vspace{6pt}\fi
  {\centering
   \begin{minipage}{0.84\textwidth}\small\normalfont\setlength{\parindent}{1.5em}%
     \noindent\ignorespaces#2\end{minipage}\par}}
\NewEnviron{authorsnote}{%
  \clearpage\thispagestyle{empty}%
  \vspace*{\stretch{2}}%
  \npp@noteblock{Authors' note}{\BODY}%
  \vspace*{\stretch{3}}%
  \clearpage}
\NewEnviron{abstractpage}{%
  \clearpage\thispagestyle{empty}%
  \vspace*{\stretch{2}}%
  \npp@noteblock{Abstract}{\BODY}%
  \vspace*{\stretch{3}}%
  \clearpage}
\newcommand{\npptitlepage}[2]{%
  \hypersetup{pageanchor=false}%
  \begin{titlepage}
  \centering\normalfont
  \vspace*{1.7in}
  {\large\strut\par}
  \vspace{1.05in}
  {\fontsize{28}{35}\selectfont #1\par}
  \vspace{1.1in}
  {\Large\def\\{\par\vspace{10pt}}#2\par}
  \vfill
  \end{titlepage}%
  \hypersetup{pageanchor=true}}
\makeatother

\newcommand{\E}{\mathbb E}
\newcommand{\Pp}{\mathbb P}
\newcommand{\R}{\mathbb R}
\newcommand{\dd}{\,\mathrm d}
\DeclareMathOperator{\Var}{Var}
\DeclareMathOperator{\Cov}{Cov}
\DeclareMathOperator{\Tr}{Tr}
\DeclareMathOperator{\supp}{supp}
\DeclareMathOperator{\sech}{sech}
\newcommand{\id}{\mathrm{id}}
\newcommand{\eps}{\varepsilon}
\newcommand{\ind}{\mathbf1}
\newcommand{\scal}{\mathrm{sc}}
\newcommand{\calF}{\mathcal F}
\newcommand{\calD}{\mathcal D}
\newcommand{\calR}{\mathcal R}
\newcommand{\calG}{\mathcal G}
\newcommand{\norm}[1]{\left\lVert#1\right\rVert}
\newcommand{\abs}[1]{\left|#1\right|}

\allowdisplaybreaks[2]

\clubpenalty=10000
\patchcmd{\thebibliography}{\clubpenalty4000}{\clubpenalty10000}{}{}
\begin{document}
\frontmatter
\npptitlepage{An $N^{-7/10}$ Bound for the\\ Finite-Size Bias of the\\ Sherrington--Kirkpatrick\\ Free Energy}{P.\ M.\ Aronow\\ Patrick Lopatto}
\begin{authorsnote}
With the exception of this note and small edits to the abstract, the entire
document is machine written. We have taken care to supervise this writing to
ensure that the proofs meet a minimum standard of readability and that
previous literature is appropriately cited. However, the resulting document
remains lacking in exposition and overall coherence. We are nonetheless
releasing this manuscript on Hexagon to disseminate the result as quickly as
possible, because we believe it may be of interest to the mathematics and
physics communities. We welcome suggestions concerning mathematical
corrections, omitted citations, and attribution of ideas.

P.L.\ was partially supported by NSF grant DMS-2450004.
\end{authorsnote}
\begin{abstractpage}
For every fixed inverse temperature $\beta>1$, we prove that the expected free
energy per spin of the zero-field Gaussian Sherrington--Kirkpatrick model
differs from its Parisi limit by at most $C_\beta N^{-7/10}$. This improves
the exponent $2/3$ of our earlier work (arXiv:2609.09103). We also compare the
bound with the finite-size correction of order $N^{-3/4}$ predicted by Franz,
Parisi, and Ricci-Tersenghi, and with the strict improvement over $2/3$ that
would follow from the ideas in OpenAI's recently released manuscripts.

\end{abstractpage}
\tableofcontents
\mainmatter

\chapter{Introduction}\label{sec:introduction}

\section{The Sherrington--Kirkpatrick model and the Parisi formula}\label{sec:intro-background}

Let $(g_{ij})_{1\le i<j\le N}$ be independent standard Gaussian
variables. The Sherrington--Kirkpatrick (SK)
model~\cite{SherringtonKirkpatrick1975} at zero external field and
inverse temperature $\beta>0$ has Hamiltonian
\[
 H_N(\sigma)=\frac1{\sqrt N}\sum_{1\le i<j\le N}g_{ij}\sigma_i\sigma_j,
 \qquad \sigma\in\{-1,1\}^N.
\]
Its logarithmic partition function and expected pressure are
\[
 F_{N,\beta}=\log\sum_{\sigma\in\{-1,1\}^N}e^{\beta H_N(\sigma)},
 \qquad
 p_{N,\beta}=\frac{\E F_{N,\beta}}N.
\]
For a probability measure $\mu$ on $[0,1]$ with distribution function
$\alpha(q)=\mu([0,q])$, let $u_\alpha$ be the solution of the Parisi
equation on $[0,1]\times\R$, and define the Parisi functional
$\mathcal P_\beta$ and the Parisi value $P_\beta$ by
\begin{equation}\label{eq:intro-parisi}
\begin{gathered}
 \partial_qu_\alpha=-\frac{\beta^2}{2}
 \bigl(\partial_{xx}u_\alpha+\alpha(q)(\partial_xu_\alpha)^2\bigr),
 \qquad u_\alpha(1,x)=\log(2\cosh x),\\
 \mathcal P_\beta(\alpha)=u_\alpha(0,0)
       -\frac{\beta^2}{2}\int_0^1q\alpha(q)\dd q,
 \qquad P_\beta=\inf_{\mu\in\mathcal P([0,1])}\mathcal P_\beta(\alpha),
\end{gathered}
\end{equation}
where $\mathcal P([0,1])$ is the set of probability measures on
$[0,1]$. For a general $\mu$, the equation is understood in the weak
sense of~\cite[Theorem~4]{JagannathTobasco2016}.

Sherrington and Kirkpatrick introduced this model and studied it by a
replica-symmetric calculation, which fails at low
temperature~\cite{SherringtonKirkpatrick1975}; a longer treatment, with
numerical comparisons, is given in~\cite{KirkpatrickSherrington1978}.
De~Almeida and Thouless showed that the replica-symmetric saddle point
becomes unstable~\cite{deAlmeidaThouless1978}. Parisi proposed a
hierarchical breaking of the replica symmetry in~\cite{Parisi1979}, and
he developed a sequence of approximate solutions and the order
parameter on $[0,1]$ that leads to \eqref{eq:intro-parisi}
in~\cite{Parisi1980Sequence,Parisi1980OrderParameter}. He later
interpreted this order parameter through the distribution of overlaps
between states~\cite{Parisi1983OrderParameter}. M\'ezard, Parisi,
Sourlas, Toulouse, and Virasoro then described the organization of the
states, with sample-dependent weights and ultrametric overlaps,
in~\cite{MezardEtAl1984Nature}.

The Parisi formula states that $p_{N,\beta}\to P_\beta$ as
$N\to\infty$. Guerra proved the upper bound $p_{N,\beta}\le P_\beta$ for
every $N$~\cite{Guerra2003}, and Talagrand proved the matching lower
bound in the limit~\cite{Talagrand}. The infimum in
\eqref{eq:intro-parisi} is attained at a unique measure
$\mu_\beta$~\cite{AuffingerChen2015Uniqueness}, the Parisi measure,
whose distribution function we denote by $\alpha_\beta$. For $\beta>1$,
its support is an interval $[0,q_\beta]$ with
$0<q_\beta<1$~\cite[Theorem~1.1]{Lopatto}; for $\beta>1$ close to $1$,
this was proved earlier by Zhou~\cite{Zhou2026FRSB}.

\section{Finite-size corrections}\label{sec:intro-finite-size}

In this paper, we study the rate of convergence in the Parisi formula at
low temperature. By Guerra's bound, the \emph{mean bias}
$P_\beta-p_{N,\beta}$ is nonnegative for every $N$. Up to a factor
$1/\beta$, it is the finite-size correction to the expected free energy
per spin (\cref{sec:intro-results}).

At high temperature, the mean bias is small. For $0<\beta<1$, Guerra
and Toninelli showed by a quadratic coupling of two replicas that it is
of order at most $N^{-1}$~\cite[Section~3]{GuerraToninelli2002Quadratic}.
With a nonzero external field, Lopatto proved that the limiting
pressure equals the replica-symmetric value up to and including the
de~Almeida--Thouless line~\cite[Theorem~1.1]{Lopatto2026ATLine}, and
Kusuoka and Nakajima proved that the pressure is within $O(N^{-1})$ of
this value on compact subsets of the strict de~Almeida--Thouless
region~\cite[Theorem~1.1]{KusuokaNakajima2026}.

Below the critical temperature, the mean bias is expected to decay more
slowly. In~\cite[Theorem~1.1, (1.5)]{AL}, we proved that
$0\le P_\beta-p_{N,\beta}\le C\beta N^{-1/2}$ for every $\beta>0$ and
$N\ge1$, with a numerical constant $C$. For fixed $\beta>1$, we used
the interval support of $\mu_\beta$ to prove the sharper upper bound
$C_\beta N^{-2/3}$ and the lower bound $c_\beta N^{-11/12}$, for all
sufficiently large $N$~\cite[Theorem~1.4]{AL}. In physics, the
correction is expected to decay like $N^{-\omega}$ with $\omega<1$ at
low temperature. Two values have been proposed for the full model,
$\omega=2/3$ and $\omega=3/4$. A calculation retaining only
longitudinal fluctuations gives
$\omega=4/5$~\cite{ParisiSarraTalamanca2019}. We review these predictions in
\cref{subsec:intro-physics}. We improve the exponent of the upper bound
from $2/3$ to $7/10$.

\section{Main result}\label{sec:intro-results}

Our main result is the following.

\begin{theorem}\label{thm:main}
For every $\beta>1$ there is a constant $C_\beta<\infty$ such that
\[
 P_\beta-p_{N,\beta}\le C_\beta N^{-7/10}
 \qquad(N\ge2).
\]
\end{theorem}

Since $P_\beta-p_{N,\beta}\ge0$ for every $N$~\cite{Guerra2003},
\cref{thm:main} is a two-sided estimate. In physical units, the
expected free energy density is $\E[-F_{N,\beta}/(\beta N)]$, and its
limit is $-P_\beta/\beta$. By \cref{thm:main}, we have
\[
 0\le \E\Bigl[-\frac{F_{N,\beta}}{\beta N}\Bigr]
     -\Bigl(-\frac{P_\beta}{\beta}\Bigr)
 =\frac{P_\beta-p_{N,\beta}}{\beta}
 \le\frac{C_\beta}{\beta}N^{-7/10}
 \qquad(N\ge2).
\]
For the expected logarithm of the partition function itself, the
corresponding bound is $0\le NP_\beta-\E F_{N,\beta}\le C_\beta N^{3/10}$. The constant
$C_\beta$ depends on $\beta$, and we do not control it as
$\beta\downarrow1$ or as $\beta\to\infty$. In particular, we do not
obtain a corresponding bound for the ground-state energy.

Together with the lower bound in~\cite[Theorem~1.4]{AL},
\cref{thm:main} implies that for every fixed $\beta>1$,
\[
 c_\beta N^{-11/12}\le P_\beta-p_{N,\beta}\le C_\beta N^{-7/10}
\]
for all sufficiently large $N$. The exponent $7/10$ lies between the
exponent $2/3$ of~\cite[Theorem~1.4]{AL} and the exponent $3/4$
predicted by Franz, Parisi, and
Ricci-Tersenghi~\cite{FranzParisiRicciTersenghi2026}, which we discuss
in \cref{subsec:intro-physics}. In \cref{subsec:intro-openai}, we
compare \cref{thm:main} with two conditional consequences of results
stated in the OpenAI
manuscripts~\cite{OpenAIFluctuationScale,OpenAILimitingLaw}. An
uncentered comparison stated there would give an unspecified strict
improvement over $2/3$, and the variance asymptotic stated there would
imply the lower bound $P_\beta-p_{N,\beta}\ge c_\beta N^{-5/6}$.

We now outline the proof; \cref{sec:overview} gives a detailed
overview. We follow the interpolation scheme of~\cite{AL}, which is a
quantitative form of the Guerra interpolation~\cite{Guerra2003} that
Talagrand used for the positive fractional moments of the partition
function~\cite{Talagrand2007FractionalMoments}. For a tilt
$0<\lambda\le\lambda_0$, let $e_{N,\beta}(\lambda)\ge0$ be the error,
in energy units, between the constrained Parisi value and the normalized
logarithmic moment of order $\lambda/\beta$ of the completed partition
function. The constraint requires mass at least $\lambda/\beta$ at zero
(\cref{sec:fnd}).
Imposing this mass raises the Parisi value by at most
$C_\beta\lambda^5$ in energy units~\cite[Lemma~7.2]{AL}. Using this
bound and the convexity of Gaussian logarithmic
moments~\cite[Theorem~1]{ChenConvexity}, as in the proof
of~\cite[Theorem~7.8, Step~1]{AL}, we find
\[
 0\le\frac{P_\beta-p_{N,\beta}}{\beta}
 \le C_\beta\lambda^5+2e_{N,\beta}(\lambda).
\]
For sufficiently large $N$, we take $\lambda=N^{-7/50}$, so it
suffices to show that
$e_{N,\beta}(N^{-7/50})\le C_\beta N^{-7/10}$.

Up to a factor of four, $e_{N,\beta}(\lambda)$ is the integral over
$\rho\in(0,1]$ of a defect $\calD_N(\rho)=\beta(A+2J+V)$ along the
interpolation, evaluated at a minimizing profile of an $N$-dimensional
variational problem. The terms $A$ and $J$ are controlled directly
(\cref{sec:fnd,sec:th}). The difficulty is $V$, the integrated variance
of a self-overlap. For $\rho$ above a cutoff $\rho_0$, we bound $V$ by
a one-site cavity argument (\cref{part:cavity}). Removing one spin leads
to an inequality $V\le(1-\rho)(1+\eps_N)V+{}$errors, where
$1+\eps_N$ bounds a scalar response operator at the minimizing profile,
and this inequality contracts once $\eps_N\le\rho/4$. To prove this
bound, we show that the response operator at the Parisi measure is
bounded by the identity, which reflects the marginal stability of
$\mu_\beta$, and that the minimizing profile is close to a scalar trial
profile (\cref{part:profile}). For $\rho\le\rho_0$, we bound $V$ by
the quantile-shell estimates of~\cite[Section~3]{AL}, refined in
\cref{sec:wi}, and by a localization of the lowest quantiles of the
minimizer (\cref{sec:lz-root}), which refines~\cite[Lemma~7.6]{AL}. The cutoff is
$\rho_0=\min\{1,K_\beta\eta^{1/7}\}$, where the trial budget
$\eta=e_{N,\beta}(\lambda)+C_\beta\lambda^5$ contains the unknown error.
The result is an implicit inequality of the form
\[
 e_{N,\beta}(\lambda)\lesssim N^{-2/3}
 \bigl(e_{N,\beta}(\lambda)+\lambda^5\bigr)^{1/21}+\text{smaller terms}.
\]
If $e_{N,\beta}(\lambda)$ and $\lambda^5$ are both of order $N^{-a}$,
this inequality is balanced when $a=2/3+a/21$, that is, $a=7/10$.

\section{Relation to physical predictions}\label{subsec:intro-physics}

\paragraph{Expectations and early calculations.}
Below the critical temperature, physicists expect the finite-size
correction to the free energy density to be nonanalytic, of order
$N^{-\omega}$ with $\omega<1$. This expectation is related to marginal
stability. In a truncated replica model near criticality,
De~Dominicis and Kondor found that the stability matrix of the Parisi
solution has nonnegative eigenvalues and many zero
modes~\cite[pp.~606--608]{DeDominicisKondor1983}. These zero modes make the
expansion in fluctuations around this solution singular;
see~\cite[Section~2]{FranzParisiRicciTersenghi2026}. Parisi, Ritort,
and Slanina studied finite-size corrections at
criticality~\cite{ParisiRitortSlanina1993Critical},
where the free energy density has a correction $(\log N)/(12N)$
(see~\cite{Billoire2006FiniteSize}), as well as slightly below
criticality and along the de~Almeida--Thouless line, including
corrections to overlap distributions and to
ultrametricity~\cite{ParisiRitortSlanina1993Corrections}.

\paragraph{The exponent $2/3$.}
At zero temperature, Bouchaud, Krzakala, and Martin found numerically,
for Gaussian couplings, that the mean ground-state energy per spin
approaches its limit like
$N^{-2/3}$~\cite[Section~VI\,B]{BouchaudKrzakalaMartin2003}. In a
study of the SK model with binary couplings, Boettcher fits the corrections to the
mean of the ground-state energy per spin and to the width of its
distribution with exponents close to $2/3$ and $3/4$,
respectively~\cite[Section~3]{Boettcher2005}. In a later study of
diluted versions of the model, the fit for the undiluted model is
again close to $2/3$~\cite[Table~I]{Boettcher2020Diluted}. At positive
temperature, Billoire estimated the correction numerically for binary
couplings, by integrating overlap observables along the interpolation of
Guerra and Toninelli~\cite{GuerraToninelli2002}, and found data
compatible with $\omega=2/3$~\cite{Billoire2006FiniteSize}. Aspelmeier,
Billoire, Marinari, and Moore derived $\omega=2/3$ from a heuristic
argument: a system of $N$ spins resolves about $N^{1/6}$ steps of
replica symmetry breaking, and truncation after $K$ steps changes the
free energy density by order
$K^{-4}$~\cite[Section~2]{AspelmeierEtAl2008FiniteSize}.

\paragraph{Longitudinal fluctuations and the exponent $4/5$.}
Parisi and Biscari obtained an earlier contribution of order
$N^{-3/4}\log N$ from the positive longitudinal eigenvalues near
criticality, neglecting interactions between
modes~\cite[Section~1 and Eq.~(5.25)]{ParisiBiscari1992}. Their
contribution from the longitudinal zero modes diverges when treated separately
from the omitted transverse modes~\cite[discussion following
Eq.~(5.27)]{ParisiBiscari1992}; see
\cite[Section~1]{ParisiSarraTalamanca2019}. Parisi, Sarra,
and Talamanca later used a resummation-based scaling argument for the
longitudinal fluctuations of a truncated free-energy functional for the SK model near
criticality and predicted
$\omega=4/5$~\cite[Section~3, Eqs.~(16)--(19)]{ParisiSarraTalamanca2019}.
They omit the transverse fluctuations, so this exponent describes their
longitudinal approximation. To determine the exponent for the full model, one
must also control the omitted modes.

\paragraph{The exponent $3/4$.}
Franz, Parisi, and Ricci-Tersenghi~\cite{FranzParisiRicciTersenghi2026}
argue that $\omega=2/3$ is incorrect and that the correct value is
$\omega=3/4$, possibly up to logarithmic
factors~\cite[Eq.~(14)]{FranzParisiRicciTersenghi2026}. Their
calculation for the Gaussian SK model matches two regimes: one in which the
overlaps of two coupled systems lock and the Guerra--Toninelli
integrand decays like $(Nt)^{-1/2}$, and a perturbative regime in which
this integrand is of order $t^{-3/2}/N$~\cite[Eqs.~(9)--(12) and
(15)--(17)]{FranzParisiRicciTersenghi2026}. They estimate the large-$z$ locking
exponent semianalytically by sampling random ultrametric
trees generated by the Bolthausen--Sznitman coalescent and extrapolating
in the number of leaves~\cite[Materials and Methods]{FranzParisiRicciTersenghi2026}.
Matching this behavior to the perturbative regime gives the scaling
prediction. The
fit of the energy uses random regular graphs of degree
four~\cite[Eq.~(18) and Fig.~2]{FranzParisiRicciTersenghi2026}, and the
simulations of the mode-locking integral use rewired random regular
graphs of degree four with binary couplings (see the subsection
``Direct simulations of the mode-locking integral''
of~\cite[Materials and Methods]{FranzParisiRicciTersenghi2026}). These
models differ from the Gaussian SK model in their connectivity, and the
latter also in the law of the couplings.

\paragraph{Consequences of \cref{thm:main}.}
\Cref{thm:main} implies that $P_\beta-p_{N,\beta}=O(N^{-7/10})$ for
every fixed $\beta>1$; in particular, the correction is $o(N^{-2/3})$. We therefore exclude a nonzero leading
correction of order $N^{-2/3}$ for the Gaussian SK model at every fixed
temperature below the critical one. This is
consistent with the prediction of Franz, Parisi, and Ricci-Tersenghi.
\Cref{thm:main} does not determine the leading exponent $\omega$. Since our
constants are not uniform as $\beta\to\infty$, it also says nothing
about the numerical results at zero temperature described above.

The estimates on which our proof rests have a similar structure with
two regimes, although they concern a different interpolation. If we
ignore logarithms and the dependence on the tilt $\lambda$, the defect
$\calD_N(\rho)$ is at most of order $(N\rho)^{-2/3}$ for small $\rho$
and of order $(N\rho^2)^{-1}$ on the interval where our cavity argument
applies. If the second bound held for all $\rho\ge N^{-1/4}$, where the
two bounds are equal, integration would give $N^{-3/4}$. In our proof,
the cavity interval begins only at $\rho_0\ge N^{-1/10}$, and the tilt
$\lambda=N^{-7/50}$ enters both bounds; these two features account for
the difference between $7/10$ and $3/4$.

Marginal stability also enters our proof, through a scalar statement:
at the Parisi measure, the multiplier that dominates our scalar response
operator equals one on $\supp\mu_\beta$ (\cref{sec:sc}). Crisanti and
Rizzo wrote the corresponding identity in physical variables, on the
part of the order parameter that
varies~\cite[Eqs.~(22)--(23)]{CrisantiRizzo2002}. This scalar statement
does not describe the full stability spectrum of the replica
calculation.

\paragraph{Mean, width, and large deviations.}
The exponents $\omega$ for the finite-size correction discussed above concern
the mean. The sample-to-sample fluctuations
of the free energy and its large deviations are governed by different
exponents. In a truncated replica model near criticality, Kondor
obtained a fifth-order dependence on $n$ along the Parisi branch for small positive
$n$~\cite[pp.~L129--L130, Eq.~(11)]{Kondor1983}. Interpreted
as a formal prediction for the SK model, this says that the
replica-normalized pressure in the thermodynamic limit,
$\lim_{N\to\infty}(Nn)^{-1}\log\E e^{nF_{N,\beta}}$, differs from
$P_\beta$ first at order $n^5$ as $n\downarrow0$ with
$n>0$ (see~\cite[pp.~1 and~3]{ParisiRizzo2008}).
Parisi and Rizzo extended this calculation throughout the
low-temperature phase in~\cite{ParisiRizzo2008}
and~\cite[Appendix~B]{ParisiRizzo2009}. Equivalently, in this formal
expansion, the limit of $N^{-1}\log\E e^{n(F_{N,\beta}-NP_\beta)}$
begins at order $n^6$, and its Legendre transform gives a rate with
power $6/5$ for upward deviations of
$F_{N,\beta}/N$~\cite{ParisiRizzo2008}.
Crisanti, Paladin, Sommers, and
Vulpiani interpreted the replica number
as a tilt over the disorder and predicted that the free energy per spin
fluctuates on the scale
$N^{-5/6}$~\cite{CrisantiEtAl1992Fluctuations}. Under an asymmetric
scaled-law assumption, the same calculation also predicts a correction to the mean
density on this scale~\cite[p.~1331, Eq.~(34) and following
text]{CrisantiEtAl1992Fluctuations}. This conclusion uses the
truncated replica theory near criticality and a matching of the disorder
tails to the typical peak.
In our convention, in which $F_{N,\beta}$ is the logarithm of the partition
function, Parisi and Rizzo's
replica predictions give upward density deviations with speed $N$ and
power $6/5$, and downward density deviations with speed $N^2$ and power
$12/5$~\cite[Eqs.~(12)--(15)]{ParisiRizzo2010}. Smoothly matching these
tails to the finite-size peak predicts the typical density scale
$N^{-5/6}$. This matching is an additional assumption; the limits
$n\downarrow0$ and $N\to\infty$ cannot simply be
interchanged~\cite[Section~VI]{ParisiRizzo2009}.
Aspelmeier, Billoire,
Marinari, and Moore also predicted a variance of order $N^{1/3}$ for
$F_{N,\beta}$~\cite{AspelmeierEtAl2008FiniteSize}.
Aspelmeier derived exact Gaussian interpolation identities relating
the variance of the free energy to two- and four-replica overlaps between
systems with correlated disorder~\cite[Eqs.~(14) and~(31)]{Aspelmeier2008Exact}.
Combining these identities with replica estimates of the bond-chaos
overlap distribution gives a predicted upper bound $1/4$ for the
extensive standard-deviation exponent~\cite[Section~4.3]{Aspelmeier2008Exact,Aspelmeier2008Bond}.
The identities are exact; this exponent estimate uses replica
approximations and concerns the width.
At zero temperature,
the width of the distribution of the ground-state energy has been
studied numerically in~\cite{Boettcher2005,Palassini2008,Boettcher2010Fluctuations}.
For Gaussian couplings, Palassini finds data consistent with a mean
correction exponent $2/3$ and, after subleading corrections, favors the
exponent $5/6$ over
$3/4$ for the standard deviation per spin~\cite[Sections~3--4 and~6]{Palassini2008}.
Boettcher discusses the difficulty of distinguishing the exponents
$3/4$ and $5/6$ after corrections to
scaling~\cite[Section~IV]{Boettcher2010Fluctuations}.
Monthus and Garel studied rare-disorder scenarios and the matching
between typical and large-deviation tails for the ground-state energy of the SK model at zero
temperature~\cite[Section~III C--F]{MonthusGarel2010}.
Their alternative width--tail relations illustrate the additional
assumptions involved in such matching.

\section{Related work}\label{sec:intro-related}

\paragraph{The Parisi formula.}
Besides the proofs of Guerra and Talagrand, the rigorous theory of the
Parisi formula includes the following results. Guerra and Toninelli
proved by interpolation that the thermodynamic limit
exists~\cite{GuerraToninelli2002}, and Aizenman, Sims, and Starr
formulated an extended variational principle in terms of random overlap
structures~\cite{AizenmanSimsStarr2003}. Panchenko proved that the
Ghirlanda--Guerra identities~\cite{GhirlandaGuerra1998} imply ultrametricity of
overlaps~\cite[Theorem~1]{Panchenko2013Ultrametricity}, and he proved
the Parisi formula for mixed $p$-spin models, including mixtures with
odd interactions~\cite{Panchenko2014Parisi}. We refer to the books of
Talagrand~\cite{TalagrandVolI2011,TalagrandVolII2011} and
Panchenko~\cite{Panchenko2013Book} for accounts of this theory;
Panchenko develops it through Ruelle probability cascades.

\paragraph{The Parisi measure.}
Toninelli proved that the replica-symmetric value is not optimal for
$\beta>1$ at zero field~\cite{Toninelli2002}. Auffinger and Chen proved
that the minimizer is unique~\cite{AuffingerChen2015Uniqueness}, that
zero belongs to its support, and that the largest point of its support
is an atom for $1<\beta\le3/2$~\cite[Theorems~1
and~4]{AuffingerChen2015Properties}. They also proved that its
distribution function is smooth on $[a,b)$ whenever $(a,b)$ is
contained in its support, for $0\le a<b\le1$~\cite[Theorem~2]{AuffingerChen2015Properties}.
They expressed the Parisi functional through a stochastic control
problem, which they used in their uniqueness proof, and they proved
martingale identities for the derivatives of
$u_\alpha$~\cite[Theorem~3 and Lemma~2]{AuffingerChen2015Uniqueness}. An earlier
Hamilton--Jacobi--Bellman representation appears
in~\cite[Section~6.2]{BovierKlimovsky2009}, and Jagannath and Tobasco
developed a dynamic programming treatment of the
functional~\cite[Lemma~18]{JagannathTobasco2016}. First-variation
formulas along segments and optimality conditions on the support are
given by Jagannath and Tobasco~\cite[Lemma~3.2 and
Proposition~1.1]{JagannathTobasco2017PhaseDiagram} and by
Chen~\cite[Theorem~2 and Proposition~1]{Chen2017Variational}. Lopatto
proved, in addition to the interval support for every $\beta>1$, that
$\mu_\beta$ is the sum of a density that is smooth up to the endpoints
of $[0,q_\beta]$ and an atom at $q_\beta$~\cite[Theorem~1.1]{Lopatto}.
Auffinger, Chen, and Zeng proved that the Parisi
minimizer at zero temperature has infinitely many support
points~\cite{AuffingerChenZeng2020}. H.-B.~Chen proved that
$q_\beta\to1$ as $\beta\to\infty$, and that the Stieltjes measure
associated with the Parisi order parameter at zero temperature has a
smooth density and support
$[0,1)$~\cite[Theorems~1.2 and~1.3]{Chen2026ZeroTemperatureSupport}. The
full support at zero temperature is also stated
in~\cite[Theorem~1.1]{OpenAIFullSupport}. We use the interval support
at fixed $\beta>1$ throughout, and the regularity of $\mu_\beta$ in
\cref{part:profile}.

\paragraph{Fluctuations of the free energy.}
The fluctuations of $F_{N,\beta}$ depend on $\beta$ and on the external
field. For $0<\beta<1$, Aizenman, Lebowitz, and Ruelle proved a
Gaussian limit for its fluctuations, which are of order
one~\cite{AizenmanLebowitzRuelle1987}, and Comets and Neveu studied this
limit by stochastic calculus~\cite{CometsNeveu1995}. With a nonzero
external field, Chen, Dey, and Panchenko proved that the variance is of
order $N$ and that a central limit theorem holds, under assumptions on
the interactions that include the SK
model~\cite[Theorems~1 and~2]{ChenDeyPanchenko2017}. At $\beta=1$, Chen
and Lam proved a variance bound of order
$\log^2N$~\cite[Theorem~1.1]{ChenLam2019}. Schertzer proved that the
variance is at most $(\log N)/4+C$ and at least of order
$\log\log\log N$~\cite[Theorems~1 and~2]{Schertzer2026Critical}, and Du
and Huang proved that $\Var(F_{N,1})=(\log N)/6+O(1)$ and that the
fluctuations have a Gaussian limit~\cite[Theorem~1.3]{DuHuang2026}. Dey and Kang obtained
variance asymptotics and a Gaussian limit as $\beta$ approaches $1$ from
below along sequences $\beta=\beta_N$ satisfying
$N^{1/3}(1-\beta_N^2)\to c\in(0,\infty]$~\cite[Theorem~1.1]{DeyKang2026}.
For fixed $\beta>1$, Chatterjee proved the superconcentration bound
$\Var(F_{N,\beta})\le C_\beta N/\log N$ and related it to disorder
chaos~\cite[Theorem~1.5]{Chatterjee2009Chaos}. In~\cite[Theorem~1.4]{AL},
we proved that $\Var(F_{N,\beta})$ lies between constant multiples of
$N^{4/15}$ and $N^{7/15}$ for large $N$. The OpenAI manuscripts state
the order $N^{1/3+o(1)}$, a positive limit for
$N^{-1/3}\Var(F_{N,\beta})$, and a limiting law for the normalized
fluctuations; see \cref{subsec:intro-openai}. A variance bound alone
does not bound the mean bias from above, but a lower bound on the
variance gives a lower bound on the mean bias
through~\cite[Lemma~7.4]{AL}. For the spherical SK model, Baik and Lee
proved that the free energy per spin has Gaussian fluctuations of order
$N^{-1}$ at high temperature and Tracy--Widom fluctuations of order
$N^{-2/3}$ at low temperature, governed by the largest eigenvalue of
the coupling matrix~\cite[Theorems~2.10 and~2.11]{BaikLee2016}.

\paragraph{Gaussian-process comparisons.}
Chu and Raginsky use Gaussian-process geometry and information-theoretic
tensorization to bound expected soft maxima~\cite{ChuRaginsky2026}.
Their application to the SK model gives a finite-volume comparison with the Parisi
variational value up to universal multiplicative
constants~\cite[Section~6, Theorem~10]{ChuRaginsky2026}.
It uses a uniform spin prior and normalizes the Hamiltonian differently.
Its conclusion is a multiplicative comparison; the present theorem
controls the additive error to the limiting value.

\paragraph{Fractional moments.}
Talagrand computed the limits of normalized logarithms of positive
moments of the partition function~\cite{Talagrand2007FractionalMoments}. For moments of order
$0<s<1$, the limit of $(Ns)^{-1}\log\E e^{sF_{N,\beta}}$ is the minimum
$P_\beta(s)$ of $\mathcal P_\beta$ over measures $\mu$ with
$\mu(\{0\})\ge s$. For the SK model at zero field, this formula is
contained in the one-level specialization of the multiscale formula of
Contucci and Mingione~\cite[Theorem~2.1]{ContucciMingione2019}. Chen,
Guionnet, Ko, Lacroix-A-Chez-Toine, and Mourrat stated it in this form
and used it to study one-sided large deviations of the ground-state
energy~\cite[Theorem~2.1]{ChenEtAl2026LargeDeviations}.
In~\cite[Lemma~7.2 and Proposition~7.3]{AL}, we proved that for fixed
$\beta>1$, the increase $P_\beta(s)-P_\beta$ is of order $s^5$ as
$s\downarrow0$. This is a rigorous counterpart of the replica
calculations of Kondor and of Parisi and Rizzo described in
\cref{subsec:intro-physics}, and Pei identified the leading coefficient
of this increase~\cite[Theorem~1(a)]{Pei2026}. The tilt
$\lambda=N^{-7/50}$ in our proof balances this cost against the
fractional-moment error.

\paragraph{Origins of the methods.}
The interpolation of~\cite[Section~2]{AL} is a finite-dimensional,
quantitative form of Guerra's interpolation, with which he proved the
upper bound in the Parisi formula~\cite{Guerra2003}. For posterior
estimates under its tilted laws, we combine the Brascamp--Lieb
inequality~\cite[Theorem~4.1]{BrascampLieb1976}, in its inverse-Hessian
form~\cite[Proposition~2.2]{AL}, with a joint-source argument from the
proof of~\cite[Lemma~2.3]{OpenAIFluctuationScale} (\cref{sec:gauss}).
Near the singular end of the interpolation, our estimates refine the
quantile-window bound of~\cite[Lemma~3.3]{AL} and the localization of
lowest quantiles in~\cite[Lemma~7.6]{AL}. We close the final implicit
inequality as in the proof of~\cite[Proposition~7.7]{AL}.

The cavity method originates in physics. An early precursor is the
single-site description of Thouless, Anderson, and
Palmer~\cite[Section~3, Eqs.~(10)--(12)]{ThoulessAndersonPalmer1977},
and M\'ezard, Parisi, and
Virasoro formulated the cavity construction for broken replica symmetry
in~\cite{MezardParisiVirasoro1986}; see also the exposition
in~\cite{MezardParisiVirasoro1987}. We use Talagrand's rigorous version
of the one-site argument~\cite[Section~1.6]{TalagrandVolI2011}. The
expansion under hierarchical Gaussian laws relies on normalized replica
trees and covariance insertions of the type developed
in~\cite[Section~2]{OpenAIFluctuationScale}. In the scalar stability
argument, the first variation of $\mathcal P_\beta$ under transports of
levels is due to Talagrand~\cite[Lemma~3.7]{Talagrand2006ParisiMeasures};
see~\cite[Section~3, (19)]{AuffingerChen2015Properties}. We compute the
second variation and compare it with a response multiplier, using
convexity of $\mathcal P_\beta$ along segments of distribution
functions~\cite{AuffingerChen2015Uniqueness}. This multiplier equals one
on $\supp\mu_\beta$, by an identity of
Talagrand~\cite{Talagrand2006ParisiMeasures}, in the form stated
in~\cite[Theorem~5]{AuffingerChen2015Properties}
and~\cite[Lemma~2.6]{Lopatto}, and by the interval
support~\cite[Theorem~1.1]{Lopatto}. Above the support, we bound it by one
using the curvature of the forward density and its heat evolution
from~\cite[Proposition~4.1 and Lemma~2.4(ii)]{Lopatto}.

\paragraph{Our related work.}
This paper continues~\cite{AL}, where we proved the quantitative Parisi
formulas recalled in \cref{sec:intro-finite-size}, the variance bounds
above, and upper-tail estimates for $F_{N,\beta}$ with exponent
$6/5$~\cite[Theorem~1.5]{AL}. The interval-support theorem used in both
papers was proved by the second author~\cite{Lopatto}.
In~\cite{AronowLopatto2026EA}, we studied a different approximation of
the SK pressure. On $d$-dimensional tori of side length $L\ge3$ and on
hypercubes, with independent Gaussian couplings of variance $1/(2d)$
and $1/d$, respectively, the expected pressure per site converges to the
SK pressure as $d\to\infty$, for every fixed inverse temperature
$\beta\ge0$ and every uniform field $h\in\R$; on tori, the convergence
is uniform in the side length. The present paper concerns the
finite-size bias of the SK model itself.

\section{Comparison with the OpenAI preprints}\label{subsec:intro-openai}

We obtained the exponent $7/10$ on September~23, 2026, following a
suggestion from Silvio Franz on September~15, 2026. The OpenAI
mathematics repository~\cite{OpenAIMathRepository} contains
\emph{The low-temperature Sherrington--Kirkpatrick fluctuation
scale}~\cite{OpenAIFluctuationScale} and its companion \emph{The
low-temperature Sherrington--Kirkpatrick free-energy limiting
law}~\cite{OpenAILimitingLaw}, both dated September~24, 2026. We refer
to their initial repository version.

\paragraph{Stated results.}
For every fixed $\beta>1$, \cite[Theorem~1.1]{OpenAIFluctuationScale}
states that $\Var(F_{N,\beta})=N^{1/3+o(1)}$. The companion paper states
that $N^{-1/3}\Var(F_{N,\beta})$ converges to a positive limit and that
$F_{N,\beta}$, standardized by its exact mean and variance, converges in
law to a nondegenerate limit~\cite[Theorems~1.1
and~1.2]{OpenAILimitingLaw}. These results concern fluctuations about
the finite-size mean $\E F_{N,\beta}$, and they do not by themselves
estimate $P_\beta-p_{N,\beta}$. We do not use these results, or the
comparison discussed next, in the proof of \cref{thm:main}.

\paragraph{A conditional upper bound.}
An improvement of the exponent $2/3$ for the mean bias is implicit in an
uncentered comparison stated in~\cite{OpenAILimitingLaw}. We use the
notation of that paper. Let $T=\beta^2$ and $x=N^{-1/6}$, and fix an
admissible choice of its auxiliary exponents and reference parameters.
For its direct, unmasked, and unseeded reference paths,
\cite[Proposition~6.1, (G1)]{OpenAILimitingLaw} states that there
exists $c_G>0$ such that
\begin{equation}\label{eq:intro-openai-G1}
 0\le F_*(\lambda)-f_N(b_\lambda,h_{*,\lambda})\le x^{4+c_G}
\end{equation}
for all sufficiently large $N$. Here $F_*$ is a scalar trial objective,
$f_N$ is a self-corrected hierarchical pressure evaluated at the clocks
$(b_\lambda,h_{*,\lambda})$, and $\lambda$ is a perturbation strength
that is unrelated to our tilt. We take terminal mass one and the
allowed strength $\lambda=x^{12}$. From the definitions
in~\cite[Section~3.1]{OpenAILimitingLaw}, one checks that the scalar
objective at $\lambda=0$ equals $P_\beta-\log2-T/4$. Using the perturbation estimates for the direct
scalar lift~\cite[Section~3.2, proof of Lemma~3.2]{OpenAILimitingLaw}
and the construction of the clocks~\cite[Section~3.4,
(P3)--(P4)]{OpenAILimitingLaw}, we find
\[
 \|b_\lambda-T\|_\infty+\|h_{*,\lambda}\|_\infty\le C_\beta\lambda,
 \qquad
 \Bigl|F_*(\lambda)-\Bigl(P_\beta-\log2-\frac T4\Bigr)\Bigr|
 \le C_\beta\lambda.
\]
With the uniform spin prior and centered Gaussian completion
of~\cite[Section~2.1]{OpenAILimitingLaw}, the unperturbed clocks satisfy
$f_N(T,0)=p_{N,\beta}-\log2-T/4$. Integrating the pressure variation
formula~\cite[Lemma~2.1, (A1)]{OpenAILimitingLaw} along the segment
between the two pairs of clocks, and using that its integrands lie in
$[0,1]$ at terminal mass one, we obtain
\[
 \bigl|f_N(b_\lambda,h_{*,\lambda})-f_N(T,0)\bigr|
 \le\frac14\|b_\lambda-T\|_1+\frac12\|h_{*,\lambda}\|_1
 \le C_\beta\lambda.
\]
We write $P_\beta-p_{N,\beta}$ as the sum of
$P_\beta-\log2-T/4-F_*(\lambda)$,
$F_*(\lambda)-f_N(b_\lambda,h_{*,\lambda})$, and
$f_N(b_\lambda,h_{*,\lambda})-f_N(T,0)$. We bound the first and last
differences by $C_\beta\lambda$ and the middle one by
\eqref{eq:intro-openai-G1}. Since $x^{4+c_G}=N^{-2/3-c_G/6}$ and
$x^{12}=N^{-2}$, we obtain
\begin{equation}\label{eq:intro-openai-bias}
 0\le P_\beta-p_{N,\beta}
 \le x^{4+c_G}+C_\beta x^{12}
 \le C_\beta N^{-2/3-\delta_\beta},\qquad
 \delta_\beta=c_G/6>0,
\end{equation}
after decreasing $c_G$ so that $c_G<8$ if necessary. The lower bound is
Guerra's, and \eqref{eq:intro-openai-bias} holds for the finitely many
remaining $N\ge2$ after we enlarge $C_\beta$. This bound follows from
\eqref{eq:intro-openai-G1} and the perturbation estimates above; it is
not a theorem on the mean bias stated in either paper. The proof
of~\cite[Proposition~6.1]{OpenAILimitingLaw} leaves $c_G$ unspecified,
and we do not extract from it the saving $c_G\ge1/5$ that the exponent
$7/10$ would require. The budgets of order $x^{5+c}$
in~\cite[Proposition~4.1]{OpenAILimitingLaw} concern increments of a
centered pressure~\cite[(A4)]{OpenAILimitingLaw}, from which the
finite-size mean has been subtracted, so they do not bound the mean
bias.

\paragraph{Conditional lower bounds.}
In the other direction, the variance--mean-deficit inequality
of~\cite[Lemma~7.4 and Remark~7.14, (7.68)]{AL} states that for every
$N\ge1$,
\[
 P_\beta-p_{N,\beta}\ge c_\beta
 \Bigl(\frac{\Var(F_{N,\beta})+\beta^2/2}{N}\Bigr)^{5/4}.
\]
Inserting $\Var(F_{N,\beta})\ge N^{1/3-4\varepsilon/5}$, we see that the
variance exponent stated in~\cite[Theorem~1.1]{OpenAIFluctuationScale}
would imply, for every fixed $\varepsilon>0$, that
$P_\beta-p_{N,\beta}\ge c_{\beta,\varepsilon}N^{-5/6-\varepsilon}$ for
all sufficiently large $N$. The positive limit of
$N^{-1/3}\Var(F_{N,\beta})$ stated
in~\cite[Theorem~1.2]{OpenAILimitingLaw} would give the stronger bound
\begin{equation}\label{eq:intro-openai-lower}
 P_\beta-p_{N,\beta}\ge c_\beta N^{-5/6}
 \qquad\text{for all sufficiently large }N.
\end{equation}
For \eqref{eq:intro-openai-lower}, an eventual lower bound
$\Var(F_{N,\beta})\ge k_\beta N^{1/3}$ with $k_\beta>0$ would suffice.
Combined with \cref{thm:main}, the variance exponent
of~\cite{OpenAIFluctuationScale} would imply that if
$P_\beta-p_{N,\beta}$ is of exact order $N^{-a}$, then $7/10\le a\le5/6$.

\paragraph{Verification status.}
We have checked the statements quoted above and the conditional
deductions from them. We have not independently reconstructed complete
proofs of the global comparison and fluctuation estimates
of~\cite{OpenAIFluctuationScale,OpenAILimitingLaw}, and we do not
assume their conclusions. The repository also contains Lean sources on
the fluctuations of the SK model. Their main theorem on the fluctuation
exponents, as we inspected it, takes as an explicit hypothesis a
two-sided bound on the centered logarithmic Laplace transform of the
completed free energy, \texttt{PositiveLaplaceWindow}; in the
fluctuation-scale paper, this bound is the conclusion
of~\cite[Proposition~7.5]{OpenAIFluctuationScale}. The conclusion of
that theorem concerns the fluctuation exponents, not a positive variance
prefactor or a limiting law.\footnote{We inspected the 151 files of the
directory \nolinkurl{lean/OAI/Probability/SKFluctuations} at commit
\nolinkurl{adc7f1241b42e322a6451854ab7e4b4c146bf78a}. The theorem
\nolinkurl{main_of_positiveLaplaceWindow}, at line~294 of
\nolinkurl{EventuallyNatPos.lean}, takes a hypothesis
\texttt{hw} of type \nolinkurl{PositiveLaplaceWindow}, which is defined at
lines~453--459 of \nolinkurl{Spin.lean}; the conclusion of the
theorem, \nolinkurl{MainFluctuationStatement}, is defined at
lines~531--540 of the same file. We found no proof supplying the
hypothesis in this directory. The directory was unchanged at commit
\nolinkurl{fd4aeeb2ee4fc729c18d98444fed42fd0529eeeb}, which we checked on
October~9, 2026. This was an inspection of the sources, not a
compilation.} Neither this conditional formal statement nor our
incomplete reconstruction shows that the stated results are false.
Our inspection did not identify an invalid step in the portions
examined.

\paragraph{Proof techniques.}
The limiting-law paper contains a one-site cavity expansion whose
structure is similar to that of \cref{part:cavity}.
In~\cite[Lemma~2.2]{OpenAILimitingLaw}, one site is chosen by
exchangeability and removed, and its Gaussian kernels are interpolated
to those of an independent scalar model on one site; the linear cavity
kernel is half of a connected four-spin correlation of the removed
spin~\cite[(A8)]{OpenAILimitingLaw}. In the proof
of~\cite[Proposition~6.1]{OpenAILimitingLaw}, the resulting linear
equation is inverted under a gap condition on the scalar
response~\cite[Proposition~5.1, (H1)]{OpenAILimitingLaw},
and~\cite[Lemma~5.2, (H2)]{OpenAILimitingLaw} writes the response on
the additive sector as a susceptibility multiplier minus a positive
semidefinite Hessian kernel. In our argument, the cavity inequality
contracts because the scalar response operator is at most $1+\eps_N$
with $\eps_N\le\rho/4$, while the interpolation contributes the factor
$1-\rho$. We derive the bound on the response operator from the second
variation of the Parisi functional under transports of levels
(\cref{lem:Jdom}); this operator acts on a space of transports, not on
the additive sector of~\cite{OpenAILimitingLaw}. The fluctuation-scale
paper also contains a one-site closure for the mean
overlap~\cite[Proposition~4.1]{OpenAIFluctuationScale}. We use two
tools from that paper. \Cref{part:cavity} relies on normalized replica
trees and covariance insertions of the type developed
in~\cite[Section~2]{OpenAIFluctuationScale} and used
in~\cite[Sections~2.1--2.3]{OpenAILimitingLaw}. In
\cref{lem:fnd-posterior}, we adapt the joint-source Brascamp--Lieb
argument of the proof of~\cite[Lemma~2.3]{OpenAIFluctuationScale}, which
is also used in the proof of~\cite[Lemma~2.1]{OpenAILimitingLaw}, to
squared conditional means on selected coordinates and to anisotropic
laws.



\section{Organization and conventions}\label{sec:intro-organization}

\Cref{sec:overview} contains an overview of the proof and a map of its
dependencies. The proof itself occupies five parts.
\begin{itemize}
\item \Cref{part:interpolation} concerns the interpolation.
\Cref{sec:fnd} recalls the interpolation of~\cite{AL}, its tilted law,
and the decomposition of the defect, and reduces the mean bias to
the fractional-moment error. \Cref{sec:gauss} proves two Gaussian tools
for tilted laws defined by arbitrary profiles, which we use in
\cref{part:profile,part:cavity}.
\item \Cref{part:overlap} contains overlap estimates at a minimizing
profile. \Cref{sec:wi} estimates moments of the self-overlap on
quantile intervals, and \cref{sec:th} bounds the mean--covariance term.
\item \Cref{part:profile} concerns the minimizing profile and the scalar
response. \Cref{sec:lz} shows that the minimizing profile is close in
$L^1$ to the scalar trial profile, and \cref{sec:lz-root} localizes its
lowest quantiles. \Cref{sec:sc} studies the scalar response operator and
proves its stability at the minimizing profile.
\item \Cref{part:cavity} contains the cavity argument. \Cref{sec:c7}
sets up the cavity interpolation and its tree laws,
\cref{sec:c7-endpoint} computes the derivative at the independent
endpoint, and \cref{sec:c7-closure} closes the cavity inequality and
bounds the defect.
\item In \cref{part:proof}, \cref{sec:b7} integrates these estimates and
proves \cref{thm:main}.
\end{itemize}

We rely on~\cite{AL} for the interpolation identities and the scalar
floor estimates, and on~\cite{Lopatto} and~\cite{AuffingerChen2015Properties}
for the structure and regularity of $\mu_\beta$. We restate these
results where they are used. Locators in~\cite{AL} refer to its arXiv
version~1, locators in~\cite{Lopatto} to its arXiv version~4, and
locators in the OpenAI manuscripts to their initial repository version.

Throughout, $\beta>1$ is fixed, and $C_\beta$ denotes a constant
depending only on $\beta$ whose value may change from line to line; a
numerical constant does not depend on any parameter. All profiles defining a tilted law are
deterministic: minimizers may depend on $N$, $\rho$, and $\lambda$ but
not on the couplings. For symmetric matrices or self-adjoint operators,
$B\preceq C$ means that $C-B$ is positive semidefinite. We write
$\log_+x=\log(\max\{1,x\})$ for $x\ge0$. Distribution functions, such
as $\alpha$ and the profiles $a$ of \cref{sec:fnd}, are
right-continuous. Atoms of measures are retained in Stieltjes
integrals, and an interval $(q,1]$ excludes an atom at $q$.

\chapter{Overview of the proof}\label{sec:overview}

In this section, we explain the proof of \cref{thm:main}, which
occupies \crefrange{sec:fnd}{sec:b7}. For each step, we state its main
estimate in the form in which it is proved, give the precise reference,
and explain briefly why the estimate holds and where it is used.
\Cref{sec:overview-map} contains a map of the proof, and
\cref{sec:overview-profiles} discusses three conventions that recur in
every part of it. Throughout, $\beta>1$ is fixed. We write
$X\lesssim Y$ if $X\le C_\beta Y$ for a constant $C_\beta$ that depends
only on $\beta$. Our estimates are uniform in the system size $N$, the interpolation
parameter $\rho$, the tilt $\lambda$, the minimizing profile, and the
positions and masses of its atoms. Much of the technical work goes into
this uniformity.

\section{The structure of the argument}\label{subsec:intro-outline}

We follow the interpolation scheme of~\cite{AL}; what changes is how
we estimate the self-overlap variance along the interpolation. The argument has four steps.

First, we reduce \cref{thm:main} to a bound on a fractional-moment
error $e_{N,\beta}(\lambda)$, which compares a normalized logarithmic
moment of order $\lambda/\beta$ of the partition function with a Parisi
problem constrained to have mass at least $\lambda/\beta$ at zero. By
\cref{prop:fnd-mean-bias}, the mean bias divided by $\beta$ is at most
$2e_{N,\beta}(\lambda)$ plus a term of order $\lambda^5$, and it
suffices to prove that $e_{N,\beta}(N^{-7/50})\lesssim N^{-7/10}$
(\cref{sec:overview-reduction}).

Second, we write $4e_{N,\beta}(\lambda)$ as the integral over
$\rho\in(0,1]$ of a defect $\calD_N(\rho)=\beta(A+2J+V)$ along an
interpolation, evaluated at a minimizing profile of a variational
problem in dimension $N$ (\cref{sec:overview-interpolation}). The terms
$A$ and $J$ are at most of order $\log^2(2+N)/(N\rho\lambda)$
(\cref{sec:overview-thermal}), and the difficulty is the self-overlap
variance $V$.

Third, we bound $V$ in two regimes. Near $\rho=0$, we combine estimates
on quantile intervals of the minimizing profile with a localization of
its lowest quantiles (\cref{sec:overview-shells}). For
$\rho$ above a cutoff $\rho_0$, we bound $V$ by a one-site cavity argument
(\cref{sec:overview-cavity}), which contracts if a scalar response
operator at the minimizing profile is at most $1+\rho/4$. We prove this
bound by showing that the operator is at most the identity at the Parisi
measure (\cref{sec:overview-stability}) and that the minimizing profile
is close to a scalar trial profile (\cref{sec:overview-localization}).
The two regimes give the following bounds. For almost every
$\rho\in[\rho_0,1]$, we have
\begin{equation}\label{eq:ov-stable}
 \calD_N(\rho)\lesssim\frac{\log^2(2+N)}{N\lambda\rho^2}
\end{equation}
by \cref{cor:c7-defect}, and for almost every $\rho\in(0,1]$, we have
\begin{equation}\label{eq:ov-small}
 \calD_N(\rho)\lesssim\frac{\log^2(2+N)}{N\rho\lambda}+(N\rho)^{-2/3}
 +(N\rho)^{-2/3}\lambda^{-1/3}\omega(\rho,\eta)^{2/3}
\end{equation}
by \cref{prop:sm-small}. Here $\eta=e_{N,\beta}(\lambda)+L_s\lambda^5$
is the trial budget of \cref{lem:fnd-trial}, with $L_s\ge1$ depending
only on $\beta$, and $\omega(\rho,\eta)=\min\{1,(\eta/\rho)^{1/6}\}$.

Fourth, the cutoff is $\rho_0=\min\{1,K\eta^{1/7}\}$, with $K$
depending only on $\beta$. Note that $\rho_0$ is defined through the
trial budget, which contains the unknown error. Integrating
\eqref{eq:ov-stable} over $[\rho_0,1]$ and \eqref{eq:ov-small} over
$(0,\rho_0]$, we obtain an implicit inequality whose leading term is
$N^{-2/3}\eta^{1/21}$. At $\lambda=N^{-7/50}$, this inequality
implies that $e_{N,\beta}(\lambda)\lesssim N^{-7/10}$
(\cref{sec:overview-closure}).

We take the first two steps from~\cite{AL}, where the same scheme
gives the rate $N^{-2/3}$~\cite[Theorem~1.4]{AL}, and our estimates
near $\rho=0$ refine those of~\cite[Sections~3 and~7]{AL}. We use the
cavity argument, and the stability and isolation estimates that it
requires, only above the cutoff.

\section{Reduction to the fractional-moment error}
\label{sec:overview-reduction}

As in~\cite[Section~1.1]{AL}, we complete the Hamiltonian by an
independent diagonal Gaussian term $D_N$, so that the completed
Hamiltonian has covariance exactly $NR^2/2$, where $R$ is the overlap
(\cref{sec:fnd-completion}). For $0<\lambda<\beta$, the
fractional-moment error is
\[
 e_{N,\beta}(\lambda)=E_\beta(\lambda)-\psi(\lambda),\qquad
 \psi(\lambda)=\frac1{N\lambda}\log\E_{\rm ref}
 e^{\lambda(F_{N,\beta}/\beta+D_N)},
\]
where $E_\beta(\lambda)=\beta^{-1}P_\beta(\lambda/\beta)$, and
$P_\beta(s)$ is the minimum of the Parisi functional over probability
measures $\mu$ on $[0,1]$ with $\mu(\{0\})\ge s$
\eqref{eq:fnd-error-definition}. Talagrand showed that the normalized
moments converge to this constrained
value~\cite{Talagrand2007FractionalMoments}; we need a rate.

The reduction is the midpoint argument of~\cite[Theorem~7.8,
Step~1]{AL}. Since $F_{N,\beta}/\beta+D_N$ is a convex Lipschitz function
of the Gaussian variables $(g_{ij},z_i)$, Chen's convexity
theorem~\cite[Theorem~1]{ChenConvexity} implies that $\psi$ extends to
a convex function on $[0,\infty)$ with
$\psi(0)=p_{N,\beta}/\beta$~\cite[Lemma~5.2]{AL}. Using
$\psi(0)\ge2\psi(\lambda)-\psi(2\lambda)$ and
$e_{N,\beta}(2\lambda)\ge0$, we find that
\[
 \frac{P_\beta-p_{N,\beta}}\beta
 \le E_\beta(0)-2E_\beta(\lambda)+E_\beta(2\lambda)
 +2e_{N,\beta}(\lambda).
\]
The second difference of $E_\beta$ is controlled by the cost of the
prescribed mass at zero. By~\cite[Lemma~7.2]{AL}, the Parisi profile
$\gamma_\beta=\beta\alpha_\beta$ grows linearly near $0$. As a consequence,
the scalar trial profile $\gamma^{\rm tr}_\lambda=\max\{\lambda,\gamma_\beta\}$
differs from it only on an interval of length of order $\lambda$, on
which the first variation vanishes at $\gamma_\beta$. The cost of the
trial profile is at most $L_s\lambda^5$
\eqref{eq:fnd-scalar-trial-cost}. Combining these facts, we arrive at
the bound of \cref{prop:fnd-mean-bias},
\[
 0\le\frac{P_\beta-p_{N,\beta}}\beta\le32L_s\lambda^5
 +2e_{N,\beta}(\lambda)\qquad(0<\lambda\le\lambda_0),
\]
with $\lambda_0>0$ depending only on $\beta$. We take
$\lambda=N^{-7/50}$ for $N$ so large that $\lambda\le\lambda_0$, and
then $\lambda^5=N^{-7/10}$. The tilt must be positive, because our
bounds on the defect degenerate as $\lambda\downarrow0$; we fix its exponent
in \cref{sec:overview-closure} by matching $\lambda^5$ with the
final bound on the error.

\section{The interpolation and the defect}
\label{sec:overview-interpolation}

The interpolation of~\cite[Section~2]{AL} is a finite-dimensional,
quantitative version of Guerra's interpolation~\cite{Guerra2003}, in the
form that Talagrand used for positive fractional
moments~\cite{Talagrand2007FractionalMoments}. For $0<\rho\le1$, let
$t=1-\rho$, $c=\beta^2\rho$, and $M=Nc$. An admissible profile is a
deterministic probability measure $\mu=s\delta_0+\nu$ on $[0,1]$, where
$s=\lambda/\beta$ and $\nu\ge0$, with distribution function
$a(q)=\mu([0,q])$. We solve the Parisi equation in $\R^N$,
\[
 \partial_qU=-\frac c2\bigl(\Delta U+a(q)|\nabla U|^2\bigr),\qquad
 U(1,x;W)=\log\sum_{\sigma\in\{-1,1\}^N}
 e^{\beta\sqrt t\,\sigma^TW\sigma/2+x\cdot\sigma},
\]
and we minimize the functional $\mathcal F_{t,s}(\mu)$ of
\eqref{eq:fnd-functional} over admissible profiles. At $t=0$, the
equation separates into $N$ scalar Parisi equations, and the minimum is
$P_\beta(s)$; at $t=1$, the minimum is the completed fractional moment
\eqref{eq:fnd-endpoints}. The fractional-moment error is
$\beta^{-1}$ times the decrease of the minimum along the interpolation.

\paragraph{The tilted law.}
Each profile defines a joint law of the couplings, an auxiliary field
$X_q\in\R^N$, $q\in[0,1]$, and a terminal spin $\sigma$. In this law, $W$ is
tilted by $e^{sU(0,0;W)}$, the field is a diffusion with drift
$ca(q)\nabla U(q,X_q;W)$ and noise coefficient $\sqrt c$, and $\sigma$ has
the Gibbs law in the external field $X_1$ (\cref{subsec:fnd-law}). We write $\calG_q$ for the
information revealed up to level $q$, $m_q$ and $C_q$ for the
conditional mean and covariance of $\sigma$ given $\calG_q$,
$H_q=\nabla^2U(q,X_q;W)$, and $S_q=|m_q|^2/N$. By
\cref{prop:fnd-contact}, $m_q$ is a martingale with
$\dd m_q=\sqrt c\,H_q\dd B_q$, and $a(q)C_q\preceq H_q\preceq C_q$. If
$\mu$ is a minimizer, every point $q$ of $\supp\nu$, called a contact
level, satisfies $\E S_q=q$ and $\E\Tr H_q^2\le N/c$, which are the
counterparts of the optimality conditions for the Parisi measure.

\paragraph{The defect.}
Let $R_q$ be the overlap of two replicas that are conditionally
independent given $\calG_q$. By \cref{prop:fnd-interpolation}, the
derivative of $\mathcal F_{t,s}(\mu)$ in $t$ is
$-(\beta^2/4)\int\E(R_q-q)^2\,\nu(\dd q)$ for every profile. It is a
square because the completed Hamiltonian has covariance $NR^2/2$, and
only $\nu$ appears because the outer moment of order $s$ cancels the
prescribed mass $s\delta_0$. At a contact level,
$R_q-S_q$ splits into three conditionally orthogonal pieces, which
gives $\E(R_q-q)^2=a^{\rm th}_q+2j_q+\Var(S_q)$ (\cref{lem:fnd-decomp}).
Integrating against $\nu$, we obtain \cref{cor:fnd-defect}, which
states that at almost every $\rho$ and at every minimizer,
\[
 \calD_N(\rho)=-\frac4\beta f_{N,s}'(1-\rho)=\beta(A+2J+V),
 \qquad 4e_{N,\beta}(\lambda)=\int_0^1\calD_N(\rho)\dd\rho,
\]
where $f_{N,s}(t)$ is the minimum of $\mathcal F_{t,s}$ and
\[
 A=\int\frac{\E\Tr C_q^2}{N^2}\,\nu(\dd q),\qquad
 J=\int\frac{\E[m_q^TC_qm_q]}{N^2}\,\nu(\dd q),\qquad
 V=\int\Var(S_q)\,\nu(\dd q).
\]
The terms $A$ and $J$ measure conditional fluctuations of the terminal
spin given $\calG_q$, while $V$ measures fluctuations of the information
that the field has revealed by level $q$. Each term may depend on the
minimizer, but their combination does not. We bound all three at an
arbitrary minimizer.

\paragraph{The posterior variance bound.}
Let $L_q=(\sigma\cdot m_q-|m_q|^2/2)/N$, whose conditional mean given
$\calG_q$ is $S_q/2$. By \cref{prop:fnd-variance}, we have, at a
minimizer and for every deterministic $k\in L^2(\nu)$,
\[
 \Var\Bigl(\int k(q)L_q\,\nu(\dd q)\Bigr)
 \le\frac1M\int\frac{k(q)^2}{a(q)}\,\nu(\dd q),
 \qquad A\le\frac1{Ms}.
\]
With the terminal spin fixed, the Gaussian coordinates that generate
$W$ and the field have a uniformly log-concave law, and the bound
follows from the Brascamp--Lieb inequality~\cite{BrascampLieb1976}, in
the form of~\cite[Proposition~2.2]{AL}, and the contact identities.
Note that the factor $a(q)^{-1}\le s^{-1}$ is where the prescribed mass
enters. We integrate such factors through the quantile function $Q$ of
$\nu$ \eqref{eq:fnd-quantile}, which pushes Lebesgue measure on
$(0,1-s)$ forward to $\nu$ and satisfies $a(Q(v))\ge s+v$.

\paragraph{The trial budget.}
In the energy normalization, we set $\gamma=\beta a$ and
$F_\rho(\gamma)=\beta^{-1}\mathcal F_{1-\rho,s}(\mu)$. By
\cref{lem:fnd-trial}, we have, for every $\rho\in(0,1]$,
\[
 0\le F_\rho(\gamma^{\rm tr}_\lambda)-\min_\gamma F_\rho(\gamma)
 \le e_{N,\beta}(\lambda)+L_s\lambda^5=\eta .
\]
Indeed, $\mathcal F_{t,s}(\mu)$ and its minimum are nonincreasing in
$t$, so that $F_\rho(\gamma^{\rm tr}_\lambda)$ is at most the scalar
value of the trial profile, which is at most $E_\beta(0)+L_s\lambda^5$,
while $\min F_\rho\ge\psi(\lambda)=E_\beta(\lambda)-e_{N,\beta}(\lambda)$
and $E_\beta(\lambda)\ge E_\beta(0)$. In words, the scalar trial profile has value at most $\eta$ above the minimum at every $\rho$. This budget controls
how far a minimizer can be from the trial profile
(\cref{sec:lz,sec:lz-root}), and it fixes the cutoff $\rho_0$
(\cref{sec:b7}). Since it contains the unknown error, the final inequality
is implicit.

\section{The thermal terms}\label{sec:overview-thermal}

At a contact level, by the covariance order, we have
$\Tr H_q^2\ge a(q)\Tr(H_qC_q)\ge a(q)^2\Tr C_q^2$, and the contact bound
on $\E\Tr H_q^2$ implies that $\E\Tr C_q^2/N^2\le(Ma(q)^2)^{-1}$. In
the quantile variable, we have
\[
 A\le\frac1M\int_0^{1-s}\frac{\dd v}{a(Q(v))^2}
 \le\frac1M\int_0^{1-s}\frac{\dd v}{(s+v)^2}\le\frac1{Ms}.
\]

The mean--covariance term needs more work, because
$j_q=N^{-2}\E[m_q^TC_qm_q]$ involves the conditional covariance in the
direction $m_q$, which is revealed at level $q$ itself. For a
direction $b$ that is measurable with respect to $\calG_{q'}$ for some
$q'<q$, we show by a scalar Riccati inequality that $\E[b^TC_qb]$ is at most
$\E|b|^2$ divided by the accumulated clock $c\int_{q'}^qa(u)^2\dd u$
(\cref{lem:th-ric}). However, this bound does not apply to $b=m_q$. To
overcome this obstacle, we write $m_{Q(v)}$ as a weighted average of
the earlier means $m_{Q(w)}$, $w<v$, plus increments. First-order
optimality implies that $\E S_r$ is $2$-Lipschitz toward every point
$q$ of $\supp\nu$, which bounds $\E|m_q-m_r|^2$ by $2N|q-r|$
(\cref{lem:th-gamma}). We control the increments by combining this
bound with \cref{lem:th-ric} in a dyadic decomposition (\cref{lem:th-res}).
The conditional covariance in the direction of the average is
controlled by the posterior variance bound. With a quadratic weight
$(s+w)^2$, both contributions are at most of order $(s+v)^{-2}$. We obtain
\[
 j_{Q(v)}\le\frac{341(1+\log_+M)^2}{M(s+v)^2},\qquad
 J\le\frac{341(1+\log_+M)^2}{Ms}
\]
(\cref{prop:th-j}).
Since $Ms=\beta N\rho\lambda$ and $M\le\beta^2N$, we conclude that
\[
 \beta(A+2J)\le\frac{1+682(1+\log_+M)^2}{N\rho\lambda}
 \lesssim\frac{\log^2(2+N)}{N\rho\lambda}.
\]
Integrated over $\rho\in[1/N,1]$, this is at most of order
$\log^3(2+N)/(N\lambda)$, which equals $N^{-43/50}\log^3(2+N)$ at
$\lambda=N^{-7/50}$ and is smaller than $N^{-7/10}$ by a power of $N$.

\section{The self-overlap variance for small \texorpdfstring{$\rho$}{rho}}
\label{sec:overview-shells}

The posterior variance bound controls averages of $L_q$ over levels.
However, $V$ involves $S_q$ at single levels, which may carry little
mass of $\nu$. Following~\cite[Lemma~3.3]{AL}, we compare $S_q$ with
averages of $L_p$ over windows of quantile levels next to $q$
(\cref{sec:wi}).

\paragraph{Two comparisons.}
For $p\ge q$, the conditional Jensen inequality for the martingale $m$
implies that $S_q\le2\E[L_p\mid\calG_q]$, and for $p\le q$, we have
$\E[L_p\mid\calG_q]=S_q/2-|m_q-m_p|^2/(2N)$. Averaging over a window
of quantile levels after or before $q$, with mean level $\bar p$, we
show in \cref{lem:wi-two} that, at a minimizer, $S_q-q$ is at most
(for a window after $q$) or at least (for a window before $q$) the sum
of $\bar p-q$ and twice the conditional expectation of a centered
average of $L_p$. If the window has quantile length $r$ and $a\ge u$ on
it, this conditional expectation has second moment at most
$(Mur)^{-1}$. Note that conditioning on $\calG_q$ removes the centered part
$N^{-1}(\sigma-m_q)\cdot m_q$ of $L_q$, which in~\cite[Lemma~3.3]{AL}
produces additional residual terms; we bound its variance $j_q$
separately (\cref{sec:overview-thermal}).

\paragraph{Variance on a quantile interval.}
Let $I\subset(0,1-s)$ be an interval of length $h$ such that
$a(Q(v))\ge u$ for $v\in I$, let $\ell$ be the diameter of $Q(I)$, and
set $\epsilon=(Mu)^{-1}$. For a threshold $\tau$, we take windows of
radius of order $\epsilon/\tau$, whose random parts have second moment
less than $\tau$. A point whose two windows lie in $I$ and for which
$\Var(S_{Q(v)})>\tau$ must then see an
increment of $Q$ at least of order $\tau^{1/2}$ across its window, and since $Q$
is nondecreasing with total increase $\ell$ on $I$, such points form a
set of measure at most of order $\epsilon\ell\tau^{-3/2}$. The remaining
points lie within distance of order $\epsilon/\tau$ of the ends of $I$. Integrating
$\min\{h,\epsilon\ell\tau^{-3/2}\}$ over $\tau$, we find the main term
in the bound of \cref{prop:wi-vfs},
\begin{equation}\label{eq:ov-vfs}
 \int_I\Var(S_{Q(v)})\dd v
 \le128\Bigl[\min\bigl\{h,\epsilon\bigl(1+\log_+(h/\epsilon)\bigr)\bigr\}
 +(\epsilon\ell)^{2/3}h^{1/3}\Bigr].
\end{equation}
The first term is at most $\epsilon$ up to a logarithm, and the
interval may contain part of the quantile interval of an atom.

\paragraph{Dyadic shells and root localization.}
In the proof of \cref{prop:sm-small}, we cover $(0,1-s)$ by dyadic
shells $I_u=\{v:u\le s+v<2u\}$, $u=2^js<1$, on which $a(Q(v))\ge u$,
$h\le u$, and $\epsilon=(Mu)^{-1}$. If we only use $\ell\le1$, the second
term in \eqref{eq:ov-vfs} is at most $M^{-2/3}u^{-1/3}$, and summing over
shells, with $\beta M^{-2/3}s^{-1/3}=(N\rho)^{-2/3}\lambda^{-1/3}$, we
would find only
\[
 \beta V\lesssim(N\rho)^{-2/3}\lambda^{-1/3}
 +\frac{\log^2(2+N)}{N\rho\lambda}.
\]
The bound $\beta V\le C(N\rho)^{-2/3}\lambda^{-1/3}$ that follows
from~\cite[(4.2)]{AL} has the same form. At $\lambda=N^{-7/50}$, the
factor $\lambda^{-1/3}=N^{7/150}$ would leave a positive power of $N$
after the normalization of \cref{sec:overview-closure}. To remove it, we
use the root localization of \cref{cor:lz-omega}, which states that at
every minimizer and for every trial budget $\eta$,
\[
 Q(v)\le C_Q\bigl(s+v+\omega(\rho,\eta)\bigr),\qquad
 \omega(\rho,\eta)=\min\{1,(\eta/\rho)^{1/6}\},
\]
with $C_Q$ depending only on $\beta$. Then
$\ell(I_u)\le2C_Q(u+\omega)$, the second term in \eqref{eq:ov-vfs} is
at most of order $M^{-2/3}(u^{1/3}+u^{-1/3}\omega^{2/3})$,
and its sum over shells is at most of order $M^{-2/3}(1+s^{-1/3}\omega^{2/3})$.
Together with the thermal terms, this proves \eqref{eq:ov-small}, that
is, \eqref{eq:sm-pointwise}. Its leading term $(N\rho)^{-2/3}$ does not
depend on $\lambda$, and the correction is integrable at $\rho=0$,
because $\omega^{2/3}\le(\eta/\rho)^{1/9}$. Below $\rho=1/N$, we rely
instead on the bound $\calD_N(\rho)\le C(N\rho)^{-1/2}$
of~\cite[Lemma~4.1]{AL}, recalled in \eqref{eq:fnd-moving}. Integrating,
we arrive at \eqref{eq:sm-integrated}, which states that for
$0<\rho_0\le1$,
\begin{equation}\label{eq:ov-small-integrated}
 \int_0^{\rho_0}\calD_N(\rho)\dd\rho
 \lesssim N^{-1}+\frac{\log^3(2+N)}{N\lambda}
 +N^{-2/3}\rho_0^{1/3}
 +N^{-2/3}\lambda^{-1/3}\eta^{1/9}\rho_0^{2/9}.
\end{equation}

\paragraph{\texorpdfstring{The barrier at $2/3$.}{The barrier at 2/3.}}
For $\rho_0=1$, the third term in \eqref{eq:ov-small-integrated} is of
order $N^{-2/3}$. This is the mechanism of~\cite[Proposition~7.7]{AL},
in which quantile estimates~\cite[Lemma~3.3]{AL} and a localization of
low quantiles~\cite[Lemma~7.6]{AL} lead to
$e_{N,\beta}(\lambda)\lesssim N^{-2/3}+N^{-3/4}\lambda^{-1/2}
+\log(2+N)/(N\lambda)$~\cite[(7.43)]{AL}, and the tilt
$\lambda=N^{-2/15}$ gives the rate $N^{-2/3}$
of~\cite[Theorem~1.4]{AL}. When $\rho$ is of order one, however,
$(N\rho)^{-2/3}$ is much larger than the cavity bound
\eqref{eq:ov-stable}, as we discuss heuristically in
\cref{subsec:intro-physics}. We therefore use
\eqref{eq:ov-small-integrated} only below the cutoff, where it
contributes $N^{-2/3}\rho_0^{1/3}$.

\section{The cavity contraction}\label{sec:overview-cavity}

Let $\mu$ be a minimizer at $(\rho,s)$ with distribution function $a$,
let $d=\beta^2(1-t/N)$, and let $J^{\scal}_{a,d}$ be the scalar response
operator of \cref{sec:overview-stability}, which acts on a subspace
$\mathcal H_\nu$ of $L^2(\nu)$. In \cref{prop:c7-closure}, we show that
if
\begin{equation}\label{eq:ov-stability}
 J^{\scal}_{a,d}\preceq(1+\eps_N)I\quad\text{on }\mathcal H_\nu,
 \qquad 0\le\eps_N\le\frac\rho4,
\end{equation}
then
\[
 \frac{3\rho}4V\le C_{\rm cav}\bigl[N^{-1}+A+J+H_3\bigr],\qquad
 H_3=\int\E|S_q-q|^3\,\nu(\dd q),
\]
with $C_{\rm cav}$ depending only on $\beta$. The argument occupies
\cref{part:cavity}, and \cref{sec:c7-mechanism} explains how the
inequality closes before the proofs begin.

\paragraph{Exchangeability.}
Let $n=N-1$, let $m_N(p)$ be the conditional magnetization of the last
spin given $\calG_p$, and let $A^-_p$ be the self-overlap of the
conditional mean of the other $n$ spins, which we call the rest
self-overlap. Then $NS_p=m_N(p)^2+nA^-_p$. Since the tilted law is
invariant under permutations of the sites, we have
$\Cov(S_p,m_N(p)^2)=\Var(S_p)$ \eqref{eq:c7-exchange}, and it follows
that
\[
 \Var(S_p)=\frac1N\Var\bigl(m_N(p)^2\bigr)
 +\frac nN\Cov\bigl(m_N(p)^2,A^-_p\bigr).
\]
Up to an error of order $N^{-1}$, the variance $V$ is the integrated
covariance between the squared magnetization of one spin and the
self-overlap of the others. Exchangeability is used in the same way in
the cavity method of~\cite[Section~1.6]{TalagrandVolI2011}.

\paragraph{The cavity interpolation.}
For $u\in[0,1]$, we give each coupling between the last spin and the
others variance $\alpha u/n$, and the last coordinate of the field
variance coefficient $b(u)=d-\alpha u$, where $\alpha=\beta^2tn/N$ and
$d=c+\alpha$ \eqref{eq:cavityparameters}. At $u=1$, we recover the
tilted law of the full system. At $u=0$, the last spin is independent
of the others: the law is a product of the tilted law of an $n$-spin
system, with the same profile, and the scalar law of \cref{sec:sc} with
coefficient $d$. In particular,
$\mathcal C_{pp}(u)=\Cov_u(m_N(p)^2,A^-_p)$ vanishes at $u=0$. By Taylor's formula, we
have
\[
 \mathcal C_{pp}(1)=\mathcal C_{pp}'(0)
 +\int_0^1(1-u)\,\mathcal C_{pp}''(u)\dd u .
\]
Note that the $n$-spin system at $u=0$ carries the profile of the full
system, which need not minimize its own functional.

\paragraph{Tree laws and the remainder.}
Derivatives in $u$ produce products of conditional expectations at
several levels. We represent these products by laws of finitely many
descendants on marked trees, in which deleting a leaf does not change
the law of the others (\cref{subsec:c7-tree}). Normalized replica trees of this
kind were developed in~\cite[Section~2]{OpenAIFluctuationScale}. By
\cref{lem:streaming}, the derivative in $u$ of a tree moment is
$\alpha/2$ times a combination of moments with one more overlap defect
$D_{\mathsf{ab}}=\epsilon_{\mathsf a}\epsilon_{\mathsf b}(R^-_{\mathsf{ab}}-q_{\mathsf{ab}})$,
plus a term from the normalization of the tilt. The row contributes the
rest overlap $R^-_{\mathsf{ab}}$ of two leaves, and the decrease of
$b(u)$ subtracts the level $q_{\mathsf{ab}}$ at which they split. It
follows that nonnegative tree moments at different values of $u$ agree
up to a factor $K_{\beta,k}$ that depends only on $\beta$ and the
number $k$ of leaves (\cref{cor:comparison}), and that every term of
$\mathcal C_{pp}''$ contains three defects. By
\cref{lem:cav-threeproducts,lem:remainder,lem:c7-thermal}, the Taylor
remainder is at most $C_\beta\sup_u\mathcal M_3(u)$, where
$\mathcal M_3(u)=\int\E_{2,q,u}|R^-_{12}-q|^3\,\nu(\dd q)$, and this
supremum is at most $K_{\beta,2}[64(A+J)+16H_3+32N^{-3}]$.

\paragraph{The linear term.}
Along the interpolation, the fields acting on the last spins of two
leaves have covariance
$dq_{\mathsf{ab}}+\alpha u(R^-_{\mathsf{ab}}-q_{\mathsf{ab}})$
(\cref{sec:c7-mechanism}). In the scalar law, if we replace the covariance $d(q\wedge r)$ of the
field by $dT(q\wedge r)$ with $T(q)=q+\theta\varphi(q)$, the mean square
magnetization at level $p$ changes at the rate $(J_{a,d}\varphi)(p)$
(\cref{lem:sc-response}). Heuristically, the row acts on the last spin
as a random transport of levels with displacement
$(\alpha u/d)(A^-_q-q)$, and we expect the main term of
$\mathcal C_{pp}'(0)$ to be $(\alpha/d)(J^{\scal}_{a,d}\Sigma^-)(p,p)$,
where $\Sigma^-(p,q)=\Cov^-(A^-_p,A^-_q)$ is the covariance of the rest
self-overlap at $u=0$. In \cref{lem:linear}, we prove the exact
identity
\[
 \int\mathcal C_{pp}'(0)\,\nu(\dd p)
 =\alpha\Bigl[\frac1d\Tr_\nu(J^{\scal}_{a,d}\Sigma^-)-\frac12\mathcal P
 +\int T_p\,\nu(\dd p)\Bigr],\qquad\mathcal P\ge0 .
\]
The term $-\mathcal P/2$ comes from the difference between conditional
and unconditional covariances of the rest self-overlap; we discard
it. The errors $T_p$ raise an obstacle: they contain covariances of
centered variables of the rest system with $A^-_p$, and if we applied the
Cauchy--Schwarz inequality at a single level, the bound would involve
$\Var^-(A^-_p)$, whose integral is what we want to control. To overcome this obstacle, we first integrate in $p$ against
the deterministic kernels of the scalar law, so that the covariance
order and the bound \eqref{eq:cl-BL} on the variance of $\nu$-averages
of $A^-$ apply. After a comparison with the full system through
\cref{cor:comparison}, we find in \cref{lem:cl-response} that
\[
 \alpha\Bigl|\int T_p\,\nu(\dd p)\Bigr|
 \le192\beta^2K_{\beta,3}\bigl[A+J+N^{-2}\bigr].
\]

\paragraph{\texorpdfstring{The factor $1-\rho$.}{The factor 1-rho.}}
Combining these identities, we obtain the trace identity of
\cref{lem:cav-trace},
\[
 V=\frac{V^{\rm site}}N+\frac nN\Bigl[
 \frac\alpha d\Tr_\nu(J^{\scal}_{a,d}\Sigma^-)
 -\frac\alpha2\mathcal P+\alpha\int T_p\,\nu(\dd p)
 +\mathcal R_{\rm cav}\Bigr],
\]
where $0\le V^{\rm site}\le1/4$ and $\mathcal R_{\rm cav}$ is the Taylor
remainder. The operator $\Sigma^-$ is positive, with range in
$\mathcal H_\nu$ and trace $V^-=\int\Var^-(A^-_p)\,\nu(\dd p)$. Under
\eqref{eq:ov-stability}, the trace $\Tr_\nu(J^{\scal}_{a,d}\Sigma^-)$ is
therefore at most $(1+\eps_N)V^-$. The contraction comes from the ratio
\[
 \frac\alpha d=t\,\frac{N-1}{N-t}\le t=1-\rho
\]
\eqref{eq:c7-alpha-d}. Indeed, the field acting on the last spin has
variance coefficient $d$, and at $u=1$, the part $c=\beta^2\rho$ of it
comes from the auxiliary field, whose covariance is deterministic; only
the part $\alpha$ that comes from the couplings feels the fluctuations
of the rest overlap. We conclude that
$V\le(1-\rho)(1+\eps_N)V^-+{}$errors.

\paragraph{Comparison with coefficient one.}
It remains to compare $V^-$ with $V$, and the comparison must have
coefficient one. If we compared the endpoints through
\cref{cor:comparison}, we would obtain only $V^-\le KV+{}$errors with
some $K>1$, and the inequality $V\le(1-\rho)K(1+\eps_N)V+{}$errors bounds $V$
only if $(1-\rho)K(1+\eps_N)<1$, which fails for $\rho\le1-1/K$. In
\cref{lem:cl-star}, we prove instead that
\[
 V^-\le V+C_\beta\sup_{0\le u\le1}\mathcal M_3(u)+\frac1{4n^2}.
\]
For each $p$, the second moment $\E_u(A^-_p-p)^2$ is the mean of a
product of two overlap defects on four leaves that split at $p$. Its
derivative in $u$ contains three defects, and after integration against
$\nu$, the values at $u=0$ and $u=1$ differ by at most
$C_\beta\sup_u\mathcal M_3(u)$. At $u=0$, it bounds $\Var^-(A^-_p)$, even
though we do not know that $\E^-A^-_p=p$. At $u=1$, exchangeability and
the contact identity $\E S_p=p$ show that it is at most
$\Var(S_p)+1/(4n^2)$.

\paragraph{Closing.}
As $(1-\rho)(1+\rho/4)\le1-3\rho/4$, these bounds imply that
$V\le(1-3\rho/4)V+C_\beta[N^{-1}+A+J+H_3]$, which is
\cref{prop:c7-closure}. By \cref{cor:wi-cub}, if $M\ge1$, then
\[
 H_3\le zV+\frac{425}{Ms}\log(2M)\log\frac2z\qquad(0<z\le1/4).
\]
We take $z$ proportional to $\rho$ and absorb $zV$ in the margin
$3\rho/4$. Inserting the bounds on $A$ and $J$ and using
$Ms=\beta N\rho\lambda$, we arrive at \cref{cor:c7-defect}, which states
that at every minimizer that satisfies \eqref{eq:ov-stability}, we have
\[
 \beta(A+2J+V)\le C_\beta\frac{\log^2(2+N)}{N\lambda\rho^2}
 \qquad(1/N\le\rho\le1).
\]
One factor $\rho^{-1}$ comes from the margin, and the other from $Ms$.

The framework of differentiation, comparison, and contraction on one
site is that of~\cite[Section~1.6]{TalagrandVolI2011}. A one-site
expansion with an independent scalar endpoint also appears
in~\cite[Lemma~2.2]{OpenAILimitingLaw}, where the resulting linear
equation is inverted under a gap condition on the scalar
response~\cite[Proposition~5.1]{OpenAILimitingLaw}; we compare the two
arguments in \cref{subsec:intro-openai}. We use no gap condition. The
multiplier that dominates our scalar response equals one on the support
of the Parisi measure, and the contraction comes from the ratio
$\alpha/d$.

\section{Stability of the scalar response}\label{sec:overview-stability}

At the independent endpoint, the last spin follows the scalar law with
the profile $a_N$ of the $N$-spin minimizer and the coefficient
$d_N=\beta^2(1-t/N)$, not with the Parisi measure. In \cref{sec:sc}, we
verify \eqref{eq:ov-stability} with an explicit $\eps_N$.

\paragraph{The response operator.}
For a profile $a$ and $d\in[\beta^2/2,\beta^2]$, let $u=u_{a,d}$ solve
the scalar Parisi equation with coefficient $d$ \eqref{eq:sc-pde}, and
let $Y$ be its diffusion. With $m_q=\partial_xu(q,Y_q)$ and
$h_q=\partial_{xx}u(q,Y_q)$, we define
\[
 \Gamma_{a,d}(q)=\E^0m_q^2,\qquad k_{a,d}(q)=d\,\E^0h_q^2,
\]
so that $\Gamma_{a,d}'=k_{a,d}$ \eqref{eq:sc-martingale}; we call
$k_{a,d}$ the scalar multiplier. By \cref{lem:sc-response}, the
operator $J_{a,d}$ of \eqref{eq:scalarJexplicit} is the derivative of
$\Gamma_{a,d}$ under transports of covariance levels: if we replace the
covariance $d(q\wedge r)$ of the field by $dT_\theta(q\wedge r)$, where
$T_\theta(q)=q+\theta\varphi(q)$ and $\varphi(0)=\varphi(1)=0$, the mean
square magnetization at level $p$ has derivative $(J_{a,d}\varphi)(p)$
at $\theta=0$. The scalar response operator $J^{\scal}_{a,d}$ is the
compression of $J_{a,d}$ to the space $\mathcal H_\nu$ of functions in
$L^2(\nu)$ that vanish at $0$ and $1$ when these are atoms of $\nu$. Its
kernel has no sign, and a direct estimate bounds its norm only by
$C_\beta$.

\paragraph{Domination by the multiplier.}
The following observation reduces the operator bound to a pointwise
bound. In \cref{lem:Jdom}, we show that, for every profile,
$J^{\scal}_{a,d}$ is self-adjoint and
\[
 J^{\scal}_{a,d}\preceq M_{k_{a,d}},\qquad
 (M_{k_{a,d}}v)(q)=k_{a,d}(q)v(q).
\]
To prove it, we follow the scalar Parisi functional $\mathcal P_d$ along
a transport of levels, $\Pi(\theta)=\mathcal P_d(a\circ T_\theta^{-1})$.
By \cref{lem:sc-response}, we have
$\Pi''(0)=(d/2)\langle\varphi,(I-J_{a,d})\varphi\rangle$
\eqref{eq:transporthess}; the first derivative is the first variation
under transports of levels due to
Talagrand~\cite{Talagrand2006ParisiMeasures}
(see~\cite[Section~3, (19)]{AuffingerChen2015Properties}). Since the functional is convex along segments of distribution
functions~\cite{AuffingerChen2015Uniqueness}, its increment along the
transport is at least its first variation at $a$, applied to
$a\circ T_\theta^{-1}-a$. We expand this lower bound to second order in
$\theta$, which produces the integral of $(1-k_{a,d})\varphi^2$ against
$\nu$. Comparing the two expansions, we find that
$\langle\varphi,J_{a,d}\varphi\rangle\le\langle\varphi,k_{a,d}\varphi\rangle$.
Note that $\Pi$ need not be convex, because
$\theta\mapsto a\circ T_\theta^{-1}$ is not affine; the term
$1-k_{a,d}$ accounts for this. It therefore suffices to show that
$\sup_qk_{a_N,d_N}(q)\le1+\eps_N$.

\paragraph{The multiplier at the Parisi measure.}
In \cref{lem:reference}, we prove that
\[
 k_{\alpha_\beta,\beta^2}(q)=1\quad(0\le q\le q_\beta),\qquad
 k_{\alpha_\beta,\beta^2}(q)\le e^{-\beta^2(q-q_\beta)/2}
 \quad(q_\beta\le q\le1).
\]
On the support, we rely on an identity of
Talagrand~\cite{Talagrand2006ParisiMeasures}, in the form
of~\cite[Theorem~5]{AuffingerChen2015Properties}
and~\cite[Lemma~2.6]{Lopatto}, which states that $\Gamma(q)=q$ on
$\supp\mu_\beta$. As the support is the interval
$[0,q_\beta]$~\cite[Theorem~1.1]{Lopatto}, we may differentiate this
identity along it, which gives $k=\Gamma'=1$. Above the support,
we have $\alpha_\beta=1$, and the solution $u$ is explicit. A
transformed density $z_q$ of $Y_q$ evolves there by the heat
equation~\cite[Lemma~2.4(ii)]{Lopatto}, and $z_{q_\beta}$ is even and
log-concave~\cite[Proposition~4.1]{Lopatto}. In particular,
$z_{q_\beta}$ is decreasing on $(0,\infty)$. The multiplier is an
explicit multiple of the integral of $z_q$ against $\sech^3$
\eqref{eq:sc-k-plateau}, and Gaussian smoothing of an even decreasing
density does not increase its integral against an even decreasing
function; this gives the exponential bound. We need it on all of
$[0,1]$, because a finite-dimensional minimizer may place mass above
$q_\beta$.

\paragraph{Marginality.}
At the Parisi measure, the multiplier equals one on the whole support,
and the domination gives no contraction there. The cavity
inequality contracts because of the ratio $\alpha/d\le1-\rho$, which
leaves a margin of order $\rho$ only. This is why we need $\eps_N\le\rho/4$ in
\eqref{eq:ov-stability}, and why stability holds only above a cutoff.
The equality $k=1$ on the support expresses, in scalar form, the
marginal stability that De~Dominicis and Kondor found in a truncated
replica model near the critical
temperature~\cite[pp.~606--608]{DeDominicisKondor1983}; see
\cref{subsec:intro-physics}.

\paragraph{Transfer to the minimizer.}
By \cref{lem:scalarcontinuity}, the multiplier depends Lipschitz
continuously on the profile and coefficient, in the sense that
\[
 \sup_{0\le q\le1}|k_{a,d}(q)-k_{\widetilde a,\widetilde d}(q)|
 \le C_\beta\bigl(\|a-\widetilde a\|_{L^1(0,1)}+|d-\widetilde d|\bigr).
\]
We have $|d_N-\beta^2|\le\beta^2/N$ and
$\|\gamma^{\rm tr}_\lambda-\gamma_\beta\|_{L^1(0,1)}\lesssim\lambda^2$
\eqref{eq:sc-floor}. Combining these facts with the isolation estimate
of \cref{sec:overview-localization}, we obtain \cref{prop:lz-stability},
which states that on the space $\mathcal H_{\nu_N}$ of the minimizer
itself,
\[
 J^{\scal}_{a_N,d_N}\preceq(1+\eps_N)I,\qquad
 \eps_N\le C_{\rm st}\bigl[\rho^{-3/4}\eta^{1/4}+\lambda^2+N^{-1}\bigr],
\]
with $C_{\rm st}$ depending only on $\beta$. We compare only the multipliers of
the two profiles, and the operator inequality comes from
\cref{lem:Jdom} at $a_N$.

\section{Localization of the minimizing profile}
\label{sec:overview-localization}

We need two bounds on the distance between a minimizer $\gamma_*$ of
$F_\rho$ and the scalar trial profile $\gamma^{\rm tr}$, in terms of a
trial budget $\eta$: a bound in $L^1(0,1)$ for the stability of the
scalar response (\cref{sec:lz}), and a localization of the lowest
quantiles for small $\rho$ (\cref{sec:lz-root}).

\paragraph{Convexity and response.}
From~\cite[Section~7.4 and Appendix~A]{AL}, we import the following
facts (\cref{prop:loc-inputs}). The functional $F_\rho$ is convex along
segments $\gamma_\theta=(1-\theta)\gamma_0+\theta\gamma_1$. Let
$f(\theta)=F_\rho(\gamma_\theta)$
and $\zeta_\theta(q)=\beta^2N^{-1}\E_{\gamma_\theta}\Tr H_q^2$, so that
$\rho\zeta_\theta$ is the derivative of the expected self-overlap. The
second variation controls how the self-overlap responds to
$\gamma_1-\gamma_0$:
\[
 \rho\left|\int_0^1(\gamma_1-\gamma_0)\zeta_\theta\varphi\dd q\right|
 \le\Bigl(\frac{f''(\theta)}\rho\Bigr)^{1/2}
 \bigl[Y^{1/2}+2(Y+Z)^{1/2}\bigr]
\]
\eqref{eq:loc-response}, for Lipschitz cutoffs $\varphi$, where $Y$ is a
weighted Dirichlet energy of $\varphi$ and $Z$ is the integral of
$\rho\varphi^2\zeta_\theta$ against $\dd\gamma_\theta$. Further,
$\zeta_\theta\ge k_\beta^2$ for a constant $k_\beta>0$
\eqref{eq:loc-in-curv}. The midpoint inequality \eqref{eq:loc-midpoint}
bounds the integral of $\min\{\theta,1-\theta\}f''(\theta)$ over
$(0,1)$ by $F_\rho(\gamma_0)+F_\rho(\gamma_1)-2\min F_\rho$. If both
endpoints are nearly minimal, the segment has little
curvature, while a difference of one sign and of definite size on some
region forces curvature (\cref{lem:loc-coercive}).

\paragraph{\texorpdfstring{Isolation in $L^1$.}{Isolation in L1.}}
A minimizer may have atoms, and $Z$ contains $\dd\gamma_\theta$ on the
support of the cutoff. To overcome this difficulty, we first spread the
free mass of $\gamma_*$ on a scale $h$ by a reflected regularization
(\cref{lem:loc-smoothing}). It creates a density at most $\beta/h$ at a
cost at most $\rho\beta h^2/12$; the cost is quadratic because the
first-variation potential is semiconcave, as in~\cite[Lemma~7.5]{AL}.
Along the segment from the regularized minimizer to $\gamma^{\rm tr}$,
we use cutoffs of one sign built from one-sided averages of
superlevel sets of the difference, away from the atom of
$\gamma^{\rm tr}$ at $q_\beta$; below $q_\beta$, the trial profile is
Lipschitz (\cref{lem:loc-lipschitz}). Using the layer-cake formula, we
find that $f''\gtrsim\rho^2\min\{\rho D^4,hD^2\}$, where
$D=\|\gamma_*-\gamma^{\rm tr}\|_{L^1(0,1)}$ \eqref{eq:loc-E2-second}.
For $h$ proportional to $\rho D^2$, we have $f''\gtrsim\rho^3D^4$, and
the smoothing cost is a small fraction of $\rho^3D^4$. The midpoint
inequality implies that $\rho^3D^4\lesssim\eta$, which is the bound of
\cref{prop:loc-isolation},
\[
 \|\gamma_*-\gamma^{\rm tr}\|_{L^1(0,1)}
 \le C\min\{1,\rho^{-3/4}\eta^{1/4}\},
\]
for every minimizer and every trial budget, with $C$ depending only on
$\beta$.

\paragraph{Root localization.}
In the isolation argument, the lower bound
$\rho\int\zeta_\theta\varphi\ge\rho k_\beta^2\int\varphi$ on the response
carries a factor $\rho$, which is too costly near $\rho=0$. To avoid it,
we use the local deficit certificate of \cref{lem:lz-certificate}. Suppose that
two nearly minimal profiles have bounded densities on an interval $I$ of
length $\ell$, that the density of the second is also bounded below,
and that the second exceeds the first by a definite amount on $I$. Then
the certificate bounds $\eta/\rho$ from below by an explicit expression
in $\ell$, the gap, and the density bounds. In
its proof, we move mass of the profile toward the middle of $I$ along a
bump. Since both endpoints are nearly minimal and $F_\rho$ is convex,
the response integrated against the bump must be at least of order
$\ell$, without the factor $\rho$. To localize
a quantile $Q_{\gamma_*}(v)$, we bound its truncation
$q=\min\{Q_{\gamma_*}(v),T_1\}$ at the level $T_1$ of
\cref{lem:lz-rootdensity}, which depends only on $\beta$; this suffices,
because $Q_{\gamma_*}(v)\le q/T_1$. If $q$ is large compared with
$\lambda+v$, we regularize the free mass of the minimizer below $q$ only. The result
is small on $(q/4,3q/4)$, where the trial profile has density at least
$k_0/2$ (\cref{lem:lz-rootdensity}) and exceeds it by an amount of order
$q$. Using the certificate, we find that
$\eta/\rho\gtrsim\min\{q^6,q^4/(m/h+C)\}$, where $m/h$ is the density
created by the regularization. With a width adapted to the moved mass,
we obtain \cref{prop:lz-quantile},
which states that
\[
 Q_{\gamma_*}(v)\le C\bigl[\lambda+v+\min\{1,(\eta/\rho)^{1/6}\}\bigr]
 \qquad(0<v<\beta-\lambda),
\]
together with its form \eqref{eq:lz-Q} in the probability
normalization. We follow the strategy of~\cite[Lemma~7.6]{AL}. Compared
with~\cite[(7.30)]{AL}, the bound depends on $\eta$ and $\rho$ only
through $\eta/\rho$, and the certificate is what gives this dependence.

\section{Closing the implicit inequality}\label{sec:overview-closure}

By \cref{prop:lz-stability}, the stability assumption
\eqref{eq:ov-stability} holds as soon as
$C_{\rm st}[\rho^{-3/4}\eta^{1/4}+\lambda^2+N^{-1}]\le\rho/4$. Recall
that $\eta\ge L_s\lambda^5\ge\lambda^5$, so that
$\eta^{1/7}\ge\lambda^{5/7}=N^{-1/10}$, which is at least $\lambda^2$ and
$N^{-1}$. The first term satisfies $\rho^{-3/4}\eta^{1/4}\lesssim\rho$
exactly when $\rho\gtrsim\eta^{1/7}$. With $K=12C_{\rm st}$, we set
\[
 \rho_0=\min\{1,K\eta^{1/7}\}
\]
\eqref{eq:b7-cutoff}. If $\rho_0<1$, then $\eps_N\le\rho/4$ for every
$\rho\in[\rho_0,1]$ and every minimizer, and \eqref{eq:ov-stable}
applies; if $\rho_0=1$, the stable interval is empty. Using
$\rho_0\ge\lambda^{5/7}$, we find that
\[
 \int_{\rho_0}^1\calD_N(\rho)\dd\rho
 \lesssim\frac{\log^2(2+N)}{N\lambda\rho_0}
 \lesssim\frac{\log^2(2+N)}{N\lambda^{12/7}}
 =N^{-19/25}\log^2(2+N).
\]
On $(0,\rho_0]$, we insert $\rho_0\le K\eta^{1/7}$ into
\eqref{eq:ov-small-integrated}, which gives
$\rho_0^{1/3}\lesssim\eta^{1/21}$ and
$\eta^{1/9}\rho_0^{2/9}\lesssim\eta^{1/7}$. Adding the two
contributions and recalling that
$4e_{N,\beta}(\lambda)=\int_0^1\calD_N(\rho)\dd\rho$, we obtain the
bound \eqref{eq:b7-normalized-input},
\begin{equation}\label{eq:ov-implicit}
 e_{N,\beta}(\lambda)\lesssim N^{-19/25}\log^3(2+N)
 +N^{-2/3}\eta^{1/21}+N^{-2/3}\lambda^{-1/3}\eta^{1/7}.
\end{equation}

But there is an apparent circularity here, since the cutoff $\rho_0$ is
defined through $\eta=e_{N,\beta}(\lambda)+L_s\lambda^5$, which contains
the quantity that we want to bound. The way out is that the error
enters \eqref{eq:ov-implicit} only through powers of $\eta$ smaller
than one. Let $r=N^{-7/10}=\lambda^5$ and $x=\eta/r$, which is finite
and at least $L_s\ge1$. Dividing \eqref{eq:ov-implicit} by $r$ and
adding $L_s$ to both sides, we find \eqref{eq:b7-normalized},
\[
 x\le L_s+C_\beta\bigl[N^{-3/50}\log^3(2+N)+x^{1/21}
 +N^{-1/50}x^{1/7}\bigr]\le C_\beta x^{1/7},
\]
which implies that $x\le C_\beta^{7/6}$ and
$e_{N,\beta}(\lambda)\le\eta=rx\lesssim N^{-7/10}$. We need no preliminary bound
on the error. The proof of~\cite[Proposition~7.7]{AL} closes
an inequality in which the unknown error enters the trial budget in the
same way. By \cref{prop:fnd-mean-bias}, we conclude that
$P_\beta-p_{N,\beta}\lesssim N^{-7/10}$ for all large $N$, and the
finitely many remaining values of $N$ only change the constant.

\paragraph{The exponent.}
If $\eta$ is of order $N^{-a}$, the leading term
$N^{-2/3}\eta^{1/21}$ in \eqref{eq:ov-implicit} is of the same order as
$\eta$ when
\[
 a=\frac23+\frac a{21},\qquad\text{that is,}\qquad a=\frac7{10},
\]
and the cost $\lambda^5$ is of the same order when $\lambda=N^{-7/50}$
(\cref{rem:b7-exponents}). \Cref{tab:overview-terms} lists the terms of
\eqref{eq:ov-implicit} after division by $N^{-7/10}$. Only the leading
term has no polynomial slack; the others carry the factors $N^{-1/50}$,
$N^{-3/50}$, $N^{-4/25}$, and $N^{-3/10}$, up to logarithms. After the
closure, we have $L_sN^{-7/10}\le\eta\lesssim N^{-7/10}$, and the
cutoff is then comparable to $N^{-1/10}$ for large $N$.
\Cref{tab:overview-exponents} collects the exponents of the proof and
the statements in which they arise.

\begin{table}[!tb]
\centering
\small
\begin{tabular}{@{}>{\raggedright\arraybackslash}p{0.27\textwidth}
 >{\raggedright\arraybackslash}p{0.42\textwidth}
 >{\raggedright\arraybackslash}p{0.23\textwidth}@{}}
\toprule
Term & Origin & After division by $N^{-7/10}$\\
\midrule
$N^{-1}$ & $\rho\le1/N$, by \eqref{eq:fnd-moving} & $N^{-3/10}$\\
$N^{-43/50}\log^3(2+N)$ & thermal and logarithmic terms on
 $[1/N,\rho_0]$, \eqref{eq:sm-pointwise} & $N^{-4/25}\log^3(2+N)$\\
$N^{-19/25}\log^2(2+N)$ & stable interval $[\rho_0,1]$,
 \cref{cor:c7-defect} & $N^{-3/50}\log^2(2+N)$\\
$N^{-2/3}\eta^{1/21}$ & quantile shells on $(0,\rho_0]$,
 \cref{prop:wi-vfs} & $x^{1/21}$\\
$N^{-2/3}\lambda^{-1/3}\eta^{1/7}$ & root localization correction,
 \cref{cor:lz-omega} & $N^{-1/50}x^{1/7}$\\
\bottomrule
\end{tabular}
\caption{The terms of the bound on $e_{N,\beta}(\lambda)$ at
$\lambda=N^{-7/50}$, with $x=\eta/N^{-7/10}$. The first three terms are
combined into one in \eqref{eq:ov-implicit}; see
\cref{subsec:b7-general}.}
\label{tab:overview-terms}
\end{table}

\begin{table}[!tb]
\centering
\small
\begin{tabular}{@{}>{\raggedright\arraybackslash}p{0.2\textwidth}
 >{\raggedright\arraybackslash}p{0.47\textwidth}
 >{\raggedright\arraybackslash}p{0.25\textwidth}@{}}
\toprule
Quantity & Where it comes from & Statement\\
\midrule
$\lambda^5$ & cost of the prescribed mass at zero, at the scalar end
 & \eqref{eq:fnd-scalar-trial-cost}, \cite[Lemma~7.2]{AL}\\
$(N\rho)^{-2/3}$ & variance on dyadic quantile shells, after root
 localization & \cref{prop:wi-vfs}, \eqref{eq:sm-pointwise}\\
$(\eta/\rho)^{1/6}$ & localization of the lowest quantiles
 & \cref{prop:lz-quantile}, \cref{cor:lz-omega}\\
$\rho^{-3/4}\eta^{1/4}$ & $L^1$ isolation of the minimizing profile
 & \cref{prop:loc-isolation}\\
$\rho/4$ & admissible error in the scalar response; it leaves the
 margin $3\rho/4$ & \eqref{eq:c7-stab}\\
$\eta^{1/7}$ & cutoff: $\rho^{-3/4}\eta^{1/4}\lesssim\rho$ when
 $\rho\gtrsim\eta^{1/7}$ & \eqref{eq:b7-cutoff}\\
$(N\lambda\rho^2)^{-1}$ & cavity bound: margin $3\rho/4$ and
 $Ms=\beta N\rho\lambda$ & \cref{cor:c7-defect}\\
$\rho_0^{1/3}$ & integral of $(N\rho)^{-2/3}$ over $(0,\rho_0]$
 & \eqref{eq:sm-integrated}\\
$\eta^{1/21}$ & $\rho_0^{1/3}$ with $\rho_0\le K\eta^{1/7}$
 & \eqref{eq:b7-normalized-input}\\
$a=7/10$ & balance $a=2/3+a/21$ & \cref{rem:b7-exponents}\\
$\lambda=N^{-7/50}$ & $\lambda^5=N^{-7/10}$ & \eqref{eq:b7-parameters}\\
\bottomrule
\end{tabular}
\caption{The exponents of the proof.}
\label{tab:overview-exponents}
\end{table}

\paragraph{What would improve the exponent.}
The exponent $1/21$ is the product of the exponent $1/3$ of $\rho_0$ in
the integrated shell estimate and the exponent $1/7$ of $\eta$ in the
cutoff, which comes from the isolation estimate. A larger exponent in
either place would improve the balance. A smaller cutoff, however,
increases the contribution $\log^2(2+N)/(N\lambda\rho_0)$ of the stable
interval, and such an improvement requires rechecking the other
terms of \cref{tab:overview-terms}. If the cavity bound
\eqref{eq:ov-stable} held down to $\rho$ of order $N^{-1/4}$, and if we
ignored logarithms and the tilt, the two regimes would give $N^{-3/4}$;
we discuss this heuristic in \cref{subsec:intro-physics}.

\section{Minimizers, arbitrary profiles, and atoms}
\label{sec:overview-profiles}

Three conventions recur throughout the proof.

\paragraph{Which statements use minimality.}
Estimates that use the contact identities \eqref{eq:in-contact} or the
trial budget hold at minimizers: the main results of
\cref{sec:wi,sec:th,sec:lz,sec:lz-root}, Step~2 of the proof of
\cref{lem:cl-star}, and everything built on them. The Gaussian tools of
\cref{sec:gauss}, the lemmas of \cref{sec:sc} on the scalar law of an
arbitrary profile, and the cavity calculus of
\cref{sec:c7,sec:c7-endpoint} hold for every deterministic profile (with
$\nu(\{1\})=0$ in \cref{sec:c7-endpoint}). The reference estimate
\cref{lem:reference} concerns the Parisi measure $\mu_\beta$ itself, and
\cref{prop:lz-stability} concerns minimizing profiles. We need this generality,
since at the independent endpoint, the $n$-spin system carries the
profile of the full system.

\paragraph{Every minimizer.}
A minimizer need not be unique, and it may depend on $N$, $\rho$, and
$\lambda$, but not on the couplings. Since \eqref{eq:in-envelope} holds at
every minimizer, and since our bounds on $A$, $J$, and $V$ hold at every
minimizer with the same constants, we can integrate in $\rho$ without a
measurable choice of minimizers.

\paragraph{Atoms.}
Minimizing profiles may have atoms in $[0,1)$, and $\nu$ may charge
$0$, but $\nu(\{1\})=0$ by \cref{prop:fnd-contact}. Distribution functions are right-continuous, so that $a(q)$
includes an atom at $q$ while $(q,1]$ excludes it. We compute integrals
against $\nu$ through quantiles, and a quantile window may contain part
of an atom (\cref{sec:wi-setup}). We handle scalar and tree laws
through exact path likelihoods (\cref{lem:sc-likelihood,lem:tree-density})
and through the heat identity of \cref{lem:fnd-path-heat}, in which
profile measures are held fixed with their atoms. In
\cref{sec:lz,sec:lz-root}, we spread the free mass before using cutoffs.

\section{A map of the proof}\label{sec:overview-map}

\Cref{tab:overview-map} lists, for each section, its main results, its
main inputs, and the places where its results are used. The proof
of \cref{thm:main} in \cref{subsec:b7-general} begins with a table of
its inputs and of the parameters at which they are used. In terms of
these results, the proof has the following structure.
\begin{itemize}
\item \Cref{prop:fnd-mean-bias}, \cref{cor:fnd-defect}, and
\cref{lem:fnd-trial} reduce the theorem to a bound on the integral of
$\calD_N$, evaluated at any minimizer, and give the trial budget $\eta$
that defines the cutoff $\rho_0$.
\item On $(0,\rho_0]$, \cref{prop:sm-small} bounds the defect, using
\cref{prop:wi-vfs}, \cref{cor:lz-omega}, \eqref{eq:in-Lvar},
\cref{prop:th-j}, and \eqref{eq:fnd-moving}.
\item On $[\rho_0,1]$, \cref{cor:c7-defect} bounds the defect, using
\cref{prop:c7-closure}, \cref{cor:wi-cub}, and \cref{prop:th-j}. The
contraction of \cref{prop:c7-closure} combines
\cref{lem:cav-trace,lem:linear,lem:remainder,lem:cl-response,lem:cl-star,lem:c7-thermal},
and its stability hypothesis \eqref{eq:c7-stab} is verified by
\cref{prop:lz-stability}, which combines
\cref{lem:Jdom,lem:reference,lem:scalarcontinuity,prop:loc-isolation}.
\item The closure in \cref{subsec:b7-general} combines the two bounds.
\end{itemize}

\begin{table}[!tb]
\centering
\small
\begin{tabular}{@{}>{\raggedright\arraybackslash}p{0.06\textwidth}
 >{\raggedright\arraybackslash}p{0.31\textwidth}
 >{\raggedright\arraybackslash}p{0.3\textwidth}
 >{\raggedright\arraybackslash}p{0.25\textwidth}@{}}
\toprule
Sec. & Main results & Main inputs & Used in\\
\midrule
\ref{sec:fnd} & defect identity \eqref{eq:in-envelope}; contact and
 variance bounds, \cref{prop:fnd-contact,prop:fnd-variance}; trial
 budget, \cref{lem:fnd-trial}; reduction, \cref{prop:fnd-mean-bias}
 & \cite[Section~2, Lemma~7.2, Theorem~7.8]{AL},
 \cite{ChenConvexity} & all later sections\\
\ref{sec:gauss} & posterior bound for arbitrary profiles,
 \cref{lem:fnd-posterior}; heat identity, \cref{lem:fnd-path-heat}
 & \cite[Section~2]{AL},
 \cite[Lemma~2.3]{OpenAIFluctuationScale}
 & \cref{sec:sc,sec:c7,sec:c7-endpoint}\\
\ref{sec:wi} & variance on quantile intervals, \cref{prop:wi-vfs};
 cubic moment, \cref{cor:wi-cub}
 & \cref{prop:fnd-contact,prop:fnd-variance}
 & \cref{prop:sm-small,cor:c7-defect}\\
\ref{sec:th} & mean--covariance term, \cref{prop:th-j}
 & \cref{prop:fnd-contact,prop:fnd-variance}
 & \cref{prop:sm-small,cor:c7-defect}\\
\ref{sec:lz} & isolation, \cref{prop:loc-isolation}
 & \cite[Section~7.4, Appendix~A]{AL}, \cite[Theorem~1.1]{Lopatto},
 \cref{lem:fnd-trial}
 & \cref{sec:lz-root}, \cref{prop:lz-stability}\\
\ref{sec:lz-root} & root localization,
 \cref{prop:lz-quantile,cor:lz-omega}
 & \cref{prop:loc-inputs,lem:loc-smoothing,lem:loc-coercive}
 & \cref{prop:sm-small}\\
\ref{sec:sc} & response identity, \cref{lem:sc-response}; domination,
 \cref{lem:Jdom}; reference multiplier, \cref{lem:reference};
 stability, \cref{prop:lz-stability}
 & \cref{sec:gauss}, \cref{prop:loc-isolation},
 \cite{Lopatto,AuffingerChen2015Uniqueness}
 & \cref{lem:linear,prop:c7-closure}; proof of \cref{thm:main}\\
\ref{sec:c7} & tree calculus, \cref{lem:streaming,cor:comparison};
 remainder, \cref{lem:cav-threeproducts,lem:remainder}
 & \cref{sec:gauss}, scalar law of \cref{sec:sc}
 & \cref{sec:c7-endpoint,sec:c7-closure}\\
\ref{sec:c7-endpoint} & linear response, \cref{lem:linear}; errors,
 \cref{lem:cl-response}
 & \cref{sec:c7}, \eqref{eq:scalarJexplicit}, \cref{lem:fnd-path-heat}
 & \cref{lem:cav-trace,prop:c7-closure}\\
\ref{sec:c7-closure} & coefficient one, \cref{lem:cl-star}; trace
 identity, \cref{lem:cav-trace}; contraction, \cref{prop:c7-closure};
 defect, \cref{cor:c7-defect}
 & \cref{sec:c7,sec:c7-endpoint}, \cref{lem:Jdom},
 \cref{cor:wi-cub,prop:th-j}
 & proof of \cref{thm:main}\\
\ref{sec:b7} & small $\rho$, \cref{prop:sm-small}; \cref{thm:main}
 & \cref{sec:fnd,sec:wi,sec:th,sec:lz-root,sec:sc,sec:c7-closure}
 & ---\\
\bottomrule
\end{tabular}
\caption{The sections of the proof, their main results and inputs, and
where the results are used.}
\label{tab:overview-map}
\end{table}


\part{The interpolation}\label{part:interpolation}
\chapter{The interpolation and the defect}\label{sec:fnd}

In this section, we recall the interpolation of \cite[Section~2]{AL},
which compares the fractional moments of the partition function with the
constrained scalar Parisi problem. We also fix the normalizations, the
tilted law, and the observables that we use in the rest of the
paper. The later sections build on two results. The first is the
defect identity \eqref{eq:in-envelope} of \cref{cor:fnd-defect}. It states
that four times the fractional-moment error $e_{N,\beta}(\lambda)$
is the integral over the interpolation of a defect
$\calD_N(\rho)=\beta(A+2J+V)$, whose three terms are overlap observables
at a minimizing profile. The second is \cref{prop:fnd-mean-bias}, where
we show that $(P_\beta-p_{N,\beta})/\beta$ is at
most $2e_{N,\beta}(\lambda)$ plus a term of order $\lambda^5$. Along the
way, we record the martingale and contact identities in
\cref{prop:fnd-contact}, the integrated variance bound in
\cref{prop:fnd-variance}, the decomposition of the overlap error in
\cref{lem:fnd-decomp}, the interpolation identity in
\cref{prop:fnd-interpolation}, and the trial budget in
\cref{lem:fnd-trial}.

Most of these facts are proved in \cite{AL}. We state each of them in
full, with its location in \cite{AL}. Where the argument is short, we
also sketch the proof, so that this section can be read without \cite{AL}
at hand. Besides \cite{AL}, this section uses only the scalar Parisi
functional \eqref{eq:intro-parisi} and the notation of
\cref{sec:intro-background}, together with the support theorem
\cite[Theorem~1.1]{Lopatto}. Its results are used throughout the
paper. \Cref{prop:fnd-contact,prop:fnd-variance,lem:fnd-decomp} are
used in \cref{sec:wi,sec:th,sec:c7,sec:c7-closure}; the energy
normalization and the trial budget in \cref{sec:lz,sec:lz-root,sec:sc};
and \eqref{eq:in-envelope}, \eqref{eq:fnd-moving}, and
\cref{prop:fnd-mean-bias} in \cref{sec:c7-closure,sec:b7}. We treat the
laws with anisotropic fields in \cref{sec:gauss}, where we
also prove the two Gaussian tools that we need for profiles that do not
minimize the functional.

Throughout this section, we fix $\beta>1$ and $N\ge1$, and we assume no
relation between the tilt $\lambda$ and $N$. The completion,
the interpolating functional, and the identities of
\cref{subsec:fnd-law,sec:fnd-contact,sec:fnd-variance,sec:fnd-defect}
hold for every $\beta>0$. We use the hypothesis $\beta>1$ for the support
of the Parisi measure in \cref{sec:fnd-completion} and for the
cost of the scalar floor in \cref{sec:fnd-trial}.

The interpolation of \cite[Section~2]{AL} is a quantitative form of the
Guerra interpolation that Talagrand used for the positive fractional
moments of the partition function, whose limits he determined in
\cite{Talagrand2007FractionalMoments}; see \cite[Section~1.2]{AL}. For
fixed $0<s<1$, the limit of the normalized logarithmic moment of order
$s$ is the minimum of the Parisi functional over probability measures
with mass at least $s$ at zero. For the SK model at zero field, this
formula is also contained in the one-level case of
\cite[Theorem~2.1]{ContucciMingione2019}, and it is stated in this form
in \cite[Theorem~2.1]{ChenEtAl2026LargeDeviations}, in a normalization
of the spins that changes the functional by the constant $\log2$. The
finite-$N$ functional of \cref{sec:fnd-functional} carries the same
constraint, in the sense that its profiles have the
prescribed mass $s$ at $0$.

\section{Completion and the constrained scalar value}
\label{sec:fnd-completion}

Let $W$ be the symmetric $N\times N$ matrix with entries
$W_{ij}=g_{ij}/\sqrt N$ for $i<j$ and $W_{ii}=\sqrt{2/N}\,z_i$, where
$z_1,\ldots,z_N$ are standard Gaussian variables, independent of each
other and of the couplings $(g_{ij})$. We write $\E_{\rm ref}$ for
expectation over $(g_{ij},z_i)$ under this Gaussian law; the subscript
distinguishes it from the tilted law of \cref{subsec:fnd-law}.
Following \cite[Section~1.1]{AL}, we set
\[
 D_N=\frac12\Tr W,\qquad
 F^c_{N,\beta}=\log\sum_{\sigma\in\{-1,1\}^N}e^{\beta\sigma^TW\sigma/2},
\]
and we recall that $F_{N,\beta}=\log\sum_\sigma e^{\beta H_N(\sigma)}$ is the
logarithmic partition function of \cref{sec:intro-background}. Note that
$\sigma^TW\sigma/2=H_N(\sigma)+D_N$, because $\sigma_i^2=1$, so that
$F^c_{N,\beta}=F_{N,\beta}+\beta D_N$. The variable
$D_N=(2N)^{-1/2}\sum_iz_i$ is centered Gaussian with variance $1/2$ and
is independent of $F_{N,\beta}$.

The diagonal term completes the covariance of the Hamiltonian to an
exact function of the overlap $R(\sigma,\sigma')=N^{-1}\sigma\cdot\sigma'$.
Indeed, writing $\tau_i=\sigma_i\sigma'_i$, we have
\[
 \E_{\rm ref}\bigl[H_N(\sigma)H_N(\sigma')\bigr]
 =\frac1N\sum_{i<j}\tau_i\tau_j
 =\frac1{2N}\biggl[\Bigl(\sum_{i=1}^N\tau_i\Bigr)^2-N\biggr]
 =\frac N2R(\sigma,\sigma')^2-\frac12 ,
\]
and adding $\Var(D_N)=1/2$, we find that the completed Hamiltonian
$\sigma^TW\sigma/2$ has covariance $NR(\sigma,\sigma')^2/2$. This is why
the derivative of the interpolation in \cref{prop:fnd-interpolation} is
exactly an integral of squared overlap deviations. Because $D_N$ is
independent of $F_{N,\beta}$, the completion changes the logarithmic
moments only by an explicit term. For $s>0$, the normalized logarithmic moment of the completed
partition function is \cite[(1.3)]{AL}
\[
 \phi^c_{N,\beta}(s)
   =\frac1{Ns}\log\E_{\rm ref}e^{sF^c_{N,\beta}}
   =\frac1{Ns}\log\E_{\rm ref}e^{sF_{N,\beta}}+\frac{\beta^2s}{4N},
\]
where we used $\log\E_{\rm ref}e^{s\beta D_N}=\beta^2s^2/4$. As $s\downarrow0$, we have $\phi^c_{N,\beta}(s)\to p_{N,\beta}$, since
$F_{N,\beta}$ has exponential moments of every order.

Recall the scalar Parisi functional $\mathcal P_\beta$ of
\eqref{eq:intro-parisi}, which we evaluate at the distribution function
of a probability measure on $[0,1]$. For $0\le s\le1$ and
$0\le\lambda\le\beta$, we define
\[
 P_\beta(s)=\min_{\mu\in\mathcal P([0,1]):\,\mu(\{0\})\ge s}
      \mathcal P_\beta\bigl(\mu([0,\cdot\,])\bigr),
 \qquad E_\beta(\lambda)=\beta^{-1}P_\beta(\lambda/\beta).
\]
The minimum is attained \cite[Section~1.1]{AL}, and $P_\beta(0)=P_\beta$.
The function $E_\beta$ is the constrained value in the energy
normalization of \cite[Section~7.1]{AL}, in which the order $s$ of the
moment is replaced by the tilt $\lambda=\beta s$. Note that $E_\beta$ is nondecreasing, because the admissible class
shrinks as $s$ grows; in particular, $E_\beta(\lambda)\ge E_\beta(0)$.

The unconstrained minimizer $\mu_\beta$ is unique
\cite{AuffingerChen2015Uniqueness}. We write
$\alpha_\beta(q)=\mu_\beta([0,q])$ and $\gamma_\beta=\beta\alpha_\beta$.
Since $\beta>1$, we have $\supp\mu_\beta=[0,q_\beta]$ with
$0<q_\beta<1$, by \cite[Theorem~1.1]{Lopatto}. We recall the further
structure of $\mu_\beta$ in \cref{sec:lz}.

For $0<\lambda<\beta$ and $s=\lambda/\beta$, we define the
\emph{fractional-moment error} by
\begin{equation}\label{eq:fnd-error-definition}
 e_{N,\beta}(\lambda)
    =\frac{P_\beta(s)-\phi^c_{N,\beta}(s)}\beta
    =E_\beta(\lambda)
       -\frac1{N\lambda}\log\E_{\rm ref}
          e^{\lambda(F_{N,\beta}/\beta+D_N)} .
\end{equation}
In the second expression, we used
$sF^c_{N,\beta}=\lambda(F_{N,\beta}/\beta+D_N)$; the completion
contributes $\lambda/(4N)$ to the logarithmic moment on the right.
As $\beta$ is fixed, we also write $e_N(\lambda)=e_{N,\beta}(\lambda)$.
From now on, we assume that $s=\lambda/\beta\le1/2$, which is the range
of the results of \cite[Section~2]{AL}. We will see in
\cref{cor:fnd-defect} that the error
$e_{N,\beta}(\lambda)$ is nonnegative.

\section{The interpolating functional}\label{sec:fnd-functional}

For $0\le\rho\le1$, we set
\[
 t=1-\rho,\qquad c=\beta^2\rho,\qquad M=Nc .
\]
The interpolation reaches the fractional moment at $\rho=0$, and
$\rho=1$ is the scalar end. We parametrize the interpolation by $\rho$
because the estimates of the later sections degenerate as
$\rho\downarrow0$. We work with the tilted law of \cref{subsec:fnd-law}
for $0<\rho\le1$,
where $c>0$.

An \emph{admissible profile} is a deterministic probability measure
$\mu=s\delta_0+\nu$ on $[0,1]$, where $\nu\ge0$ has mass $1-s$ and may
charge $0$. Equivalently, $\mu$ is a deterministic probability measure
such that $\mu(\{0\})\ge s$. We call $s\delta_0$ the prescribed root mass and
$\nu$ the free measure of $\mu$. We write $a(q)=\mu([0,q])$ for the
cumulative function. It is
nondecreasing and right-continuous, with $a\ge s$ and $a(1)=1$.
Conditional on $W$, let $U=U_{\mu,t}$ solve
\begin{equation}\label{eq:fnd-vector-pde}
 \partial_qU=-\frac c2\bigl(\Delta U+a(q)|\nabla U|^2\bigr),
 \qquad
 U(1,x;W)=\log\sum_{\sigma\in\{-1,1\}^N}
  e^{\beta\sqrt t\,\sigma^TW\sigma/2+x\cdot\sigma}.
\end{equation}
This is the Parisi equation in $\R^N$; its terminal value is the logarithm
of the completed partition function at inverse temperature $\beta\sqrt t$
in the external field $x$. We define the interpolating functional and its
minimum over admissible
profiles by
\begin{equation}\label{eq:fnd-functional}
 \mathcal F_{t,s}(\mu)
   =\frac1{Ns}\log\E_{\rm ref}e^{sU(0,0;W)}
       -\frac c2\int_0^1qa(q)\dd q,
 \qquad f_{N,s}(t)=\min_{\mu}\mathcal F_{t,s}(\mu).
\end{equation}

By \cite[Section~2.6 and Proposition~2.3]{AL}, the equation
\eqref{eq:fnd-vector-pde} has a solution for every admissible profile,
and $\mathcal F_{t,s}$ is continuous under weak convergence of profiles.
For a step profile, the solution is given by the Hopf--Cole recursion that
we recall in
\cref{subsec:fnd-law}. For a general profile, it is the
uniform limit of the solutions for step profiles whose cumulative
functions converge in $L^1([0,1])$. Indeed, if $U_a$ and $U_{\tilde a}$
are the solutions for the cumulative functions $a$ and $\tilde a$, then
\cite[Section~2.6, (2.19)]{AL}
\[
 \sup_{q,x,W}\bigl|U_a(q,x;W)-U_{\tilde a}(q,x;W)\bigr|
 \le\frac{cN}2\int_0^1|a(u)-\tilde a(u)|\dd u .
\]
The admissible profiles form a weakly compact set, since they are the
image of the compact set of measures of mass $1-s$ under
$\nu\mapsto s\delta_0+\nu$, and weak convergence of profiles implies
$L^1$ convergence of their cumulative functions. We conclude that the
minimum in \eqref{eq:fnd-functional} is attained.

We note that minimizers are deterministic. The Gaussian expectation is
taken inside
the functional, which means that $\mathcal F_{t,s}(\mu)$ is a number for
each profile. A minimizer may depend on $N$, $\rho$, and $\lambda$ but not
on the couplings. This matters below, because the tilted law of
\cref{subsec:fnd-law}
is constructed from a deterministic profile, and the estimates at a
minimizer combine that law with the variational inequality that the same
profile satisfies. A minimizer need not be unique; at each use, we say
which minimizer is meant.

The prescribed root mass $s\delta_0$ matches the order $s$ of the
moment. In the derivative of the interpolation, which we compute in
\cref{prop:fnd-interpolation}, the outer logarithmic moment
contributes a term that cancels the root mass exactly, so that only the
free measure $\nu$ remains. In the posterior densities of
\cref{sec:gauss}, the outer tilt cancels it in the same way.

The two endpoints of the interpolation are explicit. At $t=0$, the
terminal value in \eqref{eq:fnd-vector-pde} is $\sum_i\log(2\cosh x_i)$
and $c=\beta^2$. Then the equation separates into $N$ copies of the
scalar Parisi equation of \eqref{eq:intro-parisi}, and we have
$U(q,x;W)=\sum_iu_a(q,x_i)$, where $u_a$ is the solution of that
equation with $\alpha=a$. This shows that $U(0,0;W)=Nu_a(0,0)$ is
deterministic and that
$\mathcal F_{0,s}(\mu)=\mathcal P_\beta(a)$. At
$t=1$, we have $c=0$, which implies that
$U(0,0;W)=U(1,0;W)=F^c_{N,\beta}$ and that
$\mathcal F_{1,s}(\mu)=\phi^c_{N,\beta}(s)$ for every profile.
Minimizing over admissible profiles, we obtain
\cite[Proposition~2.3]{AL}
\begin{equation}\label{eq:fnd-endpoints}
 f_{N,s}(0)=P_\beta(s),\qquad f_{N,s}(1)=\phi^c_{N,\beta}(s).
\end{equation}
For a fixed profile, the map $t\mapsto\mathcal F_{t,s}(\mu)$ is
nonincreasing (\cref{prop:fnd-interpolation}). By
\eqref{eq:fnd-error-definition}, the fractional-moment error is
$(f_{N,s}(0)-f_{N,s}(1))/\beta$, which is the decrease of $f_{N,s}/\beta$
along the interpolation. The defect of \cref{sec:fnd-defect} is four
times the rate of this decrease.

\paragraph{The energy normalization.}
We vary profiles in the energy normalization of
\cite[Section~7.1]{AL}. For an admissible profile $\mu$ with cumulative
function $a$, we set
\[
 \gamma=\beta a,\qquad \nu_\gamma=\beta\nu,\qquad
 F_\rho(\gamma)=\beta^{-1}\mathcal F_{1-\rho,\lambda/\beta}(\mu).
\]
Then $\dd\gamma$ is a measure of mass $\beta$ with
$\dd\gamma\ge\lambda\delta_0$; we also call $\gamma$ admissible. We
have $\min_\gamma F_\rho(\gamma)=\beta^{-1}f_{N,s}(1-\rho)$, where the
minimum is over admissible $\gamma$. The name refers to the following
rescaling of the spatial variable. The function
$\widetilde U(q,y;W)=\beta^{-1}U(q,\beta y;W)$ solves
\[
 \partial_q\widetilde U
 =-\frac\rho2\bigl(\Delta\widetilde U+\gamma(q)|\nabla\widetilde U|^2\bigr),
\]
with terminal value $\beta^{-1}$ times the logarithm of the completed
partition function in the field $\beta y$, and we have
\[
 F_\rho(\gamma)=\frac1{N\lambda}\log\E_{\rm ref}
     e^{\lambda\widetilde U(0,0;W)}
 -\frac\rho2\int_0^1q\gamma(q)\dd q .
\]
This is the functional of \cite[(7.2)]{AL} with $r=\rho$, in which the
diagonal of $W$ has been integrated out explicitly
\cite[Section~7.1]{AL}. Note that the rescaling leaves the gradient unchanged,
$\nabla\widetilde U(q,y;W)=\nabla U(q,\beta y;W)$, and that it multiplies
the Hessian by $\beta$.

We define the \emph{scalar trial profile} by
$\gamma^{\rm tr}_\lambda(q)=\max\{\lambda,\gamma_\beta(q)\}$. It is
nondecreasing and right-continuous, and it satisfies
$\lambda\le\gamma^{\rm tr}_\lambda\le\beta$ and
$\gamma^{\rm tr}_\lambda(1)=\beta$. In other words,
$\gamma^{\rm tr}_\lambda/\beta$ is the cumulative function of an
admissible profile. When $\lambda$ is
fixed, we write $\gamma^{\rm tr}=\gamma^{\rm tr}_\lambda$. We discuss its
cost and its use as a comparison profile in
\cref{sec:fnd-trial}.

\section{The tilted joint law and overlap observables}
\label{subsec:fnd-law}

Fix $0<\rho\le1$ and an admissible profile $\mu$. We recall the joint
law of \cite[Section~2]{AL}, under which all later expectations are
taken. It is built in three stages.
\begin{enumerate}[label=(\roman*)]
\item The Gaussian law of $W$ is tilted by the density
$e^{sU(0,0;W)}/\E_{\rm ref}e^{sU(0,0;W)}$.
\item Given $W$, the auxiliary field solves
\begin{equation}\label{eq:fnd-diffusion}
 \dd X_q=ca(q)\nabla U(q,X_q;W)\dd q+\sqrt c\dd B_q,
 \qquad X_0=0,
\end{equation}
where $B$ is a standard Brownian motion in $\R^N$ under the conditional
law given $W$.
\item Given $W$ and the path $X$, a terminal spin
$\sigma\in\{-1,1\}^N$ is drawn with probabilities proportional to
$\exp\{\beta\sqrt t\,\sigma^TW\sigma/2+X_1\cdot\sigma\}$.
\end{enumerate}
We write $\Pp$ and $\E$ for this joint law, and
$\calG_q=\sigma(W,X_r:0\le r\le q)$ for the information revealed up to
level $q$; note that it does not contain the terminal spin. At fixed $N$, the
drift in \eqref{eq:fnd-diffusion} is bounded and Lipschitz in $x$
\cite[Section~2.6.1]{AL}, and the equation has a unique strong solution.
The tilt in (i) changes the law of $W$, which is no longer Gaussian under
$\Pp$, but it does not change the
conditional law of $B$ given $W$.

For a step profile, the law has a discrete description, which explains
the construction. The Gaussian Hopf--Cole recursion in the Parisi setting
is described in \cite[Section~3, (10)--(13)]{Guerra2003}, and the
finite-dimensional tilted law used here is the formulation of
\cite[Section~2.1]{AL}.
Suppose that the profile has grid points
$0=q_0<q_1<\cdots<q_K=1$ and that $a=a_\ell$ on $[q_\ell,q_{\ell+1})$. Then
\eqref{eq:fnd-vector-pde} is solved by the Hopf--Cole recursion
\[
 U(q_\ell,x;W)=\frac1{a_\ell}\log\E_g
 \exp\bigl\{a_\ell U\bigl(q_{\ell+1},x+\sqrt{c(q_{\ell+1}-q_\ell)}\,g;W\bigr)\bigr\},
\]
where $g$ is a standard Gaussian vector in $\R^N$, and we read the right
side as $\E_gU(q_{\ell+1},x+\sqrt{c(q_{\ell+1}-q_\ell)}\,g;W)$ when
$a_\ell=0$. Given $W$ and $X_{q_\ell}$, the increment
$X_{q_{\ell+1}}-X_{q_\ell}$ has the law of
$\sqrt{c(q_{\ell+1}-q_\ell)}\,g$, tilted by the density
$\exp\{a_\ell(U(q_{\ell+1},X_{q_{\ell+1}};W)-U(q_\ell,X_{q_\ell};W))\}$,
which integrates to one by the recursion. The outer tilt in (i) has the
same form, with mass $s$ at level $0$. In \cref{sec:gauss}, we multiply
these densities and use the resulting formula for the joint density.

\paragraph{Observables.}
We denote the conditional mean, the conditional covariance, and the
spatial Hessian by
\[
 m_q=\E[\sigma\mid\calG_q],\qquad
 C_q=\Cov(\sigma\mid\calG_q),\qquad
 H_q=\nabla_x^2U(q,X_q;W).
\]
We will see in \cref{prop:fnd-contact} that $m_q=\nabla U(q,X_q;W)$ and
that $a(q)C_q\preceq H_q\preceq C_q$. With the same terminal spin at every
level, we define
\[
 S_q=\frac{|m_q|^2}N,\qquad
 L_q=\frac{\sigma\cdot m_q-|m_q|^2/2}N,\qquad
 j_q=\frac{\E[m_q^TC_qm_q]}{N^2},\qquad
 a_q^{\rm th}=\frac{\E\Tr C_q^2}{N^2};
\]
we use the superscript in $a^{\rm th}_q$ to distinguish the covariance
trace from the cumulative function $a$. For each $q$, let $\sigma^1$ and $\sigma^2$
be conditionally independent given $\calG_q$, each with the conditional
law of $\sigma$ given $\calG_q$, and let $R_q=R(\sigma^1,\sigma^2)$. Note
that the self-overlap $S_q$ of the conditional mean is the conditional
expectation of $R_q$ given $\calG_q$, and that $L_q$ has conditional
expectation $S_q/2$ given $\calG_q$. In \cref{prop:fnd-variance}, we bound
the variance of
averages of $L_q$ over levels. We also have
$0\le S_q\le1$ and $|L_q|\le3/2$.

The processes $m$, $C$, and $H$ have continuous paths. Indeed, we have
$m_q=\nabla U(q,X_q;W)$, $H_q=\nabla^2U(q,X_q;W)$, and
$\E[\sigma\sigma^T\mid\calG_q]=\Sigma(q,X_q;W)$, where the conditional
second moment $\Sigma$ solves a linear backward equation with drift
$ca\nabla U$, and $\nabla U$, $\nabla^2U$, and $\Sigma$ are jointly
continuous in $(q,x)$ \cite[Section~2.6.1]{AL}; the path $X$ is
continuous. In particular, $S_q$ and $L_q$ are jointly measurable
in the level and the outcome, and Fubini's theorem applies to their
integrals against $\nu$. By bounded convergence, the functions
$q\mapsto j_q$, $a^{\rm th}_q$, $\E S_q$, $\Var(S_q)$, and $\E|S_q-q|^3$
are continuous on $[0,1]$, and so is $q\mapsto\E(R_q-q)^2$, because
$\E[R_q\mid\calG_q]=S_q$ and
$\E[R_q^2\mid\calG_q]=N^{-2}\sum_{i,j}\Sigma_{ij}(q,X_q;W)^2$.

We also define the integrated quantities
\[
 \begin{gathered}
 A=\int a_q^{\rm th}\,\nu(\dd q),\qquad
 J=\int j_q\,\nu(\dd q),\\
 V=\int\Var(S_q)\,\nu(\dd q),\qquad
 H_3=\int\E|S_q-q|^3\,\nu(\dd q).
 \end{gathered}
\]
We note that they depend on $N$, $\rho$, $\lambda$, and the profile. The trace term
$A$ and the mean--covariance term $J$ are controlled by the conditional
covariance. The integrated variance $V$ of the self-overlap is the main
term of the defect, and $H_3$ is the corresponding cubic moment, which
enters in \cref{sec:c7-closure}.

\paragraph{Quantiles and contact levels.}
The quantile function
\begin{equation}\label{eq:fnd-quantile}
 Q(v)=\inf\{q\in[0,1]:\nu([0,q])\ge v\},\qquad 0<v<1-s,
\end{equation}
pushes Lebesgue measure on $(0,1-s)$ forward to $\nu$, and
$a(Q(v))=s+\nu([0,Q(v)])\ge s+v$, by the right continuity of
$q\mapsto\nu([0,q])$. We use quantiles to integrate against $\nu$ in a
variable in which the lower bound on $a$ is explicit, whether or not
$\nu$ has a density. The image under $Q$ of Lebesgue measure on a Borel
subset of $(0,1-s)$ is a sub-measure of $\nu$, whose density with respect
to $\nu$ takes
values in $[0,1]$. For an interval of values of
$v$, this sub-measure may contain only part of an atom of $\nu$, and by
integration over such an interval, we mean integration against
this sub-measure.

If $\mu$ minimizes $\mathcal F_{t,s}$, we call a point of $\supp\nu$ a
\emph{contact level} of $\mu$. The contact levels are where the free
measure sits, and \cref{prop:fnd-contact} shows that the expected self-overlap
$\E S_q$ equals $q$ at each of them.

\section{Martingale and contact identities}\label{sec:fnd-contact}

In the next proposition, we collect the finite-dimensional martingale,
first-variation, and contact identities of \cite[Proposition~2.1]{AL}.
For the scalar Parisi functional, the stochastic-control representation
and the martingale identities are in
\cite[Theorem~3 and Lemma~2]{AuffingerChen2015Uniqueness}, with a dynamic
programming treatment in \cite{JagannathTobasco2016}. The first variation
and the optimality conditions on the support are in
\cite[Lemma~3.2 and Proposition~1.1]{JagannathTobasco2017PhaseDiagram}
and \cite[Theorem~2 and Proposition~1]{Chen2017Variational}.

\begin{proposition}[Martingale, covariance order, and contact]
\label{prop:fnd-contact}
Let $0<\rho\le1$, and let $\mu=s\delta_0+\nu$ be an admissible profile.
For every $q\in[0,1]$,
\begin{equation}\label{eq:in-martingale}
 m_q=\nabla U(q,X_q;W),\qquad m_0=0,\qquad
 \dd m_q=\sqrt c\,H_q\dd B_q,
 \qquad a(q)C_q\preceq H_q\preceq C_q ,
\end{equation}
and $q\mapsto\E S_q$ is continuously differentiable on $[0,1]$ with
derivative $c\,\E\Tr H_q^2/N$. If $\nu'\ge0$ has mass $1-s$, then the
right derivative of
$\theta\mapsto\mathcal F_{t,s}(s\delta_0+(1-\theta)\nu+\theta\nu')$ at
$\theta=0$ is
\begin{equation}\label{eq:fnd-first-variation}
 D\mathcal F_{t,s}(\mu)[\nu'-\nu]
    =\frac c2\int_0^1G(x)\,(\nu'-\nu)(\dd x),\qquad
 G(x)=\int_x^1(\E S_y-y)\dd y .
\end{equation}
If $\mu$ minimizes $\mathcal F_{t,s}$, then $\supp\nu$ is contained in
the set of minimum points of $G$ on $[0,1]$, and
\begin{equation}\label{eq:in-contact}
 \supp\nu\subset[0,1),\qquad
 \E S_q=q,\qquad \E\Tr H_q^2\le\frac Nc
 \qquad\text{for every contact level }q .
\end{equation}
\end{proposition}

\begin{remark}\label{rem:fnd-root}
A contact level is a point of the support of the free measure $\nu$. At
$q=0$, we have $\E S_0=0$ for every profile, since $m_0=0$. The bound
on $\E\Tr H_0^2$ in \eqref{eq:in-contact} holds when $0\in\supp\nu$;
the prescribed mass $s\delta_0$ alone does not give it, because the
minimization does not act on that mass.
\end{remark}

\begin{proof}[Proof of \cref{prop:fnd-contact}]
We recall the argument of \cite[Proposition~2.1 and Section~2.6]{AL}.
The proof has four steps. We first prove the martingale identities and
the covariance order for step profiles. We then derive the first
variation and deduce the support property and \eqref{eq:in-contact}
at a minimizer. Finally, we extend these statements to general profiles.

\paragraph{Step 1: martingale identities for step profiles.}
Let $\mu$ be a step profile. Then $U$ is smooth in $x$, every spatial
derivative of fixed order is bounded uniformly in $(q,x,W)$ at fixed
$N$, and $U$ is smooth in $q$ between consecutive grid points
\cite[Section~2.6.1]{AL}. Differentiating \eqref{eq:fnd-vector-pde} in
$x$ and applying It\^o's formula along \eqref{eq:fnd-diffusion} on each
interval between grid points, we obtain \cite[(2.5)]{AL}
\[
 \dd\,\nabla U(q,X_q;W)=\sqrt c\,H_q\dd B_q,\qquad
 \dd H_q=-ca(q)H_q^2\dd q+\sqrt c\,\nabla H_q\dd B_q ,
\]
where $\nabla H_q\dd B_q$ denotes the contraction of the third spatial
derivative of $U$ at $(q,X_q)$ with $\dd B_q$. The stochastic integrals are martingales, since their integrands are
bounded. At $q=1$, the
gradient $\nabla U(1,X_1;W)$ is the Gibbs mean of the terminal spin
given $(W,X_1)$, that is, $\E[\sigma\mid\calG_1]$. By the tower property,
we have $\nabla U(q,X_q;W)=\E[\sigma\mid\calG_q]=m_q$ for all $q$.
Replacing $\sigma$ by $-\sigma$ in the sum that defines $U(1,x;W)$, we
see that $U(1,\cdot\,;W)$ is even. The equation preserves this symmetry,
which implies
that $U(q,\cdot\,;W)$ is even and that $m_0=\nabla U(0,0;W)=0$.

For the covariance order, we first note that $C_1=H_1$, because the
Hessian of the logarithm of a partition function is the covariance of
its Gibbs measure. Applying the martingale isometry to $m$, and taking
conditional expectations in the equation for $H$, we find
\[
 C_q=\E[H_1\mid\calG_q]+c\int_q^1\E[H_u^2\mid\calG_q]\dd u,\qquad
 H_q=\E[H_1\mid\calG_q]+c\int_q^1a(u)\E[H_u^2\mid\calG_q]\dd u .
\]
Observe that every matrix in these expressions is positive
semidefinite, and that $a(q)\le a(u)\le1$ for $u\ge q$. Subtracting, we obtain
\[
 C_q-H_q=c\int_q^1\bigl(1-a(u)\bigr)\E[H_u^2\mid\calG_q]\dd u\succeq0
\]
and
\[
 H_q-a(q)C_q=\bigl(1-a(q)\bigr)\E[H_1\mid\calG_q]
 +c\int_q^1\bigl(a(u)-a(q)\bigr)\E[H_u^2\mid\calG_q]\dd u\succeq0,
\]
which proves $a(q)C_q\preceq H_q\preceq C_q$ \cite[(2.6)]{AL}. Finally,
It\^o's formula gives $\dd|m_q|^2=2m_q\cdot\dd m_q+c\Tr H_q^2\dd q$.
Taking expectations, we find that the derivative of $q\mapsto\E S_q$ is
$c\,\E\Tr H_q^2/N$ \cite[(2.7)]{AL}, which is continuous in $q$.

\paragraph{Step 2: the first variation.}
Let $\nu'\ge0$ have mass $1-s$, and let
$\delta a(q)=(\nu'-\nu)([0,q])$, so that the profile
$s\delta_0+(1-\theta)\nu+\theta\nu'$ has cumulative function
$a+\theta\,\delta a$. The derivative $\dot U$ of the solution in
$\theta$ at $\theta=0$ solves the linearized equation
\[
 \partial_q\dot U+\frac c2\Delta\dot U+ca(q)\nabla U\cdot\nabla\dot U
 =-\frac c2\,\delta a(q)|\nabla U|^2,\qquad \dot U(1,\cdot\,;W)=0 .
\]
Its generator is the generator of \eqref{eq:fnd-diffusion}. By the
Feynman--Kac formula, we have
\[
 \dot U(0,0;W)=\frac c2\,\E\Bigl[\int_0^1\delta a(q)|m_q|^2\dd q
 \Bigm|W\Bigr].
\]
When we differentiate the outer logarithmic moment in
\eqref{eq:fnd-functional}, this identity is averaged under the tilted law
of $W$, and the penalty
contributes $-(c/2)\int_0^1q\,\delta a(q)\dd q$. We obtain
\[
 D\mathcal F_{t,s}(\mu)[\nu'-\nu]
 =\frac c2\int_0^1(\E S_q-q)\,\delta a(q)\dd q .
\]
We write $\delta a(q)=\int\ind\{x\le q\}\,(\nu'-\nu)(\dd x)$ and
exchange the order of integration, which is allowed because the
integrand is bounded and $\nu'-\nu$ has finite total variation. This
gives \eqref{eq:fnd-first-variation}, which is \cite[(2.8)]{AL}. We note
that its derivation does not use minimality. In the energy
normalization, the
identity is stated for every admissible profile in
\cite[Section~7.4, (7.22)]{AL}.

\paragraph{Step 3: support and contact at a minimizer.}
Let $\mu$ minimize $\mathcal F_{t,s}$, and let $\nu'\ge0$ have mass
$1-s$. The profiles $s\delta_0+(1-\theta)\nu+\theta\nu'$ with
$0\le\theta\le1$ are admissible, and $\theta=0$ is a minimum point of
$\mathcal F_{t,s}$ along them. This implies that the right derivative in
\eqref{eq:fnd-first-variation} is nonnegative. Taking
$\nu'=(1-s)\delta_x$, we find
\[
 (1-s)G(x)\ge\int G\dd\nu\qquad\text{for every }x\in[0,1].
\]
The function $G$ is continuous on $[0,1]$. Choosing for $x$ a minimum
point of $G$, we see that $\int(G-\min_{[0,1]}G)\dd\nu\le0$, so that
$G=\min_{[0,1]}G$ holds $\nu$-almost everywhere. The set where
$G>\min_{[0,1]}G$ is relatively open in $[0,1]$ and $\nu$-null, and we
conclude that $\supp\nu$ is contained in the set of minimum points of
$G$.

By Step 1, $G$ is twice continuously differentiable, with
\[
 G'(q)=q-\E S_q,\qquad G''(q)=1-\frac cN\E\Tr H_q^2 .
\]
If $q\in\supp\nu\cap(0,1)$, then $q$ is an interior minimum point of
$G$, and therefore $G'(q)=0$ and $G''(q)\ge0$. These are the last two
relations in
\eqref{eq:in-contact}. At $q=0$, we have $G'(0)=-\E S_0=0$, because
$m_0=0$. If $0\in\supp\nu$, then $0$ is a minimum point of $G$, and
$G''(0)\ge0$, since otherwise $G''<0$ near $0$ and $G(h)<G(0)$ for
small $h>0$. At $q=1$, the terminal Gibbs measure gives positive
probability to every spin configuration. Then $|m_{1,i}|<1$ for all $i$,
and $G'(1)=1-\E S_1>0$, which implies that $G(1-h)<G(1)$ for small $h>0$
and that $1$ is not a minimum point of $G$. This proves
$\supp\nu\subset[0,1)$ and \eqref{eq:in-contact}, which is
\cite[(2.9)]{AL}.

\paragraph{Step 4: general profiles.}
Let $\mu$ be an arbitrary admissible profile. Following
\cite[Section~2.6]{AL}, we approximate $\nu$ by its pushforwards $\nu_j$
onto finite meshes that contain $0$, $1$, and any prescribed finite set
of levels, with mesh size tending to zero. These pushforwards are taken under the map that sends each point to
the smallest mesh point greater than or equal to it. Let
$\mu_j=s\delta_0+\nu_j$. Then $\mu_j\Rightarrow\mu$, the cumulative
functions converge in $L^1([0,1])$, and they agree with $a$ at the mesh
points. The spatial derivatives of fixed order are bounded uniformly
over step profiles at fixed $N$. We apply the $L^1$ estimate recalled in
\cref{sec:fnd-functional} to the solutions, and the same comparison
argument to their differentiated equations. Using also the uniform bound on the derivatives and
Gr\"onwall's inequality, we find that the solutions and their first and
second spatial derivatives converge uniformly. We couple the laws for $\mu_j$ and $\mu$
through the same matrix $W$ under its Gaussian reference law and the
same Brownian motion $B$. We also use the same uniform variable to
sample the terminal spin. Since the drifts are bounded and Lipschitz, Gr\"onwall's inequality
implies that $X$, $m$, and $H$ for $\mu_j$ converge to those for $\mu$
uniformly in $q$, in probability. The conditional second moments of the
terminal spin, which solve linear backward equations with drift
$ca\nabla U$, converge in the same way, and so does $C_q$. The ratios of the tilt densities for $\mu_j$ to
the density for $\mu$ converge uniformly to one, since these ratios lie
between $e^{-2sd_j}$ and $e^{2sd_j}$, where
$d_j=(cN/2)\|a_j-a\|_{L^1}\to0$. Then the densities converge in $L^1$ of
the Gaussian reference law. Since the conditional moments and spatial derivatives used here are
bounded at fixed $N$, their convergence in probability, uniformly in the
level, implies that they also converge in every finite $L^p$ norm
needed below.
Using these convergences, we pass \eqref{eq:in-martingale}, the formula for the
derivative of $\E S_q$, its continuity in $q$, and
\eqref{eq:fnd-first-variation} to the limit \cite[Section~2.6.2]{AL}.
If $\mu$ is a minimizer, we apply the limiting first variation directly
to $\mu$; the approximating step profiles need not be minimizers. Step 3
then applies to $\mu$ without change.
\end{proof}

\section{The integrated variance bound and the overlap error}
\label{sec:fnd-variance}

At a minimizer, we use the contact identities to turn the posterior
variance bound
of \cite[Proposition~2.2]{AL} into bounds that depend only on the
cumulative function $a$.

\begin{proposition}[Integrated variance and trace bounds]
\label{prop:fnd-variance}
Let $0<\rho\le1$, and let $\mu=s\delta_0+\nu$ minimize
$\mathcal F_{t,s}$. For every deterministic $k\in L^2(\nu)$,
\begin{equation}\label{eq:in-Lvar}
 \Var\left(\int k(q)L_q\,\nu(\dd q)\right)
      \le\frac1M\int\frac{k(q)^2}{a(q)}\,\nu(\dd q),
 \qquad A\le\frac1{Ms}.
\end{equation}
\end{proposition}

The first bound is \cite[Proposition~2.2, (2.16)]{AL}, and the second
is \cite[Section~2.5, (2.18)]{AL}. We briefly recall the argument. We fix
the terminal spin and apply the
Brascamp--Lieb inequality to the law of the independent Gaussian
coordinates that generate $W$ and the field. For a fixed spin, this law
is uniformly log-concave, although the
marginal law of $W$ under $\Pp$ is not Gaussian.

\begin{proof}[Proof of \cref{prop:fnd-variance}]
We first show that, for every admissible profile and every deterministic
$k\in L^2(\nu)$,
\begin{equation}\label{eq:fnd-L-general}
 \Var\left(\int k(q)L_q\,\nu(\dd q)\right)
 \le\frac1{N^2}\int k(q)^2\,\E\Tr(H_qC_q)\,\nu(\dd q),
\end{equation}
which is \cite[(2.15)]{AL}. Let $\mu$ be a step profile. For
$\varepsilon>0$, we start the field at $X_0=\sqrt\varepsilon z_0$, where
$z_0$ is a standard Gaussian vector included in the outer tilt. We
collect the independent standard Gaussian coordinates that generate $W$,
$z_0$, and the field increments in one vector $Z$. With the terminal spin
fixed, $Z$ has a density proportional to $e^{-\Phi_\sigma(Z)}$, and the
Hessian of $\Phi_\sigma$ is
$\mathcal Q=I+\int\mathcal A_q\,\nu(\dd q)\succeq I$, where
$\mathcal A_q=\nabla_Z^2U(q,X_q;W)\succeq0$ \cite[(2.1)--(2.3)]{AL}. We
derive this density for more general laws in the proof of
\cref{lem:fnd-posterior}. We work in these Gaussian reference
coordinates, in which, at the grid points $q_\ell$ of the profile,
\[
 X_{q_\ell}=\sqrt\varepsilon z_0
   +\sum_{h=1}^{\ell}\sqrt{c(q_h-q_{h-1})}\,z_h ;
\]
the tilted law is obtained by reweighting these coordinates. In this
representation, the field depends on $z_0$ only through the term
$\sqrt\varepsilon z_0$, which gives
$\nabla_{z_0}U(q,X_q;W)=\sqrt\varepsilon m_q$ and shows that
the block of
$\mathcal A_q$ for the coordinates of $z_0$ is $\varepsilon H_q$. Let $v_q$ be the vector whose $z_0$ block is
$(\sigma-m_q)/\sqrt\varepsilon$ and whose other blocks vanish.
Differentiating $L_q$ with the spin fixed, we find
\[
 \nabla_ZL_q=\frac1N\mathcal A_qv_q,\qquad
 v_q^T\mathcal A_qv_q=(\sigma-m_q)^TH_q(\sigma-m_q).
\]
For $Y_k=\int kL_q\,\dd\nu$ and $b=\int k(q)\mathcal A_qv_q\,\nu(\dd q)$,
we have $\nabla_ZY_k=b/N$. For every vector $w$, using the Cauchy--Schwarz
inequality for the forms $\mathcal A_q$ and for the integral against
$\nu$, together with $\int\mathcal A_q\,\dd\nu=\mathcal Q-I\preceq\mathcal Q$,
we have
\[
 |w^Tb|^2\le\Bigl(\int w^T\mathcal A_qw\,\nu(\dd q)\Bigr)
       \int k(q)^2v_q^T\mathcal A_qv_q\,\nu(\dd q)
 \le(w^T\mathcal Qw)\int k(q)^2v_q^T\mathcal A_qv_q\,\nu(\dd q).
\]
The choice $w=\mathcal Q^{-1}b$ yields
$b^T\mathcal Q^{-1}b\le\int k^2v_q^T\mathcal A_qv_q\,\dd\nu$. Then the
Brascamp--Lieb inequality \cite{BrascampLieb1976}, in the form
\cite[Proposition~2.2, (2.14)]{AL}, implies that the conditional variance
of $Y_k$ given $\sigma$ is at most
$N^{-2}\E[b^T\mathcal Q^{-1}b\mid\sigma]$. By the
spin-gauge symmetry of \cite[Section~2.1]{AL}, the conditional law of
$Y_k$ given $\sigma$ does not depend on $\sigma$. Using the law of total
variance, we obtain the same bound for $\Var(Y_k)$. Conditioning on $\calG_q$,
we have
$\E[(\sigma-m_q)^TH_q(\sigma-m_q)\mid\calG_q]=\Tr(H_qC_q)$. We remove the
regularization by dominated convergence as
$\varepsilon\downarrow0$, using $|L_q|\le3/2$, $\Tr(H_qC_q)\le N^2$, and
the Gaussian domination of the outer tilt. The approximation of Step~4 in
the proof of \cref{prop:fnd-contact} then extends \eqref{eq:fnd-L-general} to all
profiles, first for continuous $k$ and then by density for
$k\in L^2(\nu)$ \cite[Section~2.6.2]{AL}.

Now suppose that $\mu$ is a minimizer. Using the fact that the trace of
the product of two positive semidefinite matrices is nonnegative,
together with
$H_q-a(q)C_q\succeq0$, we have
\[
 \Tr H_q^2=\Tr\bigl(H_q(H_q-a(q)C_q)\bigr)+a(q)\Tr(H_qC_q)
 \ge a(q)\Tr(H_qC_q),
\]
and similarly
\[
 \Tr(H_qC_q)=\Tr\bigl((H_q-a(q)C_q)C_q\bigr)+a(q)\Tr C_q^2
 \ge a(q)\Tr C_q^2 .
\]
By \eqref{eq:in-contact}, $\E\Tr H_q^2\le N/c$ at every contact level,
and $\nu$ is carried by its support. Dividing by $a(q)\ge s>0$, we find that
$\E\Tr(H_qC_q)\le N/(ca(q))$ and $\E\Tr C_q^2\le N/(ca(q)^2)$ for
$\nu$-almost every $q$. Inserting the first bound into
\eqref{eq:fnd-L-general} and using $M=Nc$, we obtain the first bound in
\eqref{eq:in-Lvar}. For the second, the quantile function
\eqref{eq:fnd-quantile} and $a(Q(v))\ge s+v$ give
\[
 A\le\frac1M\int\frac{\nu(\dd q)}{a(q)^2}
  =\frac1M\int_0^{1-s}\frac{\dd v}{a(Q(v))^2}
  \le\frac1M\int_0^{1-s}\frac{\dd v}{(s+v)^2}
  =\frac1M\Bigl(\frac1s-1\Bigr)\le\frac1{Ms}. \qedhere
\]
\end{proof}

\begin{remark}\label{rem:fnd-partial-atoms}
Since $k$ in \eqref{eq:in-Lvar} is arbitrary, the bound applies to
sub-measures of $\nu$, such as the image of Lebesgue measure on an
interval of quantiles that contains only part of an atom. If $h$ is a
Borel function with $0\le h\le1$, then applying \eqref{eq:in-Lvar} with
$kh$ in place of $k$ and using $h^2\le h$, we obtain
\[
 \Var\Bigl(\int k(q)h(q)L_q\,\nu(\dd q)\Bigr)
 \le\frac1M\int\frac{k(q)^2h(q)}{a(q)}\,\nu(\dd q).
\]
\end{remark}

We now show that the overlap error splits into the three terms of the
defect.

\begin{lemma}[Decomposition of the overlap error]\label{lem:fnd-decomp}
Let $0<\rho\le1$. For every admissible profile and every $q\in[0,1]$, we
have $\Var(L_q)=j_q+\Var(S_q)/4$. If $q$ is a contact level of a
minimizer, then
\begin{equation}\label{eq:in-decomp}
 \E(R_q-q)^2=a_q^{\rm th}+2j_q+\Var(S_q).
\end{equation}
\end{lemma}

\begin{proof}
Both identities are displayed in \cite[Section~2.5]{AL} before (2.17).
Let $\xi^1=\sigma^1-m_q$ and $\xi^2=\sigma^2-m_q$. Conditionally on
$\calG_q$, these vectors are independent and centered, with covariance
$C_q$, and
\[
 R_q=S_q+\frac1N\bigl(m_q\cdot\xi^1+m_q\cdot\xi^2+\xi^1\cdot\xi^2\bigr).
\]
The three terms in parentheses have conditional mean zero, they are
pairwise orthogonal given $\calG_q$, and their conditional second moments
are $m_q^TC_qm_q$, $m_q^TC_qm_q$, and $\Tr C_q^2$. Taking expectations, we
find that for every admissible profile and every $q$,
\begin{equation}\label{eq:fnd-RS}
 \E(R_q-S_q)^2=a_q^{\rm th}+2j_q .
\end{equation}
Since $\E[R_q-S_q\mid\calG_q]=0$, this implies that
\[
 \E(R_q-q)^2=a_q^{\rm th}+2j_q+\E(S_q-q)^2
 =a_q^{\rm th}+2j_q+\Var(S_q)+(\E S_q-q)^2 .
\]
At a contact level of a minimizer, $\E S_q=q$ by \eqref{eq:in-contact},
which proves \eqref{eq:in-decomp}. Similarly,
$L_q=S_q/2+N^{-1}(\sigma-m_q)\cdot m_q$, where the second term has
conditional mean zero and conditional variance $N^{-2}m_q^TC_qm_q$ given
$\calG_q$. By the law of total variance, we have $\Var(L_q)=j_q+\Var(S_q)/4$.
\end{proof}

\section{The optimized interpolation and the defect}
\label{sec:fnd-defect}

We next recall from \cite[Proposition~2.3]{AL} that the derivative of the
interpolating functional in $t$ is an integral of squared overlap
deviations.

\begin{proposition}[Interpolation identity]\label{prop:fnd-interpolation}
Let $N\ge1$ and $0<s\le1/2$.
\begin{enumerate}[label=(\alph*)]
\item For every admissible profile $\mu$, the map
$t\mapsto\mathcal F_{t,s}(\mu)$ is Lipschitz on $[0,1]$ and continuously
differentiable on $(0,1)$, with
\begin{equation}\label{eq:fnd-fixed-derivative}
 \partial_t\mathcal F_{t,s}(\mu)
   =-\frac{\beta^2}4\int\E(R_q-q)^2\,\nu(\dd q),
\end{equation}
where the expectation is under the tilted law of $\mu$ at $t$. The
integral in \eqref{eq:fnd-fixed-derivative} lies in $[0,4(1-s)]$.
\item The function $f_{N,s}$ is Lipschitz on $[0,1]$. At every
$t\in(0,1)$ at which $f_{N,s}$ is differentiable, every minimizer $\mu$
of $\mathcal F_{t,s}$ satisfies
$f_{N,s}'(t)=\partial_t\mathcal F_{t,s}(\mu)$.
\end{enumerate}
\end{proposition}

\begin{proof}
Part (a) is \cite[Proposition~2.3 and (2.12)]{AL}. The regularity in $t$
and the extension from step profiles to general profiles are proved in
\cite[Section~2.6.3]{AL}. We sketch the computation for a step profile,
which shows where the square comes from. The parameter $t$ enters
$\mathcal F_{t,s}$ in three places: the factor $\beta\sqrt t$ of the
matrix interaction, the coefficient $c=\beta^2(1-t)$ of the field in
\eqref{eq:fnd-vector-pde}, and the penalty. We write $B(q)=\E R_q^2$
and $\Gamma(q)=\E S_q$, and note that $\E R_q=\Gamma(q)$.

\emph{Matrix interaction.} We write
$\beta\sigma^TW\sigma/2=\sum_eg_ev_e(\sigma)$, where $e$ runs over the
pairs $i<j$, with $g_e=g_{ij}$ and $v_e(\sigma)=\beta\sigma_i\sigma_j/\sqrt N$,
and over the diagonal pairs $e=(i,i)$, with $g_e=z_i$ and
$v_e=\beta/\sqrt{2N}$. By the computation of
\cref{sec:fnd-completion}, we have
$\sum_ev_e(\sigma)v_e(\sigma')=\beta^2NR(\sigma,\sigma')^2/2$. We regard
$y_e=\sqrt t\,g_e$ as Gaussian coordinates of variance $t$. By the
Gaussian heat identity, at fixed $c$, the derivative in $t$ of the outer
logarithmic moment is
\[
 \frac1{2N}\,\E\sum_e\bigl(\partial_{y_e}^2U(0,0;W)
   +s\bigl(\partial_{y_e}U(0,0;W)\bigr)^2\bigr).
\]
For a function $b(\sigma)$ of the terminal spin, we add $yb(\sigma)$ to
the terminal exponent. Differentiating the Hopf--Cole recursion twice in
$y$ at $y=0$, we find \cite[(2.10)]{AL}
\[
 \partial_yU(0,0)=\E[b(\sigma)\mid\calG_0],\qquad
 \partial_y^2U(0,0)=\E[b(\sigma)^2\mid\calG_0]
 -\int\E\bigl[\E[b(\sigma)\mid\calG_q]^2\bigm|\calG_0\bigr]\,\mu(\dd q).
\]
Applying this with $b=v_e$, summing over $e$, and using
$\E[v_e(\sigma)\mid\calG_q]^2=\E[v_e(\sigma^1)v_e(\sigma^2)\mid\calG_q]$
and $\mu=s\delta_0+\nu$, we find that the matrix interaction contributes
\[
 \frac{\beta^2}4\Bigl(1-\int B(q)\,\mu(\dd q)+sB(0)\Bigr)
 =\frac{\beta^2}4\Bigl(1-\int B(q)\,\nu(\dd q)\Bigr)
\]
\cite[(2.11)]{AL}. The term $s(\partial_{y_e}U)^2$, which comes from
the outer moment of order $s$, cancels the prescribed root mass.

\emph{Field.} As in Step~2 of the proof of \cref{prop:fnd-contact}, the
derivative of $U$ in $c$ solves the linearized equation with source
$-(\Delta U+a|\nabla U|^2)/2$, and by the Feynman--Kac formula, we have
\[
 \partial_c\Bigl(\frac1{Ns}\log\E_{\rm ref}e^{sU(0,0;W)}\Bigr)
 =\frac12\int_0^1\bigl(h(q)+a(q)\Gamma(q)\bigr)\dd q,
 \qquad h(q)=\frac1N\E\Tr H_q .
\]
By the equation for $H_q$ in Step~1 of the proof of
\cref{prop:fnd-contact} and the derivative of $\Gamma$, we have
$h'=-a\Gamma'$. Further, $h(1)=1-\Gamma(1)$, because $H_1=C_1$ and
$\Tr C_1=N-|m_1|^2$. Integrating by parts, and using that $h$ and $\Gamma$ are continuous
and that $a$ is right-continuous with $a(1)=1$, we find
\[
 h(q)+a(q)\Gamma(q)=1-\int_{(q,1]}\Gamma(r)\,\mu(\dd r)
 \qquad(0\le q\le1).
\]
The interval is open at $q$ because $a(q)$ already contains an atom of
$\mu$ at $q$. Integrating in $q$ and using $\dd c/\dd t=-\beta^2$, we
see that the field contributes
\[
 -\frac{\beta^2}2\Bigl(1-\int q\Gamma(q)\,\nu(\dd q)\Bigr);
\]
here we may replace $\mu$ by $\nu$ because the integrand vanishes at
$q=0$.

\emph{Penalty.} By Fubini's theorem, we have
$2\int_0^1qa(q)\dd q=1-\int q^2\,\mu(\dd q)$. It follows that the
penalty contributes
\[
 \frac{\beta^2}4\Bigl(1-\int q^2\,\nu(\dd q)\Bigr).
\]

Adding the three contributions, in which the constants cancel, we
obtain
\[
 \partial_t\mathcal F_{t,s}(\mu)
 =-\frac{\beta^2}4\int\bigl(B(q)-2q\Gamma(q)+q^2\bigr)\,\nu(\dd q),
\]
which is \eqref{eq:fnd-fixed-derivative}. The integral lies in
$[0,4(1-s)]$ because $|R_q-q|\le2$ and $\nu$ has mass $1-s$.

For (b), the bound in (a) implies that the functions
$t\mapsto\mathcal F_{t,s}(\mu)$ are Lipschitz with the common constant
$\beta^2(1-s)$. Their minimum $f_{N,s}$ is then Lipschitz as well. Let
$t\in(0,1)$ be a point at which $f_{N,s}$ is differentiable, and let
$\mu$ be a minimizer at $t$. For every $h$ such that $t+h\in[0,1]$, we have
$f_{N,s}(t+h)\le\mathcal F_{t+h,s}(\mu)$ and
$f_{N,s}(t)=\mathcal F_{t,s}(\mu)$. Dividing the difference by $h>0$
and letting $h\downarrow0$, we find
$f_{N,s}'(t)\le\partial_t\mathcal F_{t,s}(\mu)$; dividing by $h<0$
gives the reverse inequality. This is \cite[(2.20)]{AL}.
\end{proof}

In particular, $t\mapsto\mathcal F_{t,s}(\mu)$ and $f_{N,s}$ are
nonincreasing. We define the \emph{defect}
\[
 \calD_N(\rho)=-\frac4\beta f_{N,s}'(1-\rho)
\]
at every $\rho\in(0,1)$ such that $f_{N,s}$ is differentiable at
$1-\rho$. Since $f_{N,s}$ is Lipschitz, Rademacher's theorem implies that this
holds for almost every $\rho$, and $\calD_N$ is a measurable function
defined almost everywhere. All pointwise statements about $\calD_N$ are
understood at these values of $\rho$. The defect depends on $N$,
$\beta$, $\rho$, and $\lambda$, and not on a choice of minimizer.

\begin{corollary}[The defect identity]\label{cor:fnd-defect}
Let $N\ge1$ and $0<\lambda\le\beta/2$, and let $s=\lambda/\beta$. At
every $\rho\in(0,1)$ at which $\calD_N(\rho)$ is defined, and for every
minimizer $\mu$ of $\mathcal F_{1-\rho,s}$,
\begin{equation}\label{eq:in-envelope}
 \calD_N(\rho)=\beta(A+2J+V),\qquad
 4e_{N,\beta}(\lambda)=\int_0^1\calD_N(\rho)\dd\rho ,
\end{equation}
where $A$, $J$, and $V$ are computed under the tilted law of $\mu$. In
particular, $\calD_N\ge0$ and $e_{N,\beta}(\lambda)\ge0$.
\end{corollary}

\begin{proof}
Let $\rho$ be as in the statement, and let $\mu$ be a minimizer at
$t=1-\rho$. By \cref{prop:fnd-interpolation}(b) and
\eqref{eq:fnd-fixed-derivative}, we have
$\calD_N(\rho)=\beta\int\E(R_q-q)^2\,\nu(\dd q)$. Recall that the measure
$\nu$ is carried by the contact levels of $\mu$,
at which \eqref{eq:in-decomp}
holds. Integrating \eqref{eq:in-decomp} against $\nu$, we obtain the
first identity. For the second, since $f_{N,s}$ is Lipschitz, we have
\[
 \int_0^1\calD_N(\rho)\dd\rho=-\frac4\beta\int_0^1f_{N,s}'(t)\dd t
 =\frac4\beta\bigl(f_{N,s}(0)-f_{N,s}(1)\bigr)
 =\frac4\beta\bigl(P_\beta(s)-\phi^c_{N,\beta}(s)\bigr),
\]
where we used \eqref{eq:fnd-endpoints} in the last step. By
\eqref{eq:fnd-error-definition}, the right side is
$4e_{N,\beta}(\lambda)$; see also \cite[Proposition~2.3, (2.21)]{AL}.
\end{proof}

The individual terms $A$, $J$, and $V$ may depend on the choice of
minimizer, but by \eqref{eq:in-envelope}, their combination $A+2J+V$
does not. In the later sections, we bound $A$, $J$, and $V$ at an
arbitrary minimizer; such bounds then hold for $\calD_N(\rho)$ at almost
every $\rho$.

Near $\rho=0$, we use a cruder bound. For every $N\ge1$, $0<s\le1/2$,
$0<\rho\le1$, and every minimizer $\mu$ of $\mathcal F_{1-\rho,s}$, we
have
\[
 \int\E(R_q-q)^2\,\nu(\dd q)\le C\min\{1,M^{-1/2},M^{-2/3}s^{-1/3}\}
\]
with a numerical constant $C$. The first two terms are
\cite[Lemma~4.1, (4.1)]{AL}, and the third comes from the proof of
\cite[Theorem~1.1]{AL}; see \cite[(4.2)]{AL}. Keeping the first two
terms and using $\beta M^{-1/2}=(N\rho)^{-1/2}$, we obtain
\begin{equation}\label{eq:fnd-moving}
 \begin{gathered}
 \int\E(R_q-q)^2\,\nu(\dd q)\le C\min\{1,M^{-1/2}\},\\
 0\le\calD_N(\rho)\le C\min\{\beta,(N\rho)^{-1/2}\}
 \qquad(0<\rho\le1),
 \end{gathered}
\end{equation}
where the second bound holds at every $\rho$ at which $\calD_N(\rho)$ is
defined.

\section{The scalar trial profile and the trial budget}
\label{sec:fnd-trial}

The scalar trial profile $\gamma^{\rm tr}_\lambda=\max\{\lambda,\gamma_\beta\}$
raises the Parisi profile $\gamma_\beta$ to the floor $\lambda$. By
\cite[Lemma~7.2]{AL}, $\gamma_\beta$ grows linearly near $0$, and the two
profiles differ only on an interval of length of order $\lambda$, on
which the first variation vanishes at $\gamma_\beta$. The cost of the
floor is of order $\lambda^5$. A short lower plateau and its
fifth-order variational cost also occur in the replica calculation
of~\cite[Appendix~B, (B1)--(B3)]{ParisiRizzo2009}.
Here $\lambda=\beta s$ specifies a mass-floor constraint; we do not
identify it with the replica parameter in that calculation.
The uniform bound used here is the rigorous estimate of
\cite[Lemma~7.2]{AL}. Precisely, by
\cite[Lemma~7.2, (7.3)]{AL}, applied with $\beta_0=\beta$, which uses
$\beta>1$, there exist $\lambda_*>0$ and $C_*<\infty$, depending only on
$\beta$, such that
\[
 \beta^{-1}\mathcal P_\beta(\gamma^{\rm tr}_\lambda/\beta)
 \le E_\beta(0)+C_*\lambda^5
 \qquad(0<\lambda\le\lambda_*).
\]
We set $L_s=\max\{1,C_*\}$ and $\lambda_0=\min\{1,\lambda_*/2,\beta/4\}$.
Recall that $E_\beta$ is nondecreasing, and note that the profile with
cumulative function $\gamma^{\rm tr}_\lambda/\beta$ has mass at least
$\lambda/\beta$ at zero. We obtain
\begin{equation}\label{eq:fnd-scalar-trial-cost}
 0\le E_\beta(\lambda)-E_\beta(0)
 \le\beta^{-1}\mathcal P_\beta(\gamma_\lambda^{\rm tr}/\beta)
                      -E_\beta(0)
 \le L_s\lambda^5\qquad(0<\lambda\le2\lambda_0).
\end{equation}
The bound $\lambda_0\le\lambda_*/2$ allows us to use
\eqref{eq:fnd-scalar-trial-cost} at $2\lambda$, and $\lambda_0\le\beta/4$
ensures that $2\lambda/\beta\le1/2$ for $\lambda\le\lambda_0$, so that
the results above apply with $s=2\lambda/\beta$. The bound
$\lambda_0\le1$ is used in \cref{sec:lz,sec:b7}.

\begin{lemma}[Trial budget]\label{lem:fnd-trial}
For every $N\ge1$, $0<\rho\le1$, and $0<\lambda\le\lambda_0$,
\begin{equation}\label{eq:fnd-trial}
 0\le F_\rho(\gamma_\lambda^{\rm tr})-\min_\gamma F_\rho(\gamma)
       \le e_{N,\beta}(\lambda)+L_s\lambda^5 .
\end{equation}
\end{lemma}

\begin{proof}
The lower bound holds because $\gamma^{\rm tr}_\lambda$ is admissible.
Let $\mu^{\rm tr}$ be the admissible profile with cumulative function
$\gamma^{\rm tr}_\lambda/\beta$. By \cref{prop:fnd-interpolation}(a),
the map $t\mapsto\mathcal F_{t,s}(\mu^{\rm tr})$ is nonincreasing, and
$\mathcal F_{0,s}(\mu^{\rm tr})=\mathcal P_\beta(\gamma^{\rm tr}_\lambda/\beta)$
by \cref{sec:fnd-functional}. Hence
$F_\rho(\gamma^{\rm tr}_\lambda)\le
\beta^{-1}\mathcal P_\beta(\gamma^{\rm tr}_\lambda/\beta)$. Recall that
the optimized value $f_{N,s}$ is also nonincreasing. Using
\eqref{eq:fnd-endpoints} and \eqref{eq:fnd-error-definition}, we find
\[
 \min_\gamma F_\rho(\gamma)=\beta^{-1}f_{N,s}(1-\rho)
 \ge\beta^{-1}f_{N,s}(1)=\beta^{-1}\phi^c_{N,\beta}(s)
 =E_\beta(\lambda)-e_{N,\beta}(\lambda).
\]
Subtracting the two bounds, we obtain
\[
 F_\rho(\gamma^{\rm tr}_\lambda)-\min_\gamma F_\rho(\gamma)
 \le e_{N,\beta}(\lambda)
 +\beta^{-1}\mathcal P_\beta(\gamma^{\rm tr}_\lambda/\beta)-E_\beta(\lambda)
 \le e_{N,\beta}(\lambda)+L_s\lambda^5,
\]
where in the last step we used \eqref{eq:fnd-scalar-trial-cost} and
$E_\beta(\lambda)\ge E_\beta(0)$. This is
\cite[Proposition~7.7, (7.44)]{AL}.
\end{proof}

We call a number $\eta\ge0$ that bounds the left side of
\eqref{eq:fnd-trial} a \emph{trial budget}. By \cref{lem:fnd-trial},
$\eta=e_{N,\beta}(\lambda)+L_s\lambda^5$ is a trial budget
simultaneously for every $\rho\in(0,1]$ and every minimizer. In
\cref{sec:lz,sec:lz-root}, we use the trial budget to control the
distance between a minimizer and the scalar trial profile, and in
\cref{sec:b7}, it determines the cutoff between the two regimes of the proof.

\section{Reduction of the mean bias}\label{sec:fnd-mean-bias}

We now reduce the mean bias to the fractional-moment error by the
midpoint argument of Step~1 in the proof of \cite[Theorem~7.8]{AL}. We
define
\[
 \psi(\lambda)=\frac1{N\lambda}
       \log\E_{\rm ref}e^{\lambda(F_{N,\beta}/\beta+D_N)}\quad(\lambda>0),
 \qquad\psi(0)=\frac{p_{N,\beta}}\beta ,
\]
so that $\psi(\lambda)=E_\beta(\lambda)-e_{N,\beta}(\lambda)$ for
$0<\lambda<\beta$, by \eqref{eq:fnd-error-definition}. As a function
of the $N(N+1)/2$ independent standard Gaussian coordinates
$(g_{ij},z_i)$, the completed free energy $F^c_{N,\beta}/\beta$ is
convex, because it is $\beta^{-1}$ times the logarithm of a sum of
exponentials of linear functions. Its partial derivatives are
$N^{-1/2}\langle\sigma_i\sigma_j\rangle$ in $g_{ij}$, where
$\langle\cdot\rangle$ is the Gibbs average of $F^c_{N,\beta}$, and
$(2N)^{-1/2}$ in $z_i$. We find that the squared norm of its gradient is
at most
\[
 \frac1N\cdot\frac{N(N-1)}2+N\cdot\frac1{2N}=\frac N2,
\]
and the function is $\sqrt{N/2}$-Lipschitz. By \cite[Lemma~5.2]{AL}, which is a
special case of Chen's convexity theorem \cite[Theorem~1]{ChenConvexity},
the function
$\lambda\mapsto\lambda^{-1}\log\E_{\rm ref}e^{\lambda F^c_{N,\beta}/\beta}=N\psi(\lambda)$
extends to a convex function on $[0,\infty)$ whose value at zero is
$\E_{\rm ref}F^c_{N,\beta}/\beta=N\psi(0)$, where we used
$\E_{\rm ref}D_N=0$. In particular, $\psi$ is convex on $[0,\infty)$.

\begin{proposition}[Mean bias]\label{prop:fnd-mean-bias}
For every $N\ge1$ and $0<\lambda\le\lambda_0$,
\begin{equation}\label{eq:fnd-mean-bias}
 0\le\frac{P_\beta-p_{N,\beta}}\beta
 \le E_\beta(0)-2E_\beta(\lambda)+E_\beta(2\lambda)
                         +2e_{N,\beta}(\lambda)
 \le 32L_s\lambda^5+2e_{N,\beta}(\lambda).
\end{equation}
\end{proposition}

\begin{proof}
The lower bound is \cite[Theorem~1.1, (1.5)]{AL}; it is also Guerra's
bound \cite{Guerra2003}. Since $\psi$ is convex on $[0,\infty)$, we have
$\psi(\lambda)\le(\psi(0)+\psi(2\lambda))/2$, that is,
$\psi(0)\ge2\psi(\lambda)-\psi(2\lambda)$. Recall that
$\psi(\lambda)=E_\beta(\lambda)-e_{N,\beta}(\lambda)$. Since
$\lambda\le\lambda_0\le\beta/4$, we have $2\lambda/\beta\le1/2$. Then
\cref{cor:fnd-defect} applies at the tilt $2\lambda$ and implies that
$e_{N,\beta}(2\lambda)\ge0$, that is, $\psi(2\lambda)\le E_\beta(2\lambda)$.
Using $E_\beta(0)=P_\beta/\beta$, we obtain
\[
 \frac{P_\beta-p_{N,\beta}}\beta=E_\beta(0)-\psi(0)
 \le E_\beta(0)-2E_\beta(\lambda)+2e_{N,\beta}(\lambda)+E_\beta(2\lambda),
\]
which is the middle inequality in \eqref{eq:fnd-mean-bias}. For the last
inequality, we use $E_\beta(\lambda)\ge E_\beta(0)$ and
\eqref{eq:fnd-scalar-trial-cost} at $2\lambda\le2\lambda_0$:
\[
 E_\beta(0)-2E_\beta(\lambda)+E_\beta(2\lambda)
 \le E_\beta(2\lambda)-E_\beta(0)\le L_s(2\lambda)^5=32L_s\lambda^5 .
 \qedhere
\]
\end{proof}

In \cref{sec:b7}, we take $\lambda=N^{-7/50}$, so that
$\lambda^5=N^{-7/10}$. \Cref{prop:fnd-mean-bias} then reduces
\cref{thm:main} to the bound $e_{N,\beta}(N^{-7/50})\le C_\beta N^{-7/10}$
for all sufficiently large $N$.

\chapter{Gaussian tools for arbitrary profiles}\label{sec:gauss}

The results of \cref{sec:fnd} concern the law of \cref{subsec:fnd-law},
and the sharper ones hold only at a minimizing profile. In
\cref{sec:sc,sec:c7,sec:c7-endpoint}, we also need laws whose profile
minimizes no functional, and laws whose auxiliary field has different
variances in different coordinates. In this section, we introduce a
class of generalized laws that contains these cases
(\cref{def:gauss-law}), and we prove two Gaussian tools for them.

The first tool is \cref{lem:fnd-posterior}, which shows that the
martingale identities and the covariance order hold for every generalized
law. It also bounds the variance of averages over
levels of the squared conditional mean on a set of coordinates. We note
that this bound holds for every profile and that its proof does not use
contact identities. We use it in \cref{sec:sc} for the
scalar law, which is the generalized law with $N=1$, field variance $d$,
no couplings, and root mass $s=0$. We also use it in
\cref{sec:c7,sec:c7-endpoint} for the cavity laws, which have diffusion
matrix $\operatorname{diag}(cI_n,b(u))$ and root mass $s=\lambda/\beta$.

The second tool is \cref{lem:fnd-path-heat}, a version of the Gaussian
heat identity for normalized expectations. Here we take expectations of
functionals of finitely many Gaussian
variables and continuous Gaussian paths, and we hold the profile measure
fixed. We use the lemma in
\cref{sec:sc} to differentiate the scalar law under transports of its
levels, and in \cref{sec:c7,sec:c7-endpoint} to differentiate the cavity
laws in the cavity parameter. In this section, we use only the definitions of
\cref{sec:fnd}, the arguments of \cite[Section~2]{AL}, and the proof of
\cref{prop:fnd-contact}, which we adapt.

Note that some letters have a local meaning in this section. The letter
$\omega$ denotes Gaussian couplings, while in \cref{sec:c7}, $\omega$
denotes the
weights of a tree. In \cref{sec:gauss-heat}, the derivative measures $M$
and $Q$ are unrelated to $M=Nc$, to the quantile function
\eqref{eq:fnd-quantile}, and to the vector $M_q$ of \cref{sec:c7}.

\section{Generalized laws and a posterior variance bound}
\label{sec:gauss-posterior}

In \cref{sec:c7}, we decouple one spin from the others. The remaining
spins then evolve under a law whose profile need not minimize any
functional and whose auxiliary field has unequal variances. The
following class of laws contains both situations.

\begin{definition}[Generalized law]\label{def:gauss-law}
Let $\Lambda=\operatorname{diag}(\vartheta_1,\ldots,\vartheta_N)$ with
$\vartheta_i>0$, and let $(\kappa_{ij})_{1\le i\le j\le N}$ be real
numbers; $\Lambda$ and $\kappa$ are deterministic. Let
$\omega=(\omega_{ij})_{1\le i\le j\le N}$ be independent standard
Gaussian variables, and write $\E_{\rm ref}$ for expectation under their
law. Let $\mu=s\delta_0+\nu$ be either an admissible profile in the
sense of \cref{sec:fnd-functional}, so that $0<s\le1/2$, or, with
$s=0$, an arbitrary deterministic probability measure $\mu=\nu$ on
$[0,1]$. Let $a(q)=\mu([0,q])$, and let $U$ solve
\begin{equation}\label{eq:fnd-general-pde}
 \partial_qU=-\frac12\bigl(\Tr(\Lambda\nabla^2U)
      +a(q)\nabla U^T\Lambda\nabla U\bigr),\qquad
 U(1,x;\omega)=\log\sum_{\sigma\in\{-1,1\}^N}
     e^{\mathcal H_\sigma(\omega,x)},
\end{equation}
where
$\mathcal H_\sigma(\omega,x)=\sum_{i\le j}\kappa_{ij}\omega_{ij}\sigma_i\sigma_j+x\cdot\sigma$.
The \emph{generalized law} of $(\Lambda,\kappa,\mu)$ is the joint law of
$(\omega,X,\sigma)$ built in three stages.
\begin{enumerate}[label=(\roman*)]
\item The Gaussian law of $\omega$ is tilted by the density
$e^{sU(0,0;\omega)}/\E_{\rm ref}e^{sU(0,0;\omega)}$.
\item Given $\omega$, the field solves
$\dd X_q=a(q)\Lambda\nabla U(q,X_q;\omega)\dd q+\Lambda^{1/2}\dd B_q$
with $X_0=0$, where $B$ is a standard Brownian motion in $\R^N$.
\item Given $\omega$ and $X$, the spin $\sigma$ is drawn with
probabilities proportional to $e^{\mathcal H_\sigma(\omega,X_1)}$.
\end{enumerate}
As in \cref{subsec:fnd-law}, we write $\Pp$ and $\E$ for this law, and
we define $\calG_q=\sigma(\omega,X_r:0\le r\le q)$,
$m_q=\E[\sigma\mid\calG_q]$, $C_q=\Cov(\sigma\mid\calG_q)$, and
$H_q=\nabla_x^2U(q,X_q;\omega)$.
\end{definition}

When $s=0$, the outer density in (i) is identically one. Recall that a
Hopf--Cole step of mass zero is an
ordinary Gaussian expectation. The
fractional-moment functionals of \cref{sec:fnd} still have $s>0$. The
equation \eqref{eq:fnd-general-pde} has a solution for every such
profile, which is obtained from step profiles by the Hopf--Cole recursion and
uniform approximation, as we show in Step~6 of the proof of
\cref{lem:fnd-posterior}. With $\Lambda=cI$, with $\omega_{ij}=g_{ij}$
and $\kappa_{ij}=\beta\sqrt{t/N}$ for $i<j$, and with $\omega_{ii}=z_i$
and $\kappa_{ii}=\beta\sqrt{t/(2N)}$, we have
$\mathcal H_\sigma(\omega,x)=\beta\sqrt t\,\sigma^TW\sigma/2+x\cdot\sigma$.
In this case, the generalized law is the law of \cref{subsec:fnd-law}. The scalar
law of \cref{sec:sc} is the generalized law with $N=1$, $\Lambda=d$,
$\kappa=0$, and $s=0$. The cavity laws of \cref{sec:c7} are generalized
laws with $\Lambda=\operatorname{diag}(cI_n,b(u))$ and
independent Gaussian couplings with deterministic variances.

In the next lemma, we adapt the argument of the proof of
\cite[Lemma~2.3]{OpenAIFluctuationScale}, which is also used in the proof
of \cite[Lemma~2.1]{OpenAILimitingLaw}. That argument combines joint
convexity in the Gaussian drivers and a spin source, the inverse-Hessian
Brascamp--Lieb inequality, and a gauge symmetry that removes the
conditioning on the terminal spin. We use
it for averages of squared conditional means on selected coordinates, and
for diagonal anisotropic laws.

\begin{lemma}[Averaged self-overlap under an arbitrary profile]
\label{lem:fnd-posterior}
Let $\mu=s\delta_0+\nu$ be an admissible profile, and consider either
its joint law from \cref{subsec:fnd-law}, for some $0<\rho\le1$, or a
generalized law; a generalized law also allows $s=0$ and an arbitrary
probability measure $\mu=\nu$. Then
$m_q=\nabla U(q,X_q)=\E[\sigma\mid\calG_q]$, $m_0=0$, and
$a(q)C_q\preceq H_q\preceq C_q$ for every $q\in[0,1]$. Fix a set of
$n\ge1$ coordinates, let $m_q^-$ be the restriction of $m_q$ to these
coordinates, and let $H_q^{--}$ and $C_q^{--}$ be the corresponding
principal submatrices of $H_q$ and $C_q$. For every deterministic
$k\in L^2(\nu)$,
\begin{equation}\label{eq:fnd-posterior}
 \begin{aligned}
 \Var\left(\int k(q)\frac{|m_q^-|^2}{n}\,\nu(\dd q)\right)
 &\le\frac4{n^2}\int k(q)^2
       \E[(m_q^-)^TH_q^{--}m_q^-]\,\nu(\dd q)\\
 &\le\frac4{n^2}\int k(q)^2
       \E[(m_q^-)^TC_q^{--}m_q^-]\,\nu(\dd q).
 \end{aligned}
\end{equation}
\end{lemma}

Recall that the law of \cref{subsec:fnd-law} is a generalized law, so
that it suffices to prove the lemma for generalized laws. The proof also
gives the
martingale differential
\begin{equation}\label{eq:gauss-martingale}
 \dd m_q=H_q\Lambda^{1/2}\dd B_q\qquad(0\le q\le1)
\end{equation}
for every generalized law, as we will see in Steps~1 and~6 below.

\begin{proof}[Proof of \cref{lem:fnd-posterior}]
The proof has six steps. We first prove the martingale identities and
the covariance order for step profiles. Next, still for step profiles,
we write the law in terms of independent Gaussian coordinates. We
compare the curvature of the posterior density with the gradient of the
squared conditional mean, and we apply the Brascamp--Lieb inequality
with the terminal spin fixed. We then remove the conditioning on the
spin by a gauge symmetry. Finally, we pass to general profiles and
general $k$.

\paragraph{Step 1: martingale and covariance order.}
Let $\mu$ be a step profile. We adapt Step~1 of the proof of
\cref{prop:fnd-contact}, replacing $cI$ by $\Lambda$. The generator of
the field is
$\Tr(\Lambda\nabla^2)/2+a(q)(\Lambda\nabla U)\cdot\nabla$.
Differentiating \eqref{eq:fnd-general-pde} in $x$ and applying It\^o's
formula along the field on each interval between grid points, we find
that
\[
 \dd\,\nabla U(q,X_q;\omega)=H_q\Lambda^{1/2}\dd B_q,\qquad
 \dd H_q=-a(q)H_q\Lambda H_q\dd q+\nabla H_q\Lambda^{1/2}\dd B_q ,
\]
where the last term is the contraction of the third spatial derivative
of $U$ with $\Lambda^{1/2}\dd B_q$. The stochastic integrals are martingales,
since the spatial derivatives are bounded at fixed $N$. The
gradient $\nabla U(1,X_1;\omega)$ is the Gibbs mean of the terminal spin.
By the tower property, we have $\nabla U(q,X_q;\omega)=\E[\sigma\mid\calG_q]=m_q$,
which proves \eqref{eq:gauss-martingale} for step profiles. Using
$C_1=H_1$, the martingale isometry, and the equation for $H$, we find
\[
 \begin{gathered}
 C_q=\E[H_1\mid\calG_q]+\int_q^1\E[H_u\Lambda H_u\mid\calG_q]\dd u,
 \\
 H_q=\E[H_1\mid\calG_q]+\int_q^1a(u)\E[H_u\Lambda H_u\mid\calG_q]\dd u .
 \end{gathered}
\]
Note that the matrices $H_u\Lambda H_u$ are positive semidefinite and
that $a(q)\le a(u)\le1$ for $u\ge q$. Subtracting the two identities as in
Step~1 of the proof of \cref{prop:fnd-contact}, we obtain
$a(q)C_q\preceq H_q\preceq C_q$. Note that
$\mathcal H_{-\sigma}(\omega,x)=\mathcal H_\sigma(\omega,-x)$. Replacing
$\sigma$ by $-\sigma$ in the terminal sum, we see that
$U(1,\cdot\,;\omega)$ is even. The equation preserves this symmetry,
which implies that $U(q,\cdot\,;\omega)$ is even and that $m_0=0$.

\paragraph{Step 2: Gaussian coordinates and the posterior density.}
Let $0=q_0<q_1<\cdots<q_K=1$ and $a=a_\ell$ on $[q_\ell,q_{\ell+1})$,
with $a_\ell\ge s\ge0$. Then the cumulative function is constant between
grid points,
and the free measure $\nu$ is carried by the grid points, with
$\nu(\{q_0\})=a_0-s$, $\nu(\{q_\ell\})=a_\ell-a_{\ell-1}$ for
$1\le\ell\le K-1$, and $\nu(\{q_K\})=1-a_{K-1}$. Let
$z_1,\ldots,z_K$ be independent standard Gaussian vectors in $\R^N$, let
\[
 X_{q_\ell}=\sum_{h=1}^\ell(q_h-q_{h-1})^{1/2}\Lambda^{1/2}z_h
 \qquad(0\le\ell\le K),
\]
and let $Z$ collect the coordinates of $\omega$ and of
$z_1,\ldots,z_K$. We write $U_q(Z)=U(q,X_q;\omega)$ at the grid points.
We recall that \eqref{eq:fnd-general-pde} is solved by the Hopf--Cole
recursion
$U_{q_\ell}=a_\ell^{-1}\log\E_{z_{\ell+1}}e^{a_\ell U_{q_{\ell+1}}}$ when
$a_\ell>0$, and by $U_{q_\ell}=\E_{z_{\ell+1}}U_{q_{\ell+1}}$ when
$a_\ell=0$. We follow the construction of \cite[(2.1)--(2.3)]{AL}. We claim
that, restricted to $\omega$, to the field at the grid points, and to the
spin, the generalized law is the law of
$(\omega,(X_{q_\ell})_\ell,\sigma)$ when $(Z,\sigma)$ has the density
\[
 \frac{\exp\{\mathcal H_\sigma(\omega,X_1)
                    -\int U_q(Z)\,\nu(\dd q)\}}
      {\E_{\rm ref}e^{sU(0,0;\omega)}}
\]
with respect to the standard Gaussian law of $Z$ and counting measure on
$\{-1,1\}^N$. Indeed, the outer tilt contributes
$sU_{q_0}$ to the exponent, the transition from $q_\ell$ to
$q_{\ell+1}$ contributes $a_\ell(U_{q_{\ell+1}}-U_{q_\ell})$, and the
terminal Gibbs probability contributes $\mathcal H_\sigma(\omega,X_1)-U_{q_K}$.
In the sum, we see that the coefficients of
$U_{q_0},U_{q_1},\ldots,U_{q_K}$ are
\[
 s-a_0,\quad a_0-a_1,\quad\ldots,\quad a_{K-2}-a_{K-1},\quad a_{K-1}-1,
\]
which are the negatives of the masses of $\nu$ at the grid points.
A transition of mass zero has likelihood one, and the calculation also
holds at $s=0$, where the denominator equals one. For a fixed spin
$\sigma$, we find that the negative logarithm of the conditional density
of $Z$ is, up to a constant, the posterior potential
\[
 \Phi_\sigma(Z)=\frac12|Z|^2-\mathcal H_\sigma(\omega,X_1)
           +\int U_q(Z)\,\nu(\dd q).
\]
Note that $\mathcal H_\sigma(\omega,X_1)$ is linear in $Z$ and that each
$U_q$ is convex in $Z$. To see the convexity, observe that the terminal value
in \eqref{eq:fnd-general-pde} is a log-sum-exp of linear functions of
$(\omega,x)$, that each Hopf--Cole step of positive mass preserves
convexity by H\"older's inequality, and that Gaussian expectation at
mass zero preserves it as well. Since $X_q$ is linear in $Z$,
we obtain
\[
 \nabla_Z^2\Phi_\sigma=\mathcal Q=I+\int\mathcal A_q\,\nu(\dd q)\succeq I,
 \qquad
 \mathcal A_q=\nabla_Z^2U_q\succeq0 .
\]
The matrix $\mathcal Q$ does not depend on $\sigma$.

\paragraph{Step 3: joint source curvature.}
Let $\iota\colon\R^n\to\R^N$ embed a vector in the selected coordinates,
and let $\mathcal U_q(Z,y)=U(q,X_q(Z)+\iota y;\omega)$ for
$y\in\R^n$. This function is jointly convex in $(Z,y)$, since $U(q,x;\omega)$ is
convex in $(\omega,x)$ by Step~2 and the map
$(Z,y)\mapsto(\omega,X_q(Z)+\iota y)$ is linear. We write
$\mathcal B_q=\nabla_Z\nabla_y\mathcal U_q(Z,0)$, a matrix with $n$
columns. Since $\nabla_y\mathcal U_q(Z,0)=m_q^-$ and
$\nabla_y^2\mathcal U_q(Z,0)=H_q^{--}$, joint convexity and the chain
rule give
\begin{equation}\label{eq:fnd-source-curvature}
 \begin{pmatrix}
  \mathcal A_q&\mathcal B_q\\
  \mathcal B_q^T&H_q^{--}
 \end{pmatrix}\succeq0,\qquad
 \nabla_Z|m_q^-|^2=2\mathcal B_qm_q^- .
\end{equation}
For a vector $w$ and $r\in\R$, we evaluate this positive semidefinite
form at $(rw,m_q^-)$. The resulting quadratic polynomial in $r$ is
nonnegative. This implies that its discriminant is nonpositive:
\[
 |w^T\mathcal B_qm_q^-|^2
 \le(w^T\mathcal A_qw)\,(m_q^-)^TH_q^{--}m_q^- .
\]
The same statements hold at $q=0$, where $m_0^-=0$.

\paragraph{Step 4: the Brascamp--Lieb inequality.}
Let $Y=n^{-1}\int k(q)|m_q^-|^2\,\nu(\dd q)$, a finite sum over grid
points, and let $b=\int k(q)\mathcal B_qm_q^-\,\nu(\dd q)$, so that
$\nabla_ZY=2b/n$ by \eqref{eq:fnd-source-curvature}. For every vector $w$,
using the last display, the Cauchy--Schwarz
inequality for the integral against $\nu$, and
$\int\mathcal A_q\,\nu(\dd q)=\mathcal Q-I\preceq\mathcal Q$, we obtain
\[
 \begin{aligned}
 |w^Tb|^2&\le\Bigl(\int w^T\mathcal A_qw\,\nu(\dd q)\Bigr)
             \int k(q)^2(m_q^-)^TH_q^{--}m_q^-\,\nu(\dd q)\\
 &\le(w^T\mathcal Qw)\int k(q)^2(m_q^-)^TH_q^{--}m_q^-\,\nu(\dd q).
 \end{aligned}
\]
Taking $w=\mathcal Q^{-1}b$ and dividing by $b^T\mathcal Q^{-1}b$ when it is
positive, we find that
\[
 b^T\mathcal Q^{-1}b\le\int k(q)^2(m_q^-)^TH_q^{--}m_q^-\,\nu(\dd q);
\]
if $b^T\mathcal Q^{-1}b=0$, this inequality holds trivially. Fix $\sigma$. The
potential $\Phi_\sigma$ is smooth, its Hessian $\mathcal Q$ is at least
$I$, and it has Gaussian tails. By the Brascamp--Lieb inequality
\cite{BrascampLieb1976}, in the form
\cite[Proposition~2.2, (2.14)]{AL}, we have
$\Var(Y\mid\sigma)\le\E[\nabla_ZY^T\mathcal Q^{-1}\nabla_ZY\mid\sigma]$,
and we conclude that
\[
 \Var(Y\mid\sigma)\le\frac4{n^2}
   \E\Bigl[\int k(q)^2(m_q^-)^TH_q^{--}m_q^-\,\nu(\dd q)\Bigm|\sigma\Bigr].
\]

\paragraph{Step 5: removing the conditioning.}
For $\tau\in\{-1,1\}^N$, we consider the map
\[
 \sigma_i\mapsto\tau_i\sigma_i,\qquad
 \omega_{ij}\mapsto\tau_i\tau_j\omega_{ij},\qquad
 z_{h,i}\mapsto\tau_iz_{h,i}.
\]
This map preserves the Gaussian reference law. Because $\Lambda$ is diagonal,
it multiplies each $X_q$ coordinatewise by $\tau$. Further,
$\mathcal H_{\tau\sigma}$ evaluated at the transformed $(\omega,x)$
equals $\mathcal H_\sigma(\omega,x)$, and the map leaves every $U_q$
invariant. It therefore preserves the joint density of Step~2. It maps
$m_{q,i}$ to $\tau_im_{q,i}$ and $(H_q)_{ij}$ to $\tau_i\tau_j(H_q)_{ij}$,
so that $Y$ and $(m_q^-)^TH_q^{--}m_q^-$ are invariant. The map acts
transitively on spin configurations, which implies that the conditional
law of these quantities given $\sigma$ does not depend on $\sigma$. In
particular, $\E[Y\mid\sigma]$ is constant, and the law of total variance
implies that $\Var(Y)=\E\Var(Y\mid\sigma)$. Combining this with Step~4, we
obtain the first inequality in
\eqref{eq:fnd-posterior} for step profiles. The second inequality follows
from $H_q^{--}\preceq C_q^{--}$,
since principal submatrices preserve the order $H_q\preceq C_q$.

\paragraph{Step 6: general profiles and general $k$.}
We approximate $\mu$ by step profiles with the same root mass, as in
Step~4 of the proof of \cref{prop:fnd-contact} and
\cite[Section~2.6]{AL}. That is, we push $\nu$ forward to finite meshes that
contain $0$, $1$, and any prescribed finite set of levels, with mesh size
tending to zero. The argument of \cite[Section~2.6]{AL} uses only that
the spins are bounded, that spatial derivatives of fixed order are
bounded uniformly over step profiles at fixed $N$, that the solution is
Lipschitz in the profile in $L^1$, that the outer tilt has Gaussian
domination, and that the disorder, the Brownian motion, and the terminal
spin can be coupled. These facts carry over to a bounded diagonal
$\Lambda$ and to couplings $\kappa_{ij}\omega_{ij}$. For the Lipschitz
bound, let $U$ and $\widetilde U$ be the solutions for the cumulative
functions $a$ and $\tilde a$, and let $\varphi=U-\widetilde U$. Subtracting
the two equations, we find
\[
 \partial_q\varphi+\frac12\Tr(\Lambda\nabla^2\varphi)
 +\frac a2(\nabla U+\nabla\widetilde U)^T\Lambda\nabla\varphi
 =-\frac{a-\tilde a}{2}\,\nabla\widetilde U^T\Lambda\nabla\widetilde U,
 \qquad \varphi(1,\cdot\,;\omega)=0 .
\]
Since each coordinate of $\nabla\widetilde U$ is a conditional spin
mean, we have $\nabla\widetilde U^T\Lambda\nabla\widetilde U\le\Tr\Lambda$. By
the backward representation of $\varphi$, we have
\[
 \sup_{x,\omega}|\varphi(q,x;\omega)|\le\frac{\Tr\Lambda}2\int_q^1|a(u)-\tilde a(u)|\dd u
 \le\frac N2\max_i\vartheta_i\,\|a-\tilde a\|_{L^1([0,1])}.
\]
This replaces the constant $cN/2$ of \cite[(2.19)]{AL}. It follows that
the ratio of the outer densities for $\mu$ and for an approximating step
profile converges to one uniformly in $\omega$. At $s=0$, the outer
density is identically one, and the approximation requires no division
by a root mass.

Under the common coupling of \cite[Section~2.6.1]{AL}, the processes
$X$, $m$, $H$, and $C$ for the step profiles converge to those for $\mu$,
uniformly in $q$, in probability. Then the identities of Step~1, including
\eqref{eq:gauss-martingale},
pass to the limit, where for \eqref{eq:gauss-martingale} we use the It\^o
isometry and the bounded
convergence of $H$. For continuous $k$, the integrands
$q\mapsto k(q)|m_q^-|^2$ and $q\mapsto k(q)^2(m_q^-)^TH_q^{--}m_q^-$ are
bounded and continuous along each path. From any subsequence we may
extract a further subsequence on which the process convergence is
almost surely uniform in $q$. On that subsequence the integrands
converge uniformly, and their integrals then converge by weak
convergence of the deterministic free measures. The integrands are bounded, the ratio of the outer densities converges
to one uniformly in $\omega$, and, since $U(0,0;\omega)$ grows at most
linearly in $\omega$, the outer densities are dominated, uniformly along
the approximation, by a function that is integrable under the Gaussian
reference law. This allows passage to the expectations and variances. Thus \eqref{eq:fnd-posterior} holds for continuous $k$.
Next, we
show that both sides of \eqref{eq:fnd-posterior} are continuous
in $k\in L^2(\nu)$. Since $|m_q^-|^2\le n$, we have
\[
 \Bigl\|\frac1n\int\bigl(k_1(q)-k_2(q)\bigr)|m_q^-|^2\,\nu(\dd q)\Bigr\|_{L^2(\Pp)}
 \le\nu([0,1])^{1/2}\|k_1-k_2\|_{L^2(\nu)},
\]
which shows that the left side of \eqref{eq:fnd-posterior} is
continuous in $k$. Further,
$(m_q^-)^TH_q^{--}m_q^-\le(m_q^-)^TC_q^{--}m_q^-\le n\|C_q\|_{\rm op}\le nN$,
which implies that the right sides change by at most
$4n^{-1}N\|k_1-k_2\|_{L^2(\nu)}\|k_1+k_2\|_{L^2(\nu)}$. Since continuous
functions are dense in $L^2(\nu)$, this proves the lemma.
\end{proof}

\section{Differentiating normalized Gaussian path expectations}
\label{sec:gauss-heat}

In \cref{sec:sc,sec:c7,sec:c7-endpoint}, we differentiate normalized
expectations under Gaussian field laws while keeping their profile
measures fixed. The parameter changes the covariance of the Gaussian
variables. In
\cref{sec:sc}, it changes the clock of the field path, and in
\cref{sec:c7,sec:c7-endpoint}, it changes the variance of a row of
couplings and the clock of one field coordinate. The densities that we
encounter are exponentials of
integrals of smooth functions of the field against the profile measure,
which may have atoms. We use the following version of the Gaussian heat
identity, for finitely many Gaussian variables and continuous field
paths. Its finite-dimensional input is the classical formula for the
derivative of a Gaussian expectation with respect to the covariance. A
general version, under hypotheses
different from ours, is proved in \cite[Theorem~1]{Voigtlaender2021Price}.
In the proof, we reduce to
finite Gaussian vectors by Bernstein approximation of the paths, while we
keep the profile measure, including its atoms, exact.

\paragraph{Setting.}
Let $I\subset\R$ be a nondegenerate compact interval with interior $I^\circ$, and let
$z\in I$ be the parameter. Under a reference probability measure, with
expectation $\E_{\rm ref}$, we consider the following random
variables.
\begin{enumerate}[label=(\roman*)]
\item \emph{Finite coordinates}
$G^{\rm fin}_z=(v_k(z)^{1/2}\xi_k)_{1\le k\le n_0}$, where
$\xi_1,\ldots,\xi_{n_0}$ are independent standard Gaussian variables and
$v_k\colon I\to[0,\infty)$.
\item \emph{Field paths} on a fixed finite rooted tree. Each edge $e$
carries the level interval $[0,1]$, after an affine change of variable,
and $d$ components. For each edge $e$ and component $j\le d$, let
$B_{e,j}$ be a standard Brownian motion, and let
$T_{e,j,z}\colon[0,1]\to[0,\infty)$ be a deterministic nondecreasing clock
such that $T_{e,j,z}(0)=0$. The Brownian motions for different edges
and components are mutually independent. The path on the edge $e$ is
\[
 G_z(e,j,v)=G_z(e',j,1)+B_{e,j}\bigl(T_{e,j,z}(v)\bigr),
 \qquad 0\le v\le1,
\]
where $e'$ is the parent edge of $e$, and the first term is $0$ if $e$
starts at the root.
\item A finite family $\omega$ of independent standard Gaussian variables that do not
depend on $z$ and are independent of the variables in (i) and (ii), and
a variable $\sigma$ that ranges over a nonempty finite set and is integrated
against counting measure.
\end{enumerate}
The variables $\xi_k$ and the Brownian motions $B_{e,j}$ are mutually
independent. Thus the finite coordinates and the field paths are
independent under the reference law. This setting includes all the
scalar and cavity applications below.
We write $\mathcal I$ for the index space formed by the finite
coordinates $1\le k\le n_0$ and the field points and components
$(e,j,v)$. It is a compact metric space, the disjoint union of a finite
set and finitely many copies of $[0,1]$. The Gaussian coordinates
$G_z=(G_z^{\rm fin},G_z^{\rm path})$ form a random element of the space
$\mathcal C$ of functions on $\mathcal I$ that are continuous on each
copy of $[0,1]$, with the norm $|G^{\rm fin}|+\|G^{\rm path}\|_\infty$.
We denote their covariance by $K_z(i,i')=\E_{\rm ref}[G_z(i)G_z(i')]$ for
$i,i'\in\mathcal I$. For example, for one edge and one component,
$K_z(v,v')=T_z(\min\{v,v'\})$.

Let $\ell_z(G)=\ell_z(G,\omega,\sigma)$ and
$\Phi_z(G)=\Phi_z(G,\omega,\sigma)$ be real functions of $z\in I$,
$G\in\mathcal C$, $\omega$, and $\sigma$. We write $\dot\ell_z$ and
$\dot\Phi_z$ for their derivatives in $z$, with $(G,\omega,\sigma)$
fixed, and $\dot K_z$ for the derivative of $K_z$ in $z$. For a
continuous kernel $A$ on $\mathcal I^2$ and a finite signed measure $B$
on $\mathcal I^2$, we write $\langle A,B\rangle=\iint A(i,i')B(\dd i,\dd i')$,
and $S\otimes P$ denotes the product of two signed measures on
$\mathcal I$.

\begin{lemma}[Normalized Gaussian path heat identity]
\label{lem:fnd-path-heat}
In the setting above, assume the following.
\begin{enumerate}[label=\textup{(A\arabic*)}]
\item \emph{Variances and clocks.} Each $v_k$ and each $T_{e,j,z}(v)$ is
continuous in $z\in I$, jointly with $v$ for the clocks, and
continuously differentiable in $z\in I^\circ$. For $z\in I^\circ$, each
nonconstant $v_k$ is positive and each clock is strictly increasing in
$v$. Coordinates with $v_k\equiv0$ are omitted.
\item \emph{Covariance.} The covariance $K_z(i,i')$ is differentiable in
$z\in I^\circ$, and the maps $(z,i,i')\mapsto K_z(i,i')$ and
$(z,i,i')\mapsto\dot K_z(i,i')$ are jointly continuous on
$I\times\mathcal I^2$, where $\dot K_z$ at the endpoints of $I$ denotes
the continuous extension.
\item \emph{Functionals.} For each $\sigma$, the functions
$\ell_z(G,\omega,\sigma)$, $\Phi_z(G,\omega,\sigma)$,
$\dot\ell_z(G,\omega,\sigma)$, and $\dot\Phi_z(G,\omega,\sigma)$ are jointly
continuous in $(z,G,\omega)\in I\times\mathcal C\times\R^{n_1}$,
where $n_1$ is the number of auxiliary Gaussian coordinates, with
$\dot\ell_z$ and $\dot\Phi_z$ extended continuously to the endpoints of $I$.
The functions $\Phi_z$, $\dot\Phi_z$, and
$\dot\ell_z$ are bounded uniformly in all variables, and there exists
$C<\infty$ such that
\[
 |\ell_z(G)|\le C\bigl(1+|\omega|+|G^{\rm fin}|+\|G^{\rm path}\|_\infty\bigr).
\]
\item \emph{Derivative measures.} For fixed $(z,\omega,\sigma)$, the maps
$G\mapsto\ell_z(G)$ and $G\mapsto\Phi_z(G)$ are twice continuously
differentiable on $\mathcal C$, and their derivatives are represented by
finite signed measures
\[
 S=D\ell_z,\qquad M=D^2\ell_z,\qquad
 P=D\Phi_z,\qquad Q=D^2\Phi_z
\]
on $\mathcal I$ and $\mathcal I^2$, in the sense that, for example,
\[
 D\ell_z(G)[h]=\int h\dd S,\qquad
 D^2\ell_z(G)[h,h']=\iint h(i)h'(i')\,M(\dd i,\dd i')
 \qquad(h,h'\in\mathcal C).
\]
For fixed $(\omega,\sigma)$, the four measures depend continuously on
$(z,G)$ in the total-variation norm, and their total variations are
bounded uniformly in all variables.
\end{enumerate}
Let $Z_z=\E_{\rm ref}e^{\ell_z(G_z)}$, where $\E_{\rm ref}$ includes the
sum over $\sigma$, and let
$\E_z\Phi=Z_z^{-1}\E_{\rm ref}\bigl[e^{\ell_z(G_z)}\Phi\bigr]$. Then, for
every $z\in I$,
\begin{align}
 \frac\dd{\dd z}\E_z\Phi_z={}&
 \E_z\dot\Phi_z+\Cov_z(\Phi_z,\dot\ell_z)\nonumber\\
 &+\frac12\E_z\Bigl\langle\dot K_z,
 Q+S\otimes P+P\otimes S
 +(\Phi_z-\E_z\Phi_z)(M+S\otimes S)\Bigr\rangle,
 \label{eq:fnd-path-heat}
\end{align}
where the functionals and measures are evaluated at $G_z$. The
derivative is one-sided at the endpoints of $I$, including endpoints at
which a finite-coordinate variance vanishes. The same differentiation
gives
\begin{equation}\label{eq:fnd-path-normalizer}
 \frac\dd{\dd z}\log Z_z
 =\E_z\dot\ell_z+\frac12\E_z\langle\dot K_z,M+S\otimes S\rangle.
\end{equation}
\end{lemma}

\begin{remark}\label{rem:gauss-local}
We show in the proof that the right sides of \eqref{eq:fnd-path-heat} and
\eqref{eq:fnd-path-normalizer} are continuous in $z\in I$, so that
$z\mapsto\E_z\Phi_z$ and $z\mapsto\log Z_z$ are continuously
differentiable on $I$. The assumptions hold, in particular, for the
smooth local functionals that we use in
\cref{sec:sc,sec:c7,sec:c7-endpoint},
integrated against a fixed finite measure. For instance, let $\varpi$ be
a finite measure on the path part of $\mathcal I$, possibly with atoms,
and let $\ell_z(G)=-\int V_z(i,G(i))\,\varpi(\dd i)$, where $V_z(i,x)$
is twice continuously differentiable in $x$, and $\partial_xV_z$ and
$\partial_x^2V_z$ are bounded and jointly continuous in $(z,i,x)$. Then
\[
 S=-\partial_xV_z(\cdot,G(\cdot))\,\varpi,\qquad
 M=-\operatorname{diag}_{\#}\bigl(\partial_x^2V_z(\cdot,G(\cdot))\,\varpi\bigr),
\]
where $\operatorname{diag}_{\#}$ is the pushforward under
$i\mapsto(i,i)$. The total variation of the difference of the scores at
$G$ and $G'$ is at most
$\varpi(\mathcal I)\sup_i|\partial_xV_z(i,G(i))-\partial_xV_z(i,G'(i))|$,
and the same holds for the Hessians with $\partial_x^2V_z$. We note that
these functions are uniformly continuous on compact sets
and that $G$ is bounded, which gives the continuity in (A4) for this
term. An atom of $\varpi$ produces an atom of $M$ on
the diagonal, and its contribution to $\langle\dot K_z,M\rangle$
involves $\dot K_z$ on the diagonal.
\end{remark}

\begin{proof}[Proof of \cref{lem:fnd-path-heat}]
The proof has seven steps. We first prove the identities for finite
Gaussian vectors. We then approximate the paths by their Bernstein
polynomials and couple the approximations on one probability space.
Next, we show that the derivative measures and the covariance kernels
converge, and we establish a common integrable envelope and a lower
bound for the normalizers. Finally, we pass to the limit in the interior
of $I$ and treat the endpoints.

\paragraph{Step 1: the finite-dimensional identity.}
We begin with a finite Gaussian vector. Let $J$ be an open interval,
and let $Y_z$, $z\in J$, be a centered Gaussian vector in $\R^p$ whose
covariance $\Sigma_z$ is continuously differentiable in $z$. We allow the
matrix $\Sigma_z$ to be
singular. Let $F_z\colon\R^p\to\R$ be continuously
differentiable in $z$ and twice continuously differentiable in
$y\in\R^p$, and assume that $F_z$, $\partial_zF_z$, $\nabla F_z$, and
$\nabla^2F_z$ are bounded by $Ce^{C|y|}$, locally uniformly in $z$. We
claim that
\begin{equation}\label{eq:gauss-finite-heat}
 \frac\dd{\dd z}\E F_z(Y_z)=\E\,\partial_zF_z(Y_z)
 +\frac12\E\bigl\langle\dot\Sigma_z,\nabla^2F_z(Y_z)\bigr\rangle .
\end{equation}
Let $\varepsilon>0$, let $\zeta$ be a standard Gaussian vector in
$\R^p$, and let $A_z=(\Sigma_z+\varepsilon I)^{1/2}$. The matrix square
root is smooth on positive definite matrices, so $A_z$ is continuously
differentiable in $z$. By dominated convergence, we may differentiate
$\E F_z(A_z\zeta)$ under the expectation. The chain rule gives, besides
$\E\,\partial_zF_z(A_z\zeta)$, the term
$\E[\nabla F_z(A_z\zeta)\cdot\dot A_z\zeta]$. Using Gaussian integration
by parts, $\E[\zeta_lg(\zeta)]=\E[\partial_lg(\zeta)]$, we find that
this term equals $\E\Tr\bigl(\dot A_z\nabla^2F_z(A_z\zeta)A_z\bigr)$.
Since $\nabla^2F_z$ is symmetric and
$\dot A_zA_z+A_z\dot A_z=\dot\Sigma_z$, this is
$\E\langle\dot\Sigma_z,\nabla^2F_z(A_z\zeta)\rangle/2$. We integrate
this identity over $[z_1,z_2]\subset J$ and let $\varepsilon\downarrow0$.
The vector $A_z\zeta$ has the law of $Y_z+\sqrt\varepsilon\zeta'$, where
$\zeta'$ is a standard Gaussian vector independent of $Y_z$. Since the
growth bounds give domination, we obtain
the integrated form of
\eqref{eq:gauss-finite-heat}. Realizing $Y_z=\Sigma_z^{1/2}\zeta$, which
is continuous in $z$, we see by the same domination that the right side
of \eqref{eq:gauss-finite-heat} is continuous in $z$. The fundamental
theorem of calculus then gives
\eqref{eq:gauss-finite-heat}. This is the classical formula. Note that we
never invert $\Sigma_z$.

We will apply \eqref{eq:gauss-finite-heat} to $e^{\ell_z}$ and to
$e^{\ell_z}\Phi_z$, composed with a linear map from $\R^p$ to
$\mathcal C$, as in Step~2. By the product rule, we have
$D e^{\ell}=e^{\ell}S$, $D^2e^{\ell}=e^{\ell}(M+S\otimes S)$, and
\[
 D^2(e^{\ell}\Phi)=e^{\ell}\bigl[Q+S\otimes P+P\otimes S+\Phi(M+S\otimes S)\bigr].
\]
Inserting these derivatives into \eqref{eq:gauss-finite-heat} for the
numerator $\E_{\rm ref}[e^{\ell_z}\Phi_z]$ and the denominator
$\E_{\rm ref}e^{\ell_z}$, and using the quotient rule, we obtain
\eqref{eq:fnd-path-heat} for a finite Gaussian vector. Observe that the
centering $\Phi_z-\E_z\Phi_z$ comes from the
derivative of the denominator.
Dividing the derivative of the denominator by the denominator gives
\eqref{eq:fnd-path-normalizer}.

\paragraph{Step 2: Bernstein approximation and kernel convergence.}
We use the Bernstein approximation introduced in
\cite{Bernstein1912}. Let $m\ge1$ and $v_l=l/m$ for $0\le l\le m$. For a
continuous function $f$ on $[0,1]$, the Bernstein polynomial is
\[
 \Pi_mf(v)=\sum_{l=0}^mb_{m,l}(v)f(v_l),\qquad
 b_{m,l}(v)=\binom ml v^l(1-v)^{m-l}.
\]
The operator $\Pi_m$ is linear, and $\Pi_mf(0)=f(0)$ and
$\Pi_mf(1)=f(1)$. The weights are nonnegative and sum to one, which
gives $\|\Pi_mf\|_\infty\le\|f\|_\infty$. Further, for every $\delta>0$, we have
\[
 \|\Pi_mf-f\|_\infty\le\omega_f(\delta)+\frac{\|f\|_\infty}{2m\delta^2},
\]
where $\omega_f$ is the modulus of continuity of $f$. This follows from
Chebyshev's inequality for a binomial variable with $m$ trials, whose
variance is at most $m/4$. In particular, $\Pi_mf\to f$ uniformly
over any bounded equicontinuous family of functions $f$. The same holds for the
tensor product $\Pi_m\otimes\Pi_m$ acting on continuous functions on
$[0,1]^2$, with the bound $\omega_f(\delta)+\|f\|_\infty/(m\delta^2)$.

On each edge and component, we let $\Pi_m$ act on the path, and we
leave the finite coordinates unchanged. We denote the resulting operator
on $\mathcal C$ also by $\Pi_m$. Then $\Pi_mG_z$ is a linear function of
the finite Gaussian vector $Y_z$ formed by the finite coordinates and
the sampled values $G_z(e,j,v_l)$. A value at a vertex of the tree may
appear in several coordinates of $Y_z$. Step~1 allows such repeated
coordinates, since it does not require an invertible covariance. The
covariance of $Y_z$ consists of the values of $K_z$ at the sampled
points. By (A2), it is continuously differentiable on $I^\circ$, with
derivative given by the values of $\dot K_z$. The growth conditions of
Step~1 hold by (A3) and (A4), since $\|\Pi_my\|_{\mathcal C}\le2|y|$.

We set
\[
 \ell^{(m)}_z=\ell_z(\Pi_mG_z),\qquad
 \Phi^{(m)}_z=\Phi_z(\Pi_mG_z),\qquad
 Z^{(m)}_z=\E_{\rm ref}e^{\ell^{(m)}_z},
\]
and we normalize the approximate moment
$\E^{(m)}_z\Phi^{(m)}_z=(Z^{(m)}_z)^{-1}\E_{\rm ref}[e^{\ell^{(m)}_z}\Phi^{(m)}_z]$
by its own normalizer $Z^{(m)}_z$. By the chain rule, the derivative of
$\ell_z(\Pi_my)$ in the sampled coordinate $(e,j,v_l)$ is
$\int b_{m,l}\dd S$, where $b_{m,l}$ is placed on the edge $e$ and the
component $j$. We obtain the second derivatives in the same way from
$M$. We use Step~1 for fixed $(\omega,\sigma)$, and we then integrate
over $\omega$ and sum over $\sigma$. The domination of Step~5 justifies
this. We conclude that \eqref{eq:fnd-path-heat} and
\eqref{eq:fnd-path-normalizer} hold for the approximations at every
$z\in I^\circ$, with $\E_z$ replaced by $\E^{(m)}_z$ and with the
functionals and
measures evaluated at $\Pi_mG_z$. In these identities, $\dot K_z$ is
replaced by the kernel
\[
 \dot K_z^{(m)}(v,w)
  =\sum_{l,l'}b_{m,l}(v)b_{m,l'}(w)\dot K_z(v_l,v_{l'})
\]
on each pair of edges and components. Blocks between finite coordinates
are unchanged, and blocks between a finite coordinate $k$ and a path use
the same sampling in their path variable:
\[
 \dot K_z^{(m)}(k,(e,j,v))=\sum_lb_{m,l}(v)\dot K_z(k,(e,j,v_l)).
\] By (A2),
the kernels $K_z$ and $\dot K_z$ are uniformly continuous on the
compact set $I\times\mathcal I^2$. Hence the families
$\{K_z\}_{z\in I}$ and $\{\dot K_z\}_{z\in I}$ are bounded and
equicontinuous. Using the Bernstein bound, we obtain
\begin{equation}\label{eq:fnd-path-kernel-limit}
 \sup_{z\in I}\|\dot K_z^{(m)}-\dot K_z\|_\infty\longrightarrow0,
 \qquad
 \sup_{z\in I}\|K_z^{(m)}-K_z\|_\infty\longrightarrow0,
\end{equation}
where $K_z^{(m)}$, the covariance of $\Pi_mG_z$, is defined in the same
way.

\paragraph{Step 3: coupling.}
In Step~1, we used only the law of $Y_z$ at each fixed $z$. To pass to
the limit, we use the realization of the setting, with one Brownian motion
$B_{e,j}$ for each edge and component, evaluated at the varying clocks.
The finite coordinates are $v_k(z)^{1/2}\xi_k$ for fixed
$\xi_k$, independent of these Brownian motions as in the setting.
By (A1), the map $(z,v)\mapsto G_z(e,j,v)$ is almost surely
continuous on the compact set $I\times[0,1]$. It is therefore uniformly
continuous, and the family $\{G_z(e,j,\cdot)\}_{z\in I}$ is
equicontinuous. By Step~2, we have
$\sup_{z\in I}\|\Pi_mG_z-G_z\|_{\mathcal C}\to0$ almost surely. Let
$\bar T$ be the maximum of the clocks over $I\times[0,1]$, with
$\bar T=0$ if there are no field paths. Let
$\bar v=\max(\{0\}\cup\{v_k(z):1\le k\le n_0,\ z\in I\})$, so that
$\bar v=0$ if there are no finite coordinates. Since the value of a path
at a point of the edge $e$ is a sum of
increments of the Brownian motions on $e$ and its ancestor edges, we find
that
\[
 \sup_{z\in I}\|\Pi_mG_z\|_{\mathcal C}\le\sup_{z\in I}\|G_z\|_{\mathcal C}
 \le\bar v^{1/2}|\xi|+\sum_{e,j}\sup_{0\le\tau\le\bar T}|B_{e,j}(\tau)| .
\]

\paragraph{Step 4: convergence of the derivative measures.}
By (A4) and Step~3, we see that the measures $S$, $M$, $P$, and $Q$
evaluated at $\Pi_mG_z$ converge in total variation to the same measures evaluated
at $G_z$, almost surely, for every $z\in I$. Combining this with
\eqref{eq:fnd-path-kernel-limit} and the uniform bounds on the total
variations, we find that every contraction in the identity of Step~2
converges. For example, writing $M^{(m)}$ for $M$ evaluated at $\Pi_mG_z$,
we have
\[
 \bigl|\langle\dot K_z^{(m)},M^{(m)}\rangle-\langle\dot K_z,M\rangle\bigr|
 \le\|\dot K_z^{(m)}-\dot K_z\|_\infty\|M^{(m)}\|_{\rm TV}
 +\|\dot K_z\|_\infty\|M^{(m)}-M\|_{\rm TV},
\]
and for product measures we use
$\|S'\otimes P'-S\otimes P\|_{\rm TV}\le\|S'-S\|_{\rm TV}\|P'\|_{\rm TV}
+\|S\|_{\rm TV}\|P'-P\|_{\rm TV}$. This includes the diagonal
contributions of atoms, as in
\cref{rem:gauss-local}, since only the field values are approximated
and $\dot K^{(m)}_z(v,v)\to\dot K_z(v,v)$ uniformly.

\paragraph{Step 5: domination and the normalizers.}
By (A3) and Step~3, we have $e^{\ell^{(m)}_z}\le\mathcal W$ for all $m$
and $z$, where
\[
 \mathcal W=\exp\Bigl\{C\Bigl(1+|\omega|+\bar v^{1/2}|\xi|
 +\sum_{e,j}\sup_{0\le\tau\le\bar T}|B_{e,j}(\tau)|\Bigr)\Bigr\}
\]
and $C$ is the constant of (A3). All differentiated terms in the
identity of Step~2 are products of a weight with quantities bounded by
(A3), (A4), and \eqref{eq:fnd-path-kernel-limit}. Hence a constant
multiple of $\mathcal W$ is a common envelope. It is integrable to every fixed
positive power, because the suprema of Brownian motion over a bounded
interval have Gaussian tails, and the sum over $\sigma$ is finite. We now
show that the approximate normalizers have a common positive
lower bound. We retain one value $\sigma_0$ of the counting variable.
Using Jensen's inequality and the bound on $|\ell_z|$ in (A3), we obtain
\[
 Z^{(m)}_z\ge\E_{\rm ref}e^{\ell_z(\Pi_mG_z,\omega,\sigma_0)}
 \ge\exp\bigl(\E_{\rm ref}\,\ell_z(\Pi_mG_z,\omega,\sigma_0)\bigr)\ge e^{-C'},
\]
where the expectations are over the Gaussian variables only, and $C'$
does not depend on $m$ or $z$.

\paragraph{Step 6: passage to the limit in the interior.}
Let $\mathcal E_m(z)=\E^{(m)}_z\Phi^{(m)}_z$, and let $\Psi_m(z)$ be the
right side of the identity of Step~2. Let $[z_1,z_2]\subset I^\circ$.
By Step~2 and the fundamental theorem of calculus, we have
\[
 \mathcal E_m(z_2)-\mathcal E_m(z_1)=\int_{z_1}^{z_2}\Psi_m(z)\dd z .
\]
By Steps~3--5 and dominated convergence, we see that
$\mathcal E_m(z)\to\E_z\Phi_z$
and $\Psi_m(z)\to\Psi(z)$ for all $z\in I$, where
$\Psi(z)$ is the right side of \eqref{eq:fnd-path-heat}. Here we used the
lower bound on the normalizers and the continuity of $\ell_z$, $\Phi_z$,
$\dot\ell_z$, and $\dot\Phi_z$ in $G$. The
functions $\Psi_m$ are bounded uniformly in $m$ and $z$. Using dominated
convergence in $z$, we obtain
\[
 \E_{z_2}\Phi_{z_2}-\E_{z_1}\Phi_{z_1}=\int_{z_1}^{z_2}\Psi(z)\dd z .
\]
The same argument with $z_k\to z$ in place of $m\to\infty$ shows that
$\Psi$ and $z\mapsto\E_z\Phi_z$ are continuous on $I$. Indeed, by the
coupling, $G_{z_k}\to G_z$ uniformly, the measures converge in total
variation by (A4), $\dot K_{z_k}\to\dot K_z$ uniformly by (A2), and the
envelope of Step~5 dominates. By the fundamental theorem of calculus, we
conclude that
\eqref{eq:fnd-path-heat} holds on $I^\circ$. The same argument with
$\Phi\equiv1$, applied to $\log Z^{(m)}_z$ instead of a ratio, gives
\eqref{eq:fnd-path-normalizer} on $I^\circ$.

\paragraph{Step 7: endpoints.}
Let $z_-$ be an endpoint of $I$, say the left one. The integrated
identity of Step~6 holds for $z_-<z_1<z_2$ in $I^\circ$. As
$z_1\downarrow z_-$, the left side converges because $z\mapsto\E_z\Phi_z$
is continuous on $I$, and the integral converges because $\Psi$ is
continuous on $I$. This also holds if a finite variance vanishes at $z_-$.
In that case,
the corresponding coordinate converges to zero, while the derivative
measures and $\dot K_z$ have continuous limits, and we never invert the
covariance matrix. We obtain
$\E_{z_2}\Phi_{z_2}-\E_{z_-}\Phi_{z_-}=\int_{z_-}^{z_2}\Psi(z)\dd z$,
and the continuity of $\Psi$ gives the one-sided derivative at
$z_-$. We treat the right endpoint and \eqref{eq:fnd-path-normalizer} in the
same way.
\end{proof}


\part{Overlap estimates at a minimizer}\label{part:overlap}
\chapter{Self-overlap moments on quantile intervals}\label{sec:wi}

In this section, we bound two moments of the self-overlap at a
minimizing profile. The first is the integrated variance
$V=\int\Var(S_q)\,\nu(\dd q)$, one of the three terms of the defect in
\eqref{eq:in-envelope}. The second is the integrated cubic moment
$H_3=\int\E|S_q-q|^3\,\nu(\dd q)$, which appears in the cavity argument
of \cref{sec:c7}. We control both by comparing $S_q$ with averages of
$L_p$ over windows of quantile levels on either side of $q$; the
variances of these averages are bounded by \eqref{eq:in-Lvar}.

The section contains three results. In \cref{lem:wi-two}, we prove two
comparisons: conditionally on $\calG_q$, the self-overlap $S_q$ is at
most twice an average of $L_p$ over later levels $p\ge q$, and it equals
twice an average of $L_p$ over earlier levels $p\le q$ plus a
nonnegative term. \Cref{prop:wi-vfs} bounds the integral of
$\Var(S_{Q(v)})$ over a quantile interval in terms of its length, a
lower bound for the cumulative mass on it, and the diameter of its image
under $Q$. \Cref{cor:wi-cub} shows that $H_3$ is at most an arbitrarily
small multiple $zV$ of $V$ plus an error of order
$(Ms)^{-1}\log(2M)\log(2/z)$.

We use the martingale and contact identities \eqref{eq:in-martingale}
and \eqref{eq:in-contact} of \cref{prop:fnd-contact}, the integrated
variance bound \eqref{eq:in-Lvar} of \cref{prop:fnd-variance}, and the
quantile function \eqref{eq:fnd-quantile}. We use \cref{prop:wi-vfs} in
the proof of \cref{prop:sm-small}, which bounds the defect for small
$\rho$, and \cref{cor:wi-cub} in the proof of \cref{cor:c7-defect},
where the cavity inequality is closed. \Cref{lem:wi-two} is used only in
this section.

We refine the quantile-window argument of \cite[Lemma~3.3]{AL}, which
estimates integrals of $\Var(L_{Q(v)})$ over quantile intervals.
By \cref{lem:fnd-decomp}, we have $\Var(L_q)=j_q+\Var(S_q)/4$, where
$j_q$ is the variance of the conditionally centered part $N^{-1}(\sigma-m_q)\cdot m_q$ of $L_q$. An
estimate for $\Var(L_q)$ must therefore also control this centered part,
and in \cite[Lemma~3.3]{AL} it produces additional residual terms. Here
we compare conditional expectations given $\calG_q$, in which the
centered part vanishes. We estimate $j_q$ separately in \cref{sec:th}.

\section{Quantile windows and kernels}\label{sec:wi-setup}

\paragraph{Standing assumptions.}
Throughout this section, we fix $N\ge1$, $0<\rho\le1$, and $0<s\le1/2$,
so that $t=1-\rho\in[0,1)$ and $M=Nc>0$. We let $\mu=s\delta_0+\nu$ be a
deterministic admissible profile, and expectations refer to its joint
law from \cref{subsec:fnd-law}. Each statement says whether $\mu$ must
minimize $\mathcal F_{t,s}$; a statement about a minimizer holds for
every minimizer, with the same constants. All constants in this section
are absolute: they do not depend on $N$, $\beta$, $\rho$, $s$, or the
profile.

Recall that $\sigma\in\{-1,1\}^N$ and $m_q=\E[\sigma\mid\calG_q]$. Then
$|m_q|^2\le N$ and $|\sigma\cdot m_q|\le N$, and it follows that
\[
 0\le S_q\le1,\qquad |L_q|\le\frac32\qquad(0\le q\le1).
\]
The observables $S_q$ and $L_q$ are jointly measurable in the level and
the outcome, because the process $q\mapsto m_q$ is continuous. They are
also bounded, and Fubini's theorem applies to their integrals against deterministic
functions in $L^1(\nu)$.

\paragraph{The variance bound.}
The estimates of this section rest on the first bound in
\eqref{eq:in-Lvar}, and its factor $a(q)^{-1}$ is the reason for
working with quantiles. In the proof of \cref{prop:fnd-variance}, this
factor comes from two facts. The first is the bound
$N^{-2}\int k(q)^2\,\E\Tr(H_qC_q)\,\nu(\dd q)$ for the variance, which
holds for every admissible profile \cite[Proposition~2.2, (2.15)]{AL}
and rests on the Brascamp--Lieb inequality for the posterior law of the
Gaussian coordinates. The second is the inequality
$\Tr(H_qC_q)\le a(q)^{-1}\Tr H_q^2$ at contact levels, which follows
from $H_q-a(q)C_q\succeq0$, together with $\E\Tr H_q^2\le N/c$.

\paragraph{Quantiles.}
Recall that the quantile function $Q$ of \eqref{eq:fnd-quantile} pushes
Lebesgue measure on $(0,1-s)$ forward to $\nu$. We write
$F(p)=\nu([0,p])$ and $F_-(p)=\nu([0,p))$, so that $a(p)=s+F(p)$. The
infimum in \eqref{eq:fnd-quantile} is attained because $F$ is
right-continuous, that is, $F(Q(v))\ge v$; in particular, $a(Q(v))\ge s+v$. We
also use
\begin{equation}\label{eq:wi-supp}
 a(q)\ge s\quad(0\le q\le1),\qquad
 Q(v)\in\supp\nu\quad(0<v<1-s).
\end{equation}
The first bound holds because $\mu\ge s\delta_0$. For the second, let
$q_*=Q(v)$, so that $\nu([0,q_*])\ge v>0$. If $q_*=0$, then $0$ is an
atom of $\nu$ and lies in $\supp\nu$. If $q_*>0$ and the interval
$(q_*-r,q_*+r)$ were $\nu$-null for some $0<r<q_*$, then
$\nu([0,q_*-r/2])=\nu([0,q_*])\ge v$, which contradicts the definition
of $Q(v)$ as an infimum.

If $\mu$ is a minimizer, then every point of $\supp\nu$ is a contact
level. By \eqref{eq:wi-supp} and \eqref{eq:in-contact}, we then have
$\E S_{Q(v)}=Q(v)$ for every $v\in(0,1-s)$. In that case, we define
\[
 \xi_v=S_{Q(v)}-Q(v)\qquad(0<v<1-s).
\]
Note that $|\xi_v|\le1$, $\E\xi_v=0$, and $\E\xi_v^2=\Var(S_{Q(v)})$. Using
that $Q$ pushes Lebesgue measure forward to $\nu$, we can write
\[
 V=\int_0^{1-s}\E\xi_v^2\dd v,\qquad
 H_3=\int_0^{1-s}\E|\xi_v|^3\dd v .
\]
The integrands are Borel functions of $v$, by the joint measurability
noted above, Fubini's theorem, and the monotonicity of $Q$.

\paragraph{Windows.}
Let $I\subset(0,1-s)$ be an interval of length $|I|>0$. The image of
Lebesgue measure on $I$ under $Q$ is dominated by the image of Lebesgue
measure on $(0,1-s)$, which is $\nu$. As in \cite[Section~3]{AL}, it
has a density $f_I$ with respect to $\nu$, with
\[
 0\le f_I\le1,\qquad \int f_I\dd\nu=|I| .
\]
The density is explicit at an atom. If $p$ is an atom of $\nu$, with
mass $r_p=\nu(\{p\})$, then $\{v:Q(v)=p\}$ is an interval with
endpoints $F_-(p)$ and $F(p)$, and
\[
 f_I(p)=\frac{|I\cap(F_-(p),F(p))|}{r_p}.
\]
In particular, a window may contain only part of the quantile interval
of an atom; the atom then receives the corresponding fraction of its
mass. We define the normalized density $\kappa_I=f_I/|I|$. It satisfies
\[
 0\le\kappa_I\le|I|^{-1},\qquad \int\kappa_I\dd\nu=1,\qquad
 \int\varphi\,\kappa_I\dd\nu=\frac1{|I|}\int_I\varphi(Q(v))\dd v
\]
for every bounded Borel function $\varphi$. We say that $I$ has
\emph{threshold} $u>0$ if $a(Q(v))\ge u$ for almost every $v\in I$. Every
such interval has threshold $s$, by \eqref{eq:wi-supp}.

\paragraph{Kernels.}
More generally, we call a deterministic Borel function $\kappa\ge0$ on
$[0,1]$ with $\int\kappa\dd\nu=1$ and $\int\kappa^2\dd\nu<\infty$ a
\emph{kernel}. For a kernel $\kappa$, we define its mean level and the
average of $L_p-p/2$, which is centered when $\mu$ is a minimizer,
\[
 \bar p_\kappa=\int p\,\kappa(p)\,\nu(\dd p),\qquad
 \Lambda_\kappa=\int\kappa(p)\Bigl(L_p-\frac p2\Bigr)\nu(\dd p).
\]
We write $\bar p_I$ and $\Lambda_I$ when $\kappa=\kappa_I$. Then
\[
 \bar p_I=\frac1{|I|}\int_IQ(v)\dd v,\qquad
 \Lambda_I=\frac1{|I|}\int_I\Bigl(L_{Q(v)}-\frac{Q(v)}2\Bigr)\dd v .
\]

\paragraph{Fixed kernels.}
Each assertion of \cref{lem:wi-two} below concerns one level $q$ and one
kernel $\kappa$, both deterministic. In \cref{prop:wi-vfs,cor:wi-cub},
the level and the kernels depend on a quantile parameter $v$. We then
apply \cref{lem:wi-two} separately for each value of $v$, and we
integrate over $v$ only the resulting deterministic expectations. In
this way, we never need versions of the conditional expectations that
are chosen simultaneously for all $v$.

\section{Two-sided conditional comparison}\label{sec:wi-comparison}

For a later level $p\ge q$, we show by the conditional Jensen inequality
that $S_q$ is at most $2\E[L_p\mid\calG_q]$. For an earlier level $p\le q$,
the conditional expectation of $L_p$ given $\calG_q$ is $S_q/2$ minus the
nonnegative quantity $|m_q-m_p|^2/(2N)$. We average these relations over
$p$ against a kernel and obtain two one-sided comparisons between
$S_q-q$ and the conditional expectation of $\Lambda_\kappa$ given
$\calG_q$. At a minimizer, $\Lambda_\kappa$ is centered, and its
variance is controlled by \eqref{eq:in-Lvar}.

\begin{lemma}[Two-sided conditional comparison]\label{lem:wi-two}
Let $\mu$ be an admissible profile, let $q\in[0,1]$, and let $\kappa$ be
a kernel.
\begin{enumerate}[label=\textup{(\alph*)},series=witwo]
\item If $\kappa=0$ $\nu$-almost everywhere on $[0,q)$, then almost
surely
\[
 S_q-q\le(\bar p_\kappa-q)+2\E[\Lambda_\kappa\mid\calG_q].
\]
\item If $\kappa=0$ $\nu$-almost everywhere on $(q,1]$, then almost
surely
\[
 \begin{aligned}
 S_q-q&=(\bar p_\kappa-q)+2\E[\Lambda_\kappa\mid\calG_q]
 +\frac1N\int\kappa(p)|m_q-m_p|^2\,\nu(\dd p)\\
 &\ge(\bar p_\kappa-q)+2\E[\Lambda_\kappa\mid\calG_q].
 \end{aligned}
\]
\end{enumerate}
If $\mu$ minimizes $\mathcal F_{t,s}$, then in addition the following
hold.
\begin{enumerate}[resume*=witwo]
\item We have $\E\Lambda_\kappa=0$ and
\[
 \E\bigl[(\E[\Lambda_\kappa\mid\calG_q])^2\bigr]
 \le\Var(\Lambda_\kappa)
 \le\frac1M\int\frac{\kappa(p)^2}{a(p)}\,\nu(\dd p).
\]
\item Let $\vartheta>0$. If $\kappa$ satisfies the hypothesis of
\textup{(a)} and $\bar p_\kappa-q\le\vartheta/2$, then
\[
 \E\bigl[(S_q-q)^2\ind_{\{S_q-q>\vartheta\}}\bigr]
 \le16\Var(\Lambda_\kappa).
\]
If $\kappa$ satisfies the hypothesis of \textup{(b)} and
$q-\bar p_\kappa\le\vartheta/2$, then
\[
 \E\bigl[(S_q-q)^2\ind_{\{S_q-q<-\vartheta\}}\bigr]
 \le16\Var(\Lambda_\kappa).
\]
\item If $I\subset(0,1-s)$ is an interval of length $x>0$ with
threshold $u$, then $\Var(\Lambda_I)\le1/(Mux)$.
\end{enumerate}
For every admissible profile and every interval $I\subset(0,1-s)$ of
positive length, $\kappa_I$ satisfies the hypothesis of \textup{(a)} if
$Q\ge q$ on $I$, and that of \textup{(b)} if $Q\le q$ on $I$.
\end{lemma}

\begin{proof}
Parts (a) and (b) are statements about conditional expectations. We
prove each first for a single level $p$ and then integrate against
$\kappa(p)\,\nu(\dd p)$. Parts (c)--(e) add the contact identity and
\eqref{eq:in-Lvar}. Let
\[
 \Lambda'=\int\kappa(p)L_p\,\nu(\dd p)=\Lambda_\kappa+\frac{\bar p_\kappa}2 .
\]
Taking the conditional expectation given $\calG_p$ in the definition of
$L_p$, and recalling that $m_p=\E[\sigma\mid\calG_p]$ is
$\calG_p$-measurable, we find that
\begin{equation}\label{eq:wi-LS}
 \E[L_p\mid\calG_p]=\tfrac12S_p\qquad(0\le p\le1).
\end{equation}

\paragraph{Part (a).}
Let $p\ge q$. The process $m$ is a martingale by
\eqref{eq:in-martingale}, and $m_q=\E[m_p\mid\calG_q]$. Using the conditional Jensen
inequality for $x\mapsto|x|^2$, the tower property, and \eqref{eq:wi-LS},
we have
\[
 S_q\le\E[S_p\mid\calG_q]=2\E[L_p\mid\calG_q].
\]
Let $\zeta\ge0$ be a bounded $\calG_q$-measurable random variable. Then
$\E[\zeta S_q]\le2\E[\zeta L_p]$ for every $p\ge q$. We integrate this inequality against the probability measure
$\kappa(p)\,\nu(\dd p)$, which does not charge
$[0,q)$. By Fubini's theorem,
which applies because $|\zeta L_p|$ is bounded, this yields
$\E[\zeta S_q]\le2\E[\zeta\Lambda']$. As $\zeta$ is arbitrary and $S_q$
is $\calG_q$-measurable, we conclude that
$S_q\le2\E[\Lambda'\mid\calG_q]$ almost surely. Subtracting $q$ from
both sides and recalling that $\Lambda'=\Lambda_\kappa+\bar p_\kappa/2$,
we obtain (a).

\paragraph{Part (b).}
Let $p\le q$. Both $m_p$ and $m_q$ are $\calG_q$-measurable, and
$\E[\sigma\mid\calG_q]=m_q$. Expanding $|m_q-m_p|^2$, we find
\[
 \E[L_p\mid\calG_q]=\frac{m_q\cdot m_p-|m_p|^2/2}N
 =\frac{S_q}2-\frac{|m_q-m_p|^2}{2N}.
\]
The right side is bounded and jointly measurable in $p\in[0,q]$ and the
outcome. Its integral against $\kappa(p)\,\nu(\dd p)$, which is
concentrated on $[0,q]$, is $\calG_q$-measurable, because $m$ is
continuous and adapted. Testing against bounded $\calG_q$-measurable
variables as in (a), we find that
\[
 \E[\Lambda'\mid\calG_q]
 =\frac{S_q}2-\frac1{2N}\int\kappa(p)|m_q-m_p|^2\,\nu(\dd p).
\]
Multiplying by two, subtracting $q$, and using
$\Lambda'=\Lambda_\kappa+\bar p_\kappa/2$, we obtain the identity in
(b). The inequality holds because the last term is nonnegative.

\paragraph{Part (c).}
Let $\mu$ be a minimizer. By \eqref{eq:wi-LS} and \eqref{eq:in-contact},
we have $\E L_p=\E S_p/2=p/2$ for every $p\in\supp\nu$, and in
particular for $\nu$-almost every $p$. Then
$\E\Lambda_\kappa=0$ by Fubini's theorem. By the conditional Jensen
inequality, we have
\[
 \E\bigl[(\E[\Lambda_\kappa\mid\calG_q])^2\bigr]
 \le\E\Lambda_\kappa^2=\Var(\Lambda_\kappa).
\]
As $\Lambda_\kappa$ and $\Lambda'$ differ by a constant, the last
bound in (c) is the first bound in \eqref{eq:in-Lvar} with $k=\kappa$.

\paragraph{Part (d).}
Let $\xi=S_q-q$ and $Y=\E[\Lambda_\kappa\mid\calG_q]$. Suppose first
that $\kappa$ satisfies the hypothesis of (a) and
$\bar p_\kappa-q\le\vartheta/2$. On the event $\{\xi>\vartheta\}$, part
(a) implies that
\[
 2Y\ge\xi-(\bar p_\kappa-q)\ge\xi-\frac\vartheta2>\frac\xi2>0,
\]
so that $\xi^2<16Y^2$. Next, suppose that $\kappa$ satisfies the
hypothesis of (b) and $q-\bar p_\kappa\le\vartheta/2$. On the event
$\{\xi<-\vartheta\}$, part (b) implies that
\[
 2Y\le\xi+(q-\bar p_\kappa)\le\xi+\frac\vartheta2<\frac\xi2<0,
\]
and again $\xi^2<16Y^2$. In both cases, the expectation in (d) is at
most $16\E Y^2$, which is at most $16\Var(\Lambda_\kappa)$ by (c).

\paragraph{Part (e).}
The bound $\kappa_I\le1/x$ implies that $\kappa_I^2\le\kappa_I/x$. The function
$1/a$ is Borel and bounded by $1/s$, by \eqref{eq:wi-supp}. Using the
identity for $\int\varphi\,\kappa_I\dd\nu$ with $\varphi=1/a$, we have
\[
 \int\frac{\kappa_I(p)^2}{a(p)}\,\nu(\dd p)
 \le\frac1x\int\frac{\kappa_I(p)}{a(p)}\,\nu(\dd p)
 =\frac1{x^2}\int_I\frac{\dd v}{a(Q(v))}
 \le\frac1{ux},
\]
where we used the threshold in the last step. Inserting this bound into (c),
we find that $\Var(\Lambda_I)\le1/(Mux)$. This is the bound
\cite[Section~3, (3.6)]{AL}, because $\Lambda_I$ and
$x^{-1}\int_IL_{Q(v)}\dd v$ differ by the constant $\bar p_I/2$.

\paragraph{The hypotheses for $\kappa_I$.}
Suppose that $Q\ge q$ on $I$. Then
\[
 \int\ind_{[0,q)}\kappa_I\dd\nu=\frac1{|I|}\,|\{v\in I:Q(v)<q\}|=0,
\]
so $\kappa_I$ satisfies the hypothesis of (a). The case $Q\le q$ is the
same, with $(q,1]$ in place of $[0,q)$.
\end{proof}

\section{Variance on a quantile interval}\label{sec:wi-interval}

We now use the comparisons of \cref{lem:wi-two} to bound the
integral of $\Var(S_{Q(v)})$ over a quantile interval. For each variance
threshold, we compare every interior point $v$ of the interval with the
two windows of one common radius $r$ on either side of $v$. The random
part of each comparison is twice a conditional expectation whose second
moment is at most $(Mur)^{-1}$, which decreases as $r$ grows. The
deterministic part is at most the increment of $Q$ across the window
$(v-r,v+r)$. After we extend $Q$ by constants beyond the ends of the
interval, the integral of this increment over $\R$ equals $2r$ times the
diameter of the image of the interval, because $Q$ is nondecreasing; its
integral over the interior points is at most this quantity. This bounds the measure of the
set of points where the variance exceeds the threshold, and integrating
over the threshold gives the result. For an interval $I\subset(0,1-s)$,
we let
\[
 \ell(I)=\operatorname*{ess\,sup}_{v,w\in I}|Q(v)-Q(w)|
\]
denote the diameter of its image.

\begin{proposition}[Variance on a quantile interval]\label{prop:wi-vfs}
Let $\mu$ be a minimizer of $\mathcal F_{t,s}$, and let
$I\subset(0,1-s)$ be an interval of length $h>0$ with threshold $u>0$
and image diameter $\ell=\ell(I)$. Let $\epsilon=(Mu)^{-1}$. Then
\begin{equation}\label{eq:wi-vfs}
 \int_I\Var(S_{Q(v)})\dd v
 \le128\Bigl[\min\bigl\{h,\epsilon\bigl(1+\log_+(h/\epsilon)\bigr)\bigr\}
 +(\epsilon\ell)^{2/3}h^{1/3}\Bigr].
\end{equation}
The constant is absolute, and $I$ may contain part of the quantile
interval of an atom of $\nu$.
\end{proposition}

\begin{proof}
The proof has four steps. First, we record an identity for the
increments of $Q$. Second, we bound $\Var(S_{Q(v)})$ at an interior
point $v$ by comparing it with the two windows of radius $r$ next to
$v$. Third, for each threshold $\tau\in(0,1)$, we choose $r$ depending
on $\tau$ and bound the measure of the set where the variance exceeds
$\tau$. Fourth, we integrate over $\tau$.

\paragraph{Step 1: increments of the quantile function.}
Discarding endpoints, which changes no integral below, we write
$I=(\underline v,\overline v)$, with $h=\overline v-\underline v$. We set
\[
 \chi(v)=\E\xi_v^2=\Var(S_{Q(v)})\in[0,1]\qquad(v\in I).
\]
As $Q$ is nondecreasing, its essential infimum and supremum on $I$ are
$q_-=\lim_{v\downarrow\underline v}Q(v)$ and
$q_+=\lim_{v\uparrow\overline v}Q(v)$, and $q_+-q_-=\ell$. We extend
the restriction of $Q$ to $I$ to a nondecreasing function $\widetilde Q$
on $\R$ by setting $\widetilde Q=q_-$ on $(-\infty,\underline v]$ and
$\widetilde Q=q_+$ on $[\overline v,\infty)$. Let $\dd\widetilde Q$ be
the distributional derivative of $\widetilde Q$, a positive
Lebesgue--Stieltjes measure of total mass $q_+-q_-=\ell$. We claim that
\begin{equation}\label{eq:wi-variation}
 \int_{\R}\bigl(\widetilde Q(v+r)-\widetilde Q(v-r)\bigr)\dd v
 =2r\ell\qquad(r>0).
\end{equation}
Indeed, $\widetilde Q$ has at most countably many discontinuities. For
fixed $r>0$ and $v$ outside the two translates of this countable set by
$r$ and by $-r$, both $v-r$ and $v+r$ are continuity points, and
$\widetilde Q(v+r)-\widetilde Q(v-r)$ is the $\dd\widetilde Q$-measure of
$[v-r,v+r]$. By Tonelli's theorem,
\[
 \int_\R\dd\widetilde Q([v-r,v+r])\dd v
 =\int_\R\int_\R\ind_{\{|x-v|\le r\}}\dd v\,\dd\widetilde Q(x)
 =2r\ell .
\]
In words, each unit of mass of $\dd\widetilde Q$ is counted for the
centers $v$ within distance $r$ of it, and these form a set of length
$2r$. The identity uses only that $\widetilde Q$ is nondecreasing; in
particular, it holds when $Q$ has jumps or intervals of constancy.

\paragraph{Step 2: the two windows at a fixed radius.}
Let $0<r<h/2$ and $v\in(\underline v+r,\overline v-r)$. The windows
$I_-=(v-r,v)$ and $I_+=(v,v+r)$ lie in $I$, have length $r$, and have
threshold $u$. We set
\[
 \Delta_r(v)=Q(v+r)-Q(v-r)\ge0,\qquad
 Y_\pm=\E[\Lambda_{I_\pm}\mid\calG_{Q(v)}].
\]
Note that $Q\ge Q(v)$ on $I_+$ and $Q\le Q(v)$ on $I_-$, because $Q$ is
nondecreasing. It follows from the last assertion of \cref{lem:wi-two}
that $\kappa_{I_+}$ and $\kappa_{I_-}$ satisfy the hypotheses of
\cref{lem:wi-two}(a) and (b), respectively, with $q=Q(v)$. Further, by
monotonicity,
\[
 \begin{aligned}
 0&\le\bar p_{I_+}-Q(v)\le Q(v+r)-Q(v)\le\Delta_r(v),\\
 0&\le Q(v)-\bar p_{I_-}\le Q(v)-Q(v-r)\le\Delta_r(v).
 \end{aligned}
\]
Then \cref{lem:wi-two}(a) and (b) imply that
\[
 \xi_v\le\Delta_r(v)+2Y_+,\qquad
 -\xi_v\le\Delta_r(v)-2Y_- .
\]
On the event $\{\xi_v>0\}$, the first inequality implies that
$\xi_v^2\le(\Delta_r(v)+2|Y_+|)^2$, and this is at most
$2\Delta_r(v)^2+8Y_+^2$ because $(x+y)^2\le2x^2+2y^2$. On the event
$\{\xi_v<0\}$, we have $\xi_v^2\le2\Delta_r(v)^2+8Y_-^2$ in the same way,
by the second inequality. Since these two events
are disjoint, we have
\[
 \xi_v^2\le2\Delta_r(v)^2+8Y_+^2+8Y_-^2 .
\]
\Cref{lem:wi-two}(c) and (e) imply that
$\E Y_\pm^2\le\Var(\Lambda_{I_\pm})\le(Mur)^{-1}=\epsilon/r$. Taking
expectations, we conclude that
\begin{equation}\label{eq:wi-vpoint}
 \chi(v)\le2\Delta_r(v)^2+\frac{16\epsilon}{r}
 \qquad\text{for every }v\in(\underline v+r,\overline v-r).
\end{equation}
Note that $v$ and $r$ are fixed here. The windows and the level $Q(v)$
are then deterministic, and \cref{lem:wi-two} applies as stated.

\paragraph{Step 3: the measure of the set where the variance is large.}
Fix $\tau\in(0,1)$ and let $r=32\epsilon/\tau$, so that
$16\epsilon/r=\tau/2$. Suppose first that $r<h/2$. If
$v\in(\underline v+r,\overline v-r)$ and $\chi(v)>\tau$, then
\eqref{eq:wi-vpoint} implies that $2\Delta_r(v)^2>\tau/2$, that is,
$\Delta_r(v)>\sqrt\tau/2$. On this interior interval, we have
$\Delta_r(v)=\widetilde Q(v+r)-\widetilde Q(v-r)$. By Markov's
inequality and \eqref{eq:wi-variation}, the measure of the interior
points with $\chi(v)>\tau$ is at most
\[
 \frac2{\sqrt\tau}\int_\R\bigl(\widetilde Q(v+r)-\widetilde Q(v-r)\bigr)\dd v
 =\frac{4r\ell}{\sqrt\tau}=\frac{128\epsilon\ell}{\tau^{3/2}}.
\]
The two boundary strips $(\underline v,\underline v+r]$ and
$[\overline v-r,\overline v)$ have total length $2r=64\epsilon/\tau$. If
instead $r\ge h/2$, then $h\le2r=64\epsilon/\tau$, and the same bound
holds without any interior windows, because the whole interval has
length at most $2r$. In both cases,
using also the trivial bound $h$, we have
\[
 |\{v\in I:\chi(v)>\tau\}|
 \le\min\Bigl\{h,\frac{64\epsilon}{\tau}
                    +\frac{128\epsilon\ell}{\tau^{3/2}}\Bigr\}
 \le\min\Bigl\{h,\frac{64\epsilon}{\tau}\Bigr\}
   +\min\Bigl\{h,\frac{128\epsilon\ell}{\tau^{3/2}}\Bigr\}.
\]

\paragraph{Step 4: integration over the threshold.}
By the layer-cake formula, and since $0\le\chi\le1$, we have
\[
 \int_I\chi(v)\dd v=\int_0^1|\{v\in I:\chi(v)>\tau\}|\dd\tau .
\]
We integrate the two terms of the last bound of Step~3 separately. For
the first, we use $\min\{h,64y\}\le64\min\{h,y\}$ for $y\ge0$. If
$\epsilon\ge h$, then $\epsilon/\tau\ge h$ for $\tau\le1$, and the
integral of $\min\{h,\epsilon/\tau\}$ over $(0,1)$ equals $h$. If
$\epsilon<h$, we split the integral at $\tau=\epsilon/h$ and obtain
$\epsilon+\epsilon\log(h/\epsilon)$, which is at most $h$ because
$1+\log y\le y$ for $y\ge1$. In both cases, we have
\[
 \int_0^1\min\Bigl\{h,\frac{64\epsilon}\tau\Bigr\}\dd\tau
 \le64\min\bigl\{h,\epsilon\bigl(1+\log_+(h/\epsilon)\bigr)\bigr\}.
\]
For the second term, let $B=128\epsilon\ell$. If $\ell=0$, the term
vanishes. Otherwise, we extend the integral to $(0,\infty)$ and split
it at $\tau_0=(B/h)^{2/3}$, where $B\tau_0^{-3/2}=h$. This gives
\[
 \begin{aligned}
 \int_0^\infty\min\{h,B\tau^{-3/2}\}\dd\tau
 &=h\tau_0+2B\tau_0^{-1/2}=3B^{2/3}h^{1/3}\\
 &=3\cdot128^{2/3}(\epsilon\ell)^{2/3}h^{1/3}
 \le128(\epsilon\ell)^{2/3}h^{1/3}.
 \end{aligned}
\]
In the last inequality, we used that $3\cdot128^{2/3}\le128$, which is
equivalent to $27\le128$. Adding the two estimates, we obtain
\eqref{eq:wi-vfs}.
\end{proof}

\section{A cubic-moment bound}\label{sec:wi-cubic}

The cavity argument produces third moments of $S_q-q$. Fluctuations of
$\xi_v$ of size at most $z$ contribute at most $zV$ to $H_3$. For
fluctuations above a level $\vartheta$, we divide the range of levels
into bins of width $\vartheta/4$. Within a bin, the mean level of the
window after a point $v$ is at least $Q(v)$ and less than
$Q(v)+\vartheta/4$, and the mean level of the window before $v$ is at
most $Q(v)$ and greater than $Q(v)-\vartheta/4$. We use \cref{lem:wi-two}(d) to control both tails. Integrating
the reciprocal window lengths over $v$ produces one logarithm, and
integrating over $\vartheta$ produces the second.

\begin{corollary}[Cubic self-overlap moment]\label{cor:wi-cub}
Let $\mu$ be a minimizer of $\mathcal F_{t,s}$, and assume that $M\ge1$.
Then for every $0<z\le1/4$,
\begin{equation}\label{eq:wi-cub}
 H_3\le zV+\frac{425}{Ms}\log(2M)\log\frac2z .
\end{equation}
The constant is absolute; atoms of $\nu$, including an atom at zero, are
allowed.
\end{corollary}

\begin{proof}
The proof has three steps. We first reduce $H_3$ to integrals of the
tails of $\xi_v$, then bound these tail integrals using level bins, and
finally integrate the tail bound.

\paragraph{Step 1: reduction to tails.}
For $|y|\le1$ and $0<z\le1/4$, we have
\[
 |y|\le z+\int_z^1\ind_{\{|y|>\vartheta\}}\dd\vartheta .
\]
Indeed, the right side is at least $z$, and if $|y|>z$, the integral
equals $|y|-z$. Multiplying this inequality by $y^2$ gives
\[
 |y|^3\le zy^2+\int_z^1y^2\ind_{\{|y|>\vartheta\}}\dd\vartheta .
\]
We use this inequality with $y=\xi_v$, take expectations, and
integrate over $v\in(0,1-s)$. We define
\[
 T(\vartheta)=\int_0^{1-s}
 \E\bigl[\xi_v^2\ind_{\{|\xi_v|>\vartheta\}}\bigr]\dd v
 \qquad(0<\vartheta\le1).
\]
Using Tonelli's theorem and the identities
$V=\int_0^{1-s}\E\xi_v^2\dd v$ and $H_3=\int_0^{1-s}\E|\xi_v|^3\dd v$, we
obtain
\begin{equation}\label{eq:wi-layer}
 H_3\le zV+\int_z^1T(\vartheta)\dd\vartheta.
\end{equation}

\paragraph{Step 2: the tail bound.}
Fix $\vartheta\in(0,1]$. For $0\le k\le\lfloor4/\vartheta\rfloor$, let
$\mathcal B_k=[k\vartheta/4,(k+1)\vartheta/4)$. These level bins cover
$[0,1]$, and there are at most $4/\vartheta+1\le5/\vartheta$ of them.
The preimages
\[
 I_k=\{v\in(0,1-s):Q(v)\in\mathcal B_k\}
\]
are intervals, because $Q$ is nondecreasing. Some of them may be empty,
and we discard those. For the others, we write
$\underline v_k\le\overline v_k$ for the endpoints of $I_k$ and $h_k=\overline v_k-\underline v_k\le1$ for its
length. The intervals $I_k$ partition $(0,1-s)$, and those with
$h_k=0$ do not contribute to $T(\vartheta)$.

Fix $k$ with $h_k>0$ and $v\in(\underline v_k,\overline v_k)$. The later
window $(v,\overline v_k)$ lies in $I_k$. For $w$ in this window, we
have $Q(v)\le Q(w)<Q(v)+\vartheta/4$, since $Q(w)$ and $Q(v)$ lie in the
same bin of width $\vartheta/4$. In particular, the mean level of the
window is at least $Q(v)$ and at most $Q(v)+\vartheta/4$. Likewise, the earlier window
$(\underline v_k,v)$ satisfies $Q\le Q(v)$, and its mean level is at
least $Q(v)-\vartheta/4$. Both windows have threshold $s$ by
\eqref{eq:wi-supp}. By the last assertion of \cref{lem:wi-two}, their
normalized densities satisfy the hypotheses of (a) and (b),
respectively, with $q=Q(v)$. Since $\vartheta/4\le\vartheta/2$,
\cref{lem:wi-two}(d) and (e) now imply that
\[
 \E\bigl[\xi_v^2\ind_{\{\xi_v>\vartheta\}}\bigr]
 \le\frac{16}{Ms(\overline v_k-v)},\qquad
 \E\bigl[\xi_v^2\ind_{\{\xi_v<-\vartheta\}}\bigr]
 \le\frac{16}{Ms(v-\underline v_k)} .
\]

Suppose that $h_k>1/M$. We integrate the first bound over
$(\underline v_k,\overline v_k-1/M)$ and use $|\xi_v|\le1$ on the
remaining strip of length $1/M$. This gives
\[
 \begin{aligned}
 \int_{I_k}\E\bigl[\xi_v^2\ind_{\{\xi_v>\vartheta\}}\bigr]\dd v
 &\le\frac{16}{Ms}\int_{\underline v_k}^{\overline v_k-1/M}
       \frac{\dd v}{\overline v_k-v}+\frac1M
 =\frac{16}{Ms}\log(Mh_k)+\frac1M\\
 &\le\frac{16\log(2M)}{Ms}+\frac1M,
 \end{aligned}
\]
where we used $h_k\le1$ in the last line. The lower tail satisfies the
same bound, with the strip of length $1/M$ at the left end of $I_k$. If
$h_k\le1/M$, then both tails together contribute at most $h_k\le1/M$,
since the two events are disjoint and $|\xi_v|\le1$. Since $M\ge1$ and
$s\le1/2$, we have
\[
 \frac2M\le\frac1{Ms}\le\frac{2\log(2M)}{Ms},
\]
where we used $\log2\ge1/2$ in the second inequality. Each bin therefore
contributes at most $32\log(2M)/(Ms)+2/M\le34\log(2M)/(Ms)$ to
$T(\vartheta)$. Summing over the at most $5/\vartheta$ bins, we conclude that
\[
 T(\vartheta)\le\frac5\vartheta\cdot\frac{34\log(2M)}{Ms}
 =\frac{170\log(2M)}{Ms\,\vartheta}\qquad(0<\vartheta\le1).
\]

\paragraph{Step 3: integration.}
The tail bound holds on the whole range $(z,1]$ of the integral in
\eqref{eq:wi-layer}. Inserting it, we obtain
\[
 H_3\le zV+\frac{170\log(2M)}{Ms}\int_z^1\frac{\dd\vartheta}\vartheta
 =zV+\frac{170\log(2M)}{Ms}\log\frac1z
 \le zV+\frac{425\log(2M)}{Ms}\log\frac2z .
\]
This proves \eqref{eq:wi-cub}. In fact, we have proved the stronger
bound with $170\log(1/z)$ in place of $425\log(2/z)$.
\end{proof}

\chapter{The mean--covariance term}\label{sec:th}

In this section, we bound the mean--covariance term
$J=\int j_q\,\nu(\dd q)$ at a minimizing profile. Here
$j_q=N^{-2}\E[m_q^TC_qm_q]$ is the expected conditional variance of the
spin in the direction of its conditional mean. Our main result,
\cref{prop:th-j}, states that $J\le341(1+\log_+M)^2/(Ms)$ at every
minimizer. Together with the trace bound $A\le1/(Ms)$ of
\eqref{eq:in-Lvar}, this controls the thermal part $A+2J$ of the defect
$\calD_N(\rho)=\beta(A+2J+V)$. As
$Ms=\beta^2N\rho\cdot\lambda/\beta=\beta N\rho\lambda$, the two bounds
imply that
\[
 \beta(A+2J)\le\frac{1+682(1+\log_+M)^2}{N\rho\lambda}.
\]
We use \cref{prop:th-j} in this form in the proof of
\cref{cor:c7-defect}, for those $\rho$ where the cavity inequality
contracts, and in \eqref{eq:sm-AJ} in the proof of \cref{prop:sm-small},
for small $\rho$.

The proof has four steps.
\begin{enumerate}
\item In \cref{lem:th-ric}, we show by a scalar Riccati inequality that
conditional covariances decay in every direction that is measurable at
an earlier level: if $b$ is $\calG_{q'}$-measurable, then
$\E[b^TC_qb]$ is at most $\E|b|^2$ divided by the accumulated clock
$c\int_{q'}^qa(u)^2\dd u$. This holds for every admissible profile.
\item In \cref{lem:th-gamma}, we use first-order optimality to show that
$\mathsf S(r)=\E S_r$ is $2$-Lipschitz toward every point $q$ of
$\supp\nu$, also when $r$ lies in a gap of the support. By
\eqref{eq:th-orth} below, this bounds $\E|m_q-m_r|^2$ by $2N|q-r|$.
\item In \cref{lem:th-res}, we combine these two estimates through a
dyadic decomposition of $m_q-m_p$, for $q\in\supp\nu$ and $p\le q$. We
bound $\calR(p,q)=N^{-2}\E[(m_q-m_p)^TC_q(m_q-m_p)]$ by a multiple of
$(1+\log_+M)^2/(Ma((p+q)/2)^2)$.
\item In \cref{prop:th-j}, we average earlier means $m_{Q(w)}$, $0<w<v$,
over the whole quantile prefix with a quadratic weight, together with a
deterministic weight on $m_0=0$. We control the conditional covariance in the direction of this average by
\eqref{eq:in-Lvar}, and
its distance from $m_{Q(v)}$ by \cref{lem:th-res}. The weight produces a factor $(s+v)^{-2}$, whose
integral over $v\in(0,1-s)$ is at most $1/s$.
\end{enumerate}

We use \cref{prop:fnd-contact}, namely the martingale equation and
covariance order \eqref{eq:in-martingale}, the first variation
\eqref{eq:fnd-first-variation}, and the contact identities
\eqref{eq:in-contact}; \cref{prop:fnd-variance}; and facts about the
quantile function from \cref{sec:wi-setup}, including
\eqref{eq:wi-supp}. \Cref{lem:th-ric,lem:th-gamma,lem:th-res} are used
only in this section.

\section{Standing assumptions}\label{sec:th-setup}

Throughout this section, we fix $N\ge1$, $0<\rho\le1$, and $0<s\le1/2$,
so that $t=1-\rho\in[0,1)$, $c>0$, and $M=Nc>0$. Expectations refer to the joint law of the
deterministic admissible profile $\mu=s\delta_0+\nu$ under
consideration. \Cref{lem:th-ric} holds for every admissible profile.
The later results assume that $\mu$ minimizes $\mathcal F_{t,s}$; they
hold for every minimizer, with absolute constants.

We write $\mathsf S(q)=\E S_q$ for the expected self-overlap. Recall
from \eqref{eq:in-martingale} that $a(q)C_q\preceq H_q\preceq C_q$, and
in particular $0\preceq H_q\preceq C_q$. The diagonal entries of $C_q$
are $1-m_{q,i}^2\le1$. These two facts give
\begin{equation}\label{eq:th-op}
 \|H_q\|_{\rm op}\le\|C_q\|_{\rm op}\le\Tr C_q\le N,\qquad
 0\le\mathsf S(q)\le1,
\end{equation}
where $\|\cdot\|_{\rm op}$ is the operator norm. The process $m$ is a
bounded martingale, and $\E[m_q\mid\calG_p]=m_p$ for $p\le q$. Then
$\E|m_q|^2=\E|m_p|^2+\E|m_q-m_p|^2$. In particular, $\mathsf S$ is
nondecreasing, and
\begin{equation}\label{eq:th-orth}
 \E|m_q-m_p|^2=N\bigl(\mathsf S(q)-\mathsf S(p)\bigr)
 \qquad(0\le p\le q\le1).
\end{equation}
For $0\le p\le q\le1$, we define
\[
 \calR(p,q)=N^{-2}\E\bigl[(m_q-m_p)^TC_q(m_q-m_p)\bigr],
\]
which is the analog of $j_q$ for the increment $m_q-m_p$ in place of
$m_q$. Note that $\calR(0,q)=j_q$, since $m_0=0$ by
\eqref{eq:in-martingale}.

\section{Covariance along an earlier measurable direction}\label{sec:th-decay}

Fix a direction that is measurable at some level. As more of the
auxiliary field is revealed, the expected conditional variance of the
spin in this direction can only decrease. The martingale equation
determines the rate of decrease, and the covariance order bounds this
rate from below by a multiple of the square of the variance itself. We integrate
the resulting scalar Riccati inequality and obtain decay that is
inversely proportional to the accumulated clock.

\begin{lemma}[Directional covariance decay]\label{lem:th-ric}
Let $\mu=s\delta_0+\nu$ be any deterministic admissible profile, let
$0\le q'<q\le1$, and let
\[
 \mathcal I=c\int_{q'}^qa(u)^2\dd u .
\]
Then $\mathcal I\ge cs^2(q-q')>0$, and for every
$\calG_{q'}$-measurable random vector $b\in\R^N$ with $\E|b|^2<\infty$,
\[
 \E[b^TC_qb]\le\frac{\E|b|^2}{\mathcal I}.
\]
\end{lemma}

\begin{proof}
The lower bound for $\mathcal I$ holds because $a\ge s$. We first prove
the main bound for bounded $b$, by deriving a differential inequality for
the function $r\mapsto\E[b^TC_rb]$ on $[q',q]$, and we then remove the
boundedness by monotone convergence.

\paragraph{Step 1: the rate of decrease.}
Suppose that $b$ is bounded, and define $\varphi(r)=\E[b^TC_rb]$ for
$r\in[q',q]$. For such $r$, the vector $b$ is $\calG_r$-measurable, so
that $b^TC_rb=\E[(b\cdot(\sigma-m_r))^2\mid\calG_r]$ and
$\varphi(r)=\E(b\cdot(\sigma-m_r))^2$. We write
\[
 b\cdot(\sigma-m_{q'})=b\cdot(\sigma-m_r)+b\cdot(m_r-m_{q'}).
\]
The second term is $\calG_r$-measurable, and the first has conditional
mean zero given $\calG_r$. The cross term has expectation
zero, and
\[
 \varphi(q')=\varphi(r)+\E\bigl(b\cdot(m_r-m_{q'})\bigr)^2 .
\]
By \eqref{eq:in-martingale}, and since $b$ is $\calG_{q'}$-measurable
and bounded, we have
$b\cdot(m_r-m_{q'})=\sqrt c\int_{q'}^r(H_ub)\cdot\dd B_u$. The integrand
is progressive and bounded by $N|b|$, by \eqref{eq:th-op}. Itô's isometry
gives
\[
 \varphi(r)=\varphi(q')-c\int_{q'}^r\E|H_ub|^2\dd u\qquad(q'\le r\le q).
\]
Then $\varphi$ is absolutely continuous and nonincreasing, with
$\varphi'(r)=-c\E|H_rb|^2$ for almost every $r\in[q',q]$.

\paragraph{Step 2: the Riccati inequality.}
On the event $\{b\ne0\}$, the Cauchy--Schwarz inequality and the
covariance order $0\le a(r)b^TC_rb\le b^TH_rb$ imply that
\[
 |H_rb|^2\ge\frac{(b^TH_rb)^2}{|b|^2}\ge\frac{a(r)^2(b^TC_rb)^2}{|b|^2}.
\]
Let $Z_r=|b|^{-1}b^TC_rb$ on $\{b\ne0\}$ and $Z_r=0$ otherwise. Then
$\E|H_rb|^2\ge a(r)^2\E Z_r^2$. Further, $\varphi(r)=\E[Z_r|b|]$, and by the
Cauchy--Schwarz inequality, $\varphi(r)^2\le\E Z_r^2\,\E|b|^2$. If
$\E|b|^2>0$, we conclude that
\[
 \varphi'(r)\le-\frac{c\,a(r)^2\varphi(r)^2}{\E|b|^2}
 \qquad\text{for almost every }r\in[q',q].
\]

\paragraph{Step 3: integration.}
If $\E|b|^2=0$ or $\varphi(q)=0$, there is nothing to prove. Otherwise,
$\varphi\ge\varphi(q)>0$ on $[q',q]$, because $\varphi$ is nonincreasing.
The function $1/\varphi$ is then absolutely continuous, with
$(1/\varphi)'=-\varphi'/\varphi^2\ge c\,a(r)^2/\E|b|^2$ almost
everywhere. Integrating over $[q',q]$, we find
\[
 \frac1{\varphi(q)}\ge\frac1{\varphi(q')}+\frac{\mathcal I}{\E|b|^2}
 \ge\frac{\mathcal I}{\E|b|^2},
\]
which is the claim for bounded $b$. The profile $a$ may have jumps, but
only Lebesgue-almost every level enters this integration.

\paragraph{Step 4: general directions.}
For general $b$ with $\E|b|^2<\infty$, let $b_n=b\ind_{\{|b|\le n\}}$,
which is bounded and $\calG_{q'}$-measurable. The quantity
$b_n^TC_qb_n=(b^TC_qb)\ind_{\{|b|\le n\}}$ increases to $b^TC_qb$ as
$n\to\infty$, because $C_q\succeq0$, and $|b_n|^2$ increases to
$|b|^2$. The bound holds for each $b_n$, and monotone convergence as
$n\to\infty$ yields the claim for $b$.
\end{proof}

\section{Mean increments toward a support point}\label{sec:th-support}

In the dyadic argument below, we need $\E|m_q-m_r|^2$ to be at most of
order $N|q-r|$ when $q$ is a point of $\supp\nu$, including when $r$
lies in a gap of the support. In view of \eqref{eq:th-orth}, this is a
Lipschitz bound for $\mathsf S$ toward $q$, and it follows from the global
first-order optimality of a minimizer. Let
\[
 G(y)=\int_y^1(\mathsf S(u)-u)\dd u\qquad(0\le y\le1)
\]
be the function in the first variation \eqref{eq:fnd-first-variation}.
\Cref{prop:fnd-contact} states that if $\mu$ is a minimizer, then
$\supp\nu$ is contained in the set of minimum points of $G$ on $[0,1]$,
that is,
\begin{equation}\label{eq:th-argmin}
 G(y)\ge G(q)\qquad(q\in\supp\nu,\ 0\le y\le1).
\end{equation}
We recall the argument from the proof of \cite[Proposition~2.1]{AL},
because we will use \eqref{eq:th-argmin} at points $y$ outside
$\supp\nu$. For $y\in[0,1]$, the profiles
$s\delta_0+(1-\theta)\nu+\theta(1-s)\delta_y$, $0\le\theta\le1$, are
admissible. At a minimizer, the right derivative of $\mathcal F_{t,s}$
along them at $\theta=0$, which is \eqref{eq:fnd-first-variation} with
$\nu'=(1-s)\delta_y$, is nonnegative. Dividing by $c/2$, we
find that
\[
 (1-s)G(y)\ge\int G(p)\,\nu(\dd p)\qquad(0\le y\le1).
\]
The function $G$ is continuous. Taking $y$ to be a minimum point of $G$
and recalling that $\nu$ has mass $1-s$, we find that
$\int(G-\min G)\dd\nu\le0$. Then the open set $\{G>\min G\}$ is
$\nu$-null and therefore does not meet $\supp\nu$.

\begin{lemma}[Lipschitz bound toward the support]\label{lem:th-gamma}
Let $\mu$ be a minimizer of $\mathcal F_{t,s}$. For every $q\in\supp\nu$
and every $r\in[0,1]$,
\[
 |\mathsf S(r)-\mathsf S(q)|\le2|r-q|;\qquad\text{in particular,}\qquad
 0\le\mathsf S(q)-\mathsf S(r)\le2(q-r)\quad(0\le r\le q).
\]
\end{lemma}

\begin{proof}
Let $q\in\supp\nu$, so that $\mathsf S(q)=q$ by \eqref{eq:in-contact}.
Recall that $\mathsf S$ is nondecreasing with values in $[0,1]$, and
note that $G(y_1)-G(y_2)=\int_{y_1}^{y_2}(\mathsf S(u)-u)\dd u$ for
$y_1\le y_2$. Let $r\in[0,1]$ and $\xi=\mathsf S(r)\in[0,1]$. We compare
the values of $G$ at $\xi$ and at $q$. The point $\xi$ need not lie in
$\supp\nu$, and this is why we need \eqref{eq:th-argmin} for every
$y\in[0,1]$.

\paragraph{The case $r\le q$.}
By monotonicity, $\mathsf S(q)-\mathsf S(r)\ge0$. If $\xi\ge r$, then
$\mathsf S(q)-\mathsf S(r)=q-\xi\le q-r$. If $\xi<r$, then
$\mathsf S\le\mathsf S(r)=\xi$ on $[\xi,r]$ and
$\mathsf S\le\mathsf S(q)=q$ on $[r,q]$, by monotonicity. By
\eqref{eq:th-argmin} with $y=\xi$, we have
\[
 \begin{aligned}
 0\le G(\xi)-G(q)
 &=\int_\xi^r(\mathsf S(u)-u)\dd u+\int_r^q(\mathsf S(u)-u)\dd u\\
 &\le\int_\xi^r(\xi-u)\dd u+\int_r^q(q-u)\dd u
 =\frac{(q-r)^2-(r-\xi)^2}2 .
 \end{aligned}
\]
Then $r-\xi\le q-r$, and
$\mathsf S(q)-\mathsf S(r)=(q-r)+(r-\xi)\le2(q-r)$.

\paragraph{The case $r>q$.}
By monotonicity, $\mathsf S(r)-\mathsf S(q)\ge0$. If $\xi\le r$, then
$\mathsf S(r)-\mathsf S(q)=\xi-q\le r-q$. If $\xi>r$, then
$\mathsf S\ge q$ on $[q,r]$ and $\mathsf S\ge\xi$ on $[r,\xi]$, and in
the same way
\[
 \begin{aligned}
 0\le G(\xi)-G(q)
 &=-\int_q^r(\mathsf S(u)-u)\dd u-\int_r^\xi(\mathsf S(u)-u)\dd u\\
 &\le\int_q^r(u-q)\dd u+\int_r^\xi(u-\xi)\dd u
 =\frac{(r-q)^2-(\xi-r)^2}2 .
 \end{aligned}
\]
We obtain $\xi-r\le r-q$, and
$\mathsf S(r)-\mathsf S(q)=(r-q)+(\xi-r)\le2(r-q)$. This completes the
proof.
\end{proof}

Together, \eqref{eq:th-orth} and \cref{lem:th-gamma} imply that
\[
 \E|m_q-m_r|^2\le2N(q-r)\qquad(q\in\supp\nu,\ 0\le r\le q).
\]

\section{A dyadic decomposition of the mean increment}\label{sec:th-dyadic}

\Cref{lem:th-ric} does not apply directly to the increment $m_q-m_p$,
because this increment is not measurable at any level before $q$.
Instead, we split it into increments whose right endpoints approach $q$
geometrically. Each of them except the last is measurable strictly
before $q$, and \cref{lem:th-ric} applies to it. By \cref{lem:th-gamma},
its second moment is at most proportional to the remaining distance to $q$,
which cancels the length of the clock in \cref{lem:th-ric}. The last,
short increment is controlled by $C_q\preceq NI$.

\begin{lemma}[Covariance of a mean increment]\label{lem:th-res}
Let $\mu$ be a minimizer of $\mathcal F_{t,s}$, let $q\in\supp\nu$ and
$0\le p\le q$, and let $a_*=a((p+q)/2)$. For every integer $K\ge1$,
\begin{equation}\label{eq:th-res-K}
 \calR(p,q)\le(K+1)\Bigl[\frac{4K}{Ma_*^2}+2(q-p)2^{-K}\Bigr].
\end{equation}
Moreover,
\begin{equation}\label{eq:th-res}
 \calR(p,q)\le\frac{56\bigl(1+\log_+(M(q-p))\bigr)^2}{Ma_*^2}.
\end{equation}
\end{lemma}

\begin{proof}
We may assume that $p<q$, since for $p=q$ the left side vanishes and
both bounds hold. The proof has three steps. We split $m_q-m_p$ into
$K+1$ dyadic pieces; we bound the first $K$ pieces by
\cref{lem:th-ric,lem:th-gamma} and the last piece directly; and we then
choose $K$.

For $0\le k\le K$, let $q_k=q-(q-p)2^{-k}$, so that $q_0=p$,
$q_1=(p+q)/2$, and $q-q_k=(q-p)2^{-k}$. We set
$b_k=m_{q_{k+1}}-m_{q_k}$ for $0\le k<K$ and $b_K=m_q-m_{q_K}$. Then
$m_q-m_p=\sum_{k=0}^Kb_k$. Applying the Cauchy--Schwarz inequality to the
quadratic form of $C_q$, which is positive semidefinite, we find
\[
 (m_q-m_p)^TC_q(m_q-m_p)\le(K+1)\sum_{k=0}^Kb_k^TC_qb_k .
\]
The points $q_k$ need not be contact levels; we use \cref{lem:th-gamma}
only with the endpoint $q\in\supp\nu$.

\paragraph{Step 1: the first $K$ pieces.}
Let $0\le k<K$. The vector $b_k$ is $\calG_{q_{k+1}}$-measurable, and
$q_{k+1}<q$. Since $a$ is nondecreasing and $q_{k+1}\ge q_1$, we have
$a(u)\ge a(q_1)=a_*$ for $u\ge q_{k+1}$. The clock of \cref{lem:th-ric}
with $q'=q_{k+1}$ satisfies
\[
 \mathcal I=c\int_{q_{k+1}}^qa(u)^2\dd u\ge ca_*^2(q-q_{k+1}),
\]
and \cref{lem:th-ric} implies that
$\E[b_k^TC_qb_k]\le\E|b_k|^2/(ca_*^2(q-q_{k+1}))$. By \eqref{eq:th-orth},
the monotonicity of $\mathsf S$, and \cref{lem:th-gamma} with $r=q_k$, we
have
\[
 \E|b_k|^2=N\bigl(\mathsf S(q_{k+1})-\mathsf S(q_k)\bigr)
 \le N\bigl(\mathsf S(q)-\mathsf S(q_k)\bigr)\le2N(q-q_k)=4N(q-q_{k+1}).
\]
The last two bounds imply that
\[
 N^{-2}\E[b_k^TC_qb_k]\le\frac4{Nca_*^2}=\frac4{Ma_*^2}
 \qquad(0\le k<K).
\]

\paragraph{Step 2: the last piece.}
Using $\|C_q\|_{\rm op}\le N$ from \eqref{eq:th-op}, then
\eqref{eq:th-orth} and \cref{lem:th-gamma} with $r=q_K$, we have
\[
 N^{-2}\E[b_K^TC_qb_K]\le N^{-1}\E|b_K|^2=\mathsf S(q)-\mathsf S(q_K)
 \le2(q-q_K)=2(q-p)2^{-K}.
\]
We sum the bounds of Steps~1 and~2 over the $K+1$ pieces and insert
them into the Cauchy--Schwarz bound above. This proves
\eqref{eq:th-res-K}.

\paragraph{Step 3: the choice of $K$.}
Let $z=M(q-p)$, $\theta=\log_+z$, and $K=\lceil\log_2(2+z)\rceil$, so that
$K\ge1$. Then $2^K\ge2+z>z$, which gives
\[
 2(q-p)2^{-K}=\frac{2z}{M2^K}\le\frac2M\le\frac2{Ma_*^2},
\]
where we used $a_*\le1$. With this choice of $K$, \eqref{eq:th-res-K}
implies that
\[
 \calR(p,q)\le\frac{(K+1)(4K+2)}{Ma_*^2}.
\]
Since $2+z\le3\max\{1,z\}$, we have
$\log_2(2+z)\le\log_23+\theta/\log2<2+2\theta$, where we used
$\log_23<2$ and $1/\log2<2$. It follows that
$K<1+\log_2(2+z)<3+2\theta$, and
\[
 (K+1)(4K+2)\le(4+2\theta)(14+8\theta)=56+60\theta+16\theta^2
 \le56(1+\theta)^2 .
\]
This proves \eqref{eq:th-res}.
\end{proof}

The two powers of the logarithm in \eqref{eq:th-res} have different
sources: there are $K$ increments of comparable cost, and the
Cauchy--Schwarz inequality for their sum contributes the factor $K+1$.

\section{An integrated bound for \texorpdfstring{$j$}{j}}\label{sec:th-integrated}

With $p=0$, \cref{lem:th-res} bounds $j_q=\calR(0,q)$ by a multiple of
$(1+\log_+M)^2/(Ma(q/2)^2)$. To integrate against $\nu$, we need a bound
in terms of the cumulative mass $s+v$ at $q=Q(v)$, rather than
$a(q/2)$, which may be as small as $s$. We therefore average the earlier
means $m_{Q(w)}$, $0<w<v$, with density proportional to $(s+w)^2$, and
assign the remaining weight to $m_0=0$. The conditional covariance in the direction of
this average is controlled by the integrated variance bound
\eqref{eq:in-Lvar}, and its distance from $m_{Q(v)}$ by
\cref{lem:th-res}. With the quadratic weight, both estimates are of
order $(s+v)^{-2}$, also for $v<s$, and no separate argument is needed
near the root.

\begin{proposition}[Covariance in the mean direction]\label{prop:th-j}
Let $\mu=s\delta_0+\nu$ be a minimizer of $\mathcal F_{t,s}$. For every
$v\in(0,1-s)$,
\begin{equation}\label{eq:th-j-pointwise}
 j_{Q(v)}\le\frac{341(1+\log_+M)^2}{M(s+v)^2},
\end{equation}
and
\begin{equation}\label{eq:th-J}
 J=\int j_q\,\nu(\dd q)\le\frac{341(1+\log_+M)^2}{Ms}.
\end{equation}
The constant is absolute. The bounds hold for all $N\ge1$, $c>0$,
$0<s\le1/2$, and every minimizing profile, including profiles with atoms.
\end{proposition}

\begin{proof}
Fix $v\in(0,1-s)$, and let $q=Q(v)$ and $\mathcal L=1+\log_+M$. By
\eqref{eq:wi-supp}, we have $q\in\supp\nu$, and $Q(w)\le q$ for $w<v$
because $Q$ is nondecreasing. If $q=0$, then $m_q=m_0=0$ and $j_q=0$, so
we may assume that $q>0$. The proof has four steps. First, we write
$m_q$ as a weighted average $\bar m$ of earlier means plus a remainder
$\bar\Delta$, and we prove the bound \eqref{eq:th-split} below. Second,
we bound the conditional covariance in the direction of $\bar m$ by
\eqref{eq:in-Lvar},
using the conditional Jensen inequality on quantile fibers. Third, we
bound the remainder by \cref{lem:th-res}. Fourth, we combine the bounds
and integrate over $v$.

\paragraph{Step 1: the decomposition.}
We use the deterministic weights
\[
 r_0=\frac{s^3}{(s+v)^3},\qquad
 k_v(w)=\frac{3(s+w)^2}{(s+v)^3}\quad(0<w<v),
\]
which satisfy
\[
 r_0+\int_0^vk_v(w)\dd w=\frac{s^3+\bigl((s+v)^3-s^3\bigr)}{(s+v)^3}=1 .
\]
We use the probability measure $r_0\delta_0+k_v(w)\dd w$ on $[0,v)$ to
average the vector $m_0=0$ and the vectors $m_{Q(w)}$, $0<w<v$. We define
\[
 \begin{gathered}
 \bar m=\int_0^vk_v(w)m_{Q(w)}\dd w,\qquad
 \bar L=\int_0^vk_v(w)L_{Q(w)}\dd w,\\
 \bar\Delta=r_0m_q+\int_0^vk_v(w)\bigl(m_q-m_{Q(w)}\bigr)\dd w.
 \end{gathered}
\]
By the normalization of the weights, $m_q=\bar m+\bar\Delta$. The vector
$\bar m$ and the number $D=2^{-1}\int_0^vk_v(w)|m_{Q(w)}|^2\dd w$ are
$\calG_q$-measurable, because $Q(w)\le q$ for $w<v$, and
$N\bar L=\sigma\cdot\bar m-D$. It follows that
$\Var(\bar L\mid\calG_q)=N^{-2}\bar m^TC_q\bar m$, and by the law of
total variance,
\[
 N^{-2}\E[\bar m^TC_q\bar m]=\E\Var(\bar L\mid\calG_q)\le\Var(\bar L).
\]
Next, we write
$\bar\Delta=r_0(m_q-m_0)+\int_0^vk_v(w)(m_q-m_{Q(w)})\dd w$, which is
the average of the vectors $m_q-m_0$ and $m_q-m_{Q(w)}$ with respect to
the probability measure above. By Jensen's inequality for the
convex quadratic form of $C_q$, we have
\[
 \bar\Delta^TC_q\bar\Delta
 \le r_0(m_q-m_0)^TC_q(m_q-m_0)
 +\int_0^vk_v(w)(m_q-m_{Q(w)})^TC_q(m_q-m_{Q(w)})\dd w .
\]
The integrand is bounded by $4N^2$, by \eqref{eq:th-op}. We take
expectations, use Fubini's theorem, and divide by $N^2$. This gives
\[
 N^{-2}\E[\bar\Delta^TC_q\bar\Delta]
 \le r_0\calR(0,q)+\int_0^vk_v(w)\calR(Q(w),q)\dd w .
\]
Finally, since $C_q\succeq0$, we have
$m_q^TC_qm_q\le2\bar m^TC_q\bar m+2\bar\Delta^TC_q\bar\Delta$. Combining
the last three bounds, we obtain
\begin{equation}\label{eq:th-split}
 j_q\le2\Var(\bar L)
       +2r_0\calR(0,q)+2\int_0^v k_v(w)\calR(Q(w),q)\dd w.
\end{equation}
The quantity $\calR(0,q)=j_q$ appears on the right side. This is not
circular: in Step~3, we bound it by \cref{lem:th-res}, independently of
\eqref{eq:th-split}.

\paragraph{Step 2: the variance of the weighted prefix.}
We extend $k_v$ by zero to $(0,1-s)$, and note that $k_v\le3/(s+v)$.
The image of the measure $k_v(w)\dd w$ under $Q$ is then at most
$3/(s+v)$ times $\nu$, and it has a density $\kappa_v$ with respect to
$\nu$, with $0\le\kappa_v\le3/(s+v)$; in particular,
$\kappa_v\in L^2(\nu)$. By the definition of $\kappa_v$ and the pushforward property of $Q$, we have,
for every bounded Borel function $\varphi$,
\[
 \int_0^{1-s}k_v(w)\varphi(Q(w))\dd w=\int\kappa_v\varphi\dd\nu
 =\int_0^{1-s}\kappa_v(Q(w))\varphi(Q(w))\dd w .
\]
Applied for each outcome to the bounded function $p\mapsto L_p$, the
first identity shows that $\bar L=\int\kappa_v(p)L_p\,\nu(\dd p)$. The identity as a whole says
that, under the normalized Lebesgue measure on $(0,1-s)$, the function
$\kappa_v\circ Q$ is the conditional expectation of $k_v$ given the
$\sigma$-algebra generated by $Q$. Explicitly, if $p$ is an atom of
$\nu$ with mass $r_p$, then
\[
 \kappa_v(p)=\frac1{r_p}\int_{(F_-(p),F(p))}k_v(w)\dd w,
\]
with the notation of \cref{sec:wi-setup}. A prefix $(0,v)$ that ends
inside the quantile interval of an atom gives the atom the average
of $k_v$ over that interval. Since $1/a(Q)$ is a bounded nonnegative
function of $Q$, we may apply the conditional Jensen inequality and find
\[
 \int\frac{\kappa_v(p)^2}{a(p)}\,\nu(\dd p)
 =\int_0^{1-s}\frac{\kappa_v(Q(w))^2}{a(Q(w))}\dd w
 \le\int_0^{1-s}\frac{k_v(w)^2}{a(Q(w))}\dd w
 =\int_0^v\frac{k_v(w)^2}{a(Q(w))}\dd w .
\]
As $\mu$ is a minimizer, \cref{prop:fnd-variance} applies with
$k=\kappa_v$. Together with $a(Q(w))\ge s+w$, this yields
\[
 \begin{aligned}
 \Var(\bar L)
 &\le\frac1M\int_0^v\frac{k_v(w)^2}{s+w}\dd w
 =\frac9{M(s+v)^6}\int_0^v(s+w)^3\dd w\\
 &=\frac9{4M(s+v)^2}\Bigl[1-\frac{s^4}{(s+v)^4}\Bigr]
 \le\frac9{4M(s+v)^2}.
 \end{aligned}
\]

\paragraph{Step 3: the remainder.}
Let $0<w<v$ and $p=Q(w)$. Then $p\le q$, and since $a$ is
nondecreasing, $a((p+q)/2)\ge a(p)\ge s+w$. For the root term, we have
$a(q/2)\ge s$ by \eqref{eq:wi-supp}. All level differences are at most
one, and hence $\log_+(M(q-p))\le\log_+M$. Applying \cref{lem:th-res}
with the endpoint $q\in\supp\nu$ and the lower endpoints $0$ and $Q(w)$, we have
\[
 \calR(0,q)\le\frac{56\mathcal L^2}{Ms^2},\qquad
 \calR(Q(w),q)\le\frac{56\mathcal L^2}{M(s+w)^2}\quad(0<w<v).
\]
Using $r_0s^{-2}=s/(s+v)^3$ and $k_v(w)(s+w)^{-2}=3/(s+v)^3$, we find
that
\[
 r_0\calR(0,q)+\int_0^vk_v(w)\calR(Q(w),q)\dd w
 \le\frac{56\mathcal L^2}{M}\cdot\frac{s+3v}{(s+v)^3}
 \le\frac{168\mathcal L^2}{M(s+v)^2},
\]
where we used $s+3v\le3(s+v)$.

\paragraph{Step 4: conclusion.}
Inserting the bounds of Steps~2 and~3 into \eqref{eq:th-split} and using
$\mathcal L\ge1$, we obtain
\[
 j_{Q(v)}\le\frac9{2M(s+v)^2}+\frac{336\mathcal L^2}{M(s+v)^2}
 \le\frac{(9/2+336)\mathcal L^2}{M(s+v)^2}
 \le\frac{341\mathcal L^2}{M(s+v)^2},
\]
which is \eqref{eq:th-j-pointwise}. Finally,
$J=\int_0^{1-s}j_{Q(v)}\dd v$, because $Q$ pushes Lebesgue measure on
$(0,1-s)$ forward to $\nu$, and \eqref{eq:th-J} follows from
\[
 \int_0^{1-s}\frac{\dd v}{(s+v)^2}=\frac1s-1\le\frac1s .
\]
This completes the proof.
\end{proof}


\part{The minimizing profile and the scalar response}\label{part:profile}
\chapter{Isolation of the minimizing profile}\label{sec:lz}

In this section and the next, we study an arbitrary minimizer $\gamma_*$ of
the finite-dimensional functional $F_\rho$ in the energy normalization of
\cref{sec:fnd}. We compare $\gamma_*$ with the scalar trial profile
$\gamma^{\rm tr}_\lambda=\max\{\lambda,\gamma_\beta\}$. We state all estimates
in terms of a trial budget $\eta$, that is, a bound on the excess
$F_\rho(\gamma^{\rm tr}_\lambda)-F_\rho(\gamma_*)$. The main result of this
section is \cref{prop:loc-isolation}, which bounds the $L^1$ distance between
$\gamma_*$ and $\gamma^{\rm tr}_\lambda$ by
$C\min\{1,\rho^{-3/4}\eta^{1/4}\}$ for every minimizer and every trial budget,
with a constant $C$ depending only on $\beta$. In \cref{prop:lz-stability}, we use this bound to prove the
stability of the scalar response operator at the minimizing profile.
\Cref{sec:lz-root} uses the tools of this section to localize the lowest
quantiles of the free mass of $\gamma_*$.

We use the energy normalization, the tilted joint law, and the martingale
identities of \cref{sec:fnd}, in particular \cref{prop:fnd-contact}, and
the trial budget of \cref{lem:fnd-trial}. From \cite{AL}, we import the convexity
of $F_\rho$ along segments of profiles, its first variation, the
response inequality for the second variation, and an averaged curvature bound
(\cref{prop:loc-inputs}); the regularization of \cite[Lemma~7.5]{AL}
(\cref{lem:loc-smoothing}(a)); and the linear growth of the scalar Parisi
profile near the origin \eqref{eq:lz-gsroot}. From \cite{Lopatto}, we use the
structure of the scalar Parisi measure \eqref{eq:loc-scalar-measure}. To prove
\cref{prop:loc-isolation}, we first smooth the minimizer
(\cref{lem:loc-smoothing}(b)) and then join it to the trial profile by a
segment. With cutoffs of one sign, built from one-sided averages of the superlevel
sets of the difference of the two profiles, we turn the response inequality into
lower bounds on the second derivative of $F_\rho$ along the segment
(\cref{lem:loc-coercive}). The midpoint inequality \eqref{eq:loc-midpoint} then
compares these lower bounds with the trial budget.

\Cref{prop:loc-inputs}, \cref{lem:loc-smoothing}(a), \cref{lem:loc-coercive},
the midpoint inequality \eqref{eq:loc-midpoint}, and the facts
\eqref{eq:lz-gsroot}, \eqref{eq:loc-scalar-measure}, and
\eqref{eq:loc-trialbudget} are used again in \cref{sec:lz-root}.
\Cref{sec:sc} uses \eqref{eq:lz-gsroot}, \eqref{eq:loc-scalar-measure},
\eqref{eq:loc-trialbudget}, and \cref{prop:loc-isolation}.

\section{Energy normalization and variational facts}
\label{subsec:lz-inputs}

Throughout this section and \cref{sec:lz-root}, we fix $1<\beta<\infty$,
$N\ge1$, $0<\rho\le1$, $t=1-\rho$, and $0<\lambda\le\lambda_0$, where
$\lambda_0\le\min\{1,\beta/4\}$ is the constant fixed with
\eqref{eq:fnd-scalar-trial-cost}, and $s=\lambda/\beta$. Unless
stated otherwise, constants depend only on $\beta$. In particular, they do
not depend on $N$, $\rho$, $\lambda$, the profile, or the positions and
masses of its atoms.

\paragraph{Admissible profiles and segments.}
We use the energy normalization of \cref{sec:fnd}. An \emph{admissible
profile} is a function $\gamma=\beta a$, where $a(q)=\mu([0,q])$ is the
cumulative function of an admissible profile $\mu=s\delta_0+\nu$ of
\cref{sec:fnd}. The function $\gamma$ is right-continuous and
nondecreasing with $\gamma(1)=\beta$. With the convention $\gamma(0-)=0$, its
derivative measure $\dd\gamma$ has total mass $\beta$ and satisfies
$\dd\gamma\ge\lambda\delta_0$. We call
\[
 \nu_\gamma=\dd\gamma-\lambda\delta_0=\beta\nu
\]
the \emph{free measure} of $\gamma$. It has mass $\beta-\lambda$ and may
charge $0$. The admissible profiles form a convex set. We write $\min F_\rho$
for the minimum of $F_\rho$ over this set, which is attained
(\cref{sec:fnd}). An admissible profile $\gamma_*$ is a \emph{minimizer} if
$F_\rho(\gamma_*)=\min F_\rho$. A \emph{segment} is a family
$\gamma_\theta=(1-\theta)\gamma_0+\theta\gamma_1$, $0\le\theta\le1$, with
admissible endpoints $\gamma_0,\gamma_1$. Along a segment, the free measure
$\nu_{\gamma_\theta}=(1-\theta)\nu_{\gamma_0}+\theta\nu_{\gamma_1}$ is affine
in $\theta$.

\paragraph{Self-overlap, trace, and curvature.}
For an admissible profile $\gamma$, we write $\E_\gamma$ for expectation under
the joint law of \cref{subsec:fnd-law} with profile $a=\gamma/\beta$. We
define
\[
 S_\gamma(q)=\E_\gamma S_q,\qquad
 b_\gamma(q)=\frac\beta N\E_\gamma\Tr H_q,\qquad
 \zeta_\gamma(q)=\frac{\beta^2}N\E_\gamma\Tr H_q^2 .
\]
By \cref{prop:fnd-contact}, we have $m_0=0$, $\dd m_q=\sqrt c\,H_q\dd B_q$,
and $a(q)C_q\preceq H_q\preceq C_q$, and the function $q\mapsto\E_\gamma S_q$
is continuously differentiable with derivative $(c/N)\E_\gamma\Tr H_q^2$.
Recalling that $c=\beta^2\rho$, we see that this derivative is
$\rho\zeta_\gamma(q)$. The coordinates of $m_q$ lie in $[-1,1]$, which gives
$0\le S_q\le1$ and $\Tr C_q=\sum_i(1-m_{q,i}^2)\le N$. Further,
$0\preceq H_q\preceq C_q$ implies that $0\le\Tr H_q\le\Tr C_q$. We obtain
\begin{equation}\label{eq:loc-b}
 S_\gamma(0)=0,\qquad 0\le S_\gamma\le1,\qquad
 S_\gamma'=\rho\zeta_\gamma\ge0,\qquad
 0\le b_\gamma\le\beta .
\end{equation}
Note that $S_\gamma(q)=\rho\int_0^q\zeta_\gamma(r)\dd r$, and that $S_\gamma$ is
nondecreasing. We also note that $0\preceq H_q\preceq C_q$ implies that
$\Tr H_q^2\le\Tr C_q^2\le(\Tr C_q)^2\le N^2$, and hence that
$0\le\zeta_\gamma\le\beta^2N$. At fixed $N$, the functions $S_\gamma$,
$b_\gamma$, and $\zeta_\gamma$ are bounded.

\paragraph{The rescaled functional of \cite{AL}.}
The variational facts we need are proved in
\cite[Sections~7.1 and~7.4 and Appendix~A]{AL} for the rescaled form of
$F_\rho$ recalled in \cref{sec:fnd}, and we now check this identification in
detail. Let $\gamma$ be admissible, let $U=U_{\mu,t}$ be the solution of
\eqref{eq:fnd-vector-pde} with $a=\gamma/\beta$, and let
$\widetilde U(q,y;W)=\beta^{-1}U(q,\beta y;W)$ be the rescaled solution of
\cref{sec:fnd}; when we need to indicate the profile, we write
$\widetilde U_\gamma$. Recall from \cref{sec:fnd} that $\widetilde U$ solves
$\partial_q\widetilde U=-(\rho/2)(\Delta_y\widetilde U
+\gamma(q)|\nabla_y\widetilde U|^2)$, that
$\nabla_y\widetilde U(q,y;W)=\nabla_xU(q,\beta y;W)$ and
$\nabla_y^2\widetilde U(q,y;W)=\beta\nabla_x^2U(q,\beta y;W)$, and that
$F_\rho(\gamma)$ equals
$(N\lambda)^{-1}\log\E_{\rm ref}e^{\lambda\widetilde U(0,0;W)}
-(\rho/2)\int_0^1q\gamma(q)\dd q$. By \eqref{eq:fnd-vector-pde}, the terminal
value of $\widetilde U$ is
\[
 \widetilde U(1,y;W)=\frac1\beta\log\sum_{\sigma\in\{-1,1\}^N}
 \exp\Bigl\{\beta\Bigl(\sqrt t\,\frac{\sigma^TW\sigma}2+y\cdot\sigma\Bigr)\Bigr\}.
\]
Let $W^{\rm off}$ be the matrix $W$ with its diagonal set to zero. For
$\sigma\in\{-1,1\}^N$, we have $\sigma^TW\sigma/2=\sigma^TW^{\rm off}\sigma/2+D_N$,
where $D_N=\Tr W/2$. Since the equation for $\widetilde U$ is invariant under the addition of
constants, we find that
$\widetilde U(q,y;W)=\sqrt t\,D_N+\widetilde U^{\rm off}(q,y)$, where
$\widetilde U^{\rm off}$ is the solution of the same equation with $W^{\rm
off}$ in place of $W$ in the terminal value. This function is the solution $u$
of \cite[(7.2)]{AL} with $r=\rho$ and coupling matrix $\sqrt t\,W^{\rm off}$.
The variable $D_N$ is centered Gaussian with variance $1/2$ and independent of
$W^{\rm off}$, so that
$(N\lambda)^{-1}\log\E_{\rm ref}e^{\lambda\sqrt t D_N}=\lambda t/(4N)$.
Inserting this into the representation of $F_\rho$ recalled above, we find
\[
 F_\rho(\gamma)=\frac1{N\lambda}\log\E_{\rm ref}e^{\lambda\widetilde U^{\rm off}(0,0)}
 -\frac\rho2\int_0^1q\gamma(q)\dd q+\frac{\lambda t}{4N}.
\]
This is exactly the functional $F^c_t(\gamma)$ of \cite[(7.2)]{AL}, which is
written $F_r$ in \cite[Section~7.4]{AL}; the term $\lambda t/(4N)$ is part of
that functional and does not depend on $\gamma$.

The joint laws also agree. The tilt $e^{sU(0,0;W)}=e^{\lambda\sqrt t
D_N}e^{\lambda\widetilde U^{\rm off}(0,0)}$ of \cref{subsec:fnd-law}
factorizes, and $D_N$ enters neither the field equation
\eqref{eq:fnd-diffusion}, whose drift involves $\nabla_xU$, nor the
conditional law of the terminal spin given $W$ and the field. Under the tilted
law, the couplings $\sqrt t\,W^{\rm off}$ have the law $Q_\gamma$ of
\cite[Appendix~A]{AL}, and the rescaled field $X_q/\beta$ solves
\[
 \dd(X_q/\beta)=\rho\gamma(q)\nabla_y\widetilde U(q,X_q/\beta;W)\dd q
 +\sqrt\rho\dd B_q,\qquad X_0/\beta=0,
\]
which is the field equation of \cite[Section~7.1]{AL}. Under this law, the
magnetization $m_q=\nabla_xU(q,X_q;W)=\nabla_y\widetilde U(q,X_q/\beta;W)$ is
unchanged, while the Hessian in the rescaled variable is
$\nabla_y^2\widetilde U(q,X_q/\beta;W)=\beta H_q$. We conclude that $S_\gamma$,
$b_\gamma$, and $\zeta_\gamma$ are the functions $S$, $b$, and $\zeta$ of
\cite[Sections~7.1 and~7.4]{AL}, and that $b_\gamma$ and $\zeta_\gamma$ are the
functions of the same names in \cite[Appendix~A]{AL}. We use the results
proved there in the following form.

\begin{proposition}[Variational and response facts]
\label{prop:loc-inputs}
Let $\gamma$ be admissible and define
$G_\gamma(x)=\int_x^1(S_\gamma(q)-q)\dd q$ for $0\le x\le1$.
\begin{enumerate}
\item[(a)] Let $\nu'\ge0$ be a measure on $[0,1]$ of mass $\beta-\lambda$.
Then $F_\rho$ is convex along the segment from $\gamma$ to the profile with
free measure $\nu'$, and its one-sided derivative at $\gamma$ is
\begin{equation}\label{eq:loc-firstvariation}
 D_+F_\rho(\gamma)[\nu'-\nu_\gamma]
 =\frac\rho2\int G_\gamma\,\dd(\nu'-\nu_\gamma).
\end{equation}
The function $G_\gamma$ is continuously differentiable with
$G_\gamma'(x)=x-S_\gamma(x)$. In particular, $G_\gamma'(0)=0$, and
$x\mapsto G_\gamma'(x)-x$ is nonincreasing.
\item[(b)] Let $(\gamma_\theta)$ be a segment, $\psi=\gamma_1-\gamma_0$, and
$f(\theta)=F_\rho(\gamma_\theta)$. Then $f$ is convex and continuous on
$[0,1]$ and twice differentiable on $(0,1)$. Fix $\theta\in(0,1)$ and write
$b_\theta=b_{\gamma_\theta}$ and $\zeta_\theta=\zeta_{\gamma_\theta}$. For every
compactly supported Lipschitz function $\varphi$ on $(0,1)$, let
\[
 Y=\int_0^1 b_\theta|\varphi'|^2\dd q,\qquad
 Z=\rho\int_{(0,1)}\varphi^2\zeta_\theta\dd\gamma_\theta .
\]
Then
\begin{equation}\label{eq:loc-response}
 \rho\left|\int_0^1\psi\zeta_\theta\varphi\dd q\right|
 \le\Bigl(\frac{f''(\theta)}\rho\Bigr)^{1/2}
 \bigl[Y^{1/2}+2(Y+Z)^{1/2}\bigr].
\end{equation}
\item[(c)] There exists $k_\beta>0$, depending only on $\beta$, such that
\begin{equation}\label{eq:loc-in-curv}
 \zeta_\gamma(q)\ge k_\beta^2\qquad(0\le q<1).
\end{equation}
\end{enumerate}
\end{proposition}

The integral $Z$ is a Stieltjes integral over the open interval $(0,1)$. It
includes the atoms of $\dd\gamma_\theta$ in $(0,1)$ and excludes the
prescribed atom at $0$.

\begin{proof}
We prove the three parts in order. By the discussion preceding the
proposition, $F_\rho$ is exactly the functional $F_r$ of \cite[Section~7.4]{AL}
with $r=\rho$, and $S_\gamma,b_\gamma,
\zeta_\gamma$ are the response functions used there.

\emph{Part (a).} The segment from $\gamma$ to the profile with free measure
$\nu'$ replaces $\nu_\gamma$ by $(1-\theta)\nu_\gamma+\theta\nu'$. The
convexity of $F_r$ along every such segment and the first variation
\eqref{eq:loc-firstvariation} are stated in \cite[Section~7.4, (7.22)]{AL},
with $G_\gamma$ the function defined there. The identity
\eqref{eq:loc-firstvariation} is also \eqref{eq:fnd-first-variation} divided
by $\beta$: the free measures satisfy $\nu_\gamma=\beta\nu$, which gives
$\beta^{-1}(c/2)\int G\,\dd(\nu'/\beta-\nu)=(\rho/2)\int
G\,\dd(\nu'-\nu_\gamma)$, and the function $G$ of
\eqref{eq:fnd-first-variation} is $G_\gamma$. By \eqref{eq:loc-b},
$S_\gamma$ is continuous, nondecreasing, and vanishes at $0$. By the
definition of $G_\gamma$, we then have $G_\gamma'(x)=x-S_\gamma(x)$, which is
continuous. We also have $G_\gamma'(0)=0$, and $G_\gamma'(x)-x=-S_\gamma(x)$ is
nonincreasing. This
proves (a); see also \cite[(7.23)]{AL}.

\emph{Part (b).} Applying (a) with $\gamma=\gamma_0$ and $\nu'=\nu_{\gamma_1}$, we see
that $f$ is convex. The profile-continuity estimate
of \cite[Section~7.4]{AL} states that
\[
 \sup_{q,y}|\widetilde U^{\rm off}_\gamma(q,y)-\widetilde U^{\rm off}_{\gamma'}(q,y)|
 \le\frac{3\rho N}2\|\gamma-\gamma'\|_{L^1(0,1)}
\]
for admissible $\gamma,\gamma'$ and every value of the couplings. Inserting
this bound into the representation of $F_\rho$ above, we find
\[
 |F_\rho(\gamma)-F_\rho(\gamma')|\le2\rho\|\gamma-\gamma'\|_{L^1(0,1)},
\]
where the outer logarithmic moment contributes $3\rho/2$ and the penalty
contributes $\rho/2$. Then $f$ is Lipschitz on $[0,1]$, because
$\|\gamma_\theta-\gamma_{\theta'}\|_{L^1}=|\theta-\theta'|\,\|\gamma_1-\gamma_0\|_{L^1}$.

We now apply \cite[Proposition~A.1 and Lemma~A.2]{AL} with $r=\rho$, the
endpoints $\gamma_0,\gamma_1$, and $v=\psi$. Their hypotheses hold: $N$ is
finite, $0\le t<1$, the endpoints are profiles of \cite[Appendix~A]{AL}
($\lambda\le\gamma_i\le\beta$ on $[0,1)$, $\gamma_i(0-)=0$, and
$\gamma_i(1)=\beta$) with the same outer tilt $\lambda$, and
\cite[Proposition~A.1]{AL} allows their additional masses at $0$ to differ.
By \cite[Proposition~A.1, (A.1)]{AL}, $f$ is twice differentiable, and
$f''(\theta)$ is the sum of $\rho\int_0^1e_\theta(q)\dd q$ and a nonnegative
variance term, where $e_\theta\ge0$ is the function defined there. In
particular,
\[
 \int_0^1e_\theta(q)\dd q\le\frac{f''(\theta)}\rho .
\]
The pairing estimate \cite[(A.2)]{AL}, applied to $\varphi$ and to $-\varphi$,
and the Fisher estimate \cite[(A.3)]{AL} give a nonnegative function
$\mathcal J_\theta$, also defined in \cite[Proposition~A.1]{AL}, such that
\[
 \rho\left|\int_0^1\psi\zeta_\theta\varphi\dd q\right|
 \le\Bigl(\int_0^1e_\theta\Bigr)^{1/2}
 \Bigl[Y^{1/2}+2\Bigl(\int_0^1\mathcal J_\theta\varphi^2\Bigr)^{1/2}\Bigr],
 \qquad
 \int_0^1\mathcal J_\theta\varphi^2\le Y+Z .
\]
Combining the last three displays, we obtain \eqref{eq:loc-response}.

\emph{Part (c).} We use \cite[Lemma~A.3, (A.5)]{AL} with the level in
(A.5) equal to $\beta$; the hypothesis $\beta\ge1$ of that lemma holds. As
$\max\{1,\beta\}=\beta$, the lemma states that $\gamma(q)\le\beta$ and $q<1$ imply
$\E_\gamma[\beta H_q]\succeq k_\beta I$, where
\[
 k_\beta=\frac{e^{-16c'\beta-2\beta^2}}{2(1+c')},\qquad c'=\sqrt{\pi/2} .
\]
Every admissible profile satisfies $\gamma\le\beta$. By Jensen's inequality,
applied to the convex function $H\mapsto\Tr H^2$, we have
$\zeta_\gamma(q)=N^{-1}\E_\gamma\Tr(\beta H_q)^2\ge
N^{-1}\Tr(\E_\gamma\beta H_q)^2\ge k_\beta^2$, which is
\eqref{eq:loc-in-curv}.

In \cite{AL}, these results are first proved for step profiles, and then
extended to all admissible profiles by approximation at fixed $N$ and $\beta$
\cite[Section~7.4 and Appendix~A.3]{AL}. This includes profiles with atoms,
and the Stieltjes integral in $Z$ includes the atoms of $\dd\gamma_\theta$ in
$(0,1)$.
\end{proof}

\begin{remark}[Atoms and the Stieltjes term]\label{rem:loc-stieltjes}
The term $Z$ is where atoms of the profile enter \eqref{eq:loc-response}. In
the proof of \cite[Lemma~A.2]{AL}, the Fisher estimate (A.3) follows from the
identity of measures $\dd\mathscr P=2\mathcal J_\theta\dd q
-\rho\zeta_\theta\dd\gamma_\theta$ on $(0,1)$ \cite[(A.19)]{AL}, where
$\mathscr P(q)=N^{-1}\E_\theta\Tr\dot h$ and $\dot h$ is the derivative in
$q$, at fixed magnetization, of the Hessian of $\widetilde U_{\gamma_\theta}$ written as
a function of the magnetization. At an atom of $\dd\gamma_\theta$ at $q$, the
function $\mathscr P$ jumps by $-\rho\zeta_\theta(q)\,\dd\gamma_\theta(\{q\})$.
In \cref{prop:loc-isolation,lem:lz-certificate}, we choose the cutoffs so that
$\dd\gamma_\theta$ has a bounded density on the set where they are positive.
\end{remark}

\subsection{The scalar reference profile}
Let $\mu_\beta$ be the scalar Parisi measure, $\alpha_\beta=\mu_\beta([0,\cdot])$,
and $\gamma_\beta=\beta\alpha_\beta$, as in \cref{sec:fnd}, and let
$\gamma^{\rm tr}=\max\{\lambda,\gamma_\beta\}$ be the scalar trial profile; we
omit the subscript $\lambda$ of \cref{sec:fnd}. The function $\gamma_\beta$ is
the unconstrained minimizer of the rescaled scalar Parisi functional
$\gamma\mapsto\beta^{-1}\mathcal P_\beta(\gamma/\beta)$ of
\cite[Section~7.1]{AL}. By \cite[Lemma~7.2]{AL}, applied with $\beta_0=\beta$,
there exist constants $T_*,k_0>0$ and $C_0<\infty$, depending only on $\beta$,
such that $\gamma_\beta(0)=0$, $\gamma_\beta$ is absolutely continuous on
$[0,T_*]$, and
\[
 k_0q\le\gamma_\beta(q)\le C_0q,\qquad
 0\le\gamma_\beta'(q)\le C_0\qquad(0\le q\le T_*),
\]
the second bound for almost every $q$. We set $T_0=\min\{T_*,1/4,2/C_0\}$.
On $[0,T_0]$, we have $\gamma_\beta\le C_0T_0\le2$, and since $\lambda\le1$,
also $\gamma^{\rm tr}\le2$. The maximum of the constant $\lambda$ and the
nondecreasing absolutely continuous function $\gamma_\beta$ is absolutely
continuous on $[0,T_0]$, with derivative
$\gamma_\beta'\ind_{\{\gamma_\beta>\lambda\}}\in[0,C_0]$ almost everywhere. We
obtain
\begin{equation}\label{eq:lz-gsroot}
 \gamma_\beta(0)=0,\qquad k_0q\le\gamma_\beta(q),\qquad
 \gamma^{\rm tr}(q)\le2\qquad(0\le q\le T_0),
\end{equation}
and the restriction of $\dd\gamma^{\rm tr}$ to $(0,T_0]$ has density at most
$C_0$. At $0$, the measure $\dd\gamma^{\rm tr}$ has the prescribed atom of
mass $\gamma^{\rm tr}(0)=\lambda$.

By \cite[Theorem~1.1, (1.4) and (1.5)]{Lopatto}, there exist
$q_\beta\in(0,1)$, $c_\beta\in(0,1)$, and a nonnegative function
$\varrho_\beta\in C^\infty([0,q_\beta])$ such that
\begin{equation}\label{eq:loc-scalar-measure}
 \mu_\beta(\dd q)=\varrho_\beta(q)\ind_{[0,q_\beta)}(q)\dd q
                  +c_\beta\delta_{q_\beta}(\dd q),\qquad
 \supp\mu_\beta=[0,q_\beta].
\end{equation}
Here, as in \cite{Lopatto}, $\varrho_\beta\in C^\infty([0,q_\beta])$ means
that $\varrho_\beta$ is smooth on $(0,q_\beta)$ and that each of its
derivatives extends continuously to $[0,q_\beta]$. The only atom of
$\mu_\beta$ is at $q_\beta$, and $\alpha_\beta=1$ and $\gamma_\beta=\beta$ on
$[q_\beta,1]$. (The Parisi functional \cite[(1.3)]{Lopatto} is written with the
terminal value $\log\cosh x$ and the additional constant $\log2$; it equals the
functional in \eqref{eq:intro-parisi}, and $\mu_\beta$ is its minimizer.)

\subsection{Trial budgets}
For a minimizer $\gamma_*$ of $F_\rho$, a \emph{trial budget} is a number
$\eta\ge0$ such that
\begin{equation}\label{eq:loc-trialbudget}
 F_\rho(\gamma^{\rm tr})-F_\rho(\gamma_*)\le\eta .
\end{equation}
Since $F_\rho(\gamma_*)=\min F_\rho$, this is the notion of trial budget of
\cref{sec:fnd}, measured in energy units. By \eqref{eq:fnd-trial},
$\eta=e_{N,\beta}(\lambda)+L_s\lambda^5$ is a trial budget for every
minimizer at every $\rho\in(0,1]$. The estimates of this section and the next
hold for every trial budget.

\section{Smoothing and coercivity}

A minimizer may have atoms, and the response inequality
\eqref{eq:loc-response} involves the derivative measure $\dd\gamma_\theta$ on
the set where the cutoff is positive. We therefore first spread the free mass,
at a cost that is quadratic in the smoothing width (\cref{lem:loc-smoothing}).
Part (a) of the next lemma is \cite[Lemma~7.5]{AL}, and part (b) extends the
reflection argument of that lemma to the upper endpoint $1$. We then show
that a difference of one sign between the endpoints of a segment forces
curvature along the segment (\cref{lem:loc-coercive}).

\begin{lemma}[Reflected regularization]\label{lem:loc-smoothing}
Let $\gamma$ be admissible and $h>0$.
\begin{enumerate}
\item[(a)] Let $0<T<1$, $h\le1-T$, and $m=\nu_\gamma([0,T))$. Let
$\bar\gamma$ be the profile with $\dd\bar\gamma=\lambda\delta_0+\bar\nu$, where
$\bar\nu$ is obtained from $\nu_\gamma$ by replacing each unit of mass at a
point $q<T$ by the image of the uniform distribution on $[q-h,q+h]$ under
$x\mapsto|x|$. Then $\bar\gamma$ is
admissible, $\bar\gamma(0)=\lambda$, the restriction of $\dd\bar\gamma$ to
$(0,T)$ has density at most $m/h$, and
$F_\rho(\bar\gamma)-F_\rho(\gamma)\le\rho mh^2/12$.
\item[(b)] Let $h\le1/2$, and let $\pi(x)=|x|$ for $x\le1$ and $\pi(x)=2-x$
for $x>1$. Let $\gamma^h$ be the profile with
$\dd\gamma^h=\lambda\delta_0+\nu^h$, where $\nu^h$ is obtained from $\nu_\gamma$
by replacing each unit of mass at a point $q\in[0,1]$ by the image of the
uniform distribution on $[q-h,q+h]$ under $\pi$. Then $\gamma^h$ is admissible, $\gamma^h(0)=\lambda$, its
free measure has density at most $\beta/h$ on $[0,1]$, and
\[
 \|\gamma^h-\gamma\|_{L^1(0,1)}\le\beta h,\qquad
 F_\rho(\gamma^h)-F_\rho(\gamma)\le\frac{\rho\beta h^2}{12} .
\]
\end{enumerate}
\end{lemma}

\begin{proof}
Part (a) is \cite[Lemma~7.5]{AL}, including the cost bound
\cite[(7.24)]{AL}; by \cref{subsec:lz-inputs}, the functional $F_r$ of that
lemma is exactly $F_\rho$ with $r=\rho$. Its proof is the
argument for (b) below, with the following changes. Only the free mass on
$[0,T)$ is moved, and the reflection is $x\mapsto|x|$. The condition
$h\le1-T$ ensures that the moved mass stays in $[0,T+h)\subset[0,1]$. The semiconcavity bound
is then needed only on $[-h,T+h]\subset[-1,1]$, where the reflection at $1$
plays no role. The cost bound becomes $(\rho/2)\cdot m\cdot(h^2/6)=\rho
mh^2/12$, and the moved mass
has density at most $m/h$ on $(0,T)$, while the mass on $[T,1]$ does not
charge $(0,T)$.

For (b), we first verify the properties of $\gamma^h$, then the displacement
bound, and finally the cost bound. For $q\in[0,1]$, let $\kappa_q$ be the
image of the uniform distribution on $[q-h,q+h]$ under $\pi$. The free
measure of $\gamma^h$ is then $\nu^h=\int\kappa_q\,\nu_\gamma(\dd q)$. The map
$\pi$ sends $[-h,1+h]$ into $[0,1]$ and is piecewise affine with slopes
$\pm1$. For $x\in(0,1)$, the points of $[q-h,q+h]$ that $\pi$ sends to $x$ are
among $x$, $-x$, and $2-x$. We find that $\kappa_q$ has the density
\[
 \frac1{2h}\bigl(\ind_{\{|x-q|\le h\}}+\ind_{\{x+q\le h\}}
 +\ind_{\{x+q\ge2-h\}}\bigr),\qquad x\in(0,1).
\]
Since $h\le1/2<2-h$, the last two indicators cannot both equal one, and this
density is at most $1/h$. Thus $\kappa_q$ is a probability measure on $[0,1]$
without atoms. Integrating over $\nu_\gamma$, we see that $\nu^h$ is a measure
on $[0,1]$ of mass $\beta-\lambda$, without atoms, with density at most
$(\beta-\lambda)/h\le\beta/h$. The measure $\dd\gamma^h=\lambda\delta_0+\nu^h$
has mass $\beta$, we have $\gamma^h(0)=\lambda$, and $\gamma^h$ is
admissible.

For the displacement bound, let $\Pi$ be the measure on $[0,1]^2$ defined by
\[
 \int F\dd\Pi=\int_{[0,1]}\frac1{2h}\int_{-h}^hF\bigl(q,\pi(q+z)\bigr)\dd z\,
 \nu_\gamma(\dd q)
\]
for bounded Borel functions $F$. Its marginals are $\nu_\gamma$ and $\nu^h$.
The map $\pi$ is $1$-Lipschitz and $\pi(q)=q$ for $q\in[0,1]$. It follows
that $|\pi(q+z)-q|\le|z|\le h$, and $\Pi$ is supported on pairs $(q,r)$ with
$|r-q|\le h$. The prescribed atoms of $\dd\gamma$ and $\dd\gamma^h$ coincide,
and for $x\in[0,1]$ we have
\[
 \gamma^h(x)-\gamma(x)=\int\bigl(\ind_{\{r\le x\}}-\ind_{\{q\le x\}}\bigr)
 \Pi(\dd q,\dd r).
\]
For $q,r\in[0,1]$, we have $\int_0^1|\ind_{\{r\le x\}}-\ind_{\{q\le x\}}|\dd
x=|r-q|$. By Tonelli's theorem, we conclude that
\[
 \|\gamma^h-\gamma\|_{L^1(0,1)}\le\int|r-q|\,\Pi(\dd q,\dd r)
 \le(\beta-\lambda)h\le\beta h .
\]

To bound the cost, let $G=G_{\gamma^h}$ and $S=S_{\gamma^h}$, and define
$H(z)=G(\pi(z))-z^2/2$ for $z\in[-h,1+h]$. By \cref{prop:loc-inputs}(a), $G$
is continuously differentiable on $[0,1]$ with $G'(x)=x-S(x)$. Using
$\pi(z)=-z$ for $z<0$ and $\pi(z)=2-z$ for $z>1$, we find that $H$ is
continuous and continuously differentiable away from $0$ and $1$, with
\[
 H'(z)=
 \begin{cases}
 S(-z),&-h<z<0,\\
 -S(z),&0<z<1,\\
 S(2-z)-2,&1<z<1+h.
 \end{cases}
\]
Each piece is nonincreasing, because $S$ is nondecreasing by
\eqref{eq:loc-b}. The one-sided derivatives at $0$ agree, because $S(0)=0$. At
$1$, the derivative jumps from $-S(1)$ to $S(1)-2\le-S(1)$, where we used that
$S(1)\le1$. We see that $H'$ is nonincreasing, and $H$ is concave on
$[-h,1+h]$. For
$q\in[0,1]$, the number $G'(q)-q$ is a supergradient of $H$ at $q$; at $q=1$,
we take the left derivative $G'(1)-1=-S(1)$. By the concavity of $H$, we have
$H(q+z)\le H(q)+(G'(q)-q)z$, that is,
\[
 G\bigl(\pi(q+z)\bigr)\le G(q)+G'(q)z+\frac{z^2}2
 \qquad(q\in[0,1],\ |z|\le h).
\]
We average over $z$ in $[-h,h]$. Using that the average of $z$ is $0$ and
that the average of $z^2$ is $h^2/3$, we obtain
\begin{equation}\label{eq:loc-semiconcave}
 \int G\dd\kappa_q-G(q)\le\frac{h^2}6\qquad(q\in[0,1]).
\end{equation}
By \cref{prop:loc-inputs}(a), $F_\rho$ is convex along the segment from
$\gamma^h$ to $\gamma$, which implies that $F_\rho(\gamma)\ge
F_\rho(\gamma^h)+D_+F_\rho(\gamma^h)[\nu_\gamma-\nu^h]$. Using
\eqref{eq:loc-firstvariation} at $\gamma^h$ and \eqref{eq:loc-semiconcave}, we
conclude that
\[
 F_\rho(\gamma^h)-F_\rho(\gamma)
 \le\frac\rho2\int G\,\dd(\nu^h-\nu_\gamma)
 =\frac\rho2\int\Bigl[\int G\dd\kappa_q-G(q)\Bigr]\nu_\gamma(\dd q)
 \le\frac{\rho(\beta-\lambda)h^2}{12}. \qedhere
\]
\end{proof}

\paragraph{The midpoint inequality.}
We compare the second derivative along a segment with the values at its
endpoints. Let $f$ be convex and continuous on $[0,1]$ and twice
differentiable on $(0,1)$. Then $f'$ is nondecreasing on $(0,1)$, and $f$ is
absolutely continuous on $[0,1]$: on each interval $[\epsilon,1-\epsilon]$ it
is Lipschitz, and since $f'$ is monotone and $f$ is continuous at $0$ and $1$,
the monotone convergence theorem shows that $f(y)-f(x)=\int_x^yf'$ for all
$0\le x<y\le1$. Further, a nondecreasing function $g$ on $(0,1)$ satisfies
$g(y)-g(x)\ge\int_x^yg'$ for $0<x<y<1$, by Lebesgue's theorem on the
differentiation of monotone functions; we apply this with $g=f'$, whose
derivative $f''\ge0$ exists everywhere on $(0,1)$. We have
\[
 f(0)+f(1)-2f(\tfrac12)
 =\int_0^{1/2}\bigl(f'(x+\tfrac12)-f'(x)\bigr)\dd x
 \ge\int_0^{1/2}\!\!\int_x^{x+1/2}f''(\theta)\dd\theta\dd x .
\]
For fixed $\theta\in(0,1)$, the set of $x\in(0,1/2)$ with $x<\theta<x+1/2$
has length $\min\{1/2,\theta\}-\max\{0,\theta-1/2\}=\min\{\theta,1-\theta\}$.
By Tonelli's theorem, the right side equals
$\int_0^1\min\{\theta,1-\theta\}f''(\theta)\dd\theta$, and since $f''\ge0$,
we may restrict this integral to any Borel subset of $(0,1)$.

Now let $\gamma_*$ be a minimizer, let $(\gamma_\theta)$ be a segment, and let
$f(\theta)=F_\rho(\gamma_\theta)$. By \cref{prop:loc-inputs}(b), $f$ has the
properties just used. The midpoint $\gamma_{1/2}$ is admissible, which implies that
$f(1/2)\ge F_\rho(\gamma_*)$. Inserting this into the previous inequality, we
obtain
\begin{equation}\label{eq:loc-midpoint}
 \int_\Theta\min\{\theta,1-\theta\}f''(\theta)\dd\theta
 \le F_\rho(\gamma_0)+F_\rho(\gamma_1)-2F_\rho(\gamma_*)
 \qquad(\Theta\subset(0,1)\ \text{Borel}).
\end{equation}
The weight $\min\{\theta,1-\theta\}$ has integral $1/4$ over $(0,1)$ and
integral $2\int_{1/4}^{1/2}\theta\dd\theta=3/16$ over $[1/4,3/4]$.

\begin{lemma}[Coercivity from a difference of one sign]
\label{lem:loc-coercive}
Let $(\gamma_\theta)$ be a segment, $\psi=\gamma_1-\gamma_0$, and
$f(\theta)=F_\rho(\gamma_\theta)$, and fix $\theta\in(0,1)$. Let
$\varphi:[0,1]\to[0,1]$ be Lipschitz, with $\varphi(0)=\varphi(1)=0$, and not
identically zero, and let
\[
 \Delta=\rho\int_0^1\zeta_\theta\varphi\dd q,\qquad
 Y_0=\int_0^1|\varphi'|^2\dd q .
\]
Suppose that there exist $A>0$, $d\ge0$, and $\sigma\in\{-1,1\}$ such that
$\sigma\psi\ge A/2$ wherever $\varphi>0$, and
\begin{equation}\label{eq:loc-weighted-density}
 \rho\int_{(0,1)}\varphi^2\zeta_\theta\dd\gamma_\theta\le d\Delta.
\end{equation}
The last condition holds if the restriction of $\dd\gamma_\theta$ to the set
$\{\varphi>0\}$ has density at most $d$. Then
\begin{equation}\label{eq:loc-coercive}
 f''(\theta)\ge\rho\,\frac{A^2}{72}
 \min\left\{\frac{\Delta^2}{\beta Y_0},\ \frac\Delta d\right\},
\end{equation}
where $\Delta/0=+\infty$ if $\Delta>0$, and the right side is $0$ if
$\Delta=0$.
\end{lemma}

For fixed $Y_0$ and $d$, both entries of the minimum in
\eqref{eq:loc-coercive} are nondecreasing in $\Delta$, so we may insert any
lower bound for $\Delta$ into the right side.

\begin{proof}
We first check the sufficient condition for
\eqref{eq:loc-weighted-density}. If $\dd\gamma_\theta\le d\dd q$ on the open
set $\{\varphi>0\}$, then, since $0\le\varphi\le1$ and $\zeta_\theta\ge0$, we
have $\rho\int_{(0,1)}\varphi^2\zeta_\theta\dd\gamma_\theta\le
d\rho\int_0^1\varphi^2\zeta_\theta\dd q\le d\Delta$.

Next, we extend \eqref{eq:loc-response} to $\varphi$, which need not have
compact support in $(0,1)$. For $0<\epsilon<1/4$, let $\chi_\epsilon$ be the
piecewise linear function that vanishes within distance $\epsilon$ of
$\{0,1\}$, equals one at distance at least $2\epsilon$, and is linear in
between; note that $|\chi_\epsilon'|\le1/\epsilon$. Let $\ell_\varphi$ be the
Lipschitz constant of $\varphi$ and $\varphi_\epsilon=\chi_\epsilon\varphi$.
Then $\varphi_\epsilon$ is Lipschitz with compact support in $(0,1)$. Since
$\varphi(0)=\varphi(1)=0$, we have
$|\varphi(q)|\le\ell_\varphi\min\{q,1-q\}$, which is at most
$2\ell_\varphi\epsilon$ on the two strips where $\chi_\epsilon\ne1$; these
strips have total length $4\epsilon$. Then $\|\varphi_\epsilon-\varphi\|_\infty\le
2\ell_\varphi\epsilon$, and since
$\varphi_\epsilon'-\varphi'=\chi_\epsilon'\varphi+(\chi_\epsilon-1)\varphi'$
vanishes outside the strips and is bounded by $3\ell_\varphi$ on them, we also
have $\|\varphi_\epsilon'-\varphi'\|_{L^2}^2\le36\ell_\varphi^2\epsilon$. At
fixed $N$ and $\theta$, the functions $\psi$, $b_\theta$, and $\zeta_\theta$
are bounded by $\beta$, $\beta$, and $\beta^2N$, respectively: the first
because both profiles take values in $[\lambda,\beta]$, the second by
\eqref{eq:loc-b}, and the third because $0\preceq H_q\preceq C_q$ and
$\Tr C_q\le N$. Further,
$\dd\gamma_\theta$ is a finite measure. Using the uniform convergence of
$\varphi_\epsilon$ and the $L^2$ convergence of $\varphi_\epsilon'$, we find
that the left side of \eqref{eq:loc-response} and $Y$ converge, as
$\epsilon\to0$, to their values for $\varphi$. The term $Z$ converges as well, because
$\varphi_\epsilon^2\to\varphi^2$ uniformly and the Stieltjes integral is taken
against a finite measure; this includes the atoms of $\dd\gamma_\theta$.
Passing to the limit, we see that \eqref{eq:loc-response} holds for
$\varphi$.

If $\Delta=0$, the claim follows from the convexity of $f$. Otherwise, we have
$Y_0>0$, because $\varphi$ vanishes at $0$ and is not identically zero. In the
notation of \cref{prop:loc-inputs}(b), we have $Y\le\beta Y_0$ by
\eqref{eq:loc-b} and $Z\le d\Delta$ by \eqref{eq:loc-weighted-density}.
As $\zeta_\theta$ and $\varphi$ are nonnegative, the sign condition implies
that
$\sigma\rho\int_0^1\psi\zeta_\theta\varphi\dd q\ge A\Delta/2$. Inserting this
bound into \eqref{eq:loc-response}, we find
\[
 \frac A2\Delta\le\Bigl(\frac{f''(\theta)}\rho\Bigr)^{1/2}
 \bigl[Y^{1/2}+2(Y+Z)^{1/2}\bigr]
 \le3\Bigl(\frac{f''(\theta)}\rho\Bigr)^{1/2}(Y+Z)^{1/2}.
\]
The left side is positive, which implies that $Y+Z>0$ and
$f''(\theta)\ge\rho A^2\Delta^2/(36(Y+Z))$. Finally, we have
$Y+Z\le\beta Y_0+d\Delta\le2\max\{\beta Y_0,d\Delta\}$, which gives
\eqref{eq:loc-coercive}.
\end{proof}

\section{A density bound for the scalar measure}

By \eqref{eq:loc-scalar-measure}, the density $\varrho_\beta$ of $\mu_\beta$
on $[0,q_\beta)$ extends continuously to $[0,q_\beta]$, and in particular it
is bounded. The isolation argument uses this bound in the form of
\cref{lem:loc-lipschitz}. We include a direct derivation of the lemma from the
identity
that expresses $\alpha_\beta$ on the support in terms of spatial derivatives
of the Parisi solution. The same identity underlies the proof of
\cite[Theorem~2(ii)]{AuffingerChen2015Properties}. For the SK model, this
identity is \cite[Lemma~6.2, (6.1)]{Lopatto}, and the smoothness of the
resulting quotient up to $q_\beta$ is \cite[Lemma~6.3]{Lopatto}. In the physics
literature, Crisanti and Rizzo gave the corresponding identities in their
variables \cite[Eqs.~(22)--(23)]{CrisantiRizzo2002}.

\begin{lemma}[Density bound]\label{lem:loc-lipschitz}
There exists $L_\beta<\infty$ such that
\begin{equation}\label{eq:loc-lipschitz}
 \alpha_\beta(y)-\alpha_\beta(x)\le L_\beta(y-x)
 \qquad(0\le x<y<q_\beta).
\end{equation}
With $L=\max\{1,\beta L_\beta\}$, the trial profile satisfies
$\gamma^{\rm tr}(0)=\lambda$, $\gamma^{\rm tr}=\beta$ on $[q_\beta,1]$, and
$|\gamma^{\rm tr}(x)-\gamma^{\rm tr}(y)|\le L|x-y|$ for $x,y\in[0,q_\beta)$. In
particular, the restriction of $\dd\gamma^{\rm tr}$ to
$(0,1]\setminus\{q_\beta\}$ is absolutely continuous, with density at most $L$
on $(0,q_\beta)$ and $0$ on $(q_\beta,1]$.
\end{lemma}

The bound \eqref{eq:loc-lipschitz} is stated for $y<q_\beta$ because of the
atom of $\mu_\beta$ at $q_\beta$. The continuous function $\bar\alpha$ in the
proof below equals $\alpha_\beta$ on $[0,q_\beta)$. At $q_\beta$, it equals
the left limit $\mu_\beta([0,q_\beta))=1-c_\beta$, while
$\alpha_\beta(q_\beta)=1$.

\begin{proof}
We write $\alpha_\beta$ as a quotient of two moments along the scalar
diffusion, show that the two moments are Lipschitz and that the denominator is
bounded below, and then deduce the consequences for $\gamma^{\rm tr}$.

Let $u$ be the solution of the scalar Parisi equation in \eqref{eq:intro-parisi}
with $\alpha=\alpha_\beta$. The solution in \cite{Lopatto} has the terminal
value $\log\cosh x$ instead of $\log(2\cosh x)$; since the equation is
invariant under the addition of constants, the two solutions differ by
$\log2$, and their spatial derivatives coincide. Let $w$ be a standard
one-dimensional Brownian motion, and let $\xi$ solve
\[
 \dd\xi_r=\beta^2\alpha_\beta(r)u_x(r,\xi_r)\dd r+\beta\dd w_r,
 \qquad\xi_0=0,
\]
which is the optimal diffusion \cite[(2.4)]{Lopatto}. We write
$\hat C_r=u_{xx}(r,\xi_r)$, $\hat D_r=u_{xxx}(r,\xi_r)$, and
$\hat A_r=u_{xxxx}(r,\xi_r)$, and define
\[
 \Gamma(r)=\E u_x(r,\xi_r)^2,\qquad \Lambda(r)=\E\hat D_r^2,\qquad
 \Xi(r)=\E\hat C_r^3 \qquad(0\le r\le1).
\]
We use three facts from \cite[Section~2]{Lopatto}.
\begin{enumerate}
\item[(i)] \cite[Proposition~2.1]{Lopatto}, which combines
\cite[Propositions~1, 2, and~5]{AuffingerChen2015Properties}: every spatial
derivative of $u$ of positive order is bounded on $[0,1]\times\R$ by a
constant depending only on $\beta$ and the order, $|u_x|\le1$, and
$0<u_{xx}\le1$.
\item[(ii)] \cite[Lemma~2.6, (2.19)]{Lopatto}: $\Gamma(r)=r$ for every
$r\in\supp\mu_\beta=[0,q_\beta]$.
\item[(iii)] \cite[Lemma~2.5, (2.13) and (2.14)]{Lopatto}: the function
$\Gamma$ is continuously differentiable with $\Gamma'=\beta^2\E\hat C^2$, and
\[
 \Gamma'(y)-\Gamma'(x)=\beta^4\int_x^y
 \bigl(\Lambda(r)-2\alpha_\beta(r)\Xi(r)\bigr)\dd r
 \qquad(0\le x<y\le1).
\]
Further, by the same lemma, every polynomial moment of the spatial derivatives
along $\xi$ is continuous on $[0,1]$ and satisfies the integrated identity
\cite[(2.16)]{Lopatto}, with the polynomials of \cite[(2.18)]{Lopatto}.
\end{enumerate}
Let $K=\max_{j=3,4}\sup_{[0,1]\times\R}|\partial_x^ju|$, which is finite and
depends only on $\beta$ by (i).

We begin with the moments $\Lambda$ and $\Xi$. For finitely supported
measures, the equations \cite[(2.17)]{Lopatto} with $j=2,3$ read
$\dd\hat C_r=-\beta^2\alpha_\beta\hat C_r^2\dd r+\beta\hat D_r\dd w_r$ and
$\dd\hat D_r=-3\beta^2\alpha_\beta\hat C_r\hat D_r\dd r+\beta\hat A_r\dd w_r$,
and It\^o's formula yields
\[
 \begin{aligned}
 \dd(\hat D_r^2)&=\beta^2\bigl(\hat A_r^2-6\alpha_\beta\hat C_r\hat D_r^2\bigr)\dd r
 +2\beta\hat D_r\hat A_r\dd w_r,\\
 \dd(\hat C_r^3)&=3\beta^2\bigl(\hat C_r\hat D_r^2-\alpha_\beta\hat C_r^4\bigr)\dd r
 +3\beta\hat C_r^2\hat D_r\dd w_r .
 \end{aligned}
\]
The stochastic integrals are martingales, since their integrands are bounded
by (i). For the measure $\mu_\beta$, which has an atom, the identities we need
are the cases $f=u_{xxx}^2$ and $f=u_{xx}^3$ of \cite[(2.16)]{Lopatto}, whose
polynomials $f_0,f_1$ from \cite[(2.18)]{Lopatto} are the drifts just
computed:
\[
 \begin{aligned}
 \Lambda(y)-\Lambda(x)&=\beta^2\int_x^y
 \E\bigl[\hat A_r^2-6\alpha_\beta(r)\hat C_r\hat D_r^2\bigr]\dd r,\\
 \Xi(y)-\Xi(x)&=3\beta^2\int_x^y
 \E\bigl[\hat C_r\hat D_r^2-\alpha_\beta(r)\hat C_r^4\bigr]\dd r,
 \end{aligned}
 \qquad(0\le x<y\le1).
\]
Since $0\le\alpha_\beta\le1$, $0<\hat C\le1$, and $|\hat D|,|\hat A|\le K$, the
integrands are bounded by $7K^2$ and $K^2+1$, respectively. We find that
$\Lambda$ and $\Xi$ are Lipschitz on $[0,1]$, with constants $7\beta^2K^2$ and
$3\beta^2(K^2+1)$, and $0\le\Lambda\le K^2$. The function $\Xi$ is continuous
by (iii) and positive, since $\hat C_r>0$ by (i). We define
\[
 \Xi_0=\min_{0\le r\le1}\Xi(r)>0,
\]
which depends only on $\beta$, because $\mu_\beta$ does. (As in the proof of
\cite[Lemma~6.2]{Lopatto}, positivity of $\Xi$ comes from $u_{xx}>0$.)

The quotient $\bar\alpha=\Lambda/(2\Xi)$ is Lipschitz on $[0,1]$. Indeed, let
\[
 L_\beta=\frac{7\beta^2K^2}{2\Xi_0}+\frac{3\beta^2K^2(K^2+1)}{2\Xi_0^2}.
\]
For $x,y\in[0,1]$, using the Lipschitz bounds for $\Lambda$ and $\Xi$
together with $\Lambda\le K^2$ and $\Xi\ge\Xi_0$, we obtain
\[
 |\bar\alpha(y)-\bar\alpha(x)|
 \le\frac{|\Lambda(y)-\Lambda(x)|}{2\Xi(y)}
 +\Lambda(x)\frac{|\Xi(x)-\Xi(y)|}{2\Xi(x)\Xi(y)}
 \le L_\beta|y-x| .
\]

Next, we identify $\alpha_\beta$ with $\bar\alpha$ on $[0,q_\beta)$. By (ii),
we have $\Gamma(r)=r$ on $[0,q_\beta]$, and since $\Gamma$ is continuously
differentiable, $\Gamma'=1$ on $(0,q_\beta)$. Then (iii) implies that
$\int_x^y(\Lambda-2\alpha_\beta\Xi)\dd r=0$ for $0<x<y<q_\beta$. The integrand
is bounded, and by Lebesgue's differentiation theorem, we have
$\Lambda-2\alpha_\beta\Xi=0$, that is, $\alpha_\beta=\bar\alpha$, almost
everywhere on $(0,q_\beta)$. Let $r\in[0,q_\beta)$, and choose points
$r_n\downarrow r$ in $(r,q_\beta)$ at which $\alpha_\beta(r_n)=\bar\alpha(r_n)$.
As $n\to\infty$, the left side tends to $\alpha_\beta(r)$ by the right
continuity of $\alpha_\beta$, and the right side tends to $\bar\alpha(r)$ by
continuity. We conclude that $\alpha_\beta=\bar\alpha$ on $[0,q_\beta)$, and
\eqref{eq:loc-lipschitz} follows from the Lipschitz bound for $\bar\alpha$.

For the consequences, we have $\gamma^{\rm tr}(0)=\max\{\lambda,0\}=\lambda$ by
\eqref{eq:lz-gsroot}, and $\gamma^{\rm tr}=\max\{\lambda,\beta\}=\beta$ on
$[q_\beta,1]$ by \eqref{eq:loc-scalar-measure}. On $[0,q_\beta)$, the function
$\gamma_\beta=\beta\alpha_\beta$ is $\beta L_\beta$-Lipschitz, and so is its
maximum with the constant $\lambda$; since $L\ge\beta L_\beta$, this gives the
Lipschitz bound for $\gamma^{\rm tr}$. A nondecreasing Lipschitz function has
an absolutely continuous derivative measure, with density bounded by its
Lipschitz constant. This yields the density bound on $(0,q_\beta)$, and on
$(q_\beta,1]$ the profile $\gamma^{\rm tr}$ is constant.
\end{proof}

\section{Distance from the scalar trial profile}

We now compare a minimizer with $\gamma^{\rm tr}$ in $L^1(0,1)$. After the
regularization of \cref{lem:loc-smoothing}(b), the derivative measures along
the segment to $\gamma^{\rm tr}$ have bounded density in $(0,1)$ away from the
atom of $\gamma^{\rm tr}$ at $q_\beta$. We construct cutoffs of one sign from
one-sided averages of superlevel sets, and their masses and Dirichlet energies
lead to a single estimate for the size of every superlevel set of the
difference of the two profiles.

\begin{proposition}[Isolation against the trial]\label{prop:loc-isolation}
There exists $C=C(\beta)<\infty$ such that, for all $N\ge1$, $0<\rho\le1$,
$0<\lambda\le\lambda_0$, every minimizer $\gamma_*$ of $F_\rho$, and every
trial budget $\eta$ for $\gamma_*$,
\begin{equation}\label{eq:loc-isolation}
 \|\gamma_*-\gamma^{\rm tr}\|_{L^1(0,1)}\le C\min\{1,\rho^{-3/4}\eta^{1/4}\}.
\end{equation}
The constant depends on $\beta$ only through $\beta$, $k_\beta$, and the
constant $L$ of \cref{lem:loc-lipschitz}.
\end{proposition}

\begin{proof}
The proof has seven steps. We first fix a smoothing width $h$ and replace
$\gamma_*$ by its regularization $\gamma_*^h$ of \cref{lem:loc-smoothing}(b),
whose free measure has no atoms and which costs at most $\rho\beta h^2/12$. Along the segment
from $\gamma_*^h$ to $\gamma^{\rm tr}$, the difference
$\psi=\gamma^{\rm tr}-\gamma_*^h$ increases at rate at most $L$ below
$q_\beta$, and it is nonnegative and nonincreasing above $q_\beta$. For each
level $A>0$, we construct cutoffs of one sign from one-sided averages of the
superlevel sets of $\psi$ below $q_\beta$, and we use triangular cutoffs on the
two intervals next to $q_\beta$; none of them is positive at the atom
$q_\beta$ of $\gamma^{\rm tr}$. \Cref{lem:loc-coercive} bounds the size of
each superlevel set in terms of $f''(\theta)$, and the layer-cake formula
turns these bounds into a lower bound for $f''(\theta)$ in terms of the $L^1$
distance. Finally, we choose $h$ proportional to $\rho$ times the square of
this distance and compare with the midpoint inequality and the trial budget.

\paragraph{Step 1: Setup and constraints on $h$.}
Let $D=\|\gamma_*-\gamma^{\rm tr}\|_{L^1(0,1)}$. If $D=0$, there is nothing to
prove, so we assume from now on that $D>0$. Both profiles take values in
$[\lambda,\beta]$, which implies that $D\le\beta$. With $L\ge1$ as in
\cref{lem:loc-lipschitz}, we fix a width $h$ such that
\begin{equation}\label{eq:loc-h-constraints}
 0<h\le\tfrac12,\qquad h\le\frac{D}{2\beta},\qquad h\le\frac\beta L.
\end{equation}
The final choice, in Step 7, will be a small multiple of $\rho D^2$.

\paragraph{Step 2: The regularized segment.}
Let $\gamma_*^h$ be the regularization of $\gamma_*$ given by
\cref{lem:loc-smoothing}(b), and define
\[
 \gamma_\theta=(1-\theta)\gamma_*^h+\theta\gamma^{\rm tr},\qquad
 \psi=\gamma^{\rm tr}-\gamma_*^h,\qquad f(\theta)=F_\rho(\gamma_\theta)
 \qquad(0\le\theta\le1).
\]
The regularization is defined because $h\le1/2$. Recall that the trial
profile is admissible (\cref{sec:fnd}), so that $(\gamma_\theta)$ is a
segment. By
\cref{lem:loc-smoothing}(b) and the constraint $h\le D/(2\beta)$, we have
\[
 \bar D=\|\psi\|_{L^1(0,1)}\ge D-\|\gamma_*^h-\gamma_*\|_{L^1(0,1)}
 \ge D-\beta h\ge\frac D2 .
\]
For $\theta\in(0,1)$, the measure $\dd\gamma_\theta$ is
$(1-\theta)\dd\gamma_*^h+\theta\dd\gamma^{\rm tr}$. On
$(0,1)\setminus\{q_\beta\}$, the measure $\dd\gamma_*^h$ has density at most
$\beta/h$ by \cref{lem:loc-smoothing}(b), and $\dd\gamma^{\rm tr}$ has density
at most $L$ by \cref{lem:loc-lipschitz}. By \eqref{eq:loc-h-constraints}, we
have $L\le\beta/h$, and the restriction of $\dd\gamma_\theta$ to
$(0,1)\setminus\{q_\beta\}$ has density at most $\beta/h+L\le2\beta/h$. The
only atom of $\dd\gamma_\theta$ in $(0,1)$ is the atom of $\dd\gamma^{\rm tr}$
at $q_\beta$. Further, $\zeta_\theta\ge k_\beta^2$ by \eqref{eq:loc-in-curv}.

\paragraph{Step 3: One-sided bounds on $\psi$.}
By \cref{lem:loc-smoothing}(b) and \cref{lem:loc-lipschitz}, we have
$\gamma_*^h(0)=\gamma^{\rm tr}(0)=\lambda$. For $0\le z<x<q_\beta$, we have
$\gamma^{\rm tr}(x)-\gamma^{\rm tr}(z)\le L(x-z)$ by \cref{lem:loc-lipschitz},
and $\gamma_*^h(x)-\gamma_*^h(z)\ge0$, because $\gamma_*^h$ is nondecreasing. We
obtain
\begin{equation}\label{eq:loc-onesided}
 \psi(0)=0,\qquad
 \psi(x)-\psi(z)\le L(x-z),\qquad
 \psi(x)\le Lx\qquad(0\le z<x<q_\beta).
\end{equation}
On $[q_\beta,1]$, we have $\psi=\beta-\gamma_*^h$, which is continuous (the
free measure of $\gamma_*^h$ has no atoms), nonnegative, and nonincreasing.
Everywhere on $[0,1]$, both profiles take values in $[\lambda,\beta]$, so
$|\psi|\le\beta$.

\paragraph{Step 4: Cutoffs for the superlevel sets.}
We fix $\theta\in(0,1)$. For $A>0$, we set $r=A/(4L)$ and define the sets
\[
 \begin{aligned}
 U_A&=\{x\in(0,q_\beta):\psi(x)>A\},\\
 V_A^i&=\{x\in(0,q_\beta-r):\psi(x)<-A\},\\
 V_A^b&=\{x\in[\max\{0,q_\beta-r\},q_\beta):\psi(x)<-A\},\\
 W_A&=\{x\in[q_\beta,1]:\psi(x)>A\},
 \end{aligned}
\]
where an interval of nonpositive length is empty. These four sets are
disjoint, and since $\psi(0)=0$ and $\psi\ge0$ on $[q_\beta,1]$, their union is
$\{x\in[0,1]:|\psi(x)|>A\}$. In this step, we treat $U_A$ and $V_A^i$; Step 5
treats $V_A^b$ and $W_A$.

For a Borel set $E\subset(0,1)$, we extend its indicator by zero to $\R$ and
define
\[
 \mathcal A_+\ind_E(x)=\frac1r\int_0^r\ind_E(x+u)\dd u,
 \qquad
 \mathcal A_-\ind_E(x)=\frac1r\int_0^r\ind_E(x-u)\dd u
 \qquad(x\in\R).
\]
With $\Phi_E(x)=|E\cap(-\infty,x]|$, we have
$\mathcal A_+\ind_E(x)=(\Phi_E(x+r)-\Phi_E(x))/r$ and
$\mathcal A_-\ind_E(x)=(\Phi_E(x)-\Phi_E(x-r))/r$. The function $\Phi_E$ is
$1$-Lipschitz, and by Lebesgue's differentiation theorem, its derivative is
$\ind_E$ almost everywhere. It follows that $\mathcal A_\pm\ind_E$ is
Lipschitz, with derivative $(\ind_E(x+r)-\ind_E(x))/r$, respectively
$(\ind_E(x)-\ind_E(x-r))/r$, for almost every $x$. Using Fubini's theorem for the integral and
$|(E-r)\triangle E|\le2|E|$ for the energy, we obtain
\begin{equation}\label{eq:loc-averaging}
 0\le\mathcal A_\pm\ind_E\le1,\qquad \int_\R\mathcal A_\pm\ind_E=|E|,\qquad
 \int_\R|\partial_x\mathcal A_\pm\ind_E|^2\le\frac{2|E|}{r^2}.
\end{equation}

For $E=U_A$, we take $\varphi=\mathcal A_+\ind_E$. If $\varphi(x)>0$, then
there exists $y\in U_A$ with $x\le y\le x+r$. By \eqref{eq:loc-onesided}, we
have $A<\psi(y)\le Ly$, so that $y>A/L=4r$ and $x\ge y-r>3r>0$; also
$x\le y<q_\beta$. Then \eqref{eq:loc-onesided} implies that
$\psi(x)\ge\psi(y)-L(y-x)>A-Lr=3A/4$. We have shown that
$\{\varphi>0\}\subset(0,q_\beta)$ and that $\psi>3A/4$ on this set.

For $E=V_A^i$, we take $\varphi=\mathcal A_-\ind_E$. If $\varphi(x)>0$, then
there exists $y\in V_A^i$ with $y\le x\le y+r$, and $y+r<q_\beta$ by the
definition of $V_A^i$. Then $0<x<q_\beta$, and \eqref{eq:loc-onesided} implies
that $\psi(x)\le\psi(y)+L(x-y)<-A+Lr=-3A/4$. Again, we have
$\{\varphi>0\}\subset(0,q_\beta)$, and $\psi<-3A/4$ on this set.

In both cases, $\varphi$ is a Lipschitz function on $[0,1]$ with values in
$[0,1]$, and it vanishes at $0$ and on $[q_\beta,1]$. By
\eqref{eq:loc-averaging}, its integral over $[0,1]$ is $|E|$ and its energy
over $[0,1]$ is at most $2|E|/r^2$. The set $\{\varphi>0\}$
does not contain $q_\beta$, and by Step 2, the restriction of
$\dd\gamma_\theta$ to this set has density at most $d=2\beta/h$. If $|E|>0$,
then $\varphi$ is not identically zero, and \cref{lem:loc-coercive} applies
with $\sigma=1$ for $U_A$ and $\sigma=-1$ for $V_A^i$. Its quantities satisfy
\[
 \Delta=\rho\int_0^1\zeta_\theta\varphi\dd q\ge\rho k_\beta^2|E|,\qquad
 Y_0\le\frac{2|E|}{r^2},\qquad d=\frac{2\beta}h,
\]
by \eqref{eq:loc-in-curv} and \eqref{eq:loc-averaging}. Inserting these bounds
into \eqref{eq:loc-coercive} and using $r=A/(4L)$, we obtain
\begin{equation}\label{eq:loc-cutoff-interior}
 f''(\theta)\ge\frac{\rho^2A^2|E|}{72}
 \min\left\{\frac{\rho k_\beta^4A^2}{32\beta L^2},
 \frac{k_\beta^2h}{2\beta}\right\}
 \qquad(E=U_A,V_A^i).
\end{equation}
When $|E|=0$, the right side is zero, and the bound holds by convexity.

\paragraph{Step 5: The boundary sets.}
The set $V_A^b$ lies within distance $r$ of $q_\beta$, where an average of
$\ind_{V_A^b}$ over a window of length $r$ could be positive at $q_\beta$. We
use a single triangular cutoff instead. For an interval $(x_0,x_0+z)$ with
$z>0$, let
\[
 T_{x_0,z}(x)=\Bigl(1-\frac{4|x-x_0-z/2|}{z}\Bigr)_+ .
\]
This function takes values in $[0,1]$, is positive exactly on the middle half
$(x_0+z/4,x_0+3z/4)$, and satisfies
\[
 \int_\R T_{x_0,z}=\frac z4,\qquad\int_\R|T_{x_0,z}'|^2=\frac8z .
\]

Suppose that $V_A^b$ is nonempty, and let $x_A=\inf V_A^b$ and
$l=q_\beta-x_A$. Note that $0<l\le r$ and $|V_A^b|\le l$. For every
$z\in(x_A,q_\beta)$, there exists $x\in V_A^b$ with $x<z$, and
\eqref{eq:loc-onesided} implies that $\psi(z)\le\psi(x)+L(z-x)<-A+Lr=-3A/4$.
We take $\varphi=T_{x_A,l}$. Its positivity set is the middle half of
$(x_A,q_\beta)$, whose closure lies in $(0,q_\beta)$; on this set, $\psi<-3A/4$
and $\dd\gamma_\theta$ has density at most $2\beta/h$. \Cref{lem:loc-coercive},
applied with $\sigma=-1$, $\Delta\ge\rho k_\beta^2l/4$, $Y_0=8/l$, and
$d=2\beta/h$, yields
\begin{equation}\label{eq:loc-cutoff-boundary}
 f''(\theta)\ge\frac{\rho^2A^2l}{72}
 \min\left\{\frac{\rho k_\beta^4l^2}{128\beta},
                    \frac{k_\beta^2h}{8\beta}\right\}.
\end{equation}
Indeed, $\Delta^2/(\beta Y_0)\ge\rho^2k_\beta^4l^3/(128\beta)$ and
$\Delta/d\ge\rho k_\beta^2lh/(8\beta)$.

By Step 3, the function $\psi$ is continuous, nonnegative, and nonincreasing on
$[q_\beta,1]$. Hence $W_A$ is an interval $[q_\beta,q_\beta+w)$ of length
$w\le1$.
If $w>0$, we take $\varphi=T_{q_\beta,w}$, whose positivity set has closure in
$(q_\beta,1)$. On this set, $\psi>A$, the measure $\dd\gamma^{\rm tr}$ vanishes,
and $\dd\gamma_\theta$ has density at most $2\beta/h$. The same computation,
with $\sigma=1$, gives \eqref{eq:loc-cutoff-boundary} with $w$ in place of
$l$. If $w=0$, there is nothing to prove.

\paragraph{Step 6: Layer-cake integration.}
Let
\[
 c_1=\frac1{72}\min\left\{\frac{k_\beta^4}{32\beta L^2},
                         \frac{k_\beta^4}{128\beta},
                         \frac{k_\beta^2}{8\beta}\right\},\qquad
 K=\frac{f''(\theta)}{c_1\rho^2}.
\]
Since $\min\{ax,by\}\ge\min\{a,b\}\min\{x,y\}$ for nonnegative numbers,
\eqref{eq:loc-cutoff-interior} and \eqref{eq:loc-cutoff-boundary} imply that
\[
 A^2|E|\min\{\rho A^2,h\}\le K\quad(E=U_A,V_A^i),\qquad
 A^2z\min\{\rho z^2,h\}\le K\quad(z=l,w).
\]
For $E=U_A$ or $E=V_A^i$, one of $|E|\le K/(\rho A^4)$ and $|E|\le K/(hA^2)$
holds, and $|E|\le1$. In either case, we have
\[
 |E|\le\min\{1,K/(\rho A^4)\}+\min\{1,K/(hA^2)\}.
\]
For the boundary interval, one of $l\le(K/(\rho A^2))^{1/3}$ and
$l\le K/(hA^2)$ holds. Combining these alternatives with $|V_A^b|\le l\le
r=A/(4L)$, we obtain
\[
 |V_A^b|\le\Bigl(\frac K{\rho A^2}\Bigr)^{1/3}
 +\min\Bigl\{\frac A{4L},\frac K{hA^2}\Bigr\}.
\]
The same reasoning with $w\le1$ yields
\[
 |W_A|\le\Bigl(\frac K{\rho A^2}\Bigr)^{1/3}+\min\Bigl\{1,\frac K{hA^2}\Bigr\}.
\]
All four sets are empty for $A\ge\beta$, since $|\psi|\le\beta$. We use the
elementary integrals
\[
 \int_0^\infty\min\Bigl\{1,\frac x{A^4}\Bigr\}\dd A=\frac43x^{1/4},\qquad
 \int_0^\infty\min\Bigl\{1,\frac x{A^2}\Bigr\}\dd A=2x^{1/2},
\]
\[
 \int_0^\infty\min\Bigl\{\frac A{4L},\frac x{A^2}\Bigr\}\dd A=C_Lx^{2/3},\qquad
 \int_0^\beta\Bigl(\frac x{A^2}\Bigr)^{1/3}\dd A=3\beta^{1/3}x^{1/3},
\]
where $C_L=3/(2(4L)^{1/3})$ and $x\ge0$. The case $x=0$ is immediate. For
$x>0$, the first three follow by splitting
at the point where the two entries of the minimum agree; for the third, this
point is $A_0=(4Lx)^{1/3}$, and the integral is
$A_0^2/(8L)+x/A_0=3x/(2A_0)$. By the layer-cake formula, we have
$\bar D=\int_0^\beta|\{|\psi|>A\}|\dd A$. Integrating the bounds on the four
sets over $0<A<\beta$, we find
\begin{equation}\label{eq:loc-weaknorm}
 \bar D\le\frac83\Bigl(\frac K\rho\Bigr)^{1/4}+6\Bigl(\frac Kh\Bigr)^{1/2}
 +6\beta^{1/3}\Bigl(\frac K\rho\Bigr)^{1/3}+C_L\Bigl(\frac Kh\Bigr)^{2/3}.
\end{equation}
Let $C_w=\max\{6\beta^{1/3},C_L\}$. Each coefficient in
\eqref{eq:loc-weaknorm} is at most $C_w$, because $\beta>1$.

We now show that $K\ge c_2\min\{\rho D^4,hD^2\}$, where
$c_2=\min\{(8C_w)^{-4},\beta^{-4}\}$. Suppose first that $K\le\rho$ and $K\le
h$. Then $(K/\rho)^{1/3}\le(K/\rho)^{1/4}$ and $(K/h)^{2/3}\le(K/h)^{1/2}$,
and \eqref{eq:loc-weaknorm} gives
$\bar D\le2C_w[(K/\rho)^{1/4}+(K/h)^{1/2}]$. Using $\bar D\ge D/2$, we see that one
of the two terms in the bracket is at least $D/(8C_w)$. This implies that
$K\ge\rho D^4/(8C_w)^4$ or $K\ge hD^2/(8C_w)^2$, and since $8C_w\ge1$, in both cases
$K\ge(8C_w)^{-4}\min\{\rho D^4,hD^2\}$. If $K>\rho$, then $D\le\beta$ gives
$\min\{\rho D^4,hD^2\}\le\rho\beta^4<\beta^4K$. If $K>h$, then
$\min\{\rho D^4,hD^2\}\le h\beta^2<\beta^4K$, since $\beta>1$. This proves the
claim in all cases. Since $\theta\in(0,1)$ was arbitrary, we have proved that
\begin{equation}\label{eq:loc-E2-second}
 f''(\theta)\ge c_3\rho^2\min\{\rho D^4,hD^2\}
 \quad(0<\theta<1),\qquad c_3=c_1c_2 .
\end{equation}

\paragraph{Step 7: Choice of $h$ and conclusion.}
We now take $h=c_4\rho D^2$, where
\[
 c_4=\min\Bigl\{1,\frac1{2\beta^2},\frac1{\beta L},\frac{3c_3}{2\beta}\Bigr\}.
\]
Since $D\le\beta$ and $\rho\le1$, we have $h\le c_4\beta^2\le1/2$,
$h\le c_4\beta D\le D/(2\beta)$, and $h\le c_4\beta^2\le\beta/L$. These are the
constraints \eqref{eq:loc-h-constraints}. Because $c_4\le1$, the minimum in
\eqref{eq:loc-E2-second} equals $c_4\rho D^4$, and \eqref{eq:loc-E2-second}
gives $f''\ge c_3c_4\rho^3D^4$ on $(0,1)$. We use the midpoint inequality
\eqref{eq:loc-midpoint} with $\Theta=(0,1)$, where the weight has integral
$1/4$. Together with the regularization cost of \cref{lem:loc-smoothing}(b)
and the trial budget \eqref{eq:loc-trialbudget}, it yields
\[
 \frac{c_3c_4}{4}\rho^3D^4
 \le F_\rho(\gamma_*^h)+F_\rho(\gamma^{\rm tr})-2F_\rho(\gamma_*)
 \le\frac{\rho\beta h^2}{12}+\eta
 \le\frac{c_3c_4}{8}\rho^3D^4+\eta .
\]
In the last inequality, we used that $\rho\beta h^2/12=(\beta
c_4^2/12)\rho^3D^4$ and $c_4\le3c_3/(2\beta)$. We conclude that
$\eta\ge(c_3c_4/8)\rho^3D^4$, that is,
$D\le(8/(c_3c_4))^{1/4}\rho^{-3/4}\eta^{1/4}$. Together with $D\le\beta$, this
implies \eqref{eq:loc-isolation} with
$C=\max\{\beta,(8/(c_3c_4))^{1/4}\}$. The constants $c_1,c_2,c_3,c_4$ depend
only on $\beta$, $k_\beta$, and $L$.
\end{proof}

\chapter{Localization of the lowest quantiles}\label{sec:lz-root}

In this section, we show that the lowest quantiles of the free mass of every
minimizer of $F_\rho$ lie close to the origin. Let $\mu=s\delta_0+\nu$ be a
minimizing profile and $\eta$ a trial budget. By \cref{cor:lz-omega}, the
quantile $Q(v)$ of $\nu$ at level $v$ is at most
$C_Q(s+v+\omega(\rho,\eta))$, where
$\omega(\rho,\eta)=\min\{1,(\eta/\rho)^{1/6}\}$ and $C_Q$ depends only on
$\beta$. \Cref{prop:lz-quantile} is the same estimate in the energy
normalization. In \cref{prop:sm-small}, we use \cref{cor:lz-omega} to bound the
self-overlap variance on dyadic quantile shells for the interpolation
parameters near $\rho=0$.

We keep the standing assumptions and the notation of \cref{subsec:lz-inputs}.
From \cref{sec:lz}, we use the convexity and the first variation of $F_\rho$
(\cref{prop:loc-inputs}(a)), the identities \eqref{eq:loc-b}, the coercivity
bound of \cref{lem:loc-coercive}, and the midpoint inequality
\eqref{eq:loc-midpoint}. We also use the partial regularization of
\cref{lem:loc-smoothing}(a), the scalar root bounds \eqref{eq:lz-gsroot}, the
structure \eqref{eq:loc-scalar-measure} of $\mu_\beta$, and the trial budget
\eqref{eq:loc-trialbudget}. The isolation estimate \cref{prop:loc-isolation}
and the density bound \cref{lem:loc-lipschitz} are not used here.

The main tool is a local deficit certificate (\cref{lem:lz-certificate}). It
combines a variation that moves mass of a
profile toward the middle of an interval with the response inequality
\eqref{eq:loc-response}. If the variation violates its room or sign
condition, this failure gives a lower bound on $\eta/\rho$ as well, and the
certificate includes this bound as an alternative. In
\cref{subsec:lz-quantile}, we follow the strategy of the proof of
\cite[Lemma~7.6]{AL}. We regularize only the free mass below the selected
quantile. Then we compare the result with the scalar trial profile on an
interval where the latter has a positive density (\cref{lem:lz-rootdensity}).
The certificate gives the dependence on $\eta/\rho$ in \eqref{eq:lz-quantile}.

\section{A local deficit certificate}\label{subsec:lz-moves}

In the proof of \cref{prop:loc-isolation}, we used the uniform lower bound
\eqref{eq:loc-in-curv} on $\zeta$, which gives
$\Delta=\rho\int\zeta\varphi\ge\rho k_\beta^2\int\varphi$. For root
localization, we need $\Delta$ to be at least a constant multiple of the length
of the interval on which $\varphi$ lives, without the factor $\rho$. In the
next lemma, we achieve this by moving mass of the profile toward the middle of
the interval. If the variation does not have enough room, or if it would
destroy the profile gap, then one of the lower bounds in the conclusion already
holds.

\begin{lemma}[Local deficit certificate]\label{lem:lz-certificate}
There exists $c_{\rm cert}=c_{\rm cert}(\beta)>0$ with the following property.
Let $I\Subset(0,1)$ be an open interval of length $\ell$, let
$\gamma_0,\gamma_1$ be admissible profiles, and let $\eta>0$ be such that
$F_\rho(\gamma_i)\le\min F_\rho+\eta$ for $i=0,1$. Suppose that there exist
$k>0$, $d_0,d_1\ge0$, and $A>0$ such that, on $I$, the derivative measures and
the profile gap satisfy
\[
 \dd\gamma_0\le d_0\dd x,\qquad
 k\dd x\le\dd\gamma_1\le d_1\dd x,
 \qquad \gamma_1-\gamma_0\ge A .
\]
Then
\begin{equation}\label{eq:lz-certificate}
 \frac\eta\rho\ge c_{\rm cert}\min\left\{
 k\ell^3,\ A\ell^2,\ A^2\ell^4,
                  \frac{A^2\ell^2}{d_0+d_1+k}\right\}.
\end{equation}
\end{lemma}

\begin{proof}
The proof has four steps. In Step 1, we fix a bump function $\varphi$ on $I$
and an amplitude $\tau$. If $\tau$ is too large for the variation below, then
one of the first two alternatives in \eqref{eq:lz-certificate} holds. Otherwise, we perturb the segment from
$\gamma_0$ to $\gamma_1$ by the derivative of the bump, which moves mass of
the profile toward the middle of $I$ (Step 2). Since both endpoints are nearly
minimal and $F_\rho$ is convex, the response integrated against the bump is
at least of order $\ell$ along the perturbed segment (Step 3). Then
\cref{lem:loc-coercive} implies a lower bound on the curvature along the
perturbed segment. Comparing this curvature with the excesses at its
endpoints, we obtain the last two alternatives (Step 4). Throughout, we write $F_{\min}=\min F_\rho$ and
$I=(x_0,x_0+\ell)$.

\paragraph{Step 1: The bump and the room condition.}
Let $B(y)=16y^2(1-y)^2$ for $y\in[0,1]$, and let $\varphi(x)=B((x-x_0)/\ell)$
for $x\in I$ and $\varphi(x)=0$ for $x\notin I$. We have
\[
 B'(y)=32y(1-y)(1-2y),\qquad B''(y)=32-192y+192y^2,
\]
and $B$ and $B'$ vanish at $0$ and $1$. A direct computation shows that
$\int_0^1B=8/15$, $\max|B'|=16\sqrt3/9<4$, $\max|B''|=32$, and
$\int_0^1|B'|^2=512/105$. These properties of $B$ imply that $\varphi$ is
continuously differentiable on $\R$ with values in $[0,1]$, that $\varphi'$ is
Lipschitz, that $\varphi$ and $\varphi'$ vanish at the endpoints of $I$, and
that
\begin{equation}\label{eq:lz-bump}
 \int\varphi=\frac{8\ell}{15},\qquad \|\varphi'\|_\infty\le\frac4\ell,\qquad
 \|\varphi''\|_\infty\le\frac{32}{\ell^2},\qquad
 \int|\varphi'|^2=\frac{512}{105\,\ell},
\end{equation}
where $\varphi''$ is defined away from the endpoints of $I$. In particular,
the distributional derivative of $\varphi'$ is the function $\varphi''$, which
has no atoms at the endpoints of $I$, and $\int\varphi''=0$.

We define $\tau=30\eta/(\rho\ell)$; the factor $30$ is chosen so that
$8\eta/(\rho\tau)=4\ell/15$ is half of $\int\varphi$. Suppose first that the
room condition
\begin{equation}\label{eq:lz-room-condition}
 \tau\le\min\Bigl\{\frac{k\ell^2}{64},\frac{A\ell}8\Bigr\}
\end{equation}
fails. Then
\[
 \frac\eta\rho=\frac{\ell\tau}{30}
 >\min\Bigl\{\frac{k\ell^3}{1920},\frac{A\ell^2}{240}\Bigr\}
 \ge\frac1{1920}\min\{k\ell^3,A\ell^2\},
\]
and \eqref{eq:lz-certificate} holds whenever $c_{\rm cert}\le1/1920$. We
assume from now on that \eqref{eq:lz-room-condition} holds.

\paragraph{Step 2: The perturbed segment.}
For $0\le\theta\le1$, let
\[
 \gamma_\theta=(1-\theta)\gamma_0+\theta\gamma_1,\qquad
 \widetilde\gamma_\theta=\gamma_\theta-\theta\tau\varphi'.
\]
We claim that $\gamma_\theta-u\varphi'$ is admissible for $0\le\theta\le1$ and
$0\le u\le\theta\tau$. The function $\varphi'$ is continuous and vanishes
outside $I$, and $\overline I\subset(0,1)$. Then $\gamma_\theta-u\varphi'$ is
right-continuous and agrees with $\gamma_\theta$ near $0$ and near $1$. Its
derivative measure is $\dd\gamma_\theta-u\varphi''\dd x$. Off $I$, this is
$\dd\gamma_\theta$. On $I$, the hypotheses give $\dd\gamma_\theta\ge\theta
k\dd x$, and by \eqref{eq:lz-bump} and \eqref{eq:lz-room-condition}, we have
\[
 \dd(\gamma_\theta-u\varphi')
 \ge\Bigl(\theta k-\frac{32u}{\ell^2}\Bigr)\dd x
 \ge\Bigl(\theta k-\frac{32\theta\tau}{\ell^2}\Bigr)\dd x
 \ge\frac{\theta k}2\dd x .
\]
The signed measure $\varphi''\dd x$ has total mass zero and does not charge
$0$. We conclude that the derivative measure of $\gamma_\theta-u\varphi'$ is
nonnegative, has total mass $\beta$, and has the same atom at $0$ as
$\dd\gamma_\theta$; that is, $\gamma_\theta-u\varphi'$ is admissible. We have
$\widetilde\gamma_\theta=(1-\theta)\gamma_0+\theta(\gamma_1-\tau\varphi')$, and
$\gamma_1-\tau\varphi'$ is admissible (this is the claim with $\theta=1$ and $u=\tau$). It
follows that $(\widetilde\gamma_\theta)$ is a segment, with endpoints
$\widetilde\gamma_0=\gamma_0$ and $\widetilde\gamma_1=\gamma_1-\tau\varphi'$.

\paragraph{Step 3: The response integrated against the bump.}
Fix $\theta\in(0,1]$ and define $g_\theta(u)=F_\rho(\gamma_\theta-u\varphi')$
for $0\le u\le\theta\tau$. The profiles $\gamma_\theta-u\varphi'$,
$0\le u\le\theta\tau$, form a segment up to a linear change of parameter, so
$g_\theta$ is convex and continuous by \cref{prop:loc-inputs}. We compute its
one-sided derivatives. Let $0\le u\le\theta\tau$, and write
$\gamma'=\gamma_\theta-u\varphi'$, $G=G_{\gamma'}$, and $S=S_{\gamma'}$.
Increasing $u$ changes the free measure of $\gamma'$ in the direction
$-\varphi''\dd x$, and decreasing $u$ changes it in the direction
$+\varphi''\dd x$; in both cases, the new free measure is that of an
admissible profile on the segment. By \eqref{eq:loc-firstvariation}, the right
derivative of $g_\theta$ at $u<\theta\tau$ and the left derivative at $u>0$
are both equal to $-(\rho/2)\int G\varphi''\dd x$. We integrate by parts
twice. The function $\varphi'$ vanishes at the endpoints of $I$ and
$G'(x)=x-S(x)$, which gives $\int G\varphi''\dd x=-\int(x-S)\varphi'\dd x$. The
function $\varphi$ also vanishes at the endpoints of $I$, and $S$ is absolutely
continuous by \eqref{eq:loc-b}; integrating by parts once more, we find
\[
 -\int(x-S)\varphi'\dd x=\int\varphi\dd(x-S)=\frac{8\ell}{15}-\int\varphi\dd S .
\]
We obtain
\begin{equation}\label{eq:lz-concentrating-derivative}
 g_\theta'(u)=\frac\rho2
          \left(\int\varphi\dd S_{\gamma_\theta-u\varphi'}-\frac{8\ell}{15}\right),
\end{equation}
for the right derivative at $0\le u<\theta\tau$ and for the left derivative
at $0<u\le\theta\tau$.

By the convexity of $F_\rho$ along the segment $(\gamma_\theta)$, we have
$g_\theta(0)=F_\rho(\gamma_\theta)\le(1-\theta)F_\rho(\gamma_0)
+\theta F_\rho(\gamma_1)\le F_{\min}+\eta$, and $g_\theta(\theta\tau)\ge
F_{\min}$. For a convex function, the left derivative at the right endpoint
of an interval is at least the slope of the secant over the interval, which
gives
\[
 g_\theta'(\theta\tau)\ge\frac{g_\theta(\theta\tau)-g_\theta(0)}{\theta\tau}
 \ge-\frac\eta{\theta\tau}.
\]
Inserting \eqref{eq:lz-concentrating-derivative} at $u=\theta\tau$, where
$\gamma_\theta-\theta\tau\varphi'=\widetilde\gamma_\theta$, we find, for
$1/4\le\theta\le3/4$,
\[
 \Delta_\theta=\int\varphi\dd S_{\widetilde\gamma_\theta}
 \ge\frac{8\ell}{15}-\frac{2\eta}{\rho\theta\tau}
 \ge\frac{8\ell}{15}-\frac{8\eta}{\rho\tau}=\frac{4\ell}{15}.
\]
By \eqref{eq:loc-b}, we have $\dd S_{\widetilde\gamma_\theta}
=\rho\zeta_{\widetilde\gamma_\theta}\dd x$. In other words,
$\Delta_\theta=\rho\int\zeta_{\widetilde\gamma_\theta}\varphi\dd x$ is the
quantity $\Delta$ of \cref{lem:loc-coercive} for the segment
$(\widetilde\gamma_\theta)$ and the cutoff $\varphi$.

\paragraph{Step 4: Curvature and the endpoint excesses.}
By \eqref{eq:lz-bump} and \eqref{eq:lz-room-condition}, the direction of the
segment $(\widetilde\gamma_\theta)$ satisfies
\[
 \widetilde\gamma_1-\widetilde\gamma_0=\gamma_1-\gamma_0-\tau\varphi'
 \ge A-\frac{4\tau}\ell\ge\frac A2\qquad\text{on }I=\{\varphi>0\}.
\]
On $I$, the derivative measure of $\widetilde\gamma_\theta$ is at most
$((1-\theta)d_0+\theta d_1+32\theta\tau/\ell^2)\dd x\le\bar d\dd x$, where
\[
 \bar d=d_0+d_1+\frac{32\tau}{\ell^2}\le d_0+d_1+\frac k2 ;
\]
note that $\bar d\ge d_1\ge k>0$. Let
$\tilde f(\theta)=F_\rho(\widetilde\gamma_\theta)$. By \cref{lem:loc-coercive},
used for the segment $(\widetilde\gamma_\theta)$ with the cutoff $\varphi$,
$\sigma=1$, the level $A$, and $d=\bar d$, together with
$\Delta_\theta\ge4\ell/15$ and $Y_0=512/(105\ell)$ from \eqref{eq:lz-bump}, we
obtain, for $1/4\le\theta\le3/4$,
\[
 \tilde f''(\theta)\ge\frac{\rho A^2}{72}
 \min\left\{\Bigl(\frac{4\ell}{15}\Bigr)^2\frac{105\,\ell}{512\beta},
 \frac{4\ell}{15\bar d}\right\}
 =\frac{\rho A^2}{72}\min\left\{\frac{7\ell^3}{480\beta},
 \frac{4\ell}{15\bar d}\right\}.
\]
Because $\beta>1$, we have $4/15\ge7/(480\beta)$, which gives
\begin{equation}\label{eq:lz-certificate-curvature}
 \frac{\dd^2}{\dd\theta^2}F_\rho(\widetilde\gamma_\theta)
 \ge\frac7{34560\beta}\,\rho A^2\min\Bigl\{\ell^3,\frac\ell{\bar d}\Bigr\}
 \qquad(1/4\le\theta\le3/4).
\end{equation}

It remains to bound the excesses at the endpoints of the perturbed segment.
At $\theta=0$, we have $\tilde f(0)=F_\rho(\gamma_0)\le F_{\min}+\eta$. At
$\theta=1$, we have $\tilde f(1)=g_1(\tau)$, and the convexity of $g_1$ gives
$g_1(\tau)-g_1(0)\le\tau g_1'(\tau)$, with the left derivative. By
\eqref{eq:loc-b}, the function $S_{\widetilde\gamma_1}$ is nondecreasing with
values in $[0,1]$, and since $0\le\varphi\le1$, we have
$0\le\int\varphi\dd S_{\widetilde\gamma_1}\le1$. Then
\eqref{eq:lz-concentrating-derivative} gives $g_1'(\tau)\le\rho/2$. We
conclude that $\tilde f(1)\le F_\rho(\gamma_1)+\rho\tau/2\le F_{\min}+\eta+15\eta/\ell$.
Using $\ell\le1$, we obtain
\[
 \tilde f(0)+\tilde f(1)-2F_{\min}\le2\eta+\frac{15\eta}\ell
 \le\frac{17\eta}\ell .
\]
We now use the midpoint inequality \eqref{eq:loc-midpoint} for the segment
$(\widetilde\gamma_\theta)$ with $\Theta=[1/4,3/4]$, on which the weight has
integral $3/16$. Together with \eqref{eq:lz-certificate-curvature}, it yields
\[
 \frac3{16}\cdot\frac7{34560\beta}\,\rho A^2
 \min\Bigl\{\ell^3,\frac\ell{\bar d}\Bigr\}\le\frac{17\eta}\ell,
\]
that is,
\[
 \frac\eta\rho\ge\frac7{3133440\beta}A^2\min\Bigl\{\ell^4,\frac{\ell^2}{\bar d}\Bigr\}
 \ge\frac7{3133440\beta}\min\Bigl\{A^2\ell^4,\frac{A^2\ell^2}{d_0+d_1+k}\Bigr\},
\]
where we used $\bar d\le d_0+d_1+k$. Together with Step 1, this proves
\eqref{eq:lz-certificate} with $c_{\rm cert}=7/(3133440\beta)$, which is
smaller than $1/1920$ because $\beta>1$.
\end{proof}

\section{Root localization}\label{subsec:lz-quantile}

For an admissible profile $\gamma$, we define the quantile of its free measure
in the energy normalization by
\[
 Q_\gamma(v)=\inf\{q\in[0,1]:\nu_\gamma([0,q])\ge v\},
 \qquad 0<v<\beta-\lambda .
\]
If a quantile of a minimizer is large compared with the prescribed mass plus
its level, then the minimizer is small on an interval where the scalar trial
profile has a positive density. We control the density of the partially
regularized minimizer there, and then compare the two profiles on that
interval using \cref{lem:lz-certificate}. We follow the strategy of the proof
of \cite[Lemma~7.6]{AL}, which regularizes only the free mass below the
selected quantile and compares the result with the scalar trial profile. We
combine this strategy with \cref{lem:lz-certificate}, which gives the
dependence on $\eta/\rho$ in \eqref{eq:lz-quantile}. The positive density of
the trial profile near the origin comes from the next lemma.

\begin{lemma}[Trial density near the root]\label{lem:lz-rootdensity}
There exists $T_1=T_1(\beta)\in(0,T_0]$ such that
\begin{equation}\label{eq:lz-rootdensity}
 \dd\gamma_\beta\ge\frac{k_0}2\dd x\qquad\text{as measures on }[0,T_1].
\end{equation}
\end{lemma}

\begin{proof}
Recall from \eqref{eq:loc-scalar-measure} that $\mu_\beta$ has no atom in
$[0,q_\beta)$, and that its density $\varrho_\beta$ is continuous on
$[0,q_\beta]$. We have $\gamma_\beta(x)=\beta\int_0^x\varrho_\beta$ for
$0\le x<q_\beta$, and $\gamma_\beta(x)/x\to\beta\varrho_\beta(0)$ as
$x\downarrow0$. By \eqref{eq:lz-gsroot}, we have $\gamma_\beta(x)/x\ge k_0$ for
$0<x\le T_0$, which implies that $\beta\varrho_\beta(0)\ge k_0$. By the
continuity of $\varrho_\beta$ at $0$, there exists
$T_1\in(0,\min\{T_0,q_\beta\})$ such that $\beta\varrho_\beta\ge k_0/2$ on
$[0,T_1]$. The measure $\dd\gamma_\beta$ has no atom at $0$, because
$\gamma_\beta(0)=0$, and on $[0,T_1]$ it equals $\beta\varrho_\beta\dd x$; this gives
\eqref{eq:lz-rootdensity}. The scalar Parisi measure is unique
\cite[Corollary~1]{AuffingerChen2015Uniqueness}, and $q_\beta$ and
$\varrho_\beta$ depend only on $\beta$. The constants $k_0$ and $T_0$ also
depend only on $\beta$, and we may therefore choose $T_1$ depending only on
$\beta$.
\end{proof}

The width $T_1$ is not explicit, because the continuity of $\varrho_\beta$ at
$0$ is a qualitative statement.

\begin{proposition}[Root localization]\label{prop:lz-quantile}
There exists $C=C(\beta)<\infty$ with the following property. Let $N\ge1$,
$0<\rho\le1$, and $0<\lambda\le\lambda_0$, let $\gamma_*$ be a minimizer
of $F_\rho$, and let $\eta$ be a trial budget for $\gamma_*$. Then
\begin{equation}\label{eq:lz-quantile}
 Q_{\gamma_*}(v)\le C\left[\lambda+v+\min\{1,(\eta/\rho)^{1/6}\}\right]
 \qquad(0<v<\beta-\lambda).
\end{equation}
The constant depends on $\beta$ only through $\beta$, $k_0$, $C_0$, and
$T_1$.
\end{proposition}

\begin{proof}
The proof has five steps. After reducing to $0<\eta<\rho$, we replace the
quantile by its truncation $q=\min\{Q_{\gamma_*}(v),T_1\}$ and dispose of the
case where $q$ is small compared with $\lambda+v$ (Step 1). Otherwise, we
regularize the free mass of $\gamma_*$ below $q$, with a width adapted to the
moved mass (Step 2), compare the result with $\gamma^{\rm tr}$ on the middle
interval $(q/4,3q/4)$ (Step 3), and apply \cref{lem:lz-certificate} (Step 4).
The resulting inequality bounds $q$ either by a power of $\eta/\rho$ or in
terms of the density of the regularized profile, which we control with the
choice of the width and Young's inequality (Step 5).

\paragraph{Step 1: Reductions.}
Fix $v\in(0,\beta-\lambda)$, and let $\bar v=\lambda+v$. The constant $C$
constructed below is at least one. If $\eta\ge\rho$, the claim follows from
$Q_{\gamma_*}(v)\le1=\min\{1,(\eta/\rho)^{1/6}\}$. If $\eta=0$, every
$\eta'\in(0,\rho)$ is also a trial budget for $\gamma_*$; we apply the
argument below for positive budgets to $\eta'$ and let $\eta'\downarrow0$. We
may therefore assume that $0<\eta<\rho$, and we write
$\varepsilon=\eta/\rho\in(0,1)$.

We set $q=\min\{Q_{\gamma_*}(v),T_1\}$. Then $Q_{\gamma_*}(v)\le q/T_1$:
if $Q_{\gamma_*}(v)\le T_1$, then $q=Q_{\gamma_*}(v)$ and $T_1\le1$; otherwise
$q=T_1$ and $Q_{\gamma_*}(v)\le1=q/T_1$. It suffices to bound $q$.
If $q\le16\bar v/k_0$, then $Q_{\gamma_*}(v)\le16\bar v/(k_0T_1)$, which is of
the required form. Suppose instead that
\begin{equation}\label{eq:lz-q-large}
 q>16\bar v/k_0 .
\end{equation}

\paragraph{Step 2: Regularization below $q$.}
Let $m=\nu_{\gamma_*}([0,q))$. For every $x<q$, we have $x<Q_{\gamma_*}(v)$,
and the definition of the quantile gives $\nu_{\gamma_*}([0,x])<v$. Since the
interval $[0,q)$ is the increasing union of the intervals $[0,x]$, $x<q$,
continuity from below implies that $m\le v$. An atom of $\nu_{\gamma_*}$ at $q$ is not counted
in $m$: it is not moved by the regularization below, and it lies outside the
closure of the interval of Step 3.

We take
\[
 h=\min\{q/8,\sqrt{\varepsilon/m}\}\quad\text{if }m>0,
 \qquad h=q/8\quad\text{if }m=0 .
\]
Then $h\le q/8$ and $mh^2\le\varepsilon$, and $h\le1-q$ because $q\le
T_1\le T_0\le1/4$. Let $\bar\gamma$ be the partial regularization of
\cref{lem:loc-smoothing}(a) with $T=q$. By that lemma, we have
$F_\rho(\bar\gamma)-F_\rho(\gamma_*)\le\rho mh^2/12\le\eta/12$, and the
restriction of $\dd\bar\gamma$ to $(0,q)$ has density at most $m/h$. For
$0\le x<q$, the free mass of $\bar\gamma$ in $[0,x]$ comes from the moved
mass. Using \eqref{eq:lz-q-large} in the last step, we find that
\[
 \bar\gamma(x)\le\lambda+m\le\bar v<\frac{k_0q}{16}
 \qquad(0\le x<q).
\]

\paragraph{Step 3: The middle interval.}
Let $I=(q/4,3q/4)$. For $x\in I$, \eqref{eq:lz-gsroot} and
\eqref{eq:lz-q-large} give $\gamma_\beta(x)\ge k_0x>k_0q/4>4\bar v\ge\lambda$,
which implies that $\gamma^{\rm tr}=\gamma_\beta$ on $I$. Using
$I\subset(0,T_1]$, \cref{lem:lz-rootdensity}, and the density bound for
$\dd\gamma^{\rm tr}$ on $(0,T_0]$ stated after \eqref{eq:lz-gsroot}, we find
that
\[
 \frac{k_0}2\dd x\le\dd\gamma^{\rm tr}\le C_0\dd x\quad\text{on }I .
\]
By Step 2, we also have $\dd\bar\gamma\le(m/h)\dd x$ on $I$ and
\[
 \gamma^{\rm tr}-\bar\gamma>\frac{k_0q}4-\frac{k_0q}{16}\ge\frac{k_0q}8
 \quad\text{on }I .
\]

\paragraph{Step 4: The certificate.}
We apply \cref{lem:lz-certificate} with the interval
$I=(q/4,3q/4)\Subset(0,1)$, the endpoints $\gamma_0=\bar\gamma$ and
$\gamma_1=\gamma^{\rm tr}$, and the parameters
\[
 \ell=\frac q2,\qquad A=\frac{k_0q}8,\qquad k=\frac{k_0}2,
 \qquad d_0=\frac mh,\qquad d_1=C_0 .
\]
Both endpoints have excess at most $\eta$ over $\min F_\rho=F_\rho(\gamma_*)$:
for $\bar\gamma$ by Step 2, and for $\gamma^{\rm tr}$ by
\eqref{eq:loc-trialbudget}. The four entries of the minimum in
\eqref{eq:lz-certificate} are
\[
 \begin{gathered}
 k\ell^3=\frac{k_0q^3}{16},\qquad A\ell^2=\frac{k_0q^3}{32},\qquad
 A^2\ell^4=\frac{k_0^2q^6}{1024},\\
 \frac{A^2\ell^2}{d_0+d_1+k}=\frac{k_0^2q^4}{256(m/h+C_0+k_0/2)}.
 \end{gathered}
\]
Since $q\le1$, the first two entries are at least $k_0q^6/32$. With
$c'=c_{\rm cert}\min\{k_0/32,k_0^2/1024\}$, we find
\begin{equation}\label{eq:lz-balance}
 \varepsilon\ge c'\min\Bigl\{q^6,\frac{q^4}{m/h+C_0+k_0/2}\Bigr\}.
\end{equation}

\paragraph{Step 5: Conclusion.}
If the minimum in \eqref{eq:lz-balance} is $q^6$, then
$q\le(\varepsilon/c')^{1/6}$. Otherwise, we have
$q^4\le\varepsilon(m/h+C_0+k_0/2)/c'$. By the choice of $h$, we have
\[
 \frac mh=\max\Bigl\{\frac{8m}q,\varepsilon^{-1/2}m^{3/2}\Bigr\}
 \le\frac{8m}q+\varepsilon^{-1/2}m^{3/2}
\]
if $m>0$, and $m/h=0$ if $m=0$. Using $m\le v\le\bar v$
and \eqref{eq:lz-q-large}, we have $8m/q<k_0/2$. When the minimum is the second entry, we therefore
obtain
\[
 q^4\le\frac{(C_0+k_0)\varepsilon+\varepsilon^{1/2}m^{3/2}}{c'},
 \qquad
 q\le(c')^{-1/4}\bigl((C_0+k_0)^{1/4}\varepsilon^{1/4}
                           +\varepsilon^{1/8}m^{3/8}\bigr),
\]
where we used $(x+y)^{1/4}\le x^{1/4}+y^{1/4}$ for $x,y\ge0$. By Young's
inequality with exponents $8/5$ and $8/3$, we have
\[
 \varepsilon^{1/8}m^{3/8}
 \le\tfrac58\varepsilon^{1/5}+\tfrac38m
 \le\varepsilon^{1/6}+m,
\]
where in the last inequality we used that $0<\varepsilon<1$. Together with
$\varepsilon^{1/4}\le\varepsilon^{1/6}$ and $m\le\bar v$, this shows that
$q\le(c')^{-1/4}((C_0+k_0)^{1/4}+1)(\bar v+\varepsilon^{1/6})$. Combining the
three cases (Step 1, and the two cases of \eqref{eq:lz-balance}), we conclude
that
\[
 q\le C_1(\bar v+\varepsilon^{1/6}),\qquad
 C_1=\max\bigl\{16/k_0,(c')^{-1/6},(c')^{-1/4}((C_0+k_0)^{1/4}+1)\bigr\}.
\]
Recalling that $Q_{\gamma_*}(v)\le q/T_1$ and that
$\varepsilon^{1/6}=\min\{1,(\eta/\rho)^{1/6}\}$ for $\eta<\rho$, we obtain
\eqref{eq:lz-quantile} with $C=\max\{1,C_1/T_1\}$. This constant depends only
on $\beta$ (through $c_{\rm cert}$), $k_0$, $C_0$, and $T_1$.
\end{proof}

We comment on the exponent $1/6$ in \eqref{eq:lz-quantile}. In the setting of
the proof above, the variation in the proof of \cref{lem:lz-certificate}
moves mass toward the middle of $I$ at a cost of at most $\rho\tau/2$, which
is of order $\eta/q$, and the excesses at the endpoints of the perturbed
segment are at most of order $\eta/q$. The response integrated against the
bump is at least $4\ell/15=2q/15$. By \eqref{eq:lz-certificate-curvature}, the
second derivative along the perturbed segment is bounded below by a multiple
of $\rho\min\{q^5,q^3/\bar d\}$, where $\bar d$ is the density bound in the
proof of \cref{lem:lz-certificate}. Comparing this lower bound with the excesses at the endpoints, we obtain
$\eta/\rho\gtrsim_\beta\min\{q^6,q^4/\bar d\}$, which leads to
\eqref{eq:lz-balance}. Its first alternative gives the exponent $1/6$. In the
other alternative, the choice of the width $h$ controls the density $m/h$
without losing the linear dependence of the bound on $v$.

\begin{corollary}[Quantile bound in probability normalization]
\label{cor:lz-omega}
Let $N\ge1$, $0<\rho\le1$, and $0<\lambda\le\lambda_0$, and let
$s=\lambda/\beta$. Let $\mu=s\delta_0+\nu$ be a minimizer of
$\calF_{1-\rho,s}$ in \eqref{eq:fnd-functional}, so that
$\gamma_*=\beta\mu([0,\cdot])$ minimizes $F_\rho$, and let $\eta\ge0$ be a
trial budget for $\gamma_*$. Let
\[
 Q(v)=\inf\{q\in[0,1]:\nu([0,q])\ge v\}\qquad(0<v<1-s)
\]
be the quantile function \eqref{eq:fnd-quantile} of $\nu$, and define
\begin{equation}\label{eq:lz-omega}
 \omega(\rho,\eta)=\min\{1,(\eta/\rho)^{1/6}\}.
\end{equation}
Then, with $C_Q=\max\{1,\beta C\}$ and $C$ the constant of
\cref{prop:lz-quantile},
\begin{equation}\label{eq:lz-Q}
 Q(v)\le C_Q\bigl(s+v+\omega(\rho,\eta)\bigr)
 \qquad(0<v<1-s).
\end{equation}
\end{corollary}

\begin{proof}
The free measure of $\gamma_*$ is $\nu_{\gamma_*}=\beta\nu$, and
$\nu_{\gamma_*}([0,q])\ge\beta v$ if and only if $\nu([0,q])\ge v$. Hence
$Q(v)=Q_{\gamma_*}(\beta v)$, and $0<\beta v<\beta-\lambda$ for
$0<v<1-s$. By
\cref{prop:lz-quantile} and $\lambda=\beta s$, we have
\[
 Q(v)\le C\bigl[\beta(s+v)+\omega(\rho,\eta)\bigr]\le\beta C\bigl(s+v+\omega(\rho,\eta)\bigr),
\]
where we used $\beta>1$. We obtain \eqref{eq:lz-Q}, because $\beta C\le C_Q$.
\end{proof}

\chapter{Scalar response and stability}\label{sec:sc}\label{sec:scalar-stability}

In this section, we study the scalar Parisi equation of a single spin
under an arbitrary profile and the scalar response operator
$J^{\scal}_{a,d}$ built from it. The cavity interpolation of
\cref{sec:c7} decouples one spin from the others. At the independent
endpoint of that interpolation, the decoupled spin follows the scalar
law of \cref{sec:sc-law} with coefficient $d=\beta^2(1-t/N)$, and in
\cref{sec:c7-endpoint}, we express the derivative of the cavity covariance
at this endpoint through $J^{\scal}_{a,d}$, applied to the
covariance of the rest self-overlap. The contraction argument of
\cref{sec:c7-closure} needs the order bound
$J^{\scal}_{a,d}\preceq(1+\varepsilon_N)I$ with
$\varepsilon_N\le\rho/4$ at a minimizing profile. The goal of this
section is to prove this bound with
$\varepsilon_N\le C_{\rm st}[\rho^{-3/4}\eta^{1/4}+\lambda^2+N^{-1}]$,
where $\eta$ is a trial budget (\cref{prop:lz-stability}).

The argument has four steps.
\begin{enumerate}
\item We show in \cref{lem:sc-response} that, for every profile,
 $J_{a,d}$ is the derivative of the mean square magnetization
 $\Gamma_{a,d}$ when the covariance levels are transported by
 $q\mapsto q+\theta\varphi(q)$.
\item The scalar Parisi functional is convex. We compare its second
 derivative along a transport of levels with the lower bound given by
 convexity and the first variation, and we find that $J^{\scal}_{a,d}$
 is self-adjoint and that it is dominated by multiplication by the
 multiplier $k_{a,d}(q)=d\,\E^0h_q^2$, for every profile
 (\cref{lem:Jdom}).
\item At the scalar Parisi measure, we have $k=1$ on
 $\supp\mu_\beta=[0,q_\beta]$ and $k(q)\le e^{-\beta^2(q-q_\beta)/2}$
 above it (\cref{lem:reference}). On the support, this follows from an
 identity of Talagrand. Above the support, we use the log-concavity of
 a transformed density of the scalar diffusion and its heat evolution,
 both from \cite{Lopatto}.
\item The multiplier depends Lipschitz continuously on the profile in
 $L^1$ and on $d$ (\cref{lem:scalarcontinuity}). By
 \cref{prop:loc-isolation}, a finite-dimensional minimizer is close in
 $L^1$ to the scalar trial profile, which is close to
 $\beta\alpha_\beta$. Combining these facts with the first three steps,
 we obtain \cref{prop:lz-stability}.
\end{enumerate}
The bound is marginal: at the Parisi measure, $k=1$ on the support,
where the multiplier gives no contraction. The contraction in
\cref{sec:c7-closure} comes instead from the factor $t=1-\rho$, and it
needs only $\varepsilon_N\le\rho/4$. In the physics literature, De
Dominicis and Kondor~\cite[pp.~606--608]{DeDominicisKondor1983} found
that Parisi's solution with replica symmetry breaking is marginally stable
in a truncated replica model near the critical temperature. The
scalar identity on the support that
expresses this marginality appears in
\cite[Eqs.~(22)--(23)]{CrisantiRizzo2002}.

Sommers and Dupont gave an earlier variational formulation of the
coupled equations for the local magnetization and the distribution
of frozen fields~\cite[Section~2, Eqs.~(1)--(7), and
Eq.~(33)]{SommersDupont1984}; see also
\cite[Section~III, (14)--(18), (22)--(23)]{CrisantiRizzo2002}.
Those equations use the physical field and the order-parameter function
$q(x)$. Here we use the covariance level $q$ and the integrated
formulation in terms of probability measures described below.

We use the following results from earlier sections: the generalized
law and \cref{lem:fnd-posterior} of \cref{sec:gauss} for the basic
properties of the scalar law, \cref{lem:fnd-path-heat} for the
differentiation in \cref{lem:sc-response}, and, from \cref{sec:lz}, the
trial budget \eqref{eq:loc-trialbudget}, the root growth
\eqref{eq:lz-gsroot}, the description \eqref{eq:loc-scalar-measure} of
$\mu_\beta$, and \cref{prop:loc-isolation}. The results of this section
are used later as follows. The scalar law, described by \eqref{eq:sc-pde} and
\eqref{eq:sc-joint-law}, and the kernel bounds of \cref{lem:cl-scalar}
enter the cavity expansion in \cref{sec:c7,sec:c7-endpoint}. The
operator \eqref{eq:scalarJexplicit} and its kernel form
\eqref{eq:sc-kernel-form} are used in \cref{lem:linear}, \cref{lem:Jdom}
in \cref{prop:c7-closure}, and \cref{prop:lz-stability} in the proof of
\cref{thm:main} in \cref{sec:b7}.

\paragraph{Standing assumptions and notation.}
Throughout this section, $\beta>1$ is fixed and
$d\in[\beta^2/2,\beta^2]$. A \emph{profile} is the distribution function
$a(q)=\mu([0,q])$ of a Borel probability measure $\mu$ on $[0,1]$; it is
nondecreasing and right-continuous, and $a(1)=1$. We write
$\mu=s\delta_0+\nu$ with $s\ge0$, where the measure $\nu\ge0$ may charge
$0$. The scalar law depends only on $a$. The decomposition enters only
through $\nu$, which appears in the response operator. A \emph{step
profile} is a profile with finitely many values. Integrals over
$[0,q]$ include an atom at $q$, and integrals over $(q,1]$ exclude it.
Unless stated otherwise, $C_\beta$ denotes a constant depending only on
$\beta$, whose value may change from line to line. It does not depend
on $d\in[\beta^2/2,\beta^2]$, on $s$, or on the profile, including the
positions and masses of its atoms. No profile is assumed to minimize a
functional until \cref{sec:sc-minimizer}.

\section{The scalar law}\label{sec:sc-law}

The scalar Parisi equation and its diffusion are
\begin{equation}\label{eq:sc-pde}
 \begin{aligned}
 \partial_q u&=-\frac d2\bigl(\partial_{xx}u+a(q)(\partial_xu)^2\bigr),
 &u(1,x)&=\log(2\cosh x),\\
 \dd Y_q&=d\,a(q)\partial_xu(q,Y_q)\dd q+\sqrt d\dd B_q,
 &Y_0&=0,
 \end{aligned}
\end{equation}
where $B$ is a standard Brownian motion. We write $u=u_{a,d}$ when the
parameters need to be shown. Subscripts $x$ denote spatial derivatives,
and $u_j=\partial_x^ju$. When $a=\alpha_\beta$ and
$d=\beta^2$, the value $u(0,0)$ is the first term of the scalar Parisi
functional \eqref{eq:intro-parisi} at its minimizer.

\paragraph{Solutions and regularity.}
For a step profile, \eqref{eq:sc-pde} has a classical solution on each
interval where $a$ is constant. For a general profile, $u$ is the
uniform limit of the solutions for step profiles whose measures
converge weakly to $\mu$ \cite[Section~2.1]{Lopatto}. By
\cite[Proposition~2.1]{Lopatto}, applied with $\beta^2$ replaced by $d$,
the function $u(q,\cdot)$ is even, its spatial derivatives of every
positive order are bounded and jointly continuous on $[0,1]\times\R$,
and the spatial derivatives of the approximating solutions converge
uniformly \cite[(2.2)]{Lopatto}. The terminal value in \cite{Lopatto}
is $\log\cosh x$; it differs from ours by the constant $\log2$, which
shifts $u$ by $\log2$ and does not change its spatial derivatives. The
distribution functions of the approximating measures converge to $a$
at every continuity point of $a$, that is, at all but countably many
points. Passing to the limit in the equations integrated in the level
variable, we find that
\begin{equation}\label{eq:sc-pde-integrated}
 u(q,x)-u(r,x)=\frac d2\int_q^r\bigl(u_{xx}+a\,u_x^2\bigr)(v,x)\dd v,
 \qquad 0\le q\le r\le1,\ x\in\R .
\end{equation}
In particular, $u$ is Lipschitz in $q$, since the integrand is bounded.
For the weak space--time interpretation of the Parisi equation for
general measures, see \cite[Theorem~4]{JagannathTobasco2016}.
At an atom of $\mu$, the solution need not be differentiable in $q$. We
read \eqref{eq:sc-pde} and every other parabolic equation in this
section in this integrated form, including the linearized equations
derived below, and we never differentiate $a$.

\paragraph{Linear backward equations.}
Several proofs below use linear equations of the form
\begin{equation}\label{eq:sc-linear}
 \partial_qw+\frac{c(q)}2\,w_{xx}+b\,w_x+V\,w=-F,
\end{equation}
where $c$ is measurable with values in $[\beta^2/4,2\beta^2]$, and
$b$, $V$, and $F$ are bounded functions of $(q,x)$ that are measurable
in $q$ and have bounded continuous spatial derivatives. We use
\eqref{eq:sc-linear} only for functions $w$ that are Lipschitz in $q$
and have bounded continuous spatial derivatives of every order. We
read it in integrated form, as in \eqref{eq:sc-pde-integrated}. Fix
$(q,x)$, and let $Z$ solve $\dd Z_r=b(r,Z_r)\dd r+\sqrt{c(r)}\dd B_r$
for $r\ge q$, with $Z_q=x$. Itô's formula applies to
$r\mapsto e^{\int_q^rV(v,Z_v)\dd v}w(r,Z_r)$, with the derivative of
$w$ in $q$ taken almost everywhere; to see this, one may mollify $w$ in
the level variable and pass to the limit, using that the derivatives
are bounded. Taking expectations, we obtain the Feynman--Kac
representation
\begin{equation}\label{eq:sc-backward}
 w(q,x)=\E\Bigl[e^{\int_q^1V(v,Z_v)\dd v}\,w(1,Z_1)
   +\int_q^1e^{\int_q^rV(v,Z_v)\dd v}F(r,Z_r)\dd r\Bigr].
\end{equation}
In particular, if $w(1,\cdot)=0$, then
$\sup_x|w(q,x)|\le e^{\|V\|_\infty}\int_q^1\sup_x|F(r,x)|\dd r$.

\paragraph{Notation.}
We write $\Pp^0$ and $\E^0$ for probability and expectation under the
law of $Y$, together with a terminal spin $\epsilon\in\{-1,1\}$, which
is sampled given $Y_1$ with
\[
 \Pp^0(\epsilon=\pm1\mid Y_1)=\frac{e^{\pm Y_1}}{2\cosh Y_1}
 =e^{\pm Y_1-u(1,Y_1)}.
\]
Let $\calG^0_q$ be the filtration of $Y$; the terminal spin is not
$\calG^0_q$-measurable. We define
\[
 m_q=u_x(q,Y_q),\qquad h_q=u_{xx}(q,Y_q),\qquad
 \Gamma_{a,d}(q)=\E^0m_q^2,\qquad k_{a,d}(q)=d\,\E^0h_q^2 .
\]
We call $k_{a,d}$ the \emph{scalar multiplier}. For $0\le r\le\ell\le1$,
let $P^0_{r,\ell}\psi(y)=\E^0[\psi(Y_\ell)\mid Y_r=y]$ be the transition
operator of the Markov process $Y$, and set
\[
 f_p=m_p^2,\qquad \bar f_p=f_p-\E^0f_p,\qquad
 g_r^\ell(y)=P^0_{r,\ell}\bigl(u_x(\ell,\cdot)^2\bigr)(y),
\]
so that $g_r^\ell(Y_r)=\E^0[m_\ell^2\mid\calG_r^0]$. When we need the
diffusion started from $y$ at level $r$, we write $Y^{r,y}$ and
$\E_{r,y}$.

\begin{lemma}[Basic properties of the scalar law]\label{lem:sc-basic}
Let $a$ be a profile and $d\in[\beta^2/2,\beta^2]$.
\begin{enumerate}
\item[(a)] On $[0,1]\times\R$, we have
\begin{equation}\label{eq:sc-basic-bounds}
 |\partial_xu|\le1,\qquad
 0\le\partial_{xx}u\le1-(\partial_xu)^2\le1 .
\end{equation}
\item[(b)] For every $j\ge1$, there exists a constant $C_j(\beta)$,
depending only on $\beta$ and $j$, such that $|u_j|\le C_j(\beta)$ on
$[0,1]\times\R$. Further,
$\sup_x|u_{xx}(q,x)-u_{xx}(r,x)|\le C_\beta|q-r|$ for all
$q,r\in[0,1]$.
\item[(c)] We have
\begin{equation}\label{eq:sc-martingale}
 m_q=\E^0[\epsilon\mid\calG_q^0],\qquad m_0=0,\qquad
 \dd m_q=\sqrt d\,h_q\dd B_q,\qquad \Gamma_{a,d}'(q)=k_{a,d}(q).
\end{equation}
The function $k_{a,d}$ is continuous on $[0,1]$, and $\Gamma_{a,d}$ is
continuously differentiable on $[0,1]$, with one-sided derivatives at
the endpoints.
\end{enumerate}
\end{lemma}

\begin{proof}
We first identify the scalar law as a generalized law. We consider the
generalized law of \cref{def:gauss-law} with $N=1$, $\Lambda=d$, all
coupling coefficients equal to zero, root mass $s=0$, and the
probability measure $\mu$ in place of $\nu$; this choice is allowed for
every probability measure. The terminal value in
\eqref{eq:fnd-general-pde} is then
$\log\sum_{\sigma=\pm1}e^{x\sigma}=\log(2\cosh x)$, the Gaussian
disorder does not enter, and the outer density is one. Hence the
generalized law is the law of $(Y,\epsilon)$; it depends only on $a$.
For one spin, the conditional covariance is $C_q=1-m_q^2$, and
$H_q=h_q$.
\Cref{lem:fnd-posterior} implies that $m_q=\E^0[\epsilon\mid\calG^0_q]$,
that $m_0=0$, and that, for every $q\in[0,1]$, almost surely,
\begin{equation}\label{eq:sc-posterior-as}
 a(q)\bigl(1-m_q^2\bigr)\le h_q\le1-m_q^2 .
\end{equation}
Differentiating \eqref{eq:sc-pde-integrated} in $x$, we see that $u_x$
satisfies the integrated form of
$\partial_qu_x=-(d/2)(u_{xxx}+2a\,u_xu_{xx})$. Using Itô's formula
along $Y$, we find that the drift terms cancel and that
$\dd m_q=\sqrt d\,h_q\dd B_q$. This is the case $N=1$ of
\eqref{eq:gauss-martingale}.

\emph{Proof of (a).} Fix $q\in(0,1]$. Since the drift of $Y$ is bounded
by $d$, Girsanov's theorem implies that the law of $Y_q$ is equivalent
to the centered Gaussian law with variance $dq$. In particular, it
charges every nonempty open interval. The functions
\[
 u_{xx}(q,\cdot)-a(q)\bigl(1-u_x(q,\cdot)^2\bigr),\qquad
 1-u_x(q,\cdot)^2-u_{xx}(q,\cdot),\qquad
 1-u_x(q,\cdot)^2
\]
are continuous, and by \eqref{eq:sc-posterior-as} and $|m_q|\le1$, each
of them is nonnegative at $Y_q$ almost surely. A continuous function
with this property is nonnegative at every $x$. This proves \eqref{eq:sc-basic-bounds} for
$q>0$, since $a(q)(1-u_x^2)\ge0$. Letting $q\downarrow0$ and using the
joint continuity of $u_x$ and $u_{xx}$, we see that
\eqref{eq:sc-basic-bounds} also holds at $q=0$. (With the right continuity of
$a$, the lower bound $u_{xx}\ge a(q)(1-u_x^2)$ also holds everywhere.)

\emph{Proof of (b).} We differentiate the equation and use the
Feynman--Kac representation, following the method of
\cite[Lemma~16]{JagannathTobasco2016}. By \eqref{eq:sc-basic-bounds}, we have
$|u_1|\le1$ and $|u_2|\le1$. For $j\ge3$, the Leibniz rule gives
\[
 \partial_x^j(u_x^2)=2u_1u_{j+1}+2j\,u_2u_j
 +\sum_{i=2}^{j-2}\binom ji u_{i+1}u_{j-i+1},
\]
with an empty sum for $j=3$. Differentiating
\eqref{eq:sc-pde-integrated} $j$ times in $x$, we see that $u_j$ solves
\eqref{eq:sc-linear} with
\[
 c=d,\qquad b=d\,a\,u_1,\qquad V=j\,d\,a\,u_2,\qquad
 F=\frac{d\,a}2\sum_{i=2}^{j-2}\binom ji u_{i+1}u_{j-i+1},
\]
and terminal value $\partial_x^j\log(2\cosh x)$. Recall that the
derivatives of $u$ are bounded by \cite[Proposition~2.1]{Lopatto}; in
particular, the hypotheses of \eqref{eq:sc-backward} hold, with $Z=Y$.
By \eqref{eq:sc-basic-bounds}, $V$ takes values in $[0,j\beta^2]$; note
that the exponent in \eqref{eq:sc-backward} is nonnegative. Then
\eqref{eq:sc-backward} implies that
\[
 \|u_j\|_\infty\le e^{j\beta^2}\Bigl(
 \|\partial_x^j\log(2\cosh x)\|_\infty
 +\frac{\beta^2}2\sum_{i=2}^{j-2}\binom ji
          \|u_{i+1}\|_\infty\|u_{j-i+1}\|_\infty\Bigr).
\]
Every derivative on the right has order between $3$ and $j-1$. By
induction on $j$, we conclude that $\|u_j\|_\infty\le C_j(\beta)$ for
all $j\ge1$, uniformly over profiles and $d\in[\beta^2/2,\beta^2]$.
Differentiating \eqref{eq:sc-pde-integrated} twice in $x$, we find
\[
 u_{xx}(q,x)-u_{xx}(r,x)=\frac d2\int_q^r
 \Bigl[u_4+2a\bigl(u_2^2+u_1u_3\bigr)\Bigr](v,x)\dd v,
\]
and the integrand is bounded by a constant depending only on $\beta$.
This proves the Lipschitz bound in (b).

\emph{Proof of (c).} The first three identities in
\eqref{eq:sc-martingale} were shown above. By Itô's isometry and
$m_0=0$, we have
\[
 \Gamma_{a,d}(q)=\E^0m_q^2=d\int_0^q\E^0h_r^2\dd r
 =\int_0^qk_{a,d}(r)\dd r .
\]
By (b), we have $|h_q-h_r|\le C_\beta|q-r|+C_3(\beta)|Y_q-Y_r|$. Since
$Y$ has continuous paths and $|h|\le1$, dominated convergence implies
that $k_{a,d}$ is continuous. It follows that $\Gamma_{a,d}$ is
continuously differentiable on $[0,1]$, with $\Gamma_{a,d}'=k_{a,d}$.
\end{proof}

We also use the exact path likelihood of the scalar law, for every
profile, including profiles with atoms. It is the change of measure of
\cite[Lemma~8.5 and Corollary~8.7]{JagannathTobasco2017PhaseDiagram};
we include the derivation in order to record the conventions at atoms.
It extends the likelihood used for step profiles in the proof of
\cref{lem:fnd-posterior} to every profile, without approximation.

\begin{lemma}[Exact path likelihood]\label{lem:sc-likelihood}
Let $a$ be a profile and $d\in[\beta^2/2,\beta^2]$. Under a reference
law, let $X_q=\sqrt d\,B_q$, where $B$ is a standard Brownian motion.
For every $q\in[0,1]$, the law of $(Y_r)_{0\le r\le q}$ has density
\begin{equation}\label{eq:sc-prefix-law}
 \exp\Bigl\{a(q)u(q,X_q)-\int_{[0,q]}u(r,X_r)\,\mu(\dd r)\Bigr\}
\end{equation}
with respect to the law of $(X_r)_{0\le r\le q}$. The joint law of
$(Y,\epsilon)$ has density
\begin{equation}\label{eq:sc-joint-law}
 \exp\Bigl\{\epsilon X_1-su(0,0)-\int_{[0,1]}u(q,X_q)\,\nu(\dd q)\Bigr\}
\end{equation}
with respect to the product of the law of $X$ and counting measure on
$\{-1,1\}$. The same holds if the coefficient $d$ in \eqref{eq:sc-pde}
is replaced by a bounded measurable function $c(q)>0$ almost
everywhere, and $X_q=\int_0^q\sqrt{c(r)}\dd B_r$.
\end{lemma}

\begin{proof}
We write $\Upsilon_r=u(r,X_r)$. The drift $d\,a\,u_x$ of $Y$ is
bounded, and by \eqref{eq:sc-basic-bounds}, we have
\[
 \exp\Bigl\{\frac d2\int_0^qa(r)^2u_x(r,X_r)^2\dd r\Bigr\}\le e^{\beta^2/2},
\]
which is Novikov's condition. By Girsanov's theorem, the density of the law of $(Y_r)_{r\le q}$ with
respect to the law of $(X_r)_{r\le q}$ is
\[
 \exp\Bigl\{\int_0^qa(r)u_x(r,X_r)\dd X_r
 -\frac d2\int_0^qa(r)^2u_x(r,X_r)^2\dd r\Bigr\}.
\]
Under the reference law, using Itô's formula and
\eqref{eq:sc-pde-integrated}, we have
\[
 \dd\Upsilon_r=u_x(r,X_r)\dd X_r-\frac d2a(r)u_x(r,X_r)^2\dd r .
\]
The exponent therefore equals $\int_0^qa(r)\dd\Upsilon_r$. The function $a$ is
right-continuous and of bounded variation, with $\dd a=\mu$ on
$(0,1]$, and $\Upsilon$ is a continuous semimartingale. Integrating by
parts, we have
\[
 a(q)\Upsilon_q-a(0)\Upsilon_0
 =\int_{(0,q]}a(r-)\dd\Upsilon_r+\int_{(0,q]}\Upsilon_r\,\mu(\dd r).
\]
The functions $a(r-)$ and $a(r)$ differ only at the countably many
atoms of $\mu$, which affects neither the stochastic integral nor the
Lebesgue integral in $\dd\Upsilon_r$. Since $X_0=0$ and
$a(0)=\mu(\{0\})$, we have
$a(0)\Upsilon_0=\int_{\{0\}}u(r,X_r)\,\mu(\dd r)$. Inserting these facts
into the previous display, we obtain
\[
 \int_0^qa(r)\dd\Upsilon_r=a(q)u(q,X_q)-\int_{[0,q]}u(r,X_r)\,\mu(\dd r).
\]
This proves \eqref{eq:sc-prefix-law}. At an atom of $\mu$ at $q$, the
jump of the first term is canceled by the same atom in the integral,
and at $q=0$ the density equals one.

Given the path, the terminal spin has probability
$e^{\epsilon Y_1-u(1,Y_1)}$ relative to counting measure. Multiplying
\eqref{eq:sc-prefix-law} at $q=1$ by this factor, and using $a(1)=1$
and $\mu=s\delta_0+\nu$, we obtain \eqref{eq:sc-joint-law}. No further
normalization is needed, since both factors are probability densities.

For a coefficient $c(q)$ that is positive almost everywhere, we first
choose the everywhere positive representative
$c^+(q)=c(q)$ when $c(q)>0$ and $c^+(q)=1$ otherwise. It is measurable
and bounded and equals $c$ almost everywhere. This change preserves
every prefix clock, the integrated equation and drift, and the
stochastic integral defining $X$. Thus the law and the likelihood
are independent of this choice of representative.
Let $C(q)=\int_0^qc(r)\dd r$. The substitution
$\tau=C(q)/C(1)$ turns the equation with coefficient $c(q)$ into
\eqref{eq:sc-pde} with coefficient $C(1)$ and the profile of the image
of $\mu$ under $q\mapsto C(q)/C(1)$. This substitution is made in the
integrated equation, using $\dd\tau=c(q)\dd q/C(1)$. Since $c>0$ almost
everywhere, $q\mapsto C(q)/C(1)$ is an increasing homeomorphism of
$[0,1]$; no positive lower bound on $c$ is required. By
\cite[Proposition~2.1]{Lopatto}, applied with coefficient $C(1)$, the
solution satisfies $|u_x|\le1$ and has bounded spatial
derivatives. The argument above then applies with $d\dd r$ replaced by
$c(r)\dd r$; in Novikov's condition, $\beta^2$ is replaced by an
upper bound for $c$.
\end{proof}

\section{Continuity and kernel bounds}
\label{sec:sc-kernels}

The proof of the next lemma follows the comparison argument of
\cite[Lemma~14]{JagannathTobasco2016}, which bounds the differences of
the solutions and of their first spatial derivatives by the $L^1$
distance of the profiles. Here we also allow $d$ to vary and compare
the second derivatives along the diffusion.

\begin{lemma}[Continuity of the scalar multiplier]\label{lem:scalarcontinuity}
There exists $C_\beta<\infty$ such that, for any two profiles
$a,\widetilde a$ and any $d,\widetilde d\in[\beta^2/2,\beta^2]$,
\begin{equation}\label{eq:scalarcontinuity}
 \sup_{0\le q\le1}|k_{a,d}(q)-k_{\widetilde a,\widetilde d}(q)|
 \le C_\beta\bigl(\|a-\widetilde a\|_{L^1(0,1)}+|d-\widetilde d|\bigr).
\end{equation}
\end{lemma}

\begin{proof}
Let $u$ and $\widetilde u$ solve \eqref{eq:sc-pde} for the two
parameter pairs, and let $Y$ and $\widetilde Y$ be the corresponding
diffusions. We set $\hat u=u-\widetilde u$ and
$\Delta=\|a-\widetilde a\|_{L^1(0,1)}+|d-\widetilde d|$. The proof has
three steps. We first bound $\hat u_x$ and $\hat u_{xx}$ by
$C_\beta\Delta$, then couple the two diffusions, and finally compare
the multipliers.

\paragraph{Step 1: the difference of the solutions.}
Subtracting the two equations and writing
$u_x^2-\widetilde u_x^2=(u_x+\widetilde u_x)\hat u_x$, we find that
$\hat u$ solves \eqref{eq:sc-linear} with
\[
 c=d,\qquad b=\omega,\qquad V=0,\qquad F=\Xi,\qquad
 \hat u(1,\cdot)=0,
\]
where
\[
 \omega=\frac{d\,a(q)}2\bigl(u_x+\widetilde u_x\bigr),\qquad
 \Xi=\frac{d-\widetilde d}2\,\widetilde u_{xx}
     +\delta(q)\,\widetilde u_x^2,\qquad
 \delta(q)=\frac{d\,a(q)-\widetilde d\,\widetilde a(q)}2 .
\]
Using \cref{lem:sc-basic}(a) and~(b), we see that $0\le\omega_x\le\beta^2$ and
$|\omega_{xx}|\le C_\beta$. Further,
\[
 \Xi_x=\frac{d-\widetilde d}2\,\widetilde u_3
       +2\delta\,\widetilde u_1\widetilde u_2,\qquad
 \Xi_{xx}=\frac{d-\widetilde d}2\,\widetilde u_4
       +2\delta\bigl(\widetilde u_2^2+\widetilde u_1\widetilde u_3\bigr),
\]
and $2|\delta(q)|\le|d-\widetilde d|+\beta^2|a(q)-\widetilde a(q)|$.
Using \cref{lem:sc-basic}(b) again, we find
\[
 \int_0^1\sup_x\bigl(|\Xi_x|+|\Xi_{xx}|\bigr)(q,x)\dd q
 \le C_\beta\int_0^1\bigl(|d-\widetilde d|+|a(q)-\widetilde a(q)|\bigr)\dd q
 \le C_\beta\Delta .
\]
Differentiating the equation for $\hat u$ once and twice in $x$, we see
that $w_1=\hat u_x$ and $w_2=\hat u_{xx}$ solve \eqref{eq:sc-linear}
with $c=d$, $b=\omega$, zero terminal values, and
\[
 V=\omega_x,\quad F=\Xi_x\qquad\text{and}\qquad
 V=2\omega_x,\quad F=\Xi_{xx}+\omega_{xx}w_1,
\]
respectively. Then \eqref{eq:sc-backward} implies that
\[
 \begin{aligned}
 \|w_1\|_\infty&\le e^{\beta^2}\int_0^1\|\Xi_x(q,\cdot)\|_\infty\dd q
 \le C_\beta\Delta,\\
 \|w_2\|_\infty&\le e^{2\beta^2}\int_0^1\bigl(\|\Xi_{xx}(q,\cdot)\|_\infty
 +\|\omega_{xx}\|_\infty\|w_1\|_\infty\bigr)\dd q\le C_\beta\Delta .
 \end{aligned}
\]
This shows that
\begin{equation}\label{eq:sc-pde-distance}
 \sup_{q,x}\bigl(|u_x-\widetilde u_x|+|u_{xx}-\widetilde u_{xx}|\bigr)
 \le C_\beta\Delta .
\end{equation}

\paragraph{Step 2: coupling the diffusions.}
We drive $Y$ and $\widetilde Y$ by the same Brownian motion and set
$D(q)=\sup_{r\le q}|Y_r-\widetilde Y_r|$. By
\eqref{eq:sc-basic-bounds}, the drift $d\,a(r)u_x(r,\cdot)$ is
$\beta^2$-Lipschitz in space, and at a common point it differs from
$\widetilde d\,\widetilde a(r)\widetilde u_x(r,\cdot)$ by at most
$|d-\widetilde d|+\beta^2|a(r)-\widetilde a(r)|+\beta^2C_\beta\Delta$,
where we used \eqref{eq:sc-pde-distance}. Subtracting the integral
equations for $Y$ and $\widetilde Y$, we obtain
\[
 D(q)\le\beta^2\int_0^qD(r)\dd r+C_\beta\Delta
 +|\sqrt d-\sqrt{\widetilde d}\,|\sup_{r\le1}|B_r| .
\]
Since $d,\widetilde d\ge\beta^2/2$, we have
$|\sqrt d-\sqrt{\widetilde d}\,|\le|d-\widetilde d|/(\sqrt2\,\beta)$.
By Gronwall's inequality and $\E\sup_{r\le1}|B_r|\le2$, we conclude
that
\begin{equation}\label{eq:sc-coupling}
 \E\sup_{q\le1}|Y_q-\widetilde Y_q|
 \le C_\beta\bigl(\Delta+|\sqrt d-\sqrt{\widetilde d}\,|\,
                     \E\sup_{q\le1}|B_q|\bigr)\le C_\beta\Delta .
\end{equation}
The same argument, applied to the two diffusions started from a common
point $y$ at a common level $r$, yields the same bound, uniformly in
$(r,y)$.

\paragraph{Step 3: the multipliers.}
We write $h=u_{xx}(q,Y_q)$ and $\widetilde h=\widetilde u_{xx}(q,\widetilde Y_q)$.
Note that $|h^2-\widetilde h^2|\le2|h-\widetilde h|$, because
$0\le h,\widetilde h\le1$. Then
\[
 \begin{aligned}
 |k_{a,d}(q)-k_{\widetilde a,\widetilde d}(q)|
 &\le|d-\widetilde d|+2\beta^2\,\E|h-\widetilde h|\\
 &\le|d-\widetilde d|+2\beta^2\Bigl(\|u_{xx}-\widetilde u_{xx}\|_\infty
   +\|u_3\|_\infty\,\E|Y_q-\widetilde Y_q|\Bigr).
 \end{aligned}
\]
Inserting \eqref{eq:sc-pde-distance}, \eqref{eq:sc-coupling}, and the
bound $\|u_3\|_\infty\le C_3(\beta)$ into the last line, we obtain
\eqref{eq:scalarcontinuity}.
\end{proof}

The next lemma collects the bounds on the kernels of the response
operator that the cavity estimates of \cref{sec:c7-endpoint} use. Each
bound retains a factor $\Gamma$ at the relevant levels.

\begin{lemma}[Scalar kernels and their level weight]\label{lem:cl-scalar}
For every profile $a$ and $d\in[\beta^2/2,\beta^2]$, write
$\Gamma=\Gamma_{a,d}$. Then $\Gamma$ is nondecreasing and
\begin{equation}\label{eq:cl-Gamma-bound}
 0\le\Gamma(p)\le\min\{1,dp\}\le\min\{1,\beta^2p\}.
\end{equation}
For all $p,r,\ell\in[0,1]$,
\begin{align}
 |\E^0[m_ph_pm_r]|&\le\sqrt{\Gamma(p)\Gamma(r)},&
 |\E^0[\bar f_p\epsilon m_r]|&\le\sqrt{\Gamma(p)\Gamma(r)},
                                                        \label{eq:cl-k1}\\
 |\E^0[\bar f_pm_rm_\ell]|&\le\sqrt{\Gamma(p)\Gamma(r\wedge\ell)},&
 |\E^0[\bar f_pm_pm_r]|&\le\sqrt{\Gamma(p)\Gamma(p\wedge r)},
                                                        \label{eq:cl-k2}\\
 0\le\Var^0(f_p)&\le\Gamma(p),&
 0\le\E^0[m_p^2h_p]&\le\Gamma(p).                       \label{eq:cl-k3}
\end{align}
For $\ell\ge p$,
\begin{equation}\label{eq:cl-k4}
 |\Cov^0(f_p,g_p^\ell(Y_p))|\le\sqrt{\Gamma(p)\Gamma(\ell)}.
\end{equation}
\end{lemma}

\begin{proof}
By \cref{lem:sc-basic}(c) and \eqref{eq:sc-basic-bounds}, we have
$\Gamma(p)=d\int_0^p\E^0h_r^2\dd r$, where $0\le h_r\le1$. This implies
that $\Gamma$ is nondecreasing and that $0\le\Gamma(p)\le dp$. We also
have $\Gamma\le1$, because $|m_p|\le1$. This proves \eqref{eq:cl-Gamma-bound}. Using $|m_p|\le1$
again, we have
\[
 \Var^0(f_p)=\E^0\bar f_p^2\le\E^0m_p^4\le\E^0m_p^2=\Gamma(p).
\]
The bounds below follow from the Cauchy--Schwarz inequality, together
with $|h_p|\le1$, $|\epsilon|=1$, and $|m|\le1$:
\[
 \begin{aligned}
 |\E^0[m_ph_pm_r]|^2
 &\le\E^0[m_p^2h_p^2]\,\E^0m_r^2\le\Gamma(p)\Gamma(r),\\
 |\E^0[\bar f_p\epsilon m_r]|^2
 &\le\E^0\bar f_p^2\,\E^0m_r^2\le\Gamma(p)\Gamma(r),\\
 |\E^0[\bar f_pm_rm_\ell]|^2
 &\le\E^0\bar f_p^2\,\E^0[m_r^2m_\ell^2]\le\Gamma(p)\Gamma(r\wedge\ell).
 \end{aligned}
\]
In the last line, we used $m_r^2m_\ell^2\le m_{r\wedge\ell}^2$, which
holds because the factor at the smaller level is one of the two factors
and the other is at most one. This proves \eqref{eq:cl-k1} and the first
bound in \eqref{eq:cl-k2}; the second bound in \eqref{eq:cl-k2} is the
case $\ell=p$. The bounds \eqref{eq:cl-k3} follow from the estimate on
$\Var^0(f_p)$ and from $0\le h_p\le1$, which gives
$0\le\E^0[m_p^2h_p]\le\E^0m_p^2$. Finally, let $\ell\ge p$. We have
$0\le g_p^\ell\le1$ and $\E^0g_p^\ell(Y_p)=\E^0m_\ell^2=\Gamma(\ell)$,
which imply that
\[
 \Var^0\bigl(g_p^\ell(Y_p)\bigr)\le\E^0\bigl(g_p^\ell(Y_p)\bigr)^2
 \le\E^0g_p^\ell(Y_p)=\Gamma(\ell),
\]
and \eqref{eq:cl-k4} follows from the Cauchy--Schwarz inequality and
the bound on $\Var^0(f_p)$.
\end{proof}

\section{The response operator}\label{sec:sc-response}

For every profile, the operator defined below is the response of
$\Gamma$ to a transport of covariance levels (\cref{lem:sc-response}).
In \cite[Lemma~5.2, (H2)]{OpenAILimitingLaw}, a scalar response is
written as a diagonal multiplier minus a positive semidefinite Hessian
kernel. Here we identify the second variation of the scalar Parisi
functional along transports of levels, and we compare it using the
convexity of this functional along segments of distribution functions
\cite{AuffingerChen2015Uniqueness}. In \cref{sec:sc-domination}, this
shows that the scalar response operator is dominated by multiplication by
$k_{a,d}$ at every profile. The transports fix the levels $0$ and $1$,
and the operator acts on the \emph{scalar transport space}
\[
 \mathcal H_\nu=\bigl\{v\in L^2(\nu):
                 v(0)=0\text{ if }\nu(\{0\})>0,
                 \ v(1)=0\text{ if }\nu(\{1\})>0\bigr\}.
\]
This is a closed subspace of $L^2(\nu)$, and the orthogonal projection
onto it is multiplication by $\ind_{(0,1)}$.

The operator involves the derivative of $g_r^\ell$ in its argument.
Since the noise of $Y$ is additive, the diffusion $Y^{r,y}$ started
from $y$ at level $r$ solves an ordinary differential equation driven
by the Brownian path, and it is differentiable in $y$. Its derivative
is $\xi_\ell=\exp\{d\int_r^\ell a(v)u_{xx}(v,Y^{r,y}_v)\dd v\}$, which
lies in $[1,e^{\beta^2}]$ by \eqref{eq:sc-basic-bounds}.
Differentiating $g_r^\ell(y)=\E[u_x(\ell,Y^{r,y}_\ell)^2]$ under the
expectation, we find
\begin{equation}\label{eq:sc-dg}
 \partial_yg_r^\ell(y)=\E\bigl[2u_x(\ell,Y_\ell^{r,y})
 u_{xx}(\ell,Y_\ell^{r,y})\,\xi_\ell\bigr],\qquad
 |\partial_yg_r^\ell|\le2e^{\beta^2}.
\end{equation}

\begin{definition}[Response operator]\label{def:sc-response}
Let $a$ be a profile with $\mu=s\delta_0+\nu$, and let
$d\in[\beta^2/2,\beta^2]$. For a bounded Borel function $\varphi$ on
$[0,1]$ and $p\in[0,1]$, we define
\begin{align}
 \frac1d(J_{a,d}\varphi)(p)={}&
 \E^0\bigl[h_p^2+2(1-a(p))m_p^2h_p\bigr]\varphi(p)
 -\int_{(p,1]}\E^0\bigl[m_p\partial_y g_p^\ell(Y_p)\bigr]
                                      \varphi(\ell)\,\nu(\dd\ell)\nonumber\\
 &-2\int\E^0[m_ph_pm_r]\varphi(p\wedge r)\,\nu(\dd r)
  -\int\E^0[\bar f_p\epsilon m_r]\varphi(r)\,\nu(\dd r)\nonumber\\
 &+\frac12\int a(r)\E^0[\bar f_pm_r^2]\varphi(r)\,\nu(\dd r)
  +\frac12\iint\E^0[\bar f_pm_rm_\ell]\varphi(r\wedge\ell)
                                      \,\nu(\dd r)\nu(\dd\ell)\nonumber\\
 &+\frac12\int\nu(\dd r)\int_{(r,1]}
           \E^0\bigl[\bar f_pg_r^\ell(Y_r)\bigr]\varphi(\ell)\,\nu(\dd\ell).
 \label{eq:scalarJexplicit}
\end{align}
Integrals without a domain are over $[0,1]$. The same formula defines
$J_{a,d}v$ for $v\in L^2(\nu)$. The \emph{scalar response operator}
$J^{\scal}_{a,d}$ is the compression of $J_{a,d}$ to $\mathcal H_\nu$,
that is, $J^{\scal}_{a,d}v=\ind_{(0,1)}J_{a,d}v$ for $v\in\mathcal H_\nu$.
\end{definition}

At an atom of $\nu$ at $p$, the right-continuous value $a(p)$ includes
the atom, and the strict future interval $(p,1]$ excludes it. For
$v\in L^2(\nu)$, the formula involves $v$ only through $v(p)$ and
through values at points that are integrated against $\nu$, because
$p\wedge r\in\{p,r\}$ and $r\wedge\ell\in\{r,\ell\}$. We show in the
next paragraph that it defines $J_{a,d}v$ as an element of $L^2(\nu)$.

\paragraph{Kernel form and boundedness.}
We now rewrite \eqref{eq:scalarJexplicit}. For $r>p$, conditioning on
$\calG^0_p$ and using \eqref{eq:sc-martingale}, we have
$\E^0[m_ph_pm_r]=\E^0[m_p^2h_p]$, and $\nu((p,1])=1-a(p)$. Splitting the
third term of \eqref{eq:scalarJexplicit} at $r=p$, we see that its part
over $(p,1]$ equals $-2(1-a(p))\E^0[m_p^2h_p]\varphi(p)$, which cancels
the second part of the first coefficient. In the double integral, we
split $\varphi(r\wedge\ell)$ into the regions $r<\ell$, $\ell<r$, and the
diagonal $r=\ell$, which carries mass only at atoms of $\nu$; the first
two regions give equal contributions by symmetry. In the last term, we
exchange the order of integration. All integrands are bounded, and
Fubini's theorem justifies these steps. The result is
\begin{equation}\label{eq:sc-kernel-form}
 (J_{a,d}v)(p)=k_{a,d}(p)v(p)+d\int\mathcal K_{a,d}(p,r)v(r)\,\nu(\dd r),
\end{equation}
for every bounded Borel $v$, where
\begin{equation}\label{eq:sc-kernel}
 \begin{aligned}
 \mathcal K_{a,d}(p,r)={}&
 -\ind_{\{r>p\}}\E^0\bigl[m_p\partial_yg_p^r(Y_p)\bigr]
 -2\,\ind_{\{r\le p\}}\E^0[m_ph_pm_r]
 -\E^0[\bar f_p\epsilon m_r]\\
 &+\frac12a(r)\E^0[\bar f_pm_r^2]
 +\int_{(r,1]}\E^0[\bar f_pm_rm_\ell]\,\nu(\dd\ell)
 +\frac12\nu(\{r\})\E^0[\bar f_pm_r^2]\\
 &+\frac12\int_{[0,r)}\E^0\bigl[\bar f_pg_\ell^r(Y_\ell)\bigr]\,\nu(\dd\ell).
 \end{aligned}
\end{equation}
The term with $\nu(\{r\})$ is the diagonal part of the double integral.
The kernel $\mathcal K_{a,d}$ is a bounded Borel function. Indeed, the
first two terms are never both nonzero and are bounded by
$2e^{\beta^2}$ by \eqref{eq:sc-dg}; since $|\bar f_p|\le1$, $|m|\le1$, and $0\le g\le1$, the
third and fourth terms are bounded by $1$ and $1/2$, and the last three
terms together by $\nu([0,1])\le1$. We conclude that
$|\mathcal K_{a,d}|\le2e^{\beta^2}+3$. By the Cauchy--Schwarz
inequality,
$\|\int\mathcal K_{a,d}(\cdot,r)v(r)\,\nu(\dd r)\|_{L^2(\nu)}
\le(2e^{\beta^2}+3)\nu([0,1])\|v\|_{L^2(\nu)}$. It follows that
\eqref{eq:sc-kernel-form} extends to $v\in L^2(\nu)$, where it agrees
with \eqref{eq:scalarJexplicit}. Using $0\le k_{a,d}\le d\le\beta^2$,
we find $\|J_{a,d}\|_{L^2(\nu)\to L^2(\nu)}\le\beta^2(4+2e^{\beta^2})$.
Compression does not increase the norm, which gives
\begin{equation}\label{eq:sc-J-bounded}
 \|J^{\scal}_{a,d}\|_{\mathcal H_\nu\to\mathcal H_\nu}\le C_\beta .
\end{equation}

\paragraph{Transports of levels.}
Let $\varphi$ be the restriction to $[0,1]$ of a smooth function, with
$\varphi(0)=\varphi(1)=0$, and define $T_\theta(q)=q+\theta\varphi(q)$.
We fix $\theta_0\in(0,1]$ such that $\theta_0\|\varphi'\|_\infty\le1/2$.
For $|\theta|\le\theta_0$, the map $T_\theta$ is an
increasing bijection of $[0,1]$ with $T_\theta'\in[1/2,3/2]$. We replace
the covariance $d\,(q\wedge r)$ of $X$ in \eqref{eq:sc-joint-law} by
$d\,T_\theta(q\wedge r)$ while keeping $\nu$. Equivalently, we replace
the coefficient $d$ in \eqref{eq:sc-pde} by
\[
 c_\theta(q)=d\,T_\theta'(q)=d\bigl(1+\theta\varphi'(q)\bigr)
 \in\Bigl[\frac d2,\frac{3d}2\Bigr]\subset\Bigl[\frac{\beta^2}4,2\beta^2\Bigr],
\]
and keep $a$. We write $u_\theta$ and $Y^\theta$ for the resulting
solution and diffusion, $u_{\theta,x}=\partial_xu_\theta$ and similarly
for higher derivatives, and
\[
 \widetilde\Gamma_\theta(p)=\E\bigl[\partial_xu_\theta(p,Y^\theta_p)^2\bigr]
\]
for the resulting mean square magnetization. Let
$a_\theta=a\circ T_\theta^{-1}$ be the distribution function of the
image measure $T_{\theta\#}\mu$. Under the substitution $x'=T_\theta(q)$
in the level variable, the equation for $u_\theta$ becomes
\eqref{eq:sc-pde} with coefficient $d$ and profile $a_\theta$, and the
diffusion $Y^\theta$ becomes a time change of the diffusion for
$a_\theta$. In formulas,
\begin{equation}\label{eq:sc-transported}
 u_\theta(q,x)=u_{a_\theta,d}(T_\theta(q),x),\qquad
 \widetilde\Gamma_\theta(p)=\Gamma_{a_\theta,d}(T_\theta(p)),
\end{equation}
that is, $(u_\theta,Y^\theta)$ is the scalar law of $T_{\theta\#}\mu$
with coefficient $d$, observed at the transported levels. In
particular, $\widetilde\Gamma_0=\Gamma_{a,d}$, and by
\cref{lem:sc-basic}(b), the spatial derivatives of $u_\theta$ of each
order are bounded uniformly in $|\theta|\le\theta_0$ and over profiles.

\begin{lemma}[Response identity]\label{lem:sc-response}
Let $a$ be a profile, let $d\in[\beta^2/2,\beta^2]$, and let $\varphi$
and $\widetilde\Gamma_\theta$ be as above. Then, for every $p\in[0,1]$,
the function $\theta\mapsto\widetilde\Gamma_\theta(p)$ is continuously
differentiable near $\theta=0$, and
\[
 \frac{\dd}{\dd\theta}\widetilde\Gamma_\theta(p)\Big|_{\theta=0}
 =(J_{a,d}\varphi)(p).
\]
\end{lemma}

\begin{proof}
The proof has four steps. In Step 1, we show that $u_\theta$ and its
spatial derivatives are continuously differentiable in $\theta$; the
argument works with the integrated equations and never differentiates
$a$. In Step 2, we compute the parameter derivative of $u_\theta$ and
prove \eqref{eq:scalarNodeVariation}. In Step 3, we write
$\widetilde\Gamma_\theta(p)$ as a normalized expectation on a fixed
Gaussian path space and differentiate it with
\cref{lem:fnd-path-heat}. In Step 4, we match the resulting terms with
those of \eqref{eq:scalarJexplicit}. Throughout the proof,
$|\theta|\le\theta_0$, and constants may depend on $\beta$ and
$\varphi$.

\paragraph{Step 1: differentiability in the parameter.}
Fix $\theta$, and for small $\delta\ne0$ let $v=u_{\theta+\delta}$,
$u=u_\theta$, and $D_\delta=(v-u)/\delta$. Since
$c_{\theta+\delta}-c_\theta=\delta d\varphi'$, subtracting the two
equations and dividing by $\delta$, we find
\begin{equation}\label{eq:sc-Dh}
 \partial_qD_\delta+\frac{c_\theta}2(D_\delta)_{xx}
       +\frac{c_\theta a}2(v_x+u_x)(D_\delta)_x
 =-\frac{d\varphi'}2\bigl(v_{xx}+a\,v_x^2\bigr),\qquad D_\delta(1,\cdot)=0 .
\end{equation}
This is \eqref{eq:sc-linear} with $c=c_\theta$, $V=0$, and a drift and
a source whose spatial derivatives of each order are bounded uniformly
in $\delta$ and $\theta$. Differentiating \eqref{eq:sc-Dh} $j$ times in
$x$, we see that $\partial_x^jD_\delta$ solves an equation of the same
form, with zeroth-order coefficient $j$ times the spatial derivative
of the drift, and with a source that involves only derivatives of
$D_\delta$ of order less than $j$ and spatial derivatives of the drift
and of the source of \eqref{eq:sc-Dh}. By \eqref{eq:sc-backward} and
induction on $j$, we have
$\sup_{q,x}|\partial_x^jD_\delta|\le C_j$ for every $j\ge0$, uniformly
in $\delta$ and $\theta$. In particular, $v-u=\delta D_\delta$ and all
its spatial derivatives are $O(|\delta|)$.

Next, we write the drift in \eqref{eq:sc-Dh} as
$c_\theta a\,u_x+(c_\theta a/2)(v_x-u_x)$ and move the second part
to the right side. For two increments $\delta$ and $\delta'$, the
difference $D_\delta-D_{\delta'}$ then solves \eqref{eq:sc-linear} with
drift $c_\theta a\,u_x$, $V=0$, zero terminal value, and the source
\[
 \frac{d\varphi'}2\Bigl[(v-v')_{xx}+a\bigl(v_x^2-v_x'^2\bigr)\Bigr]
 +\frac{c_\theta a}2\Bigl[(v_x-u_x)(D_\delta)_x-(v'_x-u_x)(D_{\delta'})_x\Bigr],
\]
where $v'=u_{\theta+\delta'}$. By the bounds just proved, every spatial
derivative of this source is bounded by $C_j(|\delta|+|\delta'|)$. By
the same induction, we have
$\sup_{q,x}|\partial_x^j(D_\delta-D_{\delta'})|\le C_j(|\delta|+|\delta'|)$.
It follows that, as $\delta\to0$, the functions $D_\delta$ converge
uniformly, together with all their spatial derivatives, to a function
$\dot u_\theta$. Passing to the limit in the integrated form of
\eqref{eq:sc-Dh}, we find that $\dot u_\theta$ solves
\begin{equation}\label{eq:sc-W}
 \partial_q\dot u_\theta+\frac{c_\theta}2(\dot u_\theta)_{xx}
   +c_\theta a\,u_{\theta,x}(\dot u_\theta)_x
 =-\frac{d\varphi'}2\bigl(u_{\theta,xx}+a\,u_{\theta,x}^2\bigr),
 \qquad \dot u_\theta(1,\cdot)=0 .
\end{equation}
We conclude that $u_\theta$ and each of its spatial derivatives are
differentiable in $\theta$, with
$\partial_\theta\partial_x^ju_\theta=\partial_x^j\dot u_\theta$, and that
these derivatives are bounded uniformly in $(q,x)$ and $\theta$.
To prove that these derivatives depend continuously on $\theta$, we
set $b_\theta=c_\theta a\,u_{\theta,x}$,
$F_\theta=(d\varphi'/2)(u_{\theta,xx}+a\,u_{\theta,x}^2)$, and
$R=\dot u_\theta-\dot u_{\theta'}$. Subtracting \eqref{eq:sc-W} at the
two parameters, we obtain
\[
 \Bigl(\partial_q+\frac{c_\theta}2\partial_{xx}+b_\theta\partial_x\Bigr)R
 =-(F_\theta-F_{\theta'})
  -\frac{c_\theta-c_{\theta'}}2(\dot u_{\theta'})_{xx}
  -(b_\theta-b_{\theta'})(\dot u_{\theta'})_x,
 \qquad R(1,\cdot)=0 .
\]
Since $c_\theta-c_{\theta'}=(\theta-\theta')d\varphi'$, and since
$u_\theta-u_{\theta'}$ and all its spatial derivatives are
$O(|\theta-\theta'|)$ by the bound on $\partial_\theta\partial_x^ju_\theta$,
every fixed-order spatial derivative of the right side is
$O(|\theta-\theta'|)$. Using the same backward estimates, we have
$\sup_{q,x}|\partial_x^jR|\le C_j|\theta-\theta'|$, so
$\partial_x^j\dot u_\theta$ depends continuously on $\theta$, uniformly
in $(q,x)$. The same argument applies without
change to transports $q\mapsto q+\theta_1\varphi_1(q)+\theta_2\varphi_2(q)$
with several parameters, and gives continuous partial derivatives in
each parameter.

\paragraph{Step 2: the variation of the solution.}
We write $\dot u=\dot u_0$ for the derivative at $\theta=0$. By
\eqref{eq:sc-W} and \eqref{eq:sc-backward}, with $Z=Y$, we have
\begin{equation}\label{eq:sc-udot-rep}
 \dot u(q,y)=\frac d2\int_q^1\varphi'(r)\Psi_q(r;y)\dd r,\qquad
 \Psi_q(r;y)=\E_{q,y}\bigl[h_r+a(r)m_r^2\bigr],
\end{equation}
where $\E_{q,y}$ refers to the diffusion started from $y$ at level $q$.
To evaluate $\Psi_q(r;y)$, we use Itô's formula along $Y$. Differentiating
\eqref{eq:sc-pde-integrated} twice in $x$ and using Itô's formula for
$u_{xx}(r,Y_r)$, we find
\[
 \dd h_r=-d\,a(r)h_r^2\dd r+\sqrt d\,u_{xxx}(r,Y_r)\dd B_r,
\]
and by \eqref{eq:sc-martingale},
$\dd(m_r^2)=d\,h_r^2\dd r+2\sqrt d\,m_rh_r\dd B_r$. The product rule for
the continuous semimartingale $m_r^2$ and the deterministic function
$a$ of bounded variation gives
$\dd(a(r)m_r^2)=a(r-)\dd(m_r^2)+m_r^2\,\mu(\dd r)$ on $(0,1]$, and
$a(r-)=a(r)$ except at countably many $r$. Adding the two
differentials, we see that the drift terms $\mp d\,a(r)h_r^2\dd r$
cancel, and
\[
 \dd\bigl(h_r+a(r)m_r^2\bigr)
 =\sqrt d\bigl[u_{xxx}(r,Y_r)+2a(r)m_rh_r\bigr]\dd B_r+m_r^2\,\mu(\dd r).
\]
Taking conditional expectations given $Y_q=y$, and recalling that
$\mu=\nu$ on $(q,r]\subset(0,1]$, we obtain
\begin{equation}\label{eq:sc-Fq}
 \Psi_q(r;y)-\Psi_q(q;y)=\int_{(q,r]}g_q^v(y)\,\nu(\dd v),\qquad q\le r\le1,
\end{equation}
where $\Psi_q(q;y)=u_{xx}(q,y)+a(q)u_x(q,y)^2$.

We insert \eqref{eq:sc-Fq} into \eqref{eq:sc-udot-rep} and exchange the
order of integration over $q<v\le r\le1$. Note that
$\int_q^1\varphi'(r)\dd r=-\varphi(q)$ and
$\int_v^1\varphi'(r)\dd r=-\varphi(v)$, since $\varphi(1)=0$. We obtain
\begin{equation}\label{eq:scalarNodeVariation}
 \dot u(q,y)=-\frac d2\Bigl[
  \varphi(q)\bigl(u_{xx}(q,y)+a(q)u_x(q,y)^2\bigr)
  +\int_{(q,1]}\varphi(r)g_q^r(y)\,\nu(\dd r)\Bigr].
\end{equation}
An atom at $q$ belongs to $a(q)$ and not to the integral over $(q,1]$.
An atom at $1$ belongs to the integral, but its factor $\varphi(1)$
vanishes.

The same computation holds at every base point $\theta$, with $Y^\theta$
in place of $Y$. The diffusion coefficient $d$ is then replaced by
$c_\theta(r)$; the drift terms $\mp c_\theta(r)a(r)h_r^2\dd r$ still
cancel, and the source in \eqref{eq:sc-W} still carries the factor
$d\varphi'$, which is the derivative of $c_\theta$ in $\theta$. At
$(q,y)=(0,0)$, we have $\varphi(0)=0$ and
$\E[\partial_xu_\theta(r,Y^\theta_r)^2\mid Y^\theta_0=0]
=\widetilde\Gamma_\theta(r)$. This gives
\begin{equation}\label{eq:sc-root-variation}
 \frac{\dd}{\dd\theta}u_\theta(0,0)
 =-\frac d2\int\varphi(r)\widetilde\Gamma_\theta(r)\,\nu(\dd r),
 \qquad|\theta|<\theta_0 .
\end{equation}

\paragraph{Step 3: Gaussian differentiation at the fixed profile.}
Let $B$ be a standard Brownian motion, and under the reference law let
$X^\theta_q=B_{dT_\theta(q)}$. This Gaussian path has covariance
$K_\theta(q,r)=d\,T_\theta(q\wedge r)$ and the law of
$\int_0^q\sqrt{c_\theta(r)}\dd B_r$. Then \cref{lem:sc-likelihood},
applied with the coefficient $c_\theta$, implies that the joint law of $(Y^\theta,\epsilon)$
has density \eqref{eq:sc-joint-law} with $u_\theta$ in place of $u$,
relative to the law of $X^\theta$ and counting measure. The factor
$e^{-su_\theta(0,0)}$ does not depend on the path and cancels from
normalized expectations. With
\[
 \ell_\theta(X,\epsilon)=\epsilon X_1-\int u_\theta(r,X_r)\,\nu(\dd r),
 \qquad\Phi_\theta(X)=u_{\theta,x}(p,X_p)^2,
\]
we have
$\widetilde\Gamma_\theta(p)=\E_\theta\Phi_\theta$, where $\E_\theta$ is
the normalized expectation with weight $W_\theta=e^{\ell_\theta}$ in
the notation of \cref{lem:fnd-path-heat}. At $\theta=0$, the path $X$
under $\E_0$ has the law of $Y$. As in \cref{sec:sc-law}, we write
$m_r=u_x(r,X_r)$ and $h_r=u_{xx}(r,X_r)$, and we also write
$u_{3,r}=u_{xxx}(r,X_r)$.

We check that the assumptions (A1)--(A4) of \cref{lem:fnd-path-heat}
hold, with $z=\theta$ in the compact interval
$I=[-\theta_0/2,\theta_0/2]$. The Gaussian coordinates consist of one
field path on one edge, with the deterministic clock $dT_\theta(q)$,
which is continuous in $(q,\theta)$, continuously differentiable in
$\theta$, and strictly increasing in $q$; there are no finite
coordinates, and the terminal spin $\epsilon$ is the counting variable.
This is (A1). The kernels $K_\theta$ and
$\dot K_\theta(q,r)=d\varphi(q\wedge r)$ are jointly continuous in
$(\theta,q,r)$, which is (A2). In the notation of
\cref{lem:fnd-path-heat}, where $M$ and $Q$ denote derivative measures,
the derivative measures of $\ell_\theta$ and $\Phi_\theta$ are
\[
 \begin{aligned}
 S&=\epsilon\delta_1-u_{\theta,x}(r,X_r)\,\nu(\dd r),&\qquad
 M&=-\iota_\#\bigl(u_{\theta,xx}(r,X_r)\,\nu(\dd r)\bigr),\\
 P&=2u_{\theta,x}u_{\theta,xx}(p,X_p)\,\delta_p,&\qquad
 Q&=2\bigl(u_{\theta,xx}^2+u_{\theta,x}u_{\theta,xxx}\bigr)(p,X_p)\,
   \delta_{(p,p)},
 \end{aligned}
\]
where $\iota(r)=(r,r)$. The measure $M$ is carried by the diagonal,
with an atom at $(r,r)$ for each atom $r$ of $\nu$, as in
\cref{rem:gauss-local}. Using \cref{lem:sc-basic}(b), we see that these
measures have
total variation bounded uniformly in $\theta$ and $X$, and by Step 1,
they are continuous in $(\theta,X)$ in total variation, where $X$
carries the supremum norm. The maps $X\mapsto\ell_\theta(X,\epsilon)$
and $X\mapsto\Phi_\theta(X)$ are twice continuously differentiable,
since $u_\theta$ has bounded continuous spatial derivatives. This is
(A4). The observable $\Phi_\theta$ takes values in $[0,1]$, and the
explicit parameter derivatives
\[
 \dot\Phi_\theta=2u_{\theta,x}(p,X_p)\,\partial_x\dot u_\theta(p,X_p),\qquad
 \dot\ell_\theta=-\int\dot u_\theta(r,X_r)\,\nu(\dd r)
\]
are bounded and jointly continuous in $(\theta,X)$ by Step 1. Since
$|u_{\theta,x}|\le1$ and $u_\theta(r,0)$ is bounded, we have
$|\ell_\theta|\le C+2\sup_r|X_r|$, which is the linear growth required
in (A3). We conclude that \cref{lem:fnd-path-heat} applies to
$\E_\theta\Phi_\theta$. By \cref{rem:gauss-local}, the right side of
\eqref{eq:fnd-path-heat} is continuous in $\theta$, which proves that
$\theta\mapsto\widetilde\Gamma_\theta(p)$ is continuously differentiable
on $(-\theta_0/2,\theta_0/2)$. Each term on the right side is bounded
by a constant depending only on $\beta$ and $\varphi$; in particular,
the derivative is bounded uniformly in $p\in[0,1]$ and $\theta$. We use
this uniform bound in the proof of \cref{lem:Jdom}.

\paragraph{Step 4: the terms of the identity.}
We evaluate the right side of \eqref{eq:fnd-path-heat} at $\theta=0$
and divide it by $d$. For a signed measure $R$ on $[0,1]^2$, we have
$d^{-1}\langle\dot K_0,R\rangle=\iint\varphi(q\wedge r)\,R(\dd q,\dd r)$.
Note that, under $\E_0$, the centered observable is
$\Phi_0-\E_0\Phi_0=\bar f_p$.
We group the contributions as follows.
\begin{enumerate}
\item The second derivative of the observable contributes
 \[
  \mathrm I_1=\frac1{2d}\E_0\langle\dot K_0,Q\rangle
  =\varphi(p)\E^0\bigl[h_p^2+m_pu_{3,p}\bigr].
 \]
\item The explicit derivative of the observable is
 $\dot\Phi_0=2m_p\partial_y\dot u(p,X_p)$. Differentiating
 \eqref{eq:scalarNodeVariation} in $y$ under the integral, which is
 allowed by \eqref{eq:sc-dg}, we find
 \[
  \mathrm I_2=\frac1d\E_0\dot\Phi_0
  =-\varphi(p)\E^0\bigl[m_pu_{3,p}+2a(p)m_p^2h_p\bigr]
  -\int_{(p,1]}\varphi(\ell)\E^0\bigl[m_p\partial_yg_p^\ell(Y_p)\bigr]
   \,\nu(\dd\ell).
 \]
\item Since $\dot K_0$ is symmetric and $\dot K_0(1,p)=d\varphi(p)$,
 the score--observable terms give
 \[
  \begin{aligned}
  \mathrm I_3&=\frac1{2d}\E_0\langle\dot K_0,S\otimes P+P\otimes S\rangle
  =\frac1d\E_0\langle\dot K_0,S\otimes P\rangle\\
  &=2\varphi(p)\E^0[m_ph_p\epsilon]
  -2\int\varphi(p\wedge r)\E^0[m_ph_pm_r]\,\nu(\dd r).
  \end{aligned}
 \]
\item The product of the scores contributes the following; the
 terminal--terminal term vanishes, since $\varphi(1)=0$:
 \[
  \begin{aligned}
  \mathrm I_4&=\frac1{2d}\E_0\bigl[\bar f_p\langle\dot K_0,S\otimes S\rangle\bigr]\\
  &=-\int\varphi(r)\E^0[\bar f_p\epsilon m_r]\,\nu(\dd r)
  +\frac12\iint\varphi(r\wedge\ell)\E^0[\bar f_pm_rm_\ell]
   \,\nu(\dd r)\nu(\dd\ell).
  \end{aligned}
 \]
\item The Hessian of the logarithmic weight contributes
 \[
  \mathrm I_5=\frac1{2d}\E_0\bigl[\bar f_p\langle\dot K_0,M\rangle\bigr]
  =-\frac12\int\varphi(r)\E^0[\bar f_ph_r]\,\nu(\dd r).
 \]
\item The explicit derivative of the weight is
 $\dot\ell_0=-\int\dot u(r,X_r)\,\nu(\dd r)$. Evaluating it with
 \eqref{eq:scalarNodeVariation}, we find
 \[
  \begin{aligned}
  \mathrm I_6=\frac1d\Cov_0(\Phi_0,\dot\ell_0)
  ={}&\frac12\int\varphi(r)\E^0\bigl[\bar f_p\bigl(h_r+a(r)m_r^2\bigr)\bigr]
   \,\nu(\dd r)\\
  &+\frac12\int\nu(\dd r)\int_{(r,1]}\varphi(\ell)
   \E^0\bigl[\bar f_pg_r^\ell(Y_r)\bigr]\,\nu(\dd\ell).
  \end{aligned}
 \]
\end{enumerate}
By \eqref{eq:fnd-path-heat}, we have
\[
 \frac1d\,\frac{\dd}{\dd\theta}\widetilde\Gamma_\theta(p)\Big|_{\theta=0}
 =\mathrm I_1+\cdots+\mathrm I_6 .
\]
Conditioning on $\calG^0_p$ and using \eqref{eq:sc-martingale}, we
have
$\E^0[m_ph_p\epsilon]=\E^0[m_p^2h_p]$. The following table records
where each part of $\mathrm I_1,\ldots,\mathrm I_6$ goes; ``line $i$''
refers to the $i$-th line of \eqref{eq:scalarJexplicit}.
\begin{center}
\small
\begin{tabular}{@{}lll@{}}
\toprule
Group & Part & Contribution\\
\midrule
$\mathrm I_1$ & $\varphi(p)\E^0h_p^2$ & line 1, diagonal coefficient\\
 & $\varphi(p)\E^0[m_pu_{3,p}]$ & cancels with $\mathrm I_2$\\
$\mathrm I_2$ & $-\varphi(p)\E^0[m_pu_{3,p}]$ & cancels with $\mathrm I_1$\\
 & $-2a(p)\varphi(p)\E^0[m_p^2h_p]$ & line 1, diagonal coefficient\\
 & integral over $(p,1]$ & line 1, second term\\
$\mathrm I_3$ & $2\varphi(p)\E^0[m_ph_p\epsilon]=2\varphi(p)\E^0[m_p^2h_p]$
 & line 1, diagonal coefficient\\
 & integral of $\E^0[m_ph_pm_r]$ & line 2, first term\\
$\mathrm I_4$ & integral of $\E^0[\bar f_p\epsilon m_r]$ & line 2, second term\\
 & double integral & line 3, second term\\
$\mathrm I_5$ & $-(1/2)\int\varphi(r)\E^0[\bar f_ph_r]\,\nu(\dd r)$
 & cancels with $\mathrm I_6$\\
$\mathrm I_6$ & $(1/2)\int\varphi(r)\E^0[\bar f_ph_r]\,\nu(\dd r)$
 & cancels with $\mathrm I_5$\\
 & $(1/2)\int\varphi(r)a(r)\E^0[\bar f_pm_r^2]\,\nu(\dd r)$
 & line 3, first term\\
 & iterated integral & line 4\\
\bottomrule
\end{tabular}
\end{center}
In summary, the terms with $m_pu_{3,p}$ cancel, the covariance terms with
$h_r$ cancel, and the diagonal coefficient is
$\E^0[h_p^2-2a(p)m_p^2h_p+2m_p^2h_p]=\E^0[h_p^2+2(1-a(p))m_p^2h_p]$.
The sum $\mathrm I_1+\cdots+\mathrm I_6$ is the right side of
\eqref{eq:scalarJexplicit}, including its signs, its factors $1/2$, and
its strict future intervals. This completes the proof.
\end{proof}

\section{Domination by the scalar multiplier}\label{sec:sc-domination}

We now show that the scalar response operator is dominated by its
multiplier.
We use the scalar Parisi functional with coefficient $d$,
\[
 \mathcal P_d(a)=u_{a,d}(0,0)-\frac d2\int_0^1q\,a(q)\dd q,
\]
defined on the convex set of profiles. For $d=\beta^2$, it is the
functional $\mathcal P_\beta$ of \eqref{eq:intro-parisi}.

\begin{lemma}[Scalar multiplier domination]\label{lem:Jdom}
For every profile $a$ and $d\in[\beta^2/2,\beta^2]$, the operator
$J^{\scal}_{a,d}$ is self-adjoint on $\mathcal H_\nu$, with norm at most
$C_\beta$, and
\begin{equation}\label{eq:Jdom}
 J^{\scal}_{a,d}\preceq M_{k_{a,d}},
 \qquad (M_{k_{a,d}}v)(q)=k_{a,d}(q)v(q).
\end{equation}
\end{lemma}

\begin{proof}
The bound on the norm is \eqref{eq:sc-J-bounded}. The proof of the rest
has five steps. We first recall that $\mathcal P_d$ is convex and
compute its first variation. Then we follow $\mathcal P_d$ along a
transport of levels and compute its second derivative with
\cref{lem:sc-response}. Next, we bound the same second derivative from
below, using convexity with the first variation frozen at $a$; the
comparison yields the inequality \eqref{eq:Jdom} for smooth directions.
Two-parameter transports then give the symmetry, and a density argument
extends both conclusions to $\mathcal H_\nu$. No minimizing property of
$a$ is used. We abbreviate $\Gamma=\Gamma_{a,d}$ and $k=k_{a,d}$, and
$\langle\cdot,\cdot\rangle$ denotes the inner product of $L^2(\nu)$.

\paragraph{Step 1: convexity and the first variation.}
We first note that $\mathcal P_d$ is convex on profiles. Indeed, by
\cite[Theorem~2]{AuffingerChen2015Uniqueness}, applied with
$\xi(q)=dq^2/2$ and zero external field, the value at $(0,0)$ of the
solution with terminal value $\log\cosh x$ is a strictly convex
function of the measure, and hence of its distribution function. This
value differs from $u_{a,d}(0,0)$ by $\log2$, and the penalty term of
$\mathcal P_d$ is affine in $a$.

The first variation of $\mathcal P_d$ along segments is the scalar case of
\cite[(2.8)]{AL} (see also \cite[Lemma~3.2]{JagannathTobasco2017PhaseDiagram}
and \cite[Theorem~2]{Chen2017Variational}): for profiles $a$ and
$\widetilde a$,
\begin{equation}\label{eq:sc-first-variation}
 \lim_{\vartheta\downarrow0}\frac{\mathcal P_d(a+\vartheta(\widetilde a-a))
                               -\mathcal P_d(a)}{\vartheta}
 =\frac d2\int_0^1(\Gamma(q)-q)(\widetilde a(q)-a(q))\dd q .
\end{equation}
The statement in \cite{AL} assumes a positive root mass. We note
that the argument of Step 1 of the proof of \cref{lem:sc-response},
with the coefficient fixed and the profile varied, gives a direct proof
of \eqref{eq:sc-first-variation} for all profiles. The difference
quotients of
$u_{a+\vartheta(\widetilde a-a),d}$ converge, with all spatial
derivatives, to the solution $\dot u$ of
\[
 \partial_q\dot u+\frac d2\dot u_{xx}+d\,a\,u_x\dot u_x
 =-\frac d2(\widetilde a-a)u_x^2,\qquad \dot u(1,\cdot)=0,
\]
and \eqref{eq:sc-backward} gives
\[
 \dot u(0,0)=\frac d2\int_0^1(\widetilde a(q)-a(q))\,\E^0m_q^2\dd q
 =\frac d2\int_0^1(\widetilde a(q)-a(q))\Gamma(q)\dd q .
\]
The penalty contributes $-(d/2)\int_0^1q(\widetilde a(q)-a(q))\dd q$.
Since $\mathcal P_d$ is convex, the difference quotient in
\eqref{eq:sc-first-variation} is nondecreasing in $\vartheta\in(0,1]$.
Comparing its value at $\vartheta=1$ with its limit as
$\vartheta\downarrow0$, we obtain
\begin{equation}\label{eq:sc-convexity}
 \mathcal P_d(\widetilde a)-\mathcal P_d(a)
 \ge\frac d2\int_0^1(\Gamma(q)-q)(\widetilde a(q)-a(q))\dd q .
\end{equation}

\paragraph{Step 2: the second derivative along a transport.}
Let $\varphi$, $T_\theta$, $\theta_0$, and $a_\theta=a\circ T_\theta^{-1}$
be as in \cref{sec:sc-response}, and let
$\Pi(\theta)=\mathcal P_d(a_\theta)$ for $|\theta|<\theta_0/2$. By
\eqref{eq:sc-transported} at $q=0$, we have
$u_{a_\theta,d}(0,0)=u_\theta(0,0)$. Using Fubini's theorem, we have
\[
 \int_0^1x\,a_\theta(x)\dd x
 =\int\Bigl(\int_{T_\theta(q)}^1x\dd x\Bigr)\mu(\dd q)
 =\frac12-\frac12\int T_\theta(q)^2\,\mu(\dd q),
\]
which gives
\[
 \Pi(\theta)=u_\theta(0,0)-\frac d4+\frac d4\int T_\theta^2\dd\mu .
\]
The derivative of the last term is $(d/2)\int T_\theta\varphi\dd\mu$,
in which the atom $s\delta_0$ does not contribute, since $\varphi(0)=0$.
Combining this with \eqref{eq:sc-root-variation}, we obtain
\begin{equation}\label{eq:transporthess}
 \Pi'(\theta)=\frac d2\int\varphi\bigl(T_\theta-\widetilde\Gamma_\theta\bigr)\dd\nu,
 \qquad
 \Pi''(0)=\frac d2\langle\varphi,(I-J_{a,d})\varphi\rangle .
\end{equation}
For the second identity, we used \cref{lem:sc-response} and
differentiated under the integral over $\nu$. This is justified by
dominated convergence, since by Step 3 of the proof of
\cref{lem:sc-response}, $\partial_\theta\widetilde\Gamma_\theta(p)$ is
continuous in $\theta$ and bounded uniformly in $p$ and $\theta$. By
the same argument, $\Pi''$ is continuous near $0$. At $\theta=0$, the first identity
in \eqref{eq:transporthess} is the first variation along transports in
\cite[Section~3, (19)]{AuffingerChen2015Properties}, where it is
credited to \cite[Lemma~3.7, Eq.~(3.34)]{Talagrand2006ParisiMeasures}.

\paragraph{Step 3: the lower bound from convexity.}
Let $G(y)=\int_y^1(\Gamma(q)-q)\dd q$. Recall that
$a_\theta(q)=\mu(\{y:T_\theta(y)\le q\})$. By Fubini's theorem, we have
$\int_0^1(\Gamma(q)-q)a_\theta(q)\dd q=\int G(T_\theta y)\,\mu(\dd y)$.
Applying \eqref{eq:sc-convexity} with $\widetilde a=a_\theta$, we obtain
\[
 \Pi(\theta)-\Pi(0)\ge\frac d2\int\bigl(G(T_\theta y)-G(y)\bigr)\mu(\dd y)
 =\frac d2\int\bigl(G(T_\theta y)-G(y)\bigr)\nu(\dd y),
\]
where the atom $s\delta_0$ does not contribute because $T_\theta(0)=0$.
Then \cref{lem:sc-basic}(c) implies that $G$ is twice continuously
differentiable on $[0,1]$, with
\[
 G'(y)=y-\Gamma(y),\qquad G''(y)=1-k(y).
\]
The function $G''$ is uniformly continuous and $\varphi$ is bounded, so
by Taylor's formula, we have
\[
 G(T_\theta y)-G(y)=\theta\varphi(y)G'(y)
 +\frac{\theta^2}2\varphi(y)^2\bigl(1-k(y)\bigr)+o(\theta^2),
\]
uniformly in $y\in[0,1]$. By \eqref{eq:transporthess} at $\theta=0$,
we have $\Pi'(0)=(d/2)\int\varphi G'\dd\nu$. Integrating the expansion
against $\nu$ and inserting it into the previous inequality, we obtain
\[
 \Pi(\theta)-\Pi(0)-\theta\Pi'(0)
 \ge\frac{\theta^2}2\cdot\frac d2\int(1-k)\varphi^2\dd\nu+o(\theta^2).
\]
On the other hand, $\Pi$ is twice continuously differentiable near
$0$, so the left side equals
$(\theta^2/2)\Pi''(0)+o(\theta^2)$. Dividing by $\theta^2/2$ and
letting $\theta\to0$, we find
\[
 \Pi''(0)\ge\frac d2\int(1-k)\varphi^2\dd\nu .
\]
Comparing this with
\eqref{eq:transporthess}, we conclude that
\begin{equation}\label{eq:sc-dom-smooth}
 \langle\varphi,J_{a,d}\varphi\rangle\le\langle\varphi,k\varphi\rangle .
\end{equation}
Note that $\Pi$ need not be convex, since $\theta\mapsto a_\theta$ is
not affine. The lower bound accounts for this through the term
$1-k=G''$.

\paragraph{Step 4: symmetry.}
We now take two functions $\varphi_1$ and $\varphi_2$ as above, and we
consider the transport $T(q)=q+\theta_1\varphi_1(q)+\theta_2\varphi_2(q)$, which
fixes both endpoints. Let $\widetilde\Gamma_{\theta_1,\theta_2}$ and
$\Pi(\theta_1,\theta_2)$ be defined as before, for small
$|\theta_1|,|\theta_2|$. By Step 1 of the proof of
\cref{lem:sc-response} for several parameters, and by
\cref{lem:fnd-path-heat} applied in each parameter with the other one
fixed, the function $\widetilde\Gamma_{\theta_1,\theta_2}(p)$ has
continuous first partial derivatives, bounded uniformly in $p$. The
computation of Step 2 of that proof at every base point, with the
coefficient derivative $d\varphi_i'$, yields
\[
 \partial_{\theta_i}\Pi
 =\frac d2\int\varphi_i\bigl(T-\widetilde\Gamma_{\theta_1,\theta_2}\bigr)\dd\nu,
 \qquad i=1,2 .
\]
It follows that $\Pi$ has continuous second partial derivatives, and
using \cref{lem:sc-response} in the direction $\varphi_2$ (with
$\theta_1=0$), we find that
\[
 \partial_{\theta_2}\partial_{\theta_1}\Pi(0,0)
 =\frac d2\langle\varphi_1,(I-J_{a,d})\varphi_2\rangle .
\]
By the equality of mixed partial derivatives, we conclude that this
bilinear form is symmetric, that is, $\langle\varphi_1,J_{a,d}\varphi_2\rangle
=\langle\varphi_2,J_{a,d}\varphi_1\rangle$.

\paragraph{Step 5: extension to the transport space.}
In Steps 3 and 4, we proved the domination inequality and the symmetry
for restrictions to $[0,1]$ of smooth functions vanishing at both
endpoints. Such functions lie in $\mathcal H_\nu$, and they are dense
in it. Indeed, let $v\in\mathcal H_\nu$, and let $\psi_n$ be
continuous functions with $\|\psi_n-v\|\to0$, which exist because
continuous functions are dense in $L^2(\nu)$. For $n\ge3$, let $\chi_n$ be
continuous with values in $[0,1]$, vanishing at $0$ and $1$, and equal
to one on $[1/n,1-1/n]$. Then
$\|\chi_n\psi_n-v\|\le\|\psi_n-v\|+\|(1-\chi_n)v\|$, and
$\|(1-\chi_n)v\|\to\|\ind_{\{0,1\}}v\|=0$ by dominated convergence,
because $v$ vanishes $\nu$-almost everywhere on $\{0,1\}$. Finally, let
$f$ be continuous with $f(0)=f(1)=0$. If $P$ is a polynomial with
$|P-f|\le\tau$ on $[0,1]$, then the polynomial
$P(x)-P(0)(1-x)-P(1)x$ vanishes at both endpoints and is within
$3\tau$ of $f$; by the Weierstrass theorem, $f$ is a uniform limit of
such polynomials. For $v\in\mathcal H_\nu$, we have
$\langle v,J^{\scal}_{a,d}v\rangle=\langle v,J_{a,d}v\rangle$, and both
sides of \eqref{eq:sc-dom-smooth} and of the symmetry identity are
continuous on $L^2(\nu)$, since $J_{a,d}$ is bounded on $L^2(\nu)$ (see
the proof of \eqref{eq:sc-J-bounded}) and $0\le k\le\beta^2$. Both conclusions extend to $\mathcal H_\nu$ by
continuity. Multiplication by $k$ maps $\mathcal H_\nu$ into
itself, and \eqref{eq:Jdom} follows.
\end{proof}

\section{The reference multiplier}\label{sec:sc-reference}

A finite-dimensional minimizer may place mass outside the scalar
support. For this reason, we control the multiplier at the scalar
Parisi measure on all of $[0,1]$. Recall from
\eqref{eq:loc-scalar-measure} that $\supp\mu_\beta=[0,q_\beta]$ with
$0<q_\beta<1$ \cite[Theorem~1.1]{Lopatto}, and that $\mu_\beta$ has an
atom at $q_\beta$. In particular, $\alpha_\beta=1$ on $[q_\beta,1]$.

\begin{lemma}[Reference multiplier on the whole interval]\label{lem:reference}
At the scalar Parisi measure,
\begin{equation}\label{eq:referencek}
 k_{\alpha_\beta,\beta^2}(q)=1\quad(0\le q\le q_\beta),
 \qquad
 k_{\alpha_\beta,\beta^2}(q)
       \le e^{-\beta^2(q-q_\beta)/2}\quad(q_\beta\le q\le1).
\end{equation}
In particular, $\sup_{q\in[0,1]}k_{\alpha_\beta,\beta^2}(q)\le1$.
\end{lemma}

\begin{proof}
We abbreviate $d=\beta^2$, $a=\alpha_\beta$, $k=k_{\alpha_\beta,d}$,
and $\Gamma=\Gamma_{\alpha_\beta,d}$, and we write $u$ and $Y$ for the
corresponding solution and diffusion. The proof has four steps. We
first prove that $k=1$ on the support. We then recall that a
transformed density of $Y_{q_\beta}$ is even and log-concave, show that
this density evolves by the heat equation above $q_\beta$, and finally
use a rearrangement inequality to bound $k$ there.

\paragraph{Step 1: the support.}
By an identity of Talagrand \cite{Talagrand2006ParisiMeasures}, in the
form stated in \cite[Theorem~5]{AuffingerChen2015Properties} and
\cite[Lemma~2.6]{Lopatto} (see also
\cite[Proposition~1.1]{JagannathTobasco2017PhaseDiagram} and
\cite[Proposition~1]{Chen2017Variational}), we have $\Gamma(q)=q$ for
every $q\in\supp\mu_\beta=[0,q_\beta]$. The function $\Gamma$ of
\cite{Lopatto} coincides with ours, since $u_x$ does not depend on the constant
$\log2$ in the terminal value. By \cref{lem:sc-basic}(c), we have
$\Gamma'=k$, which implies that $k=1$ on $(0,q_\beta)$. Since $k$ is
continuous, we conclude that $k=1$ on $[0,q_\beta]$.

\paragraph{Step 2: log-concavity of the transformed density.}
For $q\in(0,1]$, let $\pi_q$ be the density of $Y_q$, and set
\[
 z_q(x)=\pi_q(x)e^{-a(q)u(q,x)},\qquad
 V_q(x)=-\log z_q(x)=a(q)u(q,x)-\log\pi_q(x).
\]
By \eqref{eq:sc-prefix-law}, $z_q(x)\dd x$ is the distribution of the endpoint
$X_q$ under the finite measure obtained by weighting the law of the
Gaussian path, without normalization, by
$\exp\{-\int_{[0,q]}u(r,X_r)\,\mu_\beta(\dd r)\}$. By
\cite[Proposition~4.1]{Lopatto}, which applies to every probability
measure on $[0,1]$, including measures with atoms, the density $\pi_q$
is smooth and positive, and the potential $a(q)u^{\rm L}(q,x)-\log\pi_q(x)$
is smooth, even, and has second derivative at least $1/(dq)$; here
$u^{\rm L}=u-\log2$ is the solution with the terminal value of
\cite{Lopatto}, and $a(q)$ is the right-continuous value. The two
potentials differ by the constant $a(q)\log2$, which implies that
\[
 V_q''(x)\ge\frac1{dq}>0,\qquad V_q(-x)=V_q(x),\qquad x\in\R .
\]
Note that $V_q'(0)=0$ by evenness. Then $V_q'(x)\ge x/(dq)>0$ for
$x>0$.
This implies that $z_q$ is even, log-concave, and decreasing on
$(0,\infty)$. We use this at $q=q_\beta$, where the right-continuous
value is $a(q_\beta)=1$, including the atom of $\mu_\beta$ at
$q_\beta$; in particular, $z_{q_\beta}=\pi_{q_\beta}e^{-u(q_\beta,\cdot)}$.

\paragraph{Step 3: heat evolution above the support.}
On $[q_\beta,1]$, we have $a=1$, and the solution of \eqref{eq:sc-pde}
is
\[
 u(q,x)=\log(2\cosh x)+\frac{d(1-q)}2,\qquad q_\beta\le q\le1,
\]
by \cite[Lemma~2.4(i)]{Lopatto}; this formula can also be checked by
direct substitution, since $u_x=\tanh x$, $u_{xx}=\sech^2x$,
$u_{xx}+u_x^2=1$, and $\partial_qu=-d/2$. We write $G_v$ for the centered
Gaussian density with variance $v$. By \cite[Lemma~2.4(ii)]{Lopatto},
applied with mass $m=1$ on $[q_\beta,1)$, we have, for
$q_\beta<q\le1$,
\[
 \pi_q(x)=e^{u(q,x)}\int G_{d(q-q_\beta)}(x-y)\,e^{-u(q_\beta,y)}\pi_{q_\beta}(y)\dd y ;
\]
the constants $\log2$ in $u(q,x)$ and $u(q_\beta,y)$ cancel. This means
that
\begin{equation}\label{eq:sc-heat}
 z_q=G_{d(q-q_\beta)}*z_{q_\beta},\qquad q_\beta<q\le1 .
\end{equation}
Equivalently, $z_q$ solves the heat equation
$\partial_qz_q=(d/2)\partial_{xx}z_q$. To see this directly, we
substitute $\pi_q=e^{u}z_q$ into the forward equation of $Y$ on the
plateau, which gives
\[
 \partial_qz_q+(\partial_qu)z_q
 =e^{-u}\Bigl(\frac d2\partial_{xx}\pi_q-d\,\partial_x(u_x\pi_q)\Bigr)
 =\frac d2\partial_{xx}z_q-\frac d2\bigl(u_{xx}+u_x^2\bigr)z_q ;
\]
the zeroth-order terms cancel because $u_{xx}+u_x^2=1$ and
$\partial_qu=-d/2$.

On the plateau, we have $u_{xx}=\sech^2$ and
$\pi_q=e^{u(q,\cdot)}z_q=2\cosh(\cdot)\,e^{d(1-q)/2}z_q$. The multiplier
is therefore
\begin{equation}\label{eq:sc-k-plateau}
 k(q)=d\int\pi_q(x)\sech^4x\dd x
 =2d\,e^{d(1-q)/2}\int_\R z_q(x)\sech^3x\dd x,\qquad q_\beta\le q\le1 .
\end{equation}

\paragraph{Step 4: the rearrangement inequality.}
Let $f$ and $g$ be nonnegative, even, continuous functions that are
nonincreasing on $[0,\infty)$. We claim that
\begin{equation}\label{eq:sc-rearrangement}
 \int f(x-y)g(x)\dd x\le\int f(x)g(x)\dd x,\qquad y\in\R .
\end{equation}
This is the centered rearrangement inequality in one dimension, which
also follows from Anderson's inequality~\cite[Theorem~1]{Anderson1955}
by the layer-cake representation. We include a proof based on intersections of intervals.
By the layer-cake representation and Tonelli's theorem,
\[
 \int f(x-y)g(x)\dd x=\int_0^\infty\!\!\int_0^\infty
 \bigl|(\{f>t\}+y)\cap\{g>t'\}\bigr|\dd t\dd t',
\]
where $|\cdot|$ is Lebesgue measure. The superlevel sets $\{f>t\}$ and
$\{g>t'\}$ are open intervals centered at $0$. For two centered intervals
$I$ and $I'$, we have $|(I+y)\cap I'|\le\min\{|I|,|I'|\}=|I\cap I'|$, and
\eqref{eq:sc-rearrangement} follows.

We use \eqref{eq:sc-rearrangement} with $f=z_{q_\beta}$ and
$g=\sech^3$, which satisfy its hypotheses by Step 2. All the integrals
below are finite, since $\sech\le1$ and
$\int z_{q_\beta}=\E[e^{-u(q_\beta,Y_{q_\beta})}]\le1/2$. By
\eqref{eq:sc-heat} and Tonelli's theorem, for $q_\beta<q\le1$,
\[
 \int z_q(x)\sech^3x\dd x
 =\int G_{d(q-q_\beta)}(y)\Bigl(\int z_{q_\beta}(x-y)\sech^3x\dd x\Bigr)\dd y
 \le\int z_{q_\beta}(x)\sech^3x\dd x .
\]
Inserting this into \eqref{eq:sc-k-plateau}, and using
\eqref{eq:sc-k-plateau} at $q=q_\beta$ and $k(q_\beta)=1$ from Step 1,
we find
\[
 k(q)\le2d\,e^{d(1-q)/2}\int z_{q_\beta}(x)\sech^3x\dd x
 =e^{-d(q-q_\beta)/2}k(q_\beta)=e^{-d(q-q_\beta)/2},
 \qquad q_\beta\le q\le1 .
\]
Together with Step 1, this proves \eqref{eq:referencek}.
\end{proof}

\section{Stability at a minimizing profile}\label{sec:sc-minimizer}

The isolation estimate of \cref{prop:loc-isolation} compares a
minimizer with the scalar trial profile
$\gamma^{\rm tr}_\lambda=\max\{\lambda,\gamma_\beta\}$, where
$\gamma_\beta=\beta\alpha_\beta$. We first bound the distance between
the trial profile and $\gamma_\beta$. Recall from \eqref{eq:lz-gsroot}
that $\gamma_\beta(q)\ge k_0q$ for $0\le q\le T_0$, with constants
$k_0,T_0>0$ depending only on $\beta$. We claim that, for every
$\lambda>0$,
\begin{equation}\label{eq:sc-floor}
 \|\gamma^{\rm tr}_\lambda-\gamma_\beta\|_{L^1(0,1)}
 \le C_\beta\lambda^2 .
\end{equation}
Indeed, $\gamma^{\rm tr}_\lambda-\gamma_\beta=(\lambda-\gamma_\beta)_+$.
Suppose first that $\lambda\le k_0T_0$. For $q\ge T_0$, since
$\gamma_\beta$ is nondecreasing, we have
$\gamma_\beta(q)\ge\gamma_\beta(T_0)\ge k_0T_0\ge\lambda$, and the
difference vanishes there. This gives
\[
 \|\gamma^{\rm tr}_\lambda-\gamma_\beta\|_{L^1(0,1)}
 =\int_0^{T_0}(\lambda-\gamma_\beta(q))_+\dd q
 \le\int_0^{\lambda/k_0}(\lambda-k_0q)\dd q=\frac{\lambda^2}{2k_0}.
\]
If $\lambda>k_0T_0$, the difference is at most $\lambda$ everywhere,
and $\lambda\le\lambda^2/(k_0T_0)$. Since $T_0\le1/4<2$, both cases give
\eqref{eq:sc-floor} with $C_\beta=1/(k_0T_0)$.

\begin{proposition}[Scalar stability at the minimizing profile]
\label{prop:lz-stability}
There exists a constant $C_{\rm st}\ge1$, depending only on $\beta$ (in
particular, not on the minimizer or its atoms), with the following
property. Let $N\ge2$, $0<\rho\le1$, $t=1-\rho$, and
$0<\lambda\le\lambda_0$, and set $s=\lambda/\beta$. Let
$\mu_N=s\delta_0+\nu_N$ be a minimizer of $\mathcal F_{t,s}$ in
\eqref{eq:fnd-functional}, let $a_N(q)=\mu_N([0,q])$ be its profile,
and let $\eta$ be a trial budget for the minimizer $\beta a_N$ of
$F_\rho$, in the sense of \eqref{eq:loc-trialbudget}. Let
$d_N=\beta^2(1-t/N)$. Then
\begin{equation}\label{eq:lz-profileconvergence}
 \|a_N-\alpha_\beta\|_{L^1(0,1)}
 \le\min\bigl\{1,C_{\rm st}[\rho^{-3/4}\eta^{1/4}+\lambda^2]\bigr\}.
\end{equation}
Further, let
$\varepsilon_N=\max\{0,\sup_{q\in[0,1]}k_{a_N,d_N}(q)-1\}$. Then, on the
scalar transport space $\mathcal H_{\nu_N}$ of $\mu_N$,
\begin{equation}\label{eq:lz-stability}
 J^{\scal}_{a_N,d_N}\preceq(1+\varepsilon_N)I,
 \qquad
 \varepsilon_N\le C_{\rm st}
       \bigl[\rho^{-3/4}\eta^{1/4}+\lambda^2+N^{-1}\bigr].
\end{equation}
\end{proposition}

\begin{proof}
We use the energy normalization of \cref{sec:fnd}, in which
$\gamma=\beta a$ and $F_\rho(\gamma)=\beta^{-1}\mathcal F_{t,s}(\mu)$.
The map $\mu\mapsto\beta a$ is a bijection onto the admissible profiles
of \cref{sec:lz}, and multiplication by $\beta^{-1}$ preserves
minimizers. Hence $\gamma_N=\beta a_N$ is a minimizer of $F_\rho$, and
$\eta$ is a trial budget for it. The hypotheses of
\cref{prop:loc-isolation} hold, and \eqref{eq:loc-isolation} implies
that
$\|\gamma_N-\gamma^{\rm tr}_\lambda\|_{L^1}\le C\rho^{-3/4}\eta^{1/4}$,
with a constant $C$ depending only on $\beta$. Using the triangle
inequality and \eqref{eq:sc-floor}, we have
\[
 \|a_N-\alpha_\beta\|_{L^1}
 \le\beta^{-1}\bigl(\|\gamma_N-\gamma^{\rm tr}_\lambda\|_{L^1}
       +\|\gamma^{\rm tr}_\lambda-\gamma_\beta\|_{L^1}\bigr)
 \le C_\beta[\rho^{-3/4}\eta^{1/4}+\lambda^2].
\]
The distance is also at most one, since both functions take values in
$[0,1]$. This proves \eqref{eq:lz-profileconvergence}.

Since $N\ge2$ and $0\le t<1$, we have $d_N\in[\beta^2/2,\beta^2]$ and
$|d_N-\beta^2|=\beta^2t/N\le\beta^2/N$. Then
\cref{lem:scalarcontinuity}, applied with $(a,d)=(a_N,d_N)$ and
$(\widetilde a,\widetilde d)=(\alpha_\beta,\beta^2)$, and
\cref{lem:reference} imply that, for every $q\in[0,1]$, whether or not
$q$ lies in $\supp\mu_\beta$,
\[
 k_{a_N,d_N}(q)\le k_{\alpha_\beta,\beta^2}(q)
 +C_\beta\bigl(\|a_N-\alpha_\beta\|_{L^1}+N^{-1}\bigr)
 \le1+C_\beta\bigl(\|a_N-\alpha_\beta\|_{L^1}+N^{-1}\bigr).
\]
Together with \eqref{eq:lz-profileconvergence}, this yields the bound
on $\varepsilon_N$ in \eqref{eq:lz-stability}. Finally, since
$k_{a_N,d_N}\le1+\varepsilon_N$ pointwise, \cref{lem:Jdom}, applied to
the profile $a_N$ with $\nu=\nu_N$, implies that
$J^{\scal}_{a_N,d_N}\preceq M_{k_{a_N,d_N}}\preceq(1+\varepsilon_N)I$ on
$\mathcal H_{\nu_N}$. Note that only the multipliers of the two
profiles are compared; the operator inequality is proved on the
transport space of the minimizer itself. We choose $C_{\rm st}\ge1$
larger than the constants in the two bounds; it depends only on
$\beta$.
\end{proof}


\part{The cavity contraction}\label{part:cavity}
\chapter{The cavity interpolation and its tree laws}\label{sec:c7}

This section and the next two prove \cref{prop:c7-closure}. Suppose that
at a minimizing profile, the scalar response operator $J^{\scal}_{a,d}$
of \cref{sec:scalar-stability} satisfies
$J^{\scal}_{a,d}\preceq(1+\eps_N)I$ with $\eps_N\le\rho/4$. The
proposition states that the self-overlap variance
$V=\int\Var(S_q)\,\nu(\dd q)$ then obeys
$(3\rho/4)V\le C_{\rm cav}[N^{-1}+A+J+H_3]$, where $C_{\rm cav}$ depends
only on $\beta$. In \cref{cor:c7-defect}, we deduce from it the bound on
the defect $\calD_N(\rho)$ that is used in \cref{sec:b7}.

The argument is a cavity interpolation on one site. We single out the
last spin and move the variance of its couplings to the other $n=N-1$
spins into its auxiliary field. At the end of the interpolation, the last
spin is independent of the others. By exchangeability, $V$ equals $n/N$
times the integrated covariance between the squared conditional
magnetization of the last spin and the self-overlap of the remaining
spins, up to an error of order $N^{-1}$. This covariance vanishes at the
independent endpoint, and its value for the full system is its derivative
there plus a Taylor remainder. In
\cref{sec:c7-endpoint}, we show that this derivative is a multiple of the
scalar response operator applied to the covariance of the self-overlap of
the remaining spins, plus a nonpositive term and explicit errors
(\cref{lem:linear}). We also bound the errors by $A$ and $J$ there
(\cref{lem:cl-response}). In \cref{sec:c7-closure}, we compare the
variance of this self-overlap at the independent endpoint with $V$, with
coefficient one (\cref{lem:cl-star}). We then express $V$ through the
expansion (\cref{lem:cav-trace}), bound the cubic moment that controls
the remainder by $A$, $J$, and $H_3$ (\cref{lem:c7-thermal}), and close the
resulting inequality.

In this section, we set up the interpolation and the calculus that the
next two sections use. In \cref{subsec:c7-row}, we define the cavity law
$\Pp_u$, $0\le u\le1$, identify its two endpoints, and record a variance
bound \eqref{eq:cl-BL} for the remaining spins. In \cref{subsec:c7-tree},
we represent products of conditional expectations by laws of finitely
many descendants on marked trees. We show that the solution of the Parisi
equation depends smoothly on the incident row and on the variance of the
last field (\cref{lem:tree-regularity}). We compute the density of the
tree laws (\cref{lem:tree-density}) and the responses
\eqref{eq:nodeTheta} and \eqref{eq:nodeb} to these two parameters
(\cref{lem:tree-responses}). From them, we derive an exact formula for the
derivative of a tree moment in $u$ (\cref{lem:streaming}) and a
comparison of tree moments at different values of $u$
(\cref{cor:comparison}). In \cref{subsec:c7-three}, we bound integrated
products of three overlap defects by a cubic moment $\mathcal M_3(u)$
(\cref{lem:cav-threeproducts}), and deduce a bound on the Taylor
remainder (\cref{lem:remainder}).

We use the generalized law of \cref{sec:gauss}, with its variance bound
(\cref{lem:fnd-posterior}) and its heat identity
(\cref{lem:fnd-path-heat}). We also use the scalar law \eqref{eq:sc-pde}
with \cref{lem:cl-scalar}, and the differentiation and approximation
arguments of \cite[Sections~2.1, 2.3, and~2.6]{AL}. The results of this
section are used as follows. The identity \eqref{eq:nodeTheta} and the
density \eqref{eq:treedensity} with one leaf enter the proof of
\cref{lem:linear}. The bound \eqref{eq:cl-BL} and \cref{cor:comparison}
enter the proofs of \cref{lem:cl-response,lem:c7-thermal}. The bound of
\cref{lem:cav-threeproducts} enters the proof of \cref{lem:cl-star}, and
\cref{lem:remainder} enters the proofs of \cref{lem:cav-trace} and
\cref{prop:c7-closure}.

The cavity method originates in the physics literature
\cite{MezardParisiVirasoro1986,MezardParisiVirasoro1987}. We use its rigorous form on
one site, the framework of differentiation, comparison, and contraction in
\cite[Section~1.6]{TalagrandVolI2011} (see Lemmas~1.6.3--1.6.4 and
Propositions~1.6.5--1.6.6 there). An expansion on one site for
hierarchical Gaussian laws appears in \cite[Lemma~2.2]{OpenAILimitingLaw}.
There, a site is chosen by exchangeability and deleted, and its Gaussian
kernels are interpolated to those of an independent scalar model on one
site. The linear cavity kernel is half of a connected correlation of four
spins of the removed site \cite[(A8)]{OpenAILimitingLaw}, and a linear
equation with this kernel is inverted under a gap condition on the scalar
response in \cite[Proposition~5.1]{OpenAILimitingLaw}. The tree laws of
\cref{subsec:c7-tree} and their calculus are adapted from
\cite[Section~2]{OpenAIFluctuationScale}; we give the precise references
there.

\paragraph{Standing assumptions.}
Throughout \cref{part:cavity}, we fix $N\ge2$ and write $n=N-1$. We also
fix $0<\rho\le1$, $t=1-\rho$, $c=\beta^2\rho$, and
$0<s=\lambda/\beta\le1/2$. In this section, the profile
$\mu=s\delta_0+\nu$ is deterministic and otherwise arbitrary. The measure
$\nu$ may have atoms anywhere in $[0,1]$, including $0$ and $1$. No result
of this section uses a minimizing property. This matters because at the
independent endpoint, the remaining $n$ spins form a system whose law is
built from the profile of the full system, and this profile need not
minimize any functional of the $n$-spin system. Constants $C$ are absolute, and
constants $C_\beta$ depend only on $\beta$. Neither depends on $N$,
$\rho$, $s$, the profile, or the positions and masses of its atoms. Note
that the interpolation parameter $u$ is not the solution $u$ of
\eqref{eq:sc-pde}; in this part of the paper, we write $U^0$ for the
latter.

\section{Removing the incident row}\label{subsec:c7-row}

We write $\sigma=(\tau,\epsilon)\in\{-1,1\}^n\times\{-1,1\}$, so that
$\epsilon=\sigma_N$ is the last spin. The couplings $g_{iN}$, $i\le n$,
between $\epsilon$ and $\tau$ form the \emph{incident row}. Let
$u\in[0,1]$ be the interpolation parameter; $t$ stays fixed. We set
\begin{equation}\label{eq:cavityparameters}
 \alpha=\beta^2t\,\frac nN,\qquad
 d=c+\alpha=\beta^2\Bigl(1-\frac tN\Bigr),\qquad
 b(u)=d-\alpha u,\qquad 0\le u\le1.
\end{equation}
The second identity holds because
$c+\alpha=\beta^2(\rho+t-t/N)=\beta^2(1-t/N)$. Since $0\le t<1$ and
$N\ge2$, we have $d\in[\beta^2/2,\beta^2]$. Further, $b(u)\ge b(1)=c>0$
for every $u$. For $x\in\R^n$ and $y\in\R$, let
\begin{gather*}
 \mathcal E(\sigma;x,y)=\mathcal E^-(\tau)
 +\sqrt{\frac{\alpha u}{n}}\;\epsilon\sum_{i=1}^ng_{iN}\tau_i
 +x\cdot\tau+y\epsilon,
 \\
 \mathcal E^-(\tau)=\frac{\beta\sqrt t}{\sqrt N}
       \sum_{1\le i<j\le n}g_{ij}\tau_i\tau_j,
\end{gather*}
where the couplings are independent standard Gaussian variables under the
reference law. Let $U(q,x,y)$ solve
\begin{equation}\label{eq:cavityPDE}
 \begin{gathered}
 \partial_q U=-\frac12\Bigl\{
 c\bigl(\Delta_xU+a(q)|\nabla_xU|^2\bigr)
 +b(u)\bigl(\partial_{yy}U+a(q)(\partial_yU)^2\bigr)\Bigr\},
 \\ U(1,x,y)=\log\sum_\sigma e^{\mathcal E(\sigma;x,y)}.
 \end{gathered}
\end{equation}
At $u=1$, each coupling of the row has the variance $\alpha/n=\beta^2t/N$
of the full system, and the last field has the variance coefficient
$b(1)=c$ of the full system. As $u$ decreases to $0$, the variance
$\alpha u/n$ of each coupling in the row is removed, and the coefficient
of the last field increases to $b(0)=d=c+\alpha$.

The completion of \cref{sec:fnd} adds to the terminal exponent the
Gaussian term $G=\beta\sqrt t\sum_{i\le N}z_i/\sqrt{2N}$, which does not
depend on the spins. It changes $U$ into $U+G$ and multiplies the tilt
below by $e^{sG}$. Since $G$ is independent of the couplings $g_{ij}$,
this factor becomes, after normalization, an independent tilt of the
variables $z_i$, and these variables do not enter the field or the
spins. We omit it.

We construct the law $\Pp_u$, with expectation $\E_u$, as follows. We
tilt the reference law of the couplings by $e^{sU(0,0,0)}/Z_s(u)$,
where $Z_s(u)=\E_{\rm ref}e^{sU(0,0,0)}$. Given the couplings, the field
path $(X_q,Y_q)\in\R^n\times\R$ solves
\[
 \begin{aligned}
 \dd X_q&=c\,a(q)\nabla_xU(q,X_q,Y_q)\dd q+\sqrt c\,\dd B_q,
 &X_0&=0,\\
 \dd Y_q&=b(u)\,a(q)\partial_yU(q,X_q,Y_q)\dd q+\sqrt{b(u)}\,\dd B'_q,
 &Y_0&=0,
 \end{aligned}
\]
where $(B,B')$ is a standard Brownian motion in $\R^{n+1}$. The diffusion
matrix is $\operatorname{diag}(cI_n,b(u))$, and the drift is $a(q)$
times this matrix applied to $\nabla U(q,X_q,Y_q)$. Given the path,
we draw $\sigma$ with probabilities proportional to
$e^{\mathcal E(\sigma;X_1,Y_1)}$. Let $\calG_q$ be the $\sigma$-field
generated by the couplings and the path up to level $q$.

\paragraph{The two endpoints.}
At $u=1$, the law $\Pp_1$ is the tilted law of \cref{subsec:fnd-law} for
the full system, with the completion omitted as above. At $u=0$, the row
is absent, and the terminal value separates as
$U(1,x,y)=U^-(1,x)+\log(2\cosh y)$, where
$U^-(1,x)=\log\sum_\tau e^{\mathcal E^-(\tau)+x\cdot\tau}$. Let $U^-$ be
the solution for the $n$-spin system with Hamiltonian $\mathcal E^-$,
field coefficient $c$, and profile $\mu$, and let $U^0$ solve
\eqref{eq:sc-pde} with coefficient $d$ and profile $\mu$. We claim that
$U=U^-+U^0$ at $u=0$. For a step profile, each Hopf--Cole step of
\eqref{eq:cavityPDE} factors, because the increments of $x$ and of $y$
are independent and $b(0)=d$. Applying the estimate \cite[(2.19)]{AL} to
each of the three equations, we extend the identity to every profile
through the approximation of \cite[Section~2.6]{AL}. The outer normalizer
factors as $Z_s(0)=e^{sU^0(0,0)}\E_{\rm ref}e^{sU^-(0,0)}$, and the
deterministic factor $e^{sU^0(0,0)}$ cancels from the tilt. The drift of
$X$ involves only $\nabla_xU^-$, and that of $Y$ only $\partial_yU^0$.
The two Brownian motions are independent, and the terminal Gibbs
probability is the product of
$e^{\mathcal E^-(\tau)+X_1\cdot\tau-U^-(1,X_1)}$ and
$e^{\epsilon Y_1}/(2\cosh Y_1)$. We conclude that $\Pp_0$ is the product
of the tilted law of the $n$-spin system, with expectation $\E^-$, and
the scalar law of \cref{sec:scalar-stability} with coefficient $d$ and
profile $\mu$, with expectation $\E^0$. Since $d\in[\beta^2/2,\beta^2]$,
we may use the scalar quantities $m_p$, $h_p$, $f_p=m_p^2$, $\bar f_p=f_p-\E^0f_p$,
$g_r^\ell$, and $\Gamma(p)=\Gamma_{a,d}(p)=\E^0m_p^2$ of
\cref{sec:scalar-stability} for this scalar law, together
with the bounds of \cref{lem:cl-scalar}.

\paragraph{The remaining spins.}
For the first $n$ spins, we define
\[
 M_q=\E_u[\tau\mid\calG_q],\qquad A_q^-=\frac{|M_q|^2}{n},\qquad
 C_q^-=\Cov_u(\tau\mid\calG_q),\qquad
 H_q^-=\nabla_x^2U(q,X_q,Y_q).
\]
We let $m_N(q)=\E_u[\epsilon\mid\calG_q]$ be the conditional
magnetization of the last spin, and we set
\begin{equation}\label{eq:cav-rest-errors}
 a_q^-(u)=\frac{\E_u\Tr\bigl((C_q^-)^2\bigr)}{n^2},\qquad
 j_q^-(u)=\frac{\E_u[M_q^TC_q^-M_q]}{n^2}.
\end{equation}
The \emph{rest self-overlap} $A_q^-$ is a random variable. It should not
be confused with the trace term $A$ of the full system.

We note that $\Pp_u$ is a generalized law in the sense of \cref{sec:gauss}.
It has $N$ coordinates, the diffusion matrix
$\Lambda=\operatorname{diag}(cI_n,b(u))$, and independent Gaussian
couplings with deterministic coefficients: $\beta\sqrt{t/N}$ for the
pairs $i<j\le n$, $(\alpha u/n)^{1/2}$ for the pairs $(i,N)$, and $0$ on
the diagonal. The matrix $\Lambda$ is positive definite because $c>0$
and $b(u)\ge c$. The coefficient $(\alpha u/n)^{1/2}$ vanishes at $u=0$,
which the generalized law allows. We use \cref{lem:fnd-posterior} for this
law, with the first $n$ coordinates selected. It implies that
$M_q=\nabla_xU(q,X_q,Y_q)$ and $m_N(q)=\partial_yU(q,X_q,Y_q)$. Since
principal submatrices preserve the order $a(q)C_q\preceq H_q\preceq C_q$,
it also implies that $a(q)C_q^-\preceq H_q^-\preceq C_q^-$, and in
particular that $0\preceq H_q^-\preceq C_q^-$. Taking $k=\psi$ in the two
inequalities in \eqref{eq:fnd-posterior}, we find that for every
$u\in[0,1]$ and every deterministic $\psi\in L^2(\nu)$,
\begin{equation}\label{eq:cl-BL}
 \Var_u\Bigl(\int\psi(q)A_q^-\,\nu(\dd q)\Bigr)
 \le4\int\psi(q)^2j_q^-(u)\,\nu(\dd q).
\end{equation}

\section{Laws of finitely many descendants}\label{subsec:c7-tree}

Derivatives in $u$ produce products of conditional expectations at
several levels. We represent them by additional terminal spins, called
leaves, under a joint law in which deleting a leaf does not change the
law of the others. The hierarchical background of such laws is Ruelle's
probability cascades \cite{Ruelle1987}; see
\cite[Chapter~2]{Panchenko2013Book}. An earlier hierarchical Gaussian
change of measure appears in
\cite[Theorem~4.3, (4.25)--(4.26)]{BolthausenSznitman1998}. Normalized
finite replica trees of the present kind were developed in
\cite[Section~2]{OpenAIFluctuationScale}, together with the
differentiation of their moments by attaching new leaves and the
resulting Taylor and comparison bounds. They are also used in
\cite[Sections~2.1--2.3]{OpenAILimitingLaw}. In
\cite{OpenAIFluctuationScale}, the covariance differentiation formula is
Proposition~2.5, the attachment rules and their uniform variation bounds
are Lemma~2.6, the bounds on the levels at which new pairs split are
Lemma~2.8, and the Taylor and comparison bounds are Lemma~2.9. We derive
below the versions that we need for the cavity law, with its root tilt
and an arbitrary profile, including its atoms.

\subsection{Marked trees}

\begin{definition}[Marked trees and tree laws]\label{def:tree-marked}
A \emph{marked tree} $\mathcal T$ with $k\ge1$ leaves is a finite rooted
tree whose vertices $v$ carry deterministic levels $q_v\in[0,1]$, with
the following properties. The root has level $0$, and levels increase
strictly along every edge. Every vertex at a level in $(0,1)$ has at
least two outgoing edges. The vertices at level $1$ have no outgoing
edges; we call them terminal points. Finally, $k$ labeled leaves
$\mathsf a=1,\ldots,k$ are attached to the terminal points, and every
terminal point carries at least one leaf.

Given $u\in[0,1]$, the \emph{tree law} $\Pp_{\mathcal T,u}$, with
expectation $\E_{\mathcal T,u}$, is constructed as follows. The couplings
are drawn once from their tilted law under $\Pp_u$. Given the couplings,
the field starts at the origin at the root and follows the diffusion of
\cref{subsec:c7-row} along each edge. At each vertex, it continues along
the outgoing edges in conditionally independent ways. Given the couplings
and the field, the leaves carry conditionally independent spins
$\sigma^{\mathsf a}=(\tau^{\mathsf a},\epsilon_{\mathsf a})$. The spin
$\sigma^{\mathsf a}$ is drawn with probabilities proportional to
$e^{\mathcal E(\sigma;X^{\mathsf a}_1,Y^{\mathsf a}_1)}$, where
$(X^{\mathsf a}_1,Y^{\mathsf a}_1)$ is the field at the terminal point of
$\mathsf a$. We write $\E_{k,q,u}$ for the tree law of $k$ leaves that
all split at the level $q$.
\end{definition}

The field is defined at every point $v$ of $\mathcal T$, that is, at every
vertex and at every point of an edge. We write $(X_v,Y_v)$ for it and
$q_v$ for the level of $v$. Two points on different branches at the same
level are different points, with different fields. If $k\ge2$ and
$0<q<1$, the tree of $k$ leaves that split at $q$ has one edge from the
root to a vertex at level $q$, followed by $k$ edges to $k$ terminal
points. If $k\ge2$ and $q=0$, the root has $k$ outgoing edges, and if
$k\ge2$ and $q=1$, the $k$ leaves share one terminal point. For $k=1$,
the tree is the single path from level $0$ to level $1$, independently of
$q$. After deleting leaves, we remove the branches that carry no
remaining leaf and suppress the interior vertices with one outgoing edge.

Given the couplings and the field at a vertex, the continuations along
different outgoing edges are conditionally independent, and their laws are
normalized. Integrating out a leaf yields the tree law of the
remaining leaves. We call this the \emph{deletion property}. In
particular, every leaf has law $\Pp_u$. We represent products of conditional
expectations given $\calG_q$ by leaves that split at $q$.
For leaves $\mathsf a\ne\mathsf b$, let $q_{\mathsf{ab}}$ be the level of
their last common ancestor, so that $q_{\mathsf{ab}}=1$ if they share a
terminal point. We define their \emph{rest overlap}, their \emph{overlap
defect}, and the product of the overlap defect with their last spins by
\begin{equation}\label{eq:defectdefinition}
 R^-_{\mathsf{ab}}=\frac{\tau^{\mathsf a}\cdot\tau^{\mathsf b}}{n},\qquad
 \Delta_{\mathsf{ab}}=R^-_{\mathsf{ab}}-q_{\mathsf{ab}},\qquad
 D_{\mathsf{ab}}=\epsilon_{\mathsf a}\epsilon_{\mathsf b}
                          \Delta_{\mathsf{ab}}.
\end{equation}
Since $|R^-_{\mathsf{ab}}|\le1$, we have
$|D_{\mathsf{ab}}|=|\Delta_{\mathsf{ab}}|\le2$. For instance, the four
spins of $\E_{4,p,u}$ are conditionally independent given $\calG_p$, each
with the conditional law of $\sigma$ given $\calG_p$. Therefore
$\E_u[m_N(p)^2A_p^-]=\E_{4,p,u}[\epsilon_1\epsilon_2R^-_{34}]$.

\subsection{Regularity in the row and in the last-field variance}

To differentiate tree moments, we temporarily replace the row by a
deterministic coupling $\epsilon\sum_i\theta_i\tau_i$ with
$\theta\in\R^n$, and the coefficient $b(u)$ by a number $b>0$. We write
$\Lambda_b=\operatorname{diag}(cI_n,b)$. We keep the other couplings $g_{ij}$,
$i<j\le n$, fixed. Under $\Pp_u$ and under the tree laws,
conditionally on the couplings, the field and the spins have the law of
this fixed system with $\theta=(\alpha u/n)^{1/2}(g_{1N},\ldots,g_{nN})$
and $b=b(u)$.

\begin{lemma}[Regularity in the parameters]\label{lem:tree-regularity}
Let $\mu$ be a deterministic profile, and let $U=U(q,x,y)$ solve
\eqref{eq:cavityPDE} with the row replaced by
$\epsilon\sum_i\theta_i\tau_i$ and $b(u)$ replaced by $b>0$. Then $U$ has
derivatives of every order in $(\theta,x,y)$. Each derivative of positive
order is bounded by a constant depending only on $N$ and its order. This
bound is uniform in the couplings, $\theta$, $b$, $(q,x,y)$, and the
profile, and each first derivative is bounded by one in absolute value.
The function $U$ and all these derivatives are jointly continuous in
$(b,q,\theta,x,y)\in(0,\infty)\times[0,1]\times\R^n\times\R^n\times\R$.
Finally, writing $\nabla$ for the gradient in $(x,y)$, we have for
$0\le r\le h\le1$ that
\begin{equation}\label{eq:tree-integrated-pde}
 U(r,x,y)-U(h,x,y)=\frac12\int_r^h\bigl[\Tr(\Lambda_b\nabla^2U)
   +a(q)\,\nabla U^T\Lambda_b\nabla U\bigr](q,x,y)\dd q.
\end{equation}
\end{lemma}

\begin{proof}
The proof has three steps. We first prove the bounds for step profiles,
using the density of the future of the field. We then pass to arbitrary
profiles with the approximation of \cite[Section~2.6]{AL}, and finally
we prove the continuity in $q$ and $b$.

\paragraph{Step 1: step profiles.}
Let $\mu$ be a step profile. The coefficients of \eqref{eq:cavityPDE} do
not depend on $(x,y)$. Translating the starting point of the field by
$(x',y')$ has the same effect as adding
$x'\cdot\tau+y'\epsilon$ to the terminal exponent. It follows that every
derivative in $(\theta,x,y)$ is a derivative at $z=0$ with respect to
parameters $z_j$ of the terminal exponent
\[
 \mathcal E_z(\sigma;x,y)=\mathcal E(\sigma;x,y)+\sum_jz_j\eta_j(\sigma),
\]
where the features $\eta_j$ range over $\tau_i\epsilon$, $\tau_i$
($i\le n$), and $\epsilon$. Note that $|\eta_j|\le1$. Let $U_z$ be the
corresponding solution. We fix a level $q$, refine the grid of the profile
so that it contains $q$, and fix the field $(x,y)$ at level $q$. On a
grid interval $[q_\ell,q_{\ell+1})$ with $a=a_\ell>0$, the Hopf--Cole form
of the equation is $e^{a_\ell U_z(q_\ell)}=\E e^{a_\ell U_z(q_{\ell+1})}$,
where the expectation is over a Gaussian increment of covariance
$(q_{\ell+1}-q_\ell)\Lambda_b$. The transition of the diffusion is this
Gaussian transition tilted by $e^{a_\ell(U_z(q_{\ell+1})-U_z(q_\ell))}$.
If $a_\ell=0$, it is the Gaussian transition. We multiply these factors
and the Gibbs probability $e^{\mathcal E_z(\sigma)-U_z(1)}$, as in
\cite[Section~2.1, (2.1)]{AL}. The coefficient of $U_z$ at a grid point
$q_\ell>q$ is $a_{\ell-1}-a_\ell=-\nu(\{q_\ell\})$, and at level $1$ it is
$a(1-)-1=-\nu(\{1\})$. We find that given the field at level $q$, the
future increments of the field and the spin have the density
\begin{equation}\label{eq:tree-future-density}
 L_z=\exp\Bigl\{\mathcal E_z(\sigma;Z_1)-a(q)U_z(q,Z_q)
     -\int_{(q,1]}U_z(r,Z_r)\,\nu(\dd r)\Bigr\}
\end{equation}
with respect to driftless Gaussian increments of covariance $\Lambda_b$
and counting measure on $\sigma$. Here $Z_r$ is the field at level $r$.

We write $\E_z$ for the expectation under this density and
$\ell_z=\log L_z$. On a grid interval with $a_\ell>0$, we differentiate
$e^{a_\ell U_z(q_\ell)}=\E e^{a_\ell U_z(q_{\ell+1})}$. This shows that
$\partial_{z_j}U_z(q_\ell)$ is the expectation of
$\partial_{z_j}U_z(q_{\ell+1})$ under the transition of the diffusion,
and the same holds if $a_\ell=0$. At level $1$, $\partial_{z_j}U_z(1)$ is
the Gibbs mean of $\eta_j$. Iterating, we obtain
\begin{equation}\label{eq:tree-first-derivative}
 \partial_{z_j}U_z(q)=\E_z\eta_j,\qquad\text{so that}\qquad
 |\partial_{z_j}U_z|\le1 .
\end{equation}
Since $a(q)+\nu((q,1])=1$, this implies that $|\partial_{z_j}\ell_z|\le2$.
Further, if all derivatives of $U_z$ of some order $m'\ge2$ are bounded
by $C_{m'}$, then all derivatives of $\ell_z$ of order $m'$ are bounded
by $C_{m'}$. Since the reference law in \eqref{eq:tree-future-density} does not
depend on $z$, we may differentiate \eqref{eq:tree-first-derivative}
under the expectation. For directions $j,j_1,\ldots,j_m$, we have
\[
 \partial_{z_{j_1}}\cdots\partial_{z_{j_m}}\partial_{z_j}U_z(q)
 =\E_z\Bigl[\eta_j\sum_\pi\prod_{B\in\pi}\partial_B\ell_z\Bigr],
\]
where $\pi$ runs over the partitions of $\{1,\ldots,m\}$ and $\partial_B$
is the derivative in the directions $j_i$, $i\in B$. By induction on $m$,
we find that every derivative of order $m+1$ in the coordinates $z_j$ is
bounded by a constant depending only on $m$. There are $2n+1$
coordinates, and this proves the bounds of the lemma for step profiles.

\paragraph{Step 2: arbitrary profiles.}
As in \cite[Section~2.6.2]{AL}, we approximate $\mu$ by step profiles
$\mu_m=s\delta_0+\nu_m$ that converge weakly to $\mu$ and whose
cumulative functions $a_m$ converge to $a$ in $L^1(0,1)$. Let $U_m$ be
the corresponding solutions. Following the proof of \cite[(2.19)]{AL}, with $cI_N$
replaced by $\Lambda_b$ and with the first derivatives bounded by one,
we obtain
\[
 \sup|U_m-U|\le\frac{\Tr\Lambda_b}2\,\|a_m-a\|_{L^1(0,1)},
\]
where the supremum is over $q,x,y$, the couplings, and $\theta$. For a
function $f$ of one real variable and every $h>0$, we have
$\|f'\|_\infty\le2\|f\|_\infty/h+(h/2)\|f''\|_\infty$. Using this
inequality in one direction at a time for $U_m-U_{m'}$ and then for its
derivatives, together with the bounds of Step 1, we find that every derivative of $U_m$ in
$(\theta,x,y)$ converges uniformly. This implies that $U$ has
derivatives of every order in $(\theta,x,y)$, that they are the uniform
limits of those of $U_m$, and that they satisfy the same bounds. Each $U_m$
satisfies \eqref{eq:tree-integrated-pde} with $a_m$ in place of $a$. As
$a_m\to a$ in $L^1$, the identity passes to the limit.

\paragraph{Step 3: continuity.}
For a step profile, the derivatives of $U$ in $(\theta,x,y)$ satisfy, on
each grid interval, the equations obtained by differentiating
\eqref{eq:cavityPDE}. By Step 1, their right sides are bounded by
constants depending only on $N$, $c$, $b$, and the order. More precisely,
for every multi-index $\kappa$ in $(\theta,x,y)$, including $\kappa=0$,
differentiation and the product rule give
\[
 \sup\bigl|\partial_qD^\kappa U_m\bigr|\le C_{N,|\kappa|}\,(nc+b)
\]
on each grid interval, where $C_{N,|\kappa|}$ depends only on $N$ and
$|\kappa|$. The derivatives are continuous across the grid points, so the
same Lipschitz bound in $q$ holds on $[0,1]$, independently of the step
approximation, and it passes to the limit. We also find that
$U$ and its derivatives in $(\theta,x,y)$ are Lipschitz in $q$,
with constants depending only on $N$, the order, and $nc+b$. For the continuity in
$b$, let $U_b$ and $U_{b'}$ be the solutions for $b,b'>0$ and the same
profile. The difference $U_b-U_{b'}$ vanishes at $q=1$, because both solutions
have the same terminal value. Subtracting their equations, we see that it
satisfies a linear backward equation with diffusion matrix
$\Lambda_{b'}$, a bounded drift, and the source term
$\frac12(b-b')\bigl[\partial_{yy}U_b+a(\partial_yU_b)^2\bigr]$. By Step 1,
we have $|\partial_{yy}U_b|\le2$ and $|\partial_yU_b|\le1$. The backward
representation then gives $\sup|U_b-U_{b'}|\le\frac32|b-b'|$, first for
step profiles and then, by Step 2, for $\mu$. Using the interpolation
inequality of Step 2 again, we find that every derivative in
$(\theta,x,y)$ is continuous in $b$, uniformly in the other variables.
These derivatives are also Lipschitz in $(\theta,x,y)$ and in $q$, so
they are jointly continuous.
\end{proof}

\subsection{The density of a tree law}

Let $\mathcal T$ be a marked tree. For a vertex $v$, let $l_v$ be the
number of its outgoing edges if $q_v<1$, and the number of leaves at $v$
if $v$ is a terminal point. We define the measure $\omega$ on the points
of $\mathcal T$ as follows. On the interior of an edge whose endpoints
have levels $r<h$, it is the copy of the restriction of $\nu$ to $(r,h)$.
A vertex $v$ carries the atom
\[
 \omega_v=\omega(\{v\})=\nu(\{q_v\})+(l_v-1)\,a(q_v).
\]
For example, the root carries $\nu(\{0\})+(l-1)a(0)$ if it has $l$
outgoing edges, and a terminal point with $l$ leaves carries
$\nu(\{1\})+l-1$. Below, the notation $\sum_v\omega_vF(v)$ means
$\int F\dd\omega$; for an atomic profile, it is an ordinary sum. The
\emph{driftless Gaussian field} on $\mathcal T$ starts at the origin at
the root and has independent increments along the edges. On each edge,
its increments have covariance $\Lambda$ times the length of the level
interval.

\begin{lemma}[Density of a tree law]\label{lem:tree-density}
Let $\mu$ be a deterministic profile, $u\in[0,1]$, and $\mathcal T$ a
marked tree with $k$ leaves. Consider the product of the original
Gaussian law of the couplings, the driftless Gaussian field on
$\mathcal T$, and counting measure on the spins. With respect to this
reference measure, the joint law of the couplings, the field on
$\mathcal T$, and the spins of the leaves under $\Pp_{\mathcal T,u}$ has
density $Z_s(u)^{-1}\mathscr D_{\mathcal T}$, where
\begin{equation}\label{eq:treedensity}
 \mathscr D_{\mathcal T}=\exp\Bigl\{
       \sum_{\mathsf a=1}^k\mathcal E(\sigma^{\mathsf a};
                  X^{\mathsf a}_1,Y^{\mathsf a}_1)
       -\int_{\mathcal T}U(q_v,X_v,Y_v)\,\omega(\dd v)\Bigr\},
 \qquad \omega\ge0,\qquad \omega(\mathcal T)=k-s.
\end{equation}
For fixed couplings, the conditional density of the field and the spins
is $e^{-sU(0,0,0)}\mathscr D_{\mathcal T}$.
\end{lemma}

\begin{proof}
The proof has three steps: we compute the likelihood of one edge, collect
the coefficients of $U$ at the vertices, and compute the mass of
$\omega$. We fix the couplings. By \cref{lem:tree-regularity}, the first
derivatives of $U$ in $(x,y)$ are bounded by one. The drift
$a(q)\Lambda\nabla U$ is bounded, with
$a(q)^2|\Lambda^{1/2}\nabla U|^2\le nc+b(u)$, and Novikov's condition
holds.

\paragraph{Step 1: one edge.}
Let $Z$ be the driftless field on an edge with levels $r<h$, so that
$\dd Z_q=\Lambda^{1/2}\dd B_q$. By Girsanov's formula, the law of the
diffusion on this edge, started from $Z_r$, has density
\[
 \exp\Bigl\{\int_r^ha(q)\nabla U(q,Z_q)\cdot\dd Z_q
   -\frac12\int_r^ha(q)^2|\Lambda^{1/2}\nabla U(q,Z_q)|^2\dd q\Bigr\}
\]
with respect to the law of $Z$. By \cref{lem:tree-regularity}, $U$ is
smooth in $(x,y)$ with bounded derivatives and satisfies
\eqref{eq:tree-integrated-pde}, whose integrand is bounded; in
particular, $U$ is Lipschitz in $q$. We may therefore apply Itô's formula
(for instance after smoothing $U$ in $q$), and by
\eqref{eq:tree-integrated-pde}, we have
\[
 \dd U(q,Z_q)=\nabla U(q,Z_q)\cdot\dd Z_q
       -\frac{a(q)}2|\Lambda^{1/2}\nabla U(q,Z_q)|^2\dd q.
\]
In the last identity, we used that the time derivative of $U$ is given by
\eqref{eq:cavityPDE} for almost every $q$. We conclude that the exponent of the density is
$\int_r^ha(q)\dd U(q,Z_q)$. The process $U_q=U(q,Z_q)$ is continuous,
and $a$ is right-continuous with Stieltjes measure $\nu$ on $(r,h]$. We
integrate by parts. The countably many jumps of $a$ do not affect the
integral against $U_q$, and we obtain
$\int_r^ha\dd U_q=a(h)U_h-a(r)U_r-\int_{(r,h]}U_q\,\nu(\dd q)$. As
$a(h)-\nu(\{h\})=a(h-)$, the edge likelihood is
\begin{equation}\label{eq:cav-edge-likelihood}
 \exp\Bigl\{a(h-)U(h,X_h,Y_h)-a(r)U(r,X_r,Y_r)
                    -\int_{(r,h)}U(q,X_q,Y_q)\,\nu(\dd q)\Bigr\},
\end{equation}
where we wrote the field as $(X,Y)$. We write the values at the endpoints
separately, so that this formula also covers atoms of $\nu$ at $r$ and
$h$.

\paragraph{Step 2: the vertices.}
We multiply the root tilt $e^{sU(0,0,0)}/Z_s(u)$, the edge likelihoods
\eqref{eq:cav-edge-likelihood}, and the Gibbs probabilities
$e^{\mathcal E(\sigma^{\mathsf a};X^{\mathsf a}_1,Y^{\mathsf a}_1)
-U(1,X^{\mathsf a}_1,Y^{\mathsf a}_1)}$ of the leaves. The open edges
contribute $-\int U\dd\nu$ along each edge. At a vertex $v$ with
$0<q_v<1$, the incoming edge contributes $a(q_v-)U(v)$, and the $l_v$
outgoing edges contribute $-l_va(q_v)U(v)$. The coefficient of $U(v)$ is
\[
 a(q_v-)-l_va(q_v)=-\bigl[\nu(\{q_v\})+(l_v-1)a(q_v)\bigr]=-\omega_v .
\]
At the root, the tilt and the $l$ outgoing edges give
$s-la(0)=-[\nu(\{0\})+(l-1)a(0)]$, since $a(0)=s+\nu(\{0\})$. At a
terminal point with $l$ leaves, the incoming edge and the $l$ Gibbs
probabilities give $a(1-)-l=-[\nu(\{1\})+l-1]$, since $a(1)=1$. All these
coefficients are nonpositive, which shows that $\omega\ge0$, and the
product is $Z_s(u)^{-1}\mathscr D_{\mathcal T}$ with
$\mathscr D_{\mathcal T}$ as in \eqref{eq:treedensity}. The edge
likelihoods and the Gibbs probabilities are normalized. Integrating
$\mathscr D_{\mathcal T}$ over the driftless field and the spins thus
gives $e^{sU(0,0,0)}$, and this proves the claim about fixed couplings.
Averaging over the original law of the couplings, we obtain $Z_s(u)$; in
particular, the normalizer is $Z_s(u)$ for every tree.

\paragraph{Step 3: the mass.}
For a single leaf, $\omega$ is a copy of $\nu$ on the path from the root
to the terminal point, and $\omega(\mathcal T)=\nu([0,1])=1-s$. Every
marked tree is obtained from a single path by successively attaching
$l-1\ge1$ new continuations at a point at some level $q$, each running to
level $1$. Such a step adds $l-1$ leaves. It adds $(l-1)a(q)$ to the
weight of the point, which is or becomes a vertex, and $(l-1)\nu((q,1])$
on the new continuations. Since $a(q)+\nu((q,1])=1$, the mass of $\omega$
increases by $l-1$, and $\omega(\mathcal T)=1-s+(k-1)=k-s$.
\end{proof}

We also see from the density \eqref{eq:treedensity} that an \emph{extended}
tree, obtained by attaching new leaves to $\mathcal T$, carries the
weights of its own shape, and that by deletion it induces
$\Pp_{\mathcal T,u}$ on the original leaves. The same proof applies when
the row is a deterministic coupling $\epsilon\sum_i\theta_i\tau_i$ and
$b(u)$ is replaced by a number $b>0$, as in \cref{lem:tree-regularity}.
We then take the density with respect to the Gaussian law of the
couplings $g_{ij}$, $i<j\le n$, and the normalizer is the expectation of
$e^{sU(0,0,0)}$ over these couplings.

\subsection{Responses to the row and to the last-field variance}

Let the field $(X,Y)$ and the spin $\sigma$ have the law of the fixed
system of \cref{lem:tree-regularity}, and let $\calG_q$ be generated by
the field up to level $q$. We set
\[
 \pi_i(q)=\E[\tau_i\epsilon\mid\calG_q],\qquad
 m_N(q)=\E[\epsilon\mid\calG_q],
\]
and we evaluate the derivatives of $U$ at $(q,X_q,Y_q)$, writing for
instance $\partial_{\theta_i}U(q)$.

\begin{lemma}[Parameter responses]\label{lem:tree-responses}
Let $\mu$ be a deterministic profile, and let $U$, $\pi_i$, and $m_N$ be
as above, for fixed couplings, $\theta\in\R^n$, and $b>0$. Then for every
$q\in[0,1]$ and $i\le n$,
\begin{equation}\label{eq:nodeTheta}
 \partial_{\theta_i}U(q)=\pi_i(q),\qquad
 \partial_{\theta_i}^2U(q)
 =1-a(q)\pi_i(q)^2-\int_{(q,1]}\E[\pi_i(r)^2\mid\calG_q]\,\nu(\dd r),
\end{equation}
\begin{equation}\label{eq:tree-node-y}
 \partial_yU(q)=m_N(q),\qquad
 \partial_{yy}U(q)
 =1-a(q)m_N(q)^2-\int_{(q,1]}\E[m_N(r)^2\mid\calG_q]\,\nu(\dd r).
\end{equation}
Further, the derivative of $U$ in $b$ exists and, for each fixed $q$, is
jointly continuous in $(b,\theta,x,y)$. It satisfies
\begin{equation}\label{eq:nodeb}
 \partial_bU(q)=\frac12\Bigl[(1-q)-\int_{(q,1]}(r-q)
                   \E[m_N(r)^2\mid\calG_q]\,\nu(\dd r)\Bigr].
\end{equation}
In particular, $0\le\partial_bU(q)\le(1-q)/2$.
\end{lemma}

For a single row, \eqref{eq:nodeTheta} and \eqref{eq:tree-node-y} are
the conditional form of \cite[(2.10)]{AL}. The identity \eqref{eq:nodeb}
is the conditional, one-coordinate form of the field-variance computation
following \cite[(2.11)]{AL}. Under $\Pp_u$ and the tree laws, the
identities hold with $\calG_q$ generated by the couplings and the field,
with $\theta=(\alpha u/n)^{1/2}(g_{1N},\ldots,g_{nN})$ and $b=b(u)$.

\begin{proof}
The proof has four steps. We prove \eqref{eq:nodeTheta} and
\eqref{eq:tree-node-y} for step profiles with the density
\eqref{eq:tree-future-density}, and we pass to arbitrary profiles with
the exact density of \cref{lem:tree-density}. We then treat the
derivative in $b$ in the same two stages.

\paragraph{Step 1: step profiles.}
Let $\mu$ be a step profile, and let $\eta$ be one of the features
$\tau_i\epsilon$ and $\epsilon$, with parameter $z$. By
\eqref{eq:tree-first-derivative}, we have
$\partial_zU(r)=\E[\eta\mid\calG_r]$ for every $r$, and we write
$\varpi(r)$ for this conditional expectation. Differentiating
\eqref{eq:tree-first-derivative} once more under the density
\eqref{eq:tree-future-density}, we find
\[
 \partial_z^2U(q)=\E\bigl[\eta\,\partial_z\ell_z\bigm|\calG_q\bigr]
 =\E[\eta^2\mid\calG_q]-a(q)\varpi(q)^2
   -\int_{(q,1]}\E\bigl[\eta\,\partial_zU(r)\bigm|\calG_q\bigr]\nu(\dd r).
\]
For $r\ge q$, by the tower property, we have
$\E[\eta\,\partial_zU(r)\mid\calG_q]=\E[\varpi(r)\E[\eta\mid\calG_r]\mid\calG_q]
=\E[\varpi(r)^2\mid\calG_q]$. Since $\eta^2=1$, this proves
\eqref{eq:nodeTheta} for $\eta=\tau_i\epsilon$ and
\eqref{eq:tree-node-y} for $\eta=\epsilon$. Note that $\mu([0,q])=a(q)$
and $\mu=\nu$ on $(q,1]$. We record the result in the form
\begin{equation}\label{eq:tree-atom-safe}
 \partial_z^2U(q)=1-\int_{[0,1]}\E\bigl[\varpi(q\vee r)^2\bigm|\calG_q\bigr]
   \,\mu(\dd r).
\end{equation}

\paragraph{Step 2: arbitrary profiles.}
Let $\mu_m\Rightarrow\mu$ be the step profiles of the proof of
\cref{lem:tree-regularity}. Applying the proof of \cref{lem:tree-density} to
the single edge from $q$ to $1$ followed by the Gibbs probability, we see
that \eqref{eq:tree-future-density} with
$z=0$ is the density of the future given the field at $q$ for every
profile. Its exponent equals
$\mathcal E(\sigma;Z_1)-\int_{[0,1]}U(q\vee r,Z_{q\vee r})\,\mu(\dd r)$,
and the integrand is continuous in $r$. The reference law does not depend
on the profile, $U_m\to U$ uniformly, and $\mu_m\Rightarrow\mu$. Then
the densities for $\mu_m$ converge pointwise to the density for $\mu$.
Writing $L_m$ for the density for $\mu_m$, we have
\[
 \log L_m=\mathcal E(\sigma;Z_1)-U_m(1,Z_1)
   +\int_{[0,1]}\bigl[U_m(1,Z_1)-U_m(q\vee r,Z_{q\vee r})\bigr]\,
   \mu_m(\dd r),
\]
since $\mu_m([0,1])=1$. The first term is nonpositive, because
$U_m(1,\cdot)$ is the logarithm of a sum that contains
$e^{\mathcal E(\sigma;\cdot)}$. Using this, the bounds on the first derivatives in
\cref{lem:tree-regularity}, and the Lipschitz bound in $q$ from Step~3 of
its proof, we find
\[
 \log L_m\le C_N\Bigl(nc+b+\sup_{q\le r\le1}|Z_r|\Bigr),
\]
where $C_N$ depends only on $N$. The same bound holds for the density for
$\mu$, and the densities have a common Gaussian integrable envelope. They
therefore converge in $L^1$ of the reference law, and expectations of
bounded functions converge. The left side of \eqref{eq:tree-first-derivative} converges by
\cref{lem:tree-regularity}, so $\partial_zU(r)=\E[\eta\mid\calG_r]$ also
for $\mu$. Next, we consider the integrand
$r\mapsto\E[\varpi(q\vee r)^2\mid\calG_q]$ in \eqref{eq:tree-atom-safe}.
Its value for $\mu_m$ differs from its value for $\mu$ by at most the
$L^1$ distance of the two densities plus twice the supremum of the
difference of the first derivatives. It therefore converges uniformly
in $r$.
The limit is continuous in $r$, because the first derivatives are
jointly continuous by \cref{lem:tree-regularity}. Together with the weak
convergence of $\mu_m$, this proves \eqref{eq:tree-atom-safe}, and hence
\eqref{eq:nodeTheta} and \eqref{eq:tree-node-y}, for $\mu$. In these
identities, an atom of $\nu$ at $q$ is counted in $a(q)$, and an atom at
$1$ is counted in the integral.

\paragraph{Step 3: the derivative in $b$ for step profiles.}
Let $\mu$ be a step profile and $q_\ell\le q<q_{\ell+1}$, with $a=a_\ell$
on $[q_\ell,q_{\ell+1})$. Suppose first that $a_\ell>0$, and let
$\phi=e^{a_\ell U}$. On this interval, $\phi$ solves the backward heat
equation $\partial_q\phi+\frac12\Tr(\Lambda_b\nabla^2\phi)=0$. In particular,
$\phi(q)$ is the Gaussian average of $\phi(q_{\ell+1})$ over an increment
of covariance $(q_{\ell+1}-q)\Lambda_b$. We differentiate in $b$, using
the Gaussian heat identity in the last coordinate, and divide by
$a_\ell\phi(q)$. We obtain
\[
 \partial_bU(q)=\E\bigl[\partial_bU(q_{\ell+1})\bigm|\calG_q\bigr]
  +\frac{q_{\ell+1}-q}2\,
  \E\bigl[(\partial_{yy}U+a_\ell(\partial_yU)^2)(q_{\ell+1})\bigm|\calG_q\bigr].
\]
The function
$\partial_{yy}\phi=a_\ell(\partial_{yy}U+a_\ell(\partial_yU)^2)\phi$ also
solves the backward heat equation. Further, the transition of the field
from $q$ to $r\le q_{\ell+1}$ is the Gaussian transition tilted by
$\phi(r)/\phi(q)$. It follows that
$\E[(\partial_{yy}U+a_\ell(\partial_yU)^2)(r)\mid\calG_q]$ does not
depend on $r\in[q,q_{\ell+1}]$, and that the last term equals
$\frac12\int_q^{q_{\ell+1}}\E[(\partial_{yy}U+a(\partial_yU)^2)(r)\mid\calG_q]\dd r$.
The case $a_\ell=0$ is the same with $\phi=U$. Iterating over the grid
intervals and using $\partial_bU(1)=0$, we find
\begin{equation}\label{eq:tree-b-response}
 \partial_bU(q)=\frac12\int_q^1
   \E\bigl[\partial_{yy}U(r)+a(r)\,m_N(r)^2\bigm|\calG_q\bigr]\dd r,
\end{equation}
where we used that $\partial_yU(r)=m_N(r)$.

\paragraph{Step 4: the derivative in $b$ for arbitrary profiles.}
By \eqref{eq:tree-node-y}, the integrand in \eqref{eq:tree-b-response}
lies in $[0,1]$. Let $R(b)$ be the right side of
\eqref{eq:tree-b-response} for $\mu$. For the step profiles $\mu_m$, the
right side of \eqref{eq:tree-b-response} converges to $R(b)$ at each
fixed $b$. Indeed, the densities converge as in Step 2, the derivatives
converge uniformly, and $a_m\to a$ in $L^1(0,1)$. These expressions are
bounded by $1/2$. By dominated convergence in the identity
$U_m(b_2)-U_m(b_1)=\int_{b_1}^{b_2}\partial_bU_m\dd b$, we obtain
$U(b_2)-U(b_1)=\int_{b_1}^{b_2}R(b)\dd b$. We write the last-field
increments of the reference law as $\sqrt b$ times a fixed Brownian
motion. Using \cref{lem:tree-regularity} and dominated convergence, we then see
that, for each fixed $q$, $R$ is jointly continuous in
$(b,\theta,x,y)$. It follows that
$\partial_bU$ exists, equals $R$, and satisfies
\eqref{eq:tree-b-response} for $\mu$. Inserting \eqref{eq:tree-node-y}
at level $r$ and using the tower property, we have
\[
 \partial_bU(q)=\frac12\int_q^1\Bigl(1-\int_{(r,1]}
   \E[m_N(\ell)^2\mid\calG_q]\,\nu(\dd\ell)\Bigr)\dd r .
\]
By Fubini's theorem, we have $\int_q^1\dd r\int_{(r,1]}F(\ell)\,\nu(\dd\ell)
=\int_{(q,1]}(\ell-q)F(\ell)\,\nu(\dd\ell)$ for bounded $F\ge0$, and
\eqref{eq:nodeb} follows. Finally, the integral in \eqref{eq:nodeb} lies
between $0$ and $(1-q)\nu((q,1])\le1-q$, which gives
$0\le\partial_bU(q)\le(1-q)/2$.
\end{proof}

\subsection{Differentiation in the interpolation parameter}

For a point $v$ of $\mathcal T$, we write $\hat v$ and $\check v$ for new
leaves whose paths leave the path through $v$ at the level $q_v$. If $v$
is an interior point of an edge, the extended tree has a new vertex at
$v$. If $v$ is a vertex, the new edge starts at $v$, and if $v$ is a
terminal point, the new leaf is attached there. Distinct new leaves are
conditionally independent, also when they are attached at the same
point. For $r>q_v$, we write $\hat r$ and $\check r$ for two new leaves
that follow a common new branch from $v$ to the level $r$ and split
there. Given the couplings and the field of $\mathcal T$, a new branch
from $v$ is a continuation of the diffusion from the field at $v$, and it
is independent of everything else. Let $F$ be a bounded function of the
spins of $\mathcal T$, and let $G,G'$ be bounded functions of one spin.
Then
\begin{equation}\label{eq:tree-replacement}
 \begin{aligned}
 \E\bigl[F\,G(\sigma^{\hat v})\bigr]
 &=\E_{\mathcal T,u}\bigl[F\,\E[G(\sigma)\mid\calG_{q_v}](v)\bigr],\\
 \E\bigl[F\,G(\sigma^{\hat r})G'(\sigma^{\check r})\bigr]
 &=\E_{\mathcal T,u}\Bigl[F\,\E\bigl[\E[G(\sigma)\mid\calG_r]\,
   \E[G'(\sigma)\mid\calG_r]\bigm|\calG_{q_v}\bigr](v)\Bigr].
 \end{aligned}
\end{equation}
Here the left sides are taken under the extended trees. A conditional
expectation under $\Pp_u$ given $\calG_{q_v}$ is a function of the
couplings and of the field at level $q_v$, and $(v)$ means that it is
evaluated at the couplings and at the field at $v$. In the same way, two
new leaves $\hat v$ and $\check w$ attached at $v$ and $w$ represent the
product of the two conditional expectations, evaluated at $v$ and at $w$.

Let $\varphi_{\mathsf{ab}}$ be a quantity attached to each pair of
distinct leaves of an extended tree. For a bounded function $F$ of the
spins of $\mathcal T$, we define the five-pattern sum
\[
 \begin{aligned}
 \mathfrak S_\varphi(F)={}&
 \sum_{\mathsf a\ne\mathsf b}\E(F\epsilon_{\mathsf a}\epsilon_{\mathsf b}
   \varphi_{\mathsf{ab}})
 -2\sum_{\mathsf a,v}\omega_v\E(F\epsilon_{\mathsf a}\epsilon_{\hat v}
   \varphi_{\mathsf a\hat v})
 +\sum_{v,w}\omega_v\omega_w\E(F\epsilon_{\hat v}\epsilon_{\check w}
   \varphi_{\hat v\check w})\\
 &+\sum_v\omega_va(q_v)\E(F\epsilon_{\hat v}\epsilon_{\check v}
   \varphi_{\hat v\check v})
 +\sum_v\omega_v\int_{(q_v,1]}\E(F\epsilon_{\hat r}\epsilon_{\check r}
   \varphi_{\hat r\check r})\,\nu(\dd r).
 \end{aligned}
\]
The sums run over leaves $\mathsf a,\mathsf b$ of $\mathcal T$ and points
$v,w$ of $\mathcal T$, and the weights $\omega$ are those of
$\mathcal T$. We take each expectation under the extended tree that
contains the leaves involved. We use this sum with
$\varphi_{\mathsf{ab}}=R^-_{\mathsf{ab}}$, with
$\varphi_{\mathsf{ab}}=q_{\mathsf{ab}}$, and with
$\varphi_{\mathsf{ab}}=\Delta_{\mathsf{ab}}$. Then
$\mathfrak S_\Delta=\mathfrak S_{R^-}-\mathfrak S_q$, and
$\epsilon_{\mathsf a}\epsilon_{\mathsf b}\Delta_{\mathsf{ab}}=D_{\mathsf{ab}}$.
The following table lists the five patterns. Its last column gives the
level $q_{\mathsf{ab}}$ of the pair.
\begin{center}
\small
\begin{tabular}{@{}llll@{}}
\toprule
pair of leaves & new leaves & coefficient & level of the pair\\
\midrule
two leaves $\mathsf a\ne\mathsf b$ of $\mathcal T$ & none & $1$
  & $q_{\mathsf{ab}}$\\
a leaf $\mathsf a$ and a leaf attached at $v$ & $\hat v$
  & $-2\,\omega_v$ & $q_{\mathsf a\hat v}$\\
leaves attached at $v$ and at $w$ & $\hat v,\check w$
  & $\omega_v\omega_w$ & $q_{\hat v\check w}$\\
two leaves attached at $v$ & $\hat v,\check v$ & $a(q_v)\,\omega_v$
  & $q_v$\\
two leaves splitting at $r>q_v$ & $\hat r,\check r$
  & $\omega_v\,\nu(\dd r)$ & $r$\\
\bottomrule
\end{tabular}
\end{center}

\begin{lemma}[Differentiation under the tree law]\label{lem:streaming}
Let $\mu$ be a deterministic profile, $\mathcal T$ a marked tree with $k$
leaves, and $F$ a bounded function of the spins of its leaves. Then
$u\mapsto\E_{\mathcal T,u}F$ belongs to $C^2([0,1])$, with one-sided
derivatives at the endpoints, and
\begin{align}
 \frac{\dd}{\dd u}\E_{\mathcal T,u}F=\frac\alpha2\bigg[&
 \sum_{\mathsf a\ne\mathsf b}\E(FD_{\mathsf{ab}})
 -2\sum_{\mathsf a,v}\omega_v\E(FD_{\mathsf a\hat v})
 +\sum_{v,w}\omega_v\omega_w\E(FD_{\hat v\check w})\nonumber\\
 &+\sum_v\omega_va(q_v)\E(FD_{\hat v\check v})
 +\sum_v\omega_v\int_{(q_v,1]}\E(FD_{\hat r\check r})\,\nu(\dd r)
 \nonumber\\
 &+s\,\E_{\mathcal T,u}F\int\E_{2,r,u}D_{12}\,\nu(\dd r)\bigg].
 \label{eq:streaming}
\end{align}
The sums run over leaves $\mathsf a,\mathsf b$ and points $v,w$ of
$\mathcal T$, with the weighted sums interpreted as integrals against
$\omega$ as above. The first five expectations are taken under extended
trees, and the last term is a product of expectations over independent
copies. The absolute coefficients in the bracket sum to at most
$4k^2+2k+1$. The second derivative is obtained by applying the same
formula to each tree moment on its right side, and the product rule to
independent products of moments.
\end{lemma}

In the notation above, \eqref{eq:streaming} reads
$\frac{\dd}{\dd u}\E_{\mathcal T,u}F=\frac\alpha2[\mathfrak S_\Delta(F)
+s\,\E_{\mathcal T,u}F\int\E_{2,r,u}D_{12}\,\nu(\dd r)]$.

\begin{proof}
If $\alpha=0$, the law does not depend on $u$, and both sides vanish. Let
$\alpha>0$. The proof has five steps. We check
that \cref{lem:fnd-path-heat} applies, and we compute the contributions
of the row and of the last field to the derivative of the unnormalized
moment. We then compute the derivative of the normalizer. Finally, we
bound the coefficients and prove the $C^2$ regularity.

\paragraph{Step 1: the heat identity.}
By \cref{lem:tree-density}, we have
$\E_{\mathcal T,u}F=Z_s(u)^{-1}\E_{\rm ref}[F\mathscr D_{\mathcal T}]$,
where $\E_{\rm ref}$ denotes integration against the reference measure of
that lemma. We write the row as $\theta=(\alpha u/n)^{1/2}(g_{iN})_{i\le n}$,
and we use \cref{lem:fnd-path-heat} on $I=[0,1]$ with the following
data. The finite coordinates are the entries of $\theta$, with variances
$\alpha u/n$. The field paths are the $n+1$ coordinates of the driftless
field on the edges of $\mathcal T$. For an edge $e$ with endpoint levels
$r_e<h_e$, we use the local coordinate $v\in[0,1]$, and the clocks are
\[
 T_{e,i,u}(v)=c\,(h_e-r_e)\,v\quad(i\le n),\qquad
 T_{e,N,u}(v)=b(u)\,(h_e-r_e)\,v .
\]
The couplings $g_{ij}$, $i<j\le n$, do not depend on $u$, and the spins
are the counting variable. The row, the Brownian motions of the
different edges and components, and the couplings $g_{ij}$ are mutually
independent. We take $\ell=\log\mathscr D_{\mathcal T}$ and
$\Phi=F$. The elements of the space $\mathcal C$ of
\cref{lem:fnd-path-heat} need not agree at the common
endpoints of edges. To define $\ell$ on the whole space, we evaluate the field at each nonroot vertex at the
endpoint of its incoming edge. At the root, we choose a fixed outgoing
edge $e_0$ and evaluate the field at the initial point of $e_0$; its
coordinates vanish on every path of the driftless field. This agrees with
\eqref{eq:treedensity} on the paths of the driftless field. The
derivative measures below, including the atom of $\omega$ at the root,
are those of this extension. We check the hypotheses \textup{(A1)}--\textup{(A4)} of
\cref{lem:fnd-path-heat}.
\begin{itemize}
\item \textup{(A1)}: the variances $\alpha u/n$ are positive for $u>0$.
The clocks are affine in $u$, and they are strictly increasing in $v$,
because $c>0$, $b(u)\ge c$, and every edge has positive length.
\item \textup{(A2)}: the covariance of the last component at two field
points $v,w$ is $b(u)h(v,w)$, where $h(v,w)$ is the level of the last
common point of the paths from the root to $v$ and to $w$. For the other
components, it is $c\,h(v,w)$. Hence $K$ is affine in $u$, with
$\dot K=(\alpha/n)I_n$ on the row block, $\dot K=-\alpha h$ on the
last-field block, and $\dot K=0$ elsewhere. These kernels are
continuous.
\item \textup{(A3)}: $\Phi=F$ is bounded and $\dot\Phi=0$. The explicit
derivative of $\ell$ comes from $b(u)$ in $U$, so that
$\dot\ell=-b'(u)\sum_v\omega_v\partial_bU(v)=\alpha\sum_v\omega_v
\partial_bU(v)$. The levels $q_v$ and the finite measure $\omega$ are
fixed. Using the fixed-level continuity and uniform bound of
\cref{lem:tree-responses}, together with dominated convergence against $\omega$, we see
that $\dot\ell$ is continuous in $(u,G)$; it is bounded by $\alpha k/2$.
If we treat each rest coupling $g_{ij}$ as an additional source parameter
with bounded feature $\beta\sqrt{t/N}\,\tau_i\tau_j$ in Step~1 of the proof of
\cref{lem:tree-regularity}, we obtain joint continuity of $U$ and its spatial
derivatives also in the rest couplings. The fixed-level response formula
of \cref{lem:tree-responses} then gives joint continuity of $\partial_bU$
in those couplings by dominated convergence, using the uniform bound of
that lemma.
Thus $\ell$ and $\dot\ell$ are jointly continuous in $(u,G,g)$, as required
by \textup{(A3)}.
The first derivatives of $U$ in
the row and in the field are bounded by one. Its derivative in $g_{ij}$,
$i<j\le n$, is the conditional mean of $\beta\sqrt{t/N}\,\tau_i\tau_j$,
by the argument of Step~1 of the proof of \cref{lem:tree-regularity}.
Integrating \eqref{eq:cavityPDE} in $q$ and using that the first spatial
derivatives of $U$ are bounded by one and the second by two (Step~1 of that
proof), we have $|U(q,0,0)|\le N\log2+\frac32(nc+b(u))$
when all couplings, the row, and the field vanish.
Since $\mathcal E$ is linear in the couplings, the row, and the field, we
conclude that $|\ell|\le C(1+|g|+|\theta|+\|G^{\rm path}\|_\infty)$, where
$C$ depends only on $N$, $k$, and $\beta$. Here $g$ collects the couplings
$g_{ij}$, $i<j\le n$.
\item \textup{(A4)}: the score of $\ell$ consists of the atoms
$\tau_i^{\mathsf a}\epsilon_{\mathsf a}$ (row), $\tau^{\mathsf a}$ and
$\epsilon_{\mathsf a}$ (field at the terminal point of $\mathsf a$), and
of $-\nabla U(v)\,\omega(\dd v)$, where $\nabla$ is the gradient in the
row and the field. Since $\mathcal E$ is linear in these variables, the
Hessian of $\ell$ is $-\nabla^2U(v)\,\omega(\dd v)$, placed at the pairs
of indices of the point $v$. It has a row--row block
$-\partial_{\theta_i\theta_j}U(v)\,\omega(\dd v)$, mixed row--field
blocks such as $-\partial_{\theta_iy}U(v)\,\omega(\dd v)$ and
$-\partial_{\theta_ix_j}U(v)\,\omega(\dd v)$ with their transposes, and
field--field blocks such as $-\partial_{yy}U(v)\,\omega(\dd v)$, placed at
the diagonal pairs of index points of $v$. By
\cref{lem:tree-regularity}, these derivatives of $U$ are bounded and
jointly continuous in $(b,q,\theta,x,y)$. As the points $(q_v,X_v,Y_v)$
stay in a compact set when $G$ varies in a bounded set, the measures are
continuous in total variation in $(u,G)$. Their total variations are
at most $C_Nk$. The measures for $\Phi$ vanish.
\end{itemize}
The constants in these hypotheses depend on $N$ and $k$. We use them
only to justify the differentiation, and the identity below has
coefficients that do not depend on $N$. Since $\Phi=F$ depends neither on
the Gaussian coordinates nor on $u$, \eqref{eq:fnd-path-heat} and
\eqref{eq:fnd-path-normalizer} imply that
\begin{equation}\label{eq:tree-quotient}
 \frac{\dd}{\dd u}\E_{\mathcal T,u}F
 =\frac{\partial_u\E_{\rm ref}[F\mathscr D_{\mathcal T}]}{Z_s(u)}
 -\E_{\mathcal T,u}F\,\frac{\dd}{\dd u}\log Z_s(u),
\end{equation}
where the derivative of the unnormalized moment is
\[
 \frac{\partial_u\E_{\rm ref}[F\mathscr D_{\mathcal T}]}{Z_s(u)}
 =\E_{\mathcal T,u}\Bigl[F\Bigl(\dot\ell+\frac\alpha{2n}\sum_{i=1}^n
   \bigl((\partial_{\theta_i}\ell)^2+\partial_{\theta_i}^2\ell\bigr)
   -\frac\alpha2\bigl\langle h,S_y\otimes S_y+M_{yy}\bigr\rangle\Bigr)\Bigr].
\]
Here $S_y$ and $M_{yy}$ are the score and Hessian measures of $\ell$ in
the last-field coordinates.

\paragraph{Step 2: the row.}
By \eqref{eq:treedensity} and \eqref{eq:nodeTheta}, we have
\[
 \partial_{\theta_i}\ell=\sum_{\mathsf a}\tau_i^{\mathsf a}\epsilon_{\mathsf a}
   -\sum_v\omega_v\pi_i(v),
 \qquad
 \partial_{\theta_i}^2\ell=-\sum_v\omega_v\partial_{\theta_i}^2U(v),
\]
where $\pi_i(v)$ and $\partial_{\theta_i}^2U(v)$ are evaluated at the
couplings and at the field at $v$. We expand the square and insert
\eqref{eq:nodeTheta}. The diagonal terms
$(\tau_i^{\mathsf a}\epsilon_{\mathsf a})^2=1$ contribute $k$, and the
constant $1$ in \eqref{eq:nodeTheta} contributes $-\omega(\mathcal T)$.
By \cref{lem:tree-density}, these constants add up to $s$. The other
terms of $(\partial_{\theta_i}\ell)^2+\partial_{\theta_i}^2\ell$ are
\[
 \begin{aligned}
 &\sum_{\mathsf a\ne\mathsf b}\tau_i^{\mathsf a}\epsilon_{\mathsf a}
   \tau_i^{\mathsf b}\epsilon_{\mathsf b}
 -2\sum_{\mathsf a,v}\omega_v\tau_i^{\mathsf a}\epsilon_{\mathsf a}\pi_i(v)
 +\sum_{v,w}\omega_v\omega_w\pi_i(v)\pi_i(w)\\
 &\quad+\sum_v\omega_va(q_v)\pi_i(v)^2
 +\sum_v\omega_v\int_{(q_v,1]}\E[\pi_i(r)^2\mid\calG_{q_v}](v)\,\nu(\dd r).
 \end{aligned}
\]
By \eqref{eq:tree-replacement}, inside $\E_{\mathcal T,u}[F\,\cdot\,]$
we may replace $\pi_i(v)$ by $\tau_i^{\hat v}\epsilon_{\hat v}$, and
$\pi_i(v)\pi_i(w)$ by the product for $\hat v$ and $\check w$. We may
also replace $\pi_i(v)^2$ by the product for $\hat v$ and $\check v$, and
$\E[\pi_i(r)^2\mid\calG_{q_v}](v)$ by the product for $\hat r$ and
$\check r$. Averaging over $i$ then replaces
$\tau_i^{\mathsf a}\tau_i^{\mathsf b}$ by $R^-_{\mathsf{ab}}$. We obtain
\begin{equation}\label{eq:tree-row-part}
 \frac\alpha{2n}\sum_{i=1}^n\E_{\mathcal T,u}\bigl[F\bigl(
   (\partial_{\theta_i}\ell)^2+\partial_{\theta_i}^2\ell\bigr)\bigr]
 =\frac\alpha2\bigl[s\,\E_{\mathcal T,u}F+\mathfrak S_{R^-}(F)\bigr].
\end{equation}

\paragraph{Step 3: the last field.}
By \eqref{eq:treedensity} and \eqref{eq:tree-node-y}, the score and
Hessian measures of $\ell$ in the last-field coordinates are
\[
 S_y=\sum_{\mathsf a}\epsilon_{\mathsf a}\delta_{\mathsf a}
                      -m_N(v)\,\omega(\dd v),\qquad
 M_{yy}=-\operatorname{diag}_{\#}
                  \bigl(\partial_{yy}U(v)\,\omega(\dd v)\bigr),
\]
where $\delta_{\mathsf a}$ is the unit mass at the terminal point of
$\mathsf a$. Here $\operatorname{diag}_{\#}$ is the image of a measure
under $v\mapsto(v,v)$, so that atoms of $\omega$ stay at exact diagonal
pairs. Let a leaf stand for its terminal point. Then the kernel $h$
satisfies $h(\mathsf a,\mathsf a)=1$,
$h(\mathsf a,\mathsf b)=q_{\mathsf{ab}}$,
$h(\mathsf a,v)=q_{\mathsf a\hat v}$, $h(v,w)=q_{\hat v\check w}$, and
$h(v,v)=q_v$. Expanding $\langle h,S_y\otimes S_y\rangle$ as in Step 2
and replacing $m_N(v)$ by $\epsilon_{\hat v}$ with
\eqref{eq:tree-replacement}, we find
\[
 \begin{aligned}
 -\frac\alpha2\E_{\mathcal T,u}\bigl[F\langle h,S_y\otimes S_y\rangle\bigr]
 =-\frac\alpha2\Bigl[&k\,\E_{\mathcal T,u}F
 +\sum_{\mathsf a\ne\mathsf b}q_{\mathsf{ab}}
   \E(F\epsilon_{\mathsf a}\epsilon_{\mathsf b})\\
 &-2\sum_{\mathsf a,v}\omega_vq_{\mathsf a\hat v}
   \E(F\epsilon_{\mathsf a}\epsilon_{\hat v})
 +\sum_{v,w}\omega_v\omega_wq_{\hat v\check w}
   \E(F\epsilon_{\hat v}\epsilon_{\check w})\Bigr].
 \end{aligned}
\]
The bracket consists of the constant $k\,\E_{\mathcal T,u}F$ and the
first three patterns of $\mathfrak S_q(F)$. The Hessian measure
contributes
$-\frac\alpha2\langle h,M_{yy}\rangle=\frac\alpha2\sum_v\omega_vq_v
\partial_{yy}U(v)$, and we recall that
$\dot\ell=\alpha\sum_v\omega_v\partial_bU(v)$. Together, these two terms
give $\frac\alpha2\sum_v\omega_v[q_v\partial_{yy}U(v)+2\partial_bU(v)]$.
By \eqref{eq:tree-node-y} and \eqref{eq:nodeb}, we have
\[
 q_v\partial_{yy}U(v)+2\partial_bU(v)
 =1-q_va(q_v)m_N(v)^2
   -\int_{(q_v,1]}r\,\E[m_N(r)^2\mid\calG_{q_v}](v)\,\nu(\dd r).
\]
In the integral, we combined $q_v\int_{(q_v,1]}$ from
\eqref{eq:tree-node-y} with $\int_{(q_v,1]}(r-q_v)$ from \eqref{eq:nodeb}.
After the replacements \eqref{eq:tree-replacement}, the constant gives
$\frac\alpha2\omega(\mathcal T)\E_{\mathcal T,u}F
=\frac\alpha2(k-s)\E_{\mathcal T,u}F$. The other two terms are the last
two patterns of $-\frac\alpha2\mathfrak S_q(F)$, with the levels
$q_{\hat v\check v}=q_v$ and $q_{\hat r\check r}=r$. Altogether, the last
field and the explicit derivative contribute
$\frac\alpha2[-k+(k-s)]\E_{\mathcal T,u}F-\frac\alpha2\mathfrak S_q(F)$.
Adding \eqref{eq:tree-row-part}, we see that the constants $\alpha s/2$,
$-\alpha k/2$, and $\alpha(k-s)/2$ cancel, and since
$\mathfrak S_{R^-}-\mathfrak S_q=\mathfrak S_\Delta$, we conclude that
\begin{equation}\label{eq:tree-unnormalized}
 \frac{\partial_u\E_{\rm ref}[F\mathscr D_{\mathcal T}]}{Z_s(u)}
 =\frac\alpha2\,\mathfrak S_\Delta(F).
\end{equation}
In words, the row inserts a rest overlap, the last field inserts the
level of the corresponding pair, and their difference is an overlap
defect.

\paragraph{Step 4: the normalizer.}
We compute the derivative of $Z_s(u)=\E_{\rm ref}e^{sU(0,0,0)}$ with
\eqref{eq:fnd-path-normalizer}; now the only active coordinates are those of the
row. The row contributes
$\frac{\alpha s}{2n}\sum_i\E_u[\partial_{\theta_i}^2U(0)
+s(\partial_{\theta_i}U(0))^2]$. As $a(0)-s=\nu(\{0\})$, we obtain from
\eqref{eq:nodeTheta} that
\[
 \partial_{\theta_i}^2U(0)+s\bigl(\partial_{\theta_i}U(0)\bigr)^2
 =1-\int_{[0,1]}\E[\pi_i(r)^2\mid\calG_0]\,\nu(\dd r).
\]
Here, as in \cite[(2.11)]{AL}, the term $s(\partial_{\theta_i}U)^2$
removes the prescribed root mass $s$, and a possible atom $\nu(\{0\})$
remains in the integral. Since
$\E_u\E[\pi_i(r)^2\mid\calG_0]=\E_{2,r,u}[\tau_i^1\epsilon_1\tau_i^2\epsilon_2]$,
the row contributes
$\frac{\alpha s}2[1-\int\E_{2,r,u}[\epsilon_1\epsilon_2R^-_{12}]\,
\nu(\dd r)]$. The coefficient $b(u)$ contributes
$sb'(u)\E_u\partial_bU(0)=-\alpha s\E_u\partial_bU(0)$. By
\eqref{eq:nodeb}, this equals
$-\frac{\alpha s}2[1-\int r\,\E_{2,r,u}[\epsilon_1\epsilon_2]\,\nu(\dd r)]$.
Adding the two contributions, we find
\begin{equation}\label{eq:normalizerderivative}
 \frac\dd{\dd u}\log Z_s(u)
 =-\frac{\alpha s}{2}\int\E_{2,r,u}D_{12}\,\nu(\dd r).
\end{equation}
Inserting \eqref{eq:tree-unnormalized} and
\eqref{eq:normalizerderivative} into \eqref{eq:tree-quotient}, we obtain
\eqref{eq:streaming}. In its last term, $\E_{2,r,u}$ is the expectation
under an independent copy of the complete tree law, including its
couplings.

\paragraph{Step 5: coefficients and regularity.}
Recall that $\omega(\mathcal T)=k-s\le k$, $a\le1$, and $\nu([0,1])\le1$. The
absolute coefficients of the six terms in \eqref{eq:streaming} are
at most $k(k-1)$, $2k^2$, $k^2$, $k$, $k$, and $1$. Their sum is
$4k^2+k+1\le4k^2+2k+1$.

We now show that each extended-tree moment on the right side of
\eqref{eq:streaming} is continuous in $u$. We write the row as
$(\alpha u/n)^{1/2}g$, and the last-field coordinates as $\sqrt{b(u)}$
times a fixed Brownian field on the tree. By \cref{lem:tree-regularity},
the density is then continuous in $u$ for fixed Gaussian variables. Its
linear growth gives a Gaussian envelope that is uniform in $u$. By
dominated convergence, the moments and their normalizers are continuous.
The weights and the measures $\omega$ and $\nu$ over which the right side
is integrated do not depend on $u$, and the integrands are bounded by
$2\|F\|_\infty$. Bounded convergence then shows that the right side is
continuous. The derivative exists at every $u\in[0,1]$, one-sided at the
endpoints, so $u\mapsto\E_{\mathcal T,u}F$ is $C^1$.
Next, we apply \eqref{eq:streaming} to each bounded spin observable $FD$
on the extended trees; these observables do not depend on $u$. The
numbers of leaves grow by at most two, and the coefficients are bounded
in terms of $k$. Each of these moments therefore has a derivative bounded
by $C_k\alpha\|F\|_\infty$, uniformly in the points $v,w$ and the level
$r$. This lets us differentiate under the integrals against $\omega$ and
$\nu$ by dominated convergence, and we use the product rule for the last
term. The resulting right side is continuous by the same argument. This
proves the $C^2$ regularity and the rule for the second derivative.
\end{proof}

\begin{corollary}[Comparison of nonnegative observables]\label{cor:comparison}
For every deterministic profile, every marked tree $\mathcal T$ with $k$
leaves, and every nonnegative bounded function $F$ of the spins of its
leaves,
\begin{equation}\label{eq:comparison}
 \E_{\mathcal T,u}F\le K_{\beta,k}\E_{\mathcal T,v}F
 \quad(0\le u,v\le1),
 \qquad K_{\beta,k}=\exp\bigl(\beta^2(4k^2+2k+1)\bigr).
\end{equation}
\end{corollary}

\begin{proof}
Since $|D_{\mathsf{ab}}|\le2$ and $F\ge0$, each expectation in
\eqref{eq:streaming} is at most twice the expectation of $F$ under the
corresponding extended tree. By deletion, the latter equals
$\E_{\mathcal T,u}F$, and the same bound holds for the product in the
last term. Using \cref{lem:streaming} and $\alpha\le\beta^2$, we find
\[
 \bigl|\partial_u\E_{\mathcal T,u}F\bigr|
 \le\beta^2(4k^2+2k+1)\,\E_{\mathcal T,u}F,\qquad0\le u\le1 .
\]
By Gronwall's inequality on the interval between $u$ and $v$, we obtain
$\E_{\mathcal T,u}F\le e^{\beta^2(4k^2+2k+1)|u-v|}\E_{\mathcal T,v}F$.
In particular, \eqref{eq:comparison} holds, also in the case
$\E_{\mathcal T,v}F=0$, in which the moment vanishes for all $u$. The
comparisons in \cite[Lemmas~2.9 and~4.3]{OpenAIFluctuationScale} are
proved by the same differential inequality. The classical analog
on one site is \cite[Lemma~1.6.4]{TalagrandVolI2011}.
\end{proof}

\section{Products of three overlap defects}\label{subsec:c7-three}

In the Taylor remainder of the cavity covariance, and in the variance
comparison of \cref{lem:cl-star}, we differentiate tree moments of
observables that already contain overlap defects. Each differentiation
adds one more defect. We show that integrated moments of products of
three defects are controlled by the integrated third absolute moment
\begin{equation}\label{eq:cav-M3}
 \mathcal M_3(u)=\int\E_{2,q,u}|R^-_{12}-q|^3\,\nu(\dd q)
\end{equation}
of the overlap defect of two descendants. In \cref{lem:c7-thermal},
we bound $\mathcal M_3$ by $A$, $J$, and $H_3$.

\paragraph{Terms and marks.}
Fix integers $k_0\ge1$ and $j\ge0$. For each $p\in[0,1]$, let
$\mathcal T_p$ be the tree of $k_0$ leaves that split at $p$, and let
$F_p$ be a bounded function of their spins, jointly measurable in $p$ and
the spins. We consider the integrated tree moment
$\int\nu(\dd p)\E_{\mathcal T_p,u}F_p$. We also consider the terms
produced from it by $j$ applications of \eqref{eq:streaming}, with the
product rule for products of moments over independent trees. Each
application multiplies by $\alpha/2$ and replaces a moment by the six
terms of \eqref{eq:streaming}. Each term produced after $j$ applications
is therefore $(\alpha/2)^j$ times an integral against a signed measure on
finitely many \emph{marks}, which are the following variables. The marks
are $p$, the points $v$ and $w$ at which the differentiations attached
leaves, the levels $r$ at which new branches split, and the split levels
$r$ of the independent trees produced by the last term of
\eqref{eq:streaming}. The integrand is a product of expectations over
independent trees, whose shapes are determined by the marks. The signed
measure collects the coefficients $1$, $-2\omega_v$, $\omega_v\omega_w$,
$a(q_v)\omega_v$, $\omega_v\nu(\dd r)$, and $s\,\nu(\dd r)$ of the
attachments, and the multiplicities of the product rule. We let
$\mathfrak m$ be its total variation. The factor $(\alpha/2)^j$ is not
part of $\mathfrak m$. The level of a mark is $p$, $q_v$, $q_w$, or $r$.

\begin{lemma}[Three-defect bound]\label{lem:cav-threeproducts}
For every deterministic profile, every such term containing three
defects $\Delta_1,\Delta_2,\Delta_3$, of pairs of leaves in the same or
in independent trees, satisfies
\begin{equation}\label{eq:threeproducts}
 \int\E\bigl|\Delta_1\Delta_2\Delta_3\bigr|\,\dd\mathfrak m
 \le C\mathcal M_3(u).
\end{equation}
Here $\E$ is the expectation under the product of the tree laws of the
term, and $C$ depends only on the numbers of leaves and of
differentiations.
\end{lemma}

The bounds on the marks in Step 1 of the proof are a counterpart, for an
arbitrary profile, of the bounds on split levels in
\cite[Lemma~2.8]{OpenAIFluctuationScale}.

\begin{proof}
The proof has three steps. We show that the level of every mark has a
marginal under $\mathfrak m$ dominated by a multiple of $\nu$, and that
the level of every pair of leaves is the level of a mark. We then use
Hölder's inequality.

\paragraph{Step 1: marginals of the marks.}
Each differentiation adds at most two leaves to one tree, or creates an
independent tree with two leaves. After $i$ differentiations, every tree
of a term has at most $K_i=k_0+2i$ leaves, and a product has at
most $k_0+i$ factors. Let $\mathfrak m_i$ be the measure $\mathfrak m$ of
a term after $i\le j$ differentiations. We show by induction on $i$ that
$\mathfrak m_i$ has total mass at most $C_i$, and that the level of each
of its marks has a marginal at most $C_i\nu$. Here $C_i$ depends only on
$k_0$ and $i$. For $i=0$, the only mark is $p$, and
$\mathfrak m_0=\nu(\dd p)$.

For the induction step, let $K=K_i$. We first bound the image of the
weights $\omega$ of a tree with at most $K$ leaves under the level map
$v\mapsto q_v$. At most $K$ points of the tree lie at any given level.
The copies of $\nu$ on the edges and the parts $\nu(\{q_v\})$ of the
vertex weights contribute at most $K\nu$. The parts
$(l_v-1)a(q_v)\le l_v-1$ contribute at most
$\sum_v(l_v-1)\delta_{q_v}$, where $v$ runs over the vertices with
$l_v\ge2$, and $\sum_v(l_v-1)\le K$. Note that every such $q_v$ is the
level of an earlier mark. We now consider one more differentiation. By
\cref{lem:streaming}, it multiplies $\mathfrak m_i$ by a kernel of total
mass at most $(k_0+i)(4K^2+2K+1)$; here $k_0+i$ bounds the
number of factors to which the product rule applies. The total mass and
the marginals of the old marks thus grow by at most the factor
$(k_0+i)(4K^2+2K+1)$. A new
point $v$ or $w$ chosen according to $\omega$ has a level whose
conditional law is at most $K[\nu+\sum_{r'}\delta_{r'}]$, where $r'$ runs
over the levels of the old marks. A new level $r$ chosen according to
$\nu(\dd r)$ on $(q_v,1]$, or according to $s\,\nu(\dd r)$, has
conditional law at most $\nu$. Factors $a(q_v)\le1$ and the restriction
$r>q_v$ only decrease these measures. For every Borel set $B$, the
marginal of the level $r_{\rm new}$ of a new mark therefore satisfies
\[
 \mathfrak m_{i+1}(r_{\rm new}\in B)
 \le C\Bigl[\mathfrak m_i(1)\,\nu(B)+\sum_{r'}\mathfrak m_i(r'\in B)\Bigr]
 \le C_{i+1}\,\nu(B),
\]
where $C$ depends only on $k_0$ and $i$. This completes the induction.

\paragraph{Step 2: levels of pairs.}
In a tree of such a term, every vertex with at least two outgoing edges or
leaves lies at the level of a mark. This level is the original split
level $p$, a level $q_v$ at which leaves were attached, a level $r$ at
which a new branch split, or the split level of an independent tree. For
two leaves $\mathsf a,\mathsf b$ of one tree, $q_{\mathsf{ab}}$ is the
level of such a vertex. In particular,
$\chi(q_{\mathsf{ab}})\le\sum_{r'}\chi(r')$ for every $\chi\ge0$, where
$r'$ runs over the levels of the marks. By Step 1, this implies that
\[
 \int\chi(q_{\mathsf{ab}})\,\dd\mathfrak m\le C\int\chi\,\dd\nu .
\]

\paragraph{Step 3: Hölder's inequality.}
As in \cite[Remark~2.4]{OpenAIFluctuationScale}, we apply Hölder's
inequality unconditionally under the positive tree laws, and we bound each
factor after deleting all leaves except one pair. On the product of the
probability space of the term and the mark space with $\mathfrak m$, we
have
\[
 \int\E|\Delta_1\Delta_2\Delta_3|\,\dd\mathfrak m
 \le\prod_{i=1}^3\Bigl(\int\E|\Delta_i|^3\,\dd\mathfrak m\Bigr)^{1/3}.
\]
If we delete all leaves except the two leaves of $\Delta_i$, two
descendants remain, splitting at their common level $q_i$. By deletion, we
have $\E|\Delta_i|^3=\E_{2,q_i,u}|R^-_{12}-q_i|^3$; for a product over
independent trees, the other trees integrate to one. We use Step 2 with
$\chi(q)=\E_{2,q,u}|R^-_{12}-q|^3$, and we find that each factor is at
most $(C\mathcal M_3(u))^{1/3}$. This proves \eqref{eq:threeproducts}.
\end{proof}

The same proof applies to products of finitely many independent initial
tree moments whose leaves split at $p$, integrated against $\nu(\dd p)$.
The constant then depends on the total number of their initial leaves and on the
number of differentiations. Indeed, the initial mark measure is still
$\nu(\dd p)$, the number of factors to which the product rule applies
is at most the total number of initial leaves plus the number of differentiations,
and Steps~2 and~3 are unchanged.

\paragraph{The cavity covariance.}
In particular, \cref{lem:cav-threeproducts} applies to the covariance
between the squared conditional magnetization of the last spin and the
rest self-overlap, and to its Taylor remainder,
\[
 \mathcal C_{pp}(u)=\Cov_u\bigl(m_N(p)^2,A_p^-\bigr),\qquad
 \mathcal R_{\rm cav}=\int\nu(\dd p)\int_0^1(1-u)
                              \mathcal C_{pp}''(u)\dd u .
\]
Representing conditional expectations by leaves, we can write
$\mathcal C_{pp}=\Psi_p-\Psi^\epsilon_p\Psi^\tau_p$, where
\[
 \Psi_p(u)=\E_{4,p,u}[\epsilon_1\epsilon_2(R^-_{34}-p)],\qquad
 \Psi^\epsilon_p(u)=\E_{2,p,u}\epsilon_1\epsilon_2,\qquad
 \Psi^\tau_p(u)=\E_{2,p,u}(R^-_{12}-p).
\]
Indeed, by conditional independence given $\calG_p$, we have
$\Psi_p=\E_u[m_N(p)^2(A_p^--p)]$, $\Psi^\epsilon_p=\E_um_N(p)^2$, and
$\Psi^\tau_p=\E_u(A_p^--p)$. The constant $p$ does not change the
covariance. Note that the observables in these tree moments, including the center
$p$, do not depend on $u$. They are bounded by $2$, and the
trees have at most four leaves. By \cref{lem:streaming}, each of the
three moments belongs to $C^2([0,1])$, with first and second
derivatives bounded by $C\alpha$ and $C\alpha^2$. Then
$\mathcal C_{pp}\in C^2([0,1])$ for every deterministic profile, and
$|\mathcal C_{pp}''|$ is bounded by a constant depending only on
$\beta$, uniformly in $p$. In particular, $\mathcal R_{\rm cav}$ is well
defined, and by Taylor's formula with integral remainder, we have
\[
 \mathcal R_{\rm cav}=\int\bigl[\mathcal C_{pp}(1)-\mathcal C_{pp}(0)
   -\mathcal C_{pp}'(0)\bigr]\,\nu(\dd p).
\]

\begin{lemma}[Taylor remainder]\label{lem:remainder}
For every deterministic profile and every $u\in[0,1]$,
\begin{equation}\label{eq:C2derivative}
 \int|\mathcal C_{pp}''(u)|\,\nu(\dd p)\le C_\beta\mathcal M_3(u),
 \qquad
 |\mathcal R_{\rm cav}|\le C_\beta\sup_{0\le w\le1}\mathcal M_3(w).
\end{equation}
\end{lemma}

\begin{proof}
We have
\[
 \mathcal C_{pp}''=\Psi_p''-(\Psi^\epsilon_p)''\Psi^\tau_p
 -2(\Psi^\epsilon_p)'(\Psi^\tau_p)'-\Psi^\epsilon_p(\Psi^\tau_p)''.
\]
Each application of \eqref{eq:streaming} adds a factor
$\alpha/2\le\beta^2/2$ and one overlap defect, either inside a factor $D$
or in the independent tree of its last term. The spin factors have
absolute value one. The moments $\Psi_p$ and $\Psi^\tau_p$ contain one
defect, and $\Psi^\epsilon_p$ contains none. We count the defects before
differentiation and those added by the two differentiations.
\begin{center}
\small
\begin{tabular}{@{}lccc@{}}
\toprule
term & defects before & added & total\\
\midrule
$\Psi_p''$ & $1$ & $2$ & $3$\\
$(\Psi^\epsilon_p)''\Psi^\tau_p$ & $0+1$ & $2+0$ & $3$\\
$(\Psi^\epsilon_p)'(\Psi^\tau_p)'$ & $0+1$ & $1+1$ & $3$\\
$\Psi^\epsilon_p(\Psi^\tau_p)''$ & $0+1$ & $0+2$ & $3$\\
\bottomrule
\end{tabular}
\end{center}
Products of moments, and the last term of \eqref{eq:streaming}, are
products over independent trees. Every term of $\mathcal C_{pp}''$ is
therefore $(\alpha/2)^2$ times a term of the kind treated in
\cref{lem:cav-threeproducts}, or in its extension to products of
independent initial moments stated after its proof, with at most four
initial leaves and $j=2$. It is multiplied by a bounded multiplicity and
by factors of absolute value at most one. Integrating $p$ against $\nu$
and using \cref{lem:cav-threeproducts} and this extension, we obtain the first bound in
\eqref{eq:C2derivative} with $C_\beta=C\beta^4$. We integrate it against
$(1-u)\dd u$ and use Tonelli's theorem. This gives
$|\mathcal R_{\rm cav}|\le C\beta^4\int_0^1(1-u)\mathcal M_3(u)\dd u
\le C\beta^4\sup_w\mathcal M_3(w)$, which is the second bound.
\end{proof}

\chapter{The derivative at the independent endpoint}\label{sec:c7-endpoint}

In \cref{sec:c7}, we constructed the cavity interpolation $\Pp_u$,
$0\le u\le1$. At $u=1$, it is the tilted law of the full system with
$N$ spins. At $u=0$, it is a product law: the last spin $\epsilon$ is
independent of the remaining $n=N-1$ spins. We also introduced the
covariance
\[
 \mathcal C_{pp}(u)=\Cov_u\bigl(m_N(p)^2,A_p^-\bigr)
\]
between the squared conditional magnetization of the last spin and the
rest self-overlap. By \cref{lem:streaming}, it is twice continuously
differentiable in $u$ (see \cref{subsec:c7-three}), and
\cref{lem:remainder} bounds its Taylor remainder. The covariance
vanishes at $u=0$. In this section, we compute its derivative there.

The main result is \cref{lem:linear}. It states that after
integration against $\nu$, the derivative $\mathcal C_{pp}'(0)$ equals
$\alpha/d$ times the trace of the scalar response operator
$J^{\scal}_{a,d}$ of \cref{sec:scalar-stability} against the
covariance operator $\Sigma^-$ of the rest self-overlap, minus a
nonnegative term $\alpha\mathcal P/2$, plus explicit error terms. The
identity before integration is \cref{lem:c7-endpoint-pointwise}. In
\cref{lem:cl-response}, we bound the integrated error terms by the
trace term $A$ and the mean--covariance term $J$ of the full system,
up to an additive $N^{-2}$. The three lemmas hold for every
deterministic profile with $\nu(\{1\})=0$; no minimizing property is
used. In
\cref{sec:c7-closure}, we insert \cref{lem:linear} into the proof of
the trace identity of \cref{lem:cav-trace} and use
\cref{lem:cl-response} to bound its error term in the proof of
\cref{prop:c7-closure}.

The computation uses the exact density \eqref{eq:treedensity} of the
cavity law and the parameter responses \eqref{eq:nodeTheta} from
\cref{sec:c7}, the Gaussian heat identity of \cref{lem:fnd-path-heat},
and the explicit form \eqref{eq:scalarJexplicit} of the scalar
response operator. The error bounds also use the kernel bounds of
\cref{lem:cl-scalar}, the posterior variance bound \eqref{eq:cl-BL},
and the comparison inequality of \cref{cor:comparison}.

\paragraph{Standing assumptions.}
We keep the assumptions of \cref{sec:c7}: $N\ge2$, $n=N-1$,
$0<\rho\le1$, $t=1-\rho$, $c=\beta^2\rho$, and $0<s=\lambda/\beta\le1/2$.
The parameters $\alpha$, $d$, and $b(u)$ are as in \eqref{eq:cavityparameters}. In
particular, $0\le\alpha\le\beta^2$ and $d\in[\beta^2/2,\beta^2]$. The
profile $\mu=s\delta_0+\nu$ is deterministic and otherwise arbitrary,
with cumulative function $a(q)=\mu([0,q])$. From now on, we also
assume that $\nu(\{1\})=0$. This holds at every minimizing profile by
\eqref{eq:in-contact}; \cref{subsec:c7-tree,subsec:c7-three} did not
use it. Under this assumption, we show below that the covariance
operator of the rest self-overlap has its range in the scalar
transport space $\mathcal H_\nu$ of \cref{sec:scalar-stability}, on
which $J^{\scal}_{a,d}$ acts.

\section{Linear response at the independent endpoint}
\label{subsec:c7-linear}

\subsection{The product law at the independent endpoint}

In this section, we take all expectations at $u=0$ unless a
subscript says otherwise. Recall from \cref{subsec:c7-row} that
$\Pp_0$ is a product law. Its first factor is the tilted law of the $n$-spin
system with Hamiltonian $\mathcal E^-$, field coefficient $c$, and profile
$\mu$, and the second is the scalar law of \cref{sec:scalar-stability}
with coefficient $d$. We
write $\E^-$, $\Var^-$, and $\Cov^-$ for the first factor (the rest
system) and $\E^0$, $\Var^0$, and $\Cov^0$ for the second;
$\calG^-_q$ and $\calG^0_q$ denote their filtrations. At $u=0$, the
solution of \eqref{eq:cavityPDE} splits as
$U(q,x,y)=U^-(q,x)+U^0(q,y)$. We define
\[
 \begin{gathered}
 M(r,x)=\nabla_xU^-(r,x),\qquad M_i(r,x)=\partial_{x_i}U^-(r,x),\\
 A^-(r,x)=\frac{|M(r,x)|^2}n,\qquad m(r,y)=\partial_yU^0(r,y),
 \end{gathered}
\]
so that $M_r=M(r,X_r)$, $A_r^-=A^-(r,X_r)$,
$H_r^-=\nabla_x^2U^-(r,X_r)$, $m_r=m(r,Y_r)$, and
$h_r=\partial_ym(r,Y_r)$. We use the scalar quantities $f_p=m_p^2$,
$\bar f_p=f_p-\E^0f_p$, $g_r^\ell(y)=P^0_{r,\ell}(m(\ell,\cdot)^2)(y)$,
and $\Gamma=\Gamma_{a,d}$, with $\Gamma(p)=\E^0m_p^2$, of
\cref{sec:scalar-stability}. Recall that
$g_r^\ell(Y_r)=\E^0[m_\ell^2\mid\calG^0_r]$ and
$m_r=\E^0[\epsilon\mid\calG^0_r]$. For $r\le\ell$, let
$P^-_{r,\ell}$ be the transition operator of the rest diffusion. Its
state includes the rest couplings, so that
$(P^-_{r,\ell}A^-_\ell)(X_r)=\E^-[A^-_\ell\mid\calG^-_r]$. For
notational convenience, we abbreviate $a_r^-=a_r^-(0)$ and
$j_r^-=j_r^-(0)$ for the quantities of \eqref{eq:cav-rest-errors} at
$u=0$.

Note that the rest Hamiltonian $\mathcal E^-$ is invariant under the
global spin reversal $\tau\mapsto-\tau$. Then $U^-(0,\cdot)$ is even,
which implies that $M_0=0$ and
\begin{equation}\label{eq:c7-root-overlap}
 A_0^-=0.
\end{equation}
Since the spatial derivatives of $U^-$ are bounded and jointly
continuous, $r\mapsto A_r^-$ is continuous with values in $[0,1]$.

\paragraph{The rest self-overlap.}
We define
\[
 \begin{gathered}
 \Sigma^-(p,q)=\Cov^-(A_p^-,A_q^-),\qquad
 \bar A_p^-=A_p^--\E^-A_p^-,\\
 \Theta(r;\ell,p)=\Cov^-\bigl(P^-_{r,\ell}A_\ell^-,A_p^-\bigr)
 \quad(r\le\ell),
 \end{gathered}
\]
where we recall that $P^-_{r,\ell}A_\ell^-$ stands for
$\E^-[A_\ell^-\mid\calG^-_r]$. Note that $\Sigma^-$ is continuous and
bounded by one, and that $\Sigma^-(0,p)=0$ by
\eqref{eq:c7-root-overlap}. If $r\ge p$, then $\bar A_p^-$ is
$\calG^-_r$-measurable, and by the tower property, we have
\begin{equation}\label{eq:c7-theta-past}
 \Theta(r;\ell,p)=\E^-\bigl[\E^-[A_\ell^-\mid\calG^-_r]\,\bar A_p^-\bigr]
 =\E^-\bigl[A_\ell^-\bar A_p^-\bigr]=\Sigma^-(\ell,p)
 \qquad(r\ge p).
\end{equation}

\paragraph{Fluctuations along the path.}
Let $\tau$ be the terminal rest spin. For levels $r,\ell\in[0,1]$, we
define
\[
 x_r=\frac{\tau\cdot M_r}n-A_r^-,\qquad
 y_{r\ell}=\frac{M_r\cdot M_\ell}n-A_{r\wedge\ell}^-,
\]
random variables of the rest system; the subscripts distinguish them
from the field variables $x$ and $y$. The vector
$M_q=\E^-[\tau\mid\calG^-_q]$ is a martingale, and by the conditional
law of total covariance, we have, for $r\le\ell$,
\begin{equation}\label{eq:c7-total-cov}
 C_r^-=\E^-[C_\ell^-\mid\calG^-_r]
 +\E^-\bigl[(M_\ell-M_r)(M_\ell-M_r)^T\bigm|\calG^-_r\bigr].
\end{equation}
In particular, $\E^-[C_\ell^-\mid\calG^-_r]\preceq C_r^-$ and
$\Cov^-(M_\ell\mid\calG^-_r)\preceq C_r^-$. We have
$x_r=M_r\cdot(\tau-M_r)/n$, and $y_{r\ell}=M_r\cdot(M_\ell-M_r)/n$ if
$r\le\ell$. By conditioning on $\calG^-_r$, we deduce from
\eqref{eq:c7-total-cov} that
\begin{equation}\label{eq:cl-xy}
 \E^-[x_r\mid\calG_r^-]=0,\quad \E^-x_r^2=j_r^-,\qquad
 \E^-y_{r\ell}=0,\quad \E^-y_{r\ell}^2\le j_{r\wedge\ell}^-.
\end{equation}
Indeed, $\E^-[x_r^2\mid\calG^-_r]=M_r^TC_r^-M_r/n^2$, and for
$r\le\ell$, the vector $M_\ell-M_r$ has conditional mean zero and
conditional covariance at most $C_r^-$ given $\calG^-_r$, so that
$\E^-[y_{r\ell}^2\mid\calG^-_r]\le M_r^TC_r^-M_r/n^2$. The case
$r>\ell$ follows by symmetry.

\paragraph{The error terms.}
We set $v_p^0=\Var^0(f_p)$ and
\begin{gather*}
 a_p^H=\frac{\E^-\Tr\bigl((H_p^-)^2\bigr)}{n^2},\qquad
 b_p^H=\frac{\E^-[M_p^TH_p^-M_p]}{n^2},\qquad
 b_{r,p}^H=\frac{\E^-[M_r^TH_p^-M_p]}{n^2},\\
 z_{p,\ell}=\frac1n\E^-\bigl[M_p\cdot\nabla(P^-_{p,\ell}A_\ell^-)(X_p)\bigr]
 \quad(\ell\ge p),
\end{gather*}
where $\nabla(P^-_{p,\ell}A_\ell^-)$ is the gradient of
$x\mapsto(P^-_{p,\ell}A_\ell^-)(x)$, with the rest couplings fixed. The
error terms are
\begin{align}
 T_p^{(1)}={}&2\E^0[m_ph_p\epsilon]\,\E^-[x_p\bar A_p^-]
 -2\int\E^0[m_ph_pm_r]\,\E^-[y_{pr}\bar A_p^-]\,\nu(\dd r),
 \label{eq:cav-T1}\\
 T_p^{(2)}={}&-\int\E^0[\bar f_p\epsilon m_r]\,
                         \E^-[x_r\bar A_p^-]\,\nu(\dd r)\nonumber\\
 &+\frac12\iint\E^0[\bar f_pm_rm_\ell]\,
                  \E^-[y_{r\ell}\bar A_p^-]\,\nu(\dd r)\nu(\dd\ell),
 \label{eq:cav-T2}\\
 T_p^{(3)}={}&v_p^0a_p^H+
       \bigl[4\E^0(m_p^2h_p)+2(1-a(p))v_p^0\bigr]b_p^H\nonumber\\
 &-\int_{(p,1]}\Cov^0\bigl(f_p,g_p^\ell(Y_p)\bigr)z_{p,\ell}\,\nu(\dd\ell)
 -2\int\E^0[\bar f_pm_pm_r]\,b_{r,p}^H\,\nu(\dd r),
 \label{eq:cav-T3}
\end{align}
and $T_p=T_p^{(1)}+T_p^{(2)}+T_p^{(3)}$. When we integrate against
$\nu$ without indicating a domain, we integrate over $[0,1]$, including
any atom of $\nu$ at $0$.
The terms $T_p^{(1)}$ and $T_p^{(2)}$ contain the fluctuations
$x_r$ and $y_{r\ell}$, and $T_p^{(3)}$ contains the derivatives of the
rest observable $A_p^-$ in the coupling vector of the incident row.

\paragraph{The conditional future correction.}
For $p,\ell>r$, we define
\[
 \Lambda_r(p,\ell)=\Cov^0\bigl(g_r^p(Y_r),g_r^\ell(Y_r)\bigr),\qquad
 \Sigma_r^-(p,\ell)=\E^-\Cov^-(A_p^-,A_\ell^-\mid\calG_r^-),
\]
and
\begin{equation}\label{eq:Ppositive}
 \mathcal P=\int\nu(\dd r)\iint_{p,\ell>r}
      \Lambda_r(p,\ell)\,\Sigma_r^-(p,\ell)\,\nu(\dd p)\nu(\dd\ell).
\end{equation}
Note that $\Lambda_r$ is a covariance of the scalar law and that
$\Sigma_r^-$ is an averaged conditional covariance of the rest system.
We show in the proof of \cref{lem:linear} that $\mathcal P\ge0$ and
that $-\alpha\mathcal P/2$ is the contribution of the difference
between $\Theta$ and $\Sigma^-$.

\paragraph{The scalar response operator applied to $\Sigma^-$.}
For a bounded Borel function $\varphi$ on $[0,1]$ and $p\in[0,1]$, the
seven terms on the right side of \eqref{eq:scalarJexplicit}, in the
order written there, are
\begin{alignat*}{2}
 &\text{(J1)}\qquad&&\E^0\bigl[h_p^2+2(1-a(p))m_p^2h_p\bigr]\varphi(p),\\
 &\text{(J2)}&&-\int_{(p,1]}\E^0[m_p\partial_y g_p^\ell(Y_p)]\,
                                      \varphi(\ell)\,\nu(\dd\ell),\\
 &\text{(J3)}&&-2\int\E^0[m_ph_pm_r]\,\varphi(p\wedge r)\,\nu(\dd r),\\
 &\text{(J4)}&&-\int\E^0[\bar f_p\epsilon m_r]\,\varphi(r)\,\nu(\dd r),\\
 &\text{(J5)}&&\frac12\int a(r)\E^0[\bar f_pm_r^2]\,\varphi(r)\,\nu(\dd r),\\
 &\text{(J6)}&&\frac12\iint\E^0[\bar f_pm_rm_\ell]\,\varphi(r\wedge\ell)
                                      \,\nu(\dd r)\nu(\dd\ell),\\
 &\text{(J7)}&&\frac12\int\nu(\dd r)\int_{(r,1]}
           \E^0[\bar f_pg_r^\ell(Y_r)]\,\varphi(\ell)\,\nu(\dd\ell),
\end{alignat*}
and their sum is $d^{-1}(J_{a,d}\varphi)(p)$. All kernels in (J1)--(J7)
are bounded by constants depending only on $\beta$
(\cref{sec:scalar-stability}). For $p\in[0,1]$, we denote by
$(J^{\scal}_{a,d}\Sigma^-)(p,p)$ the value at $p$ of the function
$J_{a,d}\Sigma^-(\cdot,p)$, that is, of $J_{a,d}$ applied to the first
variable of $\Sigma^-$. This notation agrees with the compression
$J^{\scal}_{a,d}$. Indeed, let $\varphi$ be bounded with
$\varphi(0)=0$. Each of the terms (J2)--(J7) contains a factor $m_p$
or $\bar f_p$, and $m_0=\bar f_0=0$. At $p=0$, these terms vanish, and
(J1) equals $h_0^2\varphi(0)=0$. As $\nu(\{1\})=0$, the function
$J_{a,d}\varphi$ therefore lies in $\mathcal H_\nu$,
and it coincides with $J^{\scal}_{a,d}\varphi$ when
$\varphi\in\mathcal H_\nu$. By \eqref{eq:c7-root-overlap}, this applies
to $\varphi=\Sigma^-(\cdot,p)$ for every $p$.

\subsection{The endpoint derivative}

\begin{lemma}[Endpoint derivative at a fixed level]
\label{lem:c7-endpoint-pointwise}
Let $\mu=s\delta_0+\nu$ be a deterministic profile with
$\nu(\{1\})=0$, and let $p\in[0,1]$. If $\alpha>0$, then
\begin{align}
 \frac{\mathcal C_{pp}'(0)}\alpha
 ={}&\frac1d(J^{\scal}_{a,d}\Sigma^-)(p,p)+T_p\nonumber\\
 &+\frac12\int\nu(\dd r)\int_{(r,1]}
  \Cov^0\bigl(f_p,g_r^\ell(Y_r)\bigr)
    \bigl[\Theta(r;\ell,p)-\Sigma^-(\ell,p)\bigr]\nu(\dd\ell).
 \label{eq:diagonallinear}
\end{align}
If $\alpha=0$, then $\mathcal C_{pp}'(0)=0$.
\end{lemma}

The first term on the right side of \eqref{eq:diagonallinear} is the
scalar response applied to $\Sigma^-(\cdot,p)$ in place of $\varphi$,
and $T_p$ collects the error terms. The last line has no scalar
counterpart. It comes from the conditional expectation
$\E^-[A_\ell^-\mid\calG^-_r]$ in the second derivative of the density,
which replaces $\Sigma^-(\ell,p)$ by $\Theta(r;\ell,p)$ in (J7). After
integration over $p$, it becomes the nonpositive term
$-\mathcal P/2$.

\begin{lemma}[Integrated linear response]\label{lem:linear}
Let $\mu=s\delta_0+\nu$ be a deterministic profile with
$\nu(\{1\})=0$, and let $V^-=\int\Var^-(A_p^-)\nu(\dd p)$. Then
\begin{equation}\label{eq:lineartrace}
 \int\mathcal C_{pp}'(0)\,\nu(\dd p)
 =\alpha\Bigl[\frac1d\Tr_\nu(J^{\scal}_{a,d}\Sigma^-)
                     -\frac12\mathcal P+\int T_p\,\nu(\dd p)\Bigr],
 \qquad \mathcal P\ge0,
\end{equation}
where $\mathcal P$ is given by \eqref{eq:Ppositive},
$\Tr_\nu(J^{\scal}_{a,d}\Sigma^-)
=\int(J^{\scal}_{a,d}\Sigma^-)(p,p)\nu(\dd p)$, and $T_p$ is given by
\eqref{eq:cav-T1}--\eqref{eq:cav-T3}. Further, $\Sigma^-$ is a positive
trace-class covariance operator on $L^2(\nu)$ with trace $V^-$ and range
in $\mathcal H_\nu$. With $\bar A^-$ the random function
$p\mapsto\bar A_p^-$, the number $\Tr_\nu(J^{\scal}_{a,d}\Sigma^-)$
equals $\E^-\langle\bar A^-,J^{\scal}_{a,d}\bar A^-\rangle_{L^2(\nu)}$,
the trace of the operator $J^{\scal}_{a,d}\Sigma^-$ on $\mathcal H_\nu$.
\end{lemma}

\subsection{Proof of the endpoint derivative at a fixed level}

\begin{proof}[Proof of \cref{lem:c7-endpoint-pointwise}]
If $\alpha=0$, then the row variance $\alpha u/n$ vanishes and
$b(u)=d$ for every $u$, so that $\Pp_u=\Pp_0$ for every $u$ and
$\mathcal C_{pp}$ vanishes identically. From now on, we assume that
$\alpha>0$ and fix $p\in[0,1]$. The proof has six steps. In Step~1,
we show that $\mathcal C_{pp}'(0)$ is a Laplacian, in the coupling
vector of the incident row, of a covariance under laws with
deterministic rows. In Step~2, we write this Laplacian as a sum of four
groups of terms, using the exact density of these laws. In Steps~3--5,
we compute the four groups. In Step~6, we match the result with
(J1)--(J7). At fixed $N$, the spin observables and the required
positive-order derivatives of $U$ are bounded. The reference-density
integrands have the integrable Gaussian envelopes established below, so
we may exchange expectations and finite level integrals by Fubini's
theorem.

\paragraph{Step 1: reduction to fixed couplings.}
For $w\ge0$ and $b\in[c,d]$, let $\Pp_{w,b}$ be the law of
\cref{subsec:c7-row} with an incident row of variance $w$ per
coordinate in place of $\alpha u/n$, and with last-field coefficient
$b$ in place of $b(u)$. In this notation, $\Pp_u=\Pp_{\alpha u/n,b(u)}$,
and $b(u)\in[c,d]$. We denote expectation and covariance under these laws by
$\E_{w,b}$ and $\Cov_{w,b}$, and define
\[
 C(w,b)=\Cov_{w,b}\bigl(m_N(p)^2,A_p^-\bigr),
 \qquad\text{so that}\qquad
 \mathcal C_{pp}(u)=C\bigl(\alpha u/n,b(u)\bigr).
\]
At $w=0$, the row vanishes, and $U$ splits into a function of $(q,x)$
and a function of $(q,y)$, as at $u=0$ but with $d$ replaced by $b$.
Then $\Pp_{0,b}$ is a product law, under which $m_N(p)$ and $A_p^-$
are independent, and therefore
\begin{equation}\label{eq:c7-zero-row}
 C(0,b)=0\qquad\text{for every }b\in[c,d].
\end{equation}
Applying \cref{lem:fnd-path-heat} with the row variance $w$ as its
parameter, together with \cref{lem:tree-regularity,lem:tree-responses},
we see that the function $w\mapsto C(w,b)$ is continuous on $[0,\alpha/n]$ and
differentiable on $(0,\alpha/n]$, with a right derivative at $w=0$.
The derivative $\partial_wC$ is jointly continuous in $(w,b)$. To see
this, we couple the row as $\sqrt w\,g$ and the last field as
$\sqrt b\,B$ and use the Gaussian envelope of the density, as in
the continuity part of the proof of \cref{lem:streaming}. Let
$0<u\le1$. By \eqref{eq:c7-zero-row} and the mean value theorem in
$w$, there exists $w_u\in(0,\alpha u/n)$ such that
\[
 \mathcal C_{pp}(u)=C\bigl(\alpha u/n,b(u)\bigr)-C\bigl(0,b(u)\bigr)
 =\frac{\alpha u}n\,\partial_wC\bigl(w_u,b(u)\bigr).
\]
Dividing by $u$ and letting $u\downarrow0$, we find
\begin{equation}\label{eq:c7-row-only}
 \mathcal C_{pp}'(0)=\frac\alpha n\,\partial_wC(0,d).
\end{equation}
In particular, the variation of the last-field coefficient $b(u)$
does not contribute to $\mathcal C_{pp}'(0)$.

We now compute $\partial_wC(0,d)$. For $\theta\in\R^n$, let
$\Pp^\theta$ be the law of \cref{subsec:c7-row} with last-field
coefficient $d$ and with the incident row replaced by a deterministic
coupling: the term $\sqrt{\alpha u/n}\,\epsilon\sum_ig_{iN}\tau_i$ in
$\mathcal E(\sigma;x,y)$ becomes $\epsilon\,\theta\cdot\tau$. We write
$U(q,x,y;\theta)$ for the corresponding solution of
\eqref{eq:cavityPDE}, $\E^\theta$ and $\Cov^\theta$ for expectation and
covariance under $\Pp^\theta$, and
\[
 Z(\theta)=\E_{\rm ref}\,e^{sU(0,0,0;\theta)}
\]
for its outer normalizer, which is an expectation over the couplings
$g_{ij}$, $1\le i<j\le n$, only. Under $\Pp^\theta$, observables such
as $m_N(p)=\partial_yU(p,X_p,Y_p;\theta)$ and
$M_p=\nabla_xU(p,X_p,Y_p;\theta)$ are computed from $U(\cdot;\theta)$.
Under $\Pp_{w,d}$, we may integrate first over the couplings $g_{ij}$
and then over $\theta=\sqrt w\,(g_{1N},\dots,g_{nN})$. Given $\theta$,
the conditional law is $\Pp^\theta$. Because of the outer tilt
$e^{sU(0,0,0)}$, the vector $\theta$ has the density
$Z(\theta)/\E Z(\theta)$ with respect to $N(0,wI_n)$. We obtain, for every bounded observable
$G$,
\begin{equation}\label{eq:c7-mixture}
 \E_{w,d}G=\frac{\E[Z(\theta)\,\E^\theta G]}{\E Z(\theta)},
 \qquad\theta\sim N(0,wI_n).
\end{equation}
This mixture is not the Gaussian average of the laws $\Pp^\theta$,
because $\theta$ is weighted by $Z(\theta)$. We now check that this
weight does not change the derivative at $w=0$.

Let $f=m_N(p)^2$. For an observable $G$, we write
$\Phi_G(\theta)=\E^\theta G$. Replacing the summation variable
$\epsilon$ by $-\epsilon$ in the terminal value of \eqref{eq:cavityPDE},
we find $U(1,x,y;\theta)=U(1,x,-y;-\theta)$. The invariance of
\eqref{eq:cavityPDE} under $y\mapsto-y$ then implies that
$U(q,x,y;\theta)=U(q,x,-y;-\theta)$ for every $q$. Hence the
reflection $(\epsilon,Y)\mapsto(-\epsilon,-Y)$ of the last spin and its
field maps $\Pp^\theta$ to $\Pp^{-\theta}$, preserves
$U(0,0,0;\theta)$ and $M_p$, and changes the sign of $m_N(p)$. It
follows that the functions
\[
 Z,\qquad \Phi_f,\qquad \Phi_{A_p^-},\qquad \Phi_{fA_p^-}
\]
are even in $\theta$, so that their gradients vanish at $\theta=0$. At
fixed $N$, they are twice continuously differentiable, with second
derivatives bounded by $Ce^{C|\theta|}$ for a constant $C$ that may
depend on $N$. Indeed, by \eqref{eq:nodeTheta}, we have
$|\partial_{\theta_i}U|\le1$ and $0\le\partial_{\theta_i}^2U\le1$. By
\cref{lem:tree-regularity}, all derivatives of $U$ of positive order in
$(\theta,x,y)$ are bounded at fixed $N$. Further, the density of $\Pp^\theta$ and its
first two derivatives in $\theta$ are bounded by the integrable
envelope
\[
 C\exp\Bigl\{C\Bigl(1+\sum_{1\le i<j\le n}|g_{ij}|+|\theta|
 +\sup_{q\le1}|X_q|+\sup_{q\le1}|Y_q|\Bigr)\Bigr\}
\]
under the Gaussian reference law of the couplings and the paths.
Constants that we use only to justify differentiation may depend on
$N$; they do not enter the identity we prove.

Let $F\in C^2(\R^n)$ satisfy $|\nabla^2F(\theta)|\le Ce^{C|\theta|}$, and
let $g$ be a standard Gaussian vector in $\R^n$. For $w>0$, Gaussian
integration by parts shows that the derivative of
$w\mapsto\E F(\sqrt w\,g)$ is $\E\,\Delta F(\sqrt w\,g)/2$, which
tends to $\Delta F(0)/2$ as $w\downarrow0$ by dominated convergence.
By the mean value theorem, the right derivative at $w=0$ is
$\Delta F(0)/2$. This is the heat equation for the Gaussian
semigroup. We use this fact for $F=Z\Phi_G$ and for $F=Z$ in
\eqref{eq:c7-mixture}. By the quotient rule, we have
\[
 \begin{aligned}
 \partial_w\E_{w,d}G\big|_{w=0}
 &=\frac{\Delta(Z\Phi_G)(0)-\Phi_G(0)\Delta Z(0)}{2Z(0)}\\
 &=\frac12\Delta\Phi_G(0)+\frac{\nabla Z(0)\cdot\nabla\Phi_G(0)}{Z(0)}
 =\frac12\Delta\Phi_G(0),
 \end{aligned}
\]
where in the second equality we used
$\Delta(Z\Phi_G)=Z\Delta\Phi_G+2\nabla Z\cdot\nabla\Phi_G+\Phi_G\Delta Z$,
and in the third that $\nabla Z(0)=0$. We use this for
$G\in\{f,A_p^-,fA_p^-\}$, together with the product rule. The cross term
$\nabla\Phi_f\cdot\nabla\Phi_{A_p^-}$ in the expansion of
$\Delta(\Phi_f\Phi_{A_p^-})/2$ vanishes at $\theta=0$, because
$\nabla\Phi_f(0)=0$. We conclude that
\[
 \partial_wC(0,d)=\frac12\Delta\Phi_{fA_p^-}(0)
 -\frac12\Delta\Phi_f(0)\,\Phi_{A_p^-}(0)
 -\frac12\Phi_f(0)\,\Delta\Phi_{A_p^-}(0)
 =\frac12\Delta_\theta\Cov^\theta(f,A_p^-)\Big|_{\theta=0}.
\]
Inserting this into \eqref{eq:c7-row-only}, we obtain
\begin{equation}\label{eq:c7-fixedcoupling}
 \mathcal C_{pp}'(0)=\frac{\alpha}{2n}\sum_{i=1}^n
   \partial_{\theta_i}^2\Cov^{\theta}(f,A_p^-)\Big|_{\theta=0}.
\end{equation}

\paragraph{Step 2: derivatives in one coordinate of the row.}
For $\theta=0$, the law $\Pp^\theta$ is the product law $\Pp_0$, and
$U(q,x,y;0)=U^-(q,x)+U^0(q,y)$. In this proof, $\E$ and $\Cov$ without
superscript denote expectation and covariance under $\Pp_0$. Let
\[
 \dot U_i(r,x,y)=\partial_{\theta_i}U(r,x,y;\theta)\big|_{\theta=0},
 \qquad
 \ddot U_i(r,x,y)=\partial_{\theta_i}^2U(r,x,y;\theta)\big|_{\theta=0}.
\]
Under $\Pp_0$, the conditional law of $(\tau,\epsilon)$ given $\calG_r$
is the product of the conditional laws of $\tau$ given $\calG^-_r$ and
of $\epsilon$ given $\calG^0_r$. The conditional feature
$\pi_i(r)=\E[\tau_i\epsilon\mid\calG_r]$ of \eqref{eq:nodeTheta}
therefore equals $M_i(r,X_r)m_r$, and for $\ell\ge r$, we have
$\E[\pi_i(\ell)^2\mid\calG_r]=(P^-_{r,\ell}M_i(\ell,\cdot)^2)(X_r)\,
g_r^\ell(Y_r)$. Inserting these identities into \eqref{eq:nodeTheta}, we obtain
\begin{align}
 \dot U_i(r,x,y)&=M_i(r,x)\,m(r,y),\label{eq:Vi}\\
 \ddot U_i(r,x,y)&=1-a(r)M_i(r,x)^2m(r,y)^2
 -\int_{(r,1]}\bigl(P^-_{r,\ell}M_i(\ell,\cdot)^2\bigr)(x)\,
 g_r^\ell(y)\,\nu(\dd\ell).
 \label{eq:Wi}
\end{align}
Since $\sum_iM_i(\ell,\cdot)^2=nA^-(\ell,\cdot)$, the average of
\eqref{eq:Wi} over $i$ is
\begin{equation}\label{eq:c7-average}
 \ddot U_{\rm av}(r,x,y)=\frac1n\sum_{i=1}^n\ddot U_i(r,x,y)
 =1-a(r)A^-(r,x)\,m(r,y)^2
 -\int_{(r,1]}(P^-_{r,\ell}A^-_\ell)(x)\,g_r^\ell(y)\,\nu(\dd\ell).
\end{equation}

By \cref{lem:tree-density} with one leaf, applied with the row fixed
at $\theta$, the law $\Pp^\theta$ has the density
\[
 \varpi_\theta=\frac1{Z(\theta)}\exp\Bigl\{\mathcal E^-(\tau)
 +\epsilon\,\theta\cdot\tau+X_1\cdot\tau+Y_1\epsilon
 -\int_{[0,1]}U(q,X_q,Y_q;\theta)\,\nu(\dd q)\Bigr\}
\]
with respect to the Gaussian reference law of the couplings $g_{ij}$,
the Gaussian reference paths $X$ and $Y$ with covariances
$c(q\wedge q')I_n$ and $d(q\wedge q')$, and counting measure on
$\sigma=(\tau,\epsilon)$. Recall from Step~1 that $\nabla Z(0)=0$.
Then the $\theta_i$-score of $\varpi_\theta$ at $\theta=0$ is
\[
 \mathfrak s_i=\partial_{\theta_i}\log\varpi_\theta\big|_{\theta=0}
 =\tau_i\epsilon-\int M_i(r)\,m_r\,\nu(\dd r),
 \qquad M_i(r)=M_i(r,X_r),
\]
where we used \eqref{eq:Vi}, and the second derivative satisfies
\begin{equation}\label{eq:c7-density-hessian}
 \frac{\partial_{\theta_i}^2\varpi_\theta}{\varpi_\theta}\Big|_{\theta=0}
 =\mathfrak s_i^2+\partial_{\theta_i}^2\log\varpi_\theta\big|_{\theta=0}
 =\mathfrak s_i^2-\int\ddot U_i(r,X_r,Y_r)\,\nu(\dd r)
 -\partial_{\theta_i}^2\log Z(0).
\end{equation}
This quantity has mean zero under $\Pp_0$, because $\varpi_\theta$
integrates to one for every $\theta$.

An observable $G_\theta$ under $\Pp^\theta$ depends on $\theta$
explicitly, through $U(\cdot;\theta)$, and through the law. We write
$G_i$ and $G_{ii}$ for its first and second explicit
$\theta_i$-derivatives at $\theta=0$, with the couplings, the paths, and
the spins fixed. Differentiating under the reference integral, which
the envelope of Step~1 permits, we have
\[
 \partial_{\theta_i}^2\E^\theta G_\theta\big|_{\theta=0}
 =\E\Bigl[G_{ii}+2G_i\mathfrak s_i
 +G\,\frac{\partial_{\theta_i}^2\varpi_\theta}{\varpi_\theta}
 \Big|_{\theta=0}\Bigr].
\]
This formula holds for $G=fA_p^-$, $f$, and $A_p^-$. At $\theta=0$, we
have $f=f_p$. We denote the explicit derivatives of $f$ by $f_{p,i}$ and
$f_{p,ii}$, and those of $A_p^-$ by $A_{p,i}^-$ and $A_{p,ii}^-$.
The first derivatives of $\E^\theta f$ and $\E^\theta A_p^-$ vanish at
$\theta=0$ by Step~1. Expanding
$\Cov^\theta(f,A_p^-)=\E^\theta[fA_p^-]-\E^\theta f\,\E^\theta A_p^-$
with the product rule, and using that \eqref{eq:c7-density-hessian}
has mean zero, we find that $\partial_{\theta_i}^2\Cov^\theta(f,A_p^-)$
at $\theta=0$ equals
\[
 \E\bigl[f_{p,ii}\bar A_p^-+\bar f_pA_{p,ii}^-+2f_{p,i}A_{p,i}^-\bigr]
 +2\E\bigl[\mathfrak s_i(f_{p,i}\bar A_p^-+\bar f_pA_{p,i}^-)\bigr]
 +\E\Bigl[\bar f_p\bar A_p^-\frac{\partial_{\theta_i}^2\varpi_\theta}
 {\varpi_\theta}\Big|_{\theta=0}\Bigr].
\]
Since $\Pp_0$ is a product law, we have $\E[\bar f_p\bar A_p^-]=0$,
and the constant $\partial_{\theta_i}^2\log Z(0)$ in
\eqref{eq:c7-density-hessian} drops out. Summing over $i$, we see from
\eqref{eq:c7-fixedcoupling} that
\begin{align}
 \frac{\mathcal C_{pp}'(0)}\alpha=\frac1{2n}\sum_i\Bigl\{&
 \E\bigl[f_{p,ii}\bar A_p^-+\bar f_pA_{p,ii}^-+2f_{p,i}A_{p,i}^-\bigr]
 +2\E\bigl[\mathfrak s_i(f_{p,i}\bar A_p^-+\bar f_pA_{p,i}^-)\bigr]
 \nonumber\\
 &+\E\Bigl[\bar f_p\bar A_p^-
       \Bigl(\mathfrak s_i^2-\int\ddot U_i\dd\nu\Bigr)\Bigr]\Bigr\}.
 \label{eq:productcovderivative}
\end{align}
We sort the right side into four groups, each including the factor
$1/(2n)$ and the sum over $i$:
\begin{enumerate}[label=(\alph*)]
\item the terms $f_{p,ii}\bar A_p^-$;
\item the terms $2\mathfrak s_if_{p,i}\bar A_p^-$;
\item the terms $\bar f_p\bar A_p^-(\mathfrak s_i^2-\int\ddot U_i\dd\nu)$;
\item the three terms containing a derivative of $A_p^-$, namely
$\bar f_pA_{p,ii}^-$, $2f_{p,i}A_{p,i}^-$, and
$2\mathfrak s_i\bar f_pA_{p,i}^-$.
\end{enumerate}

We now compute the explicit derivatives. We have
$f=(\partial_yU(p,X_p,Y_p;\theta))^2$ and
$A_p^-=|\nabla_xU(p,X_p,Y_p;\theta)|^2/n$, and by \eqref{eq:Vi},
$\partial_y\dot U_i(p)=M_i(p)h_p$ and $\nabla_x\dot U_i(p)=m_pH_p^-e_i$.
Differentiating once and twice in $\theta_i$, we obtain
\begin{align}
 f_{p,i}&=2m_ph_pM_i(p),&
 f_{p,ii}&=2M_i(p)^2h_p^2+2m_p\partial_y\ddot U_i(p),\nonumber\\
 A_{p,i}^-&=\frac{2m_p}{n}(H_p^-M_p)_i,&
 A_{p,ii}^-&=\frac2n\bigl\{m_p^2|H_p^-e_i|^2+
                         M_p\cdot\nabla_x\ddot U_i(p)\bigr\},
 \label{eq:cav-explicit-derivatives}
\end{align}
where $\ddot U_i(p)$ and its derivatives are evaluated at
$(p,X_p,Y_p)$ and $e_i$ is the $i$th coordinate vector of $\R^n$.

\paragraph{Step 3: the groups (a) and (b).}
The identities \eqref{eq:cav-explicit-derivatives} and
\eqref{eq:c7-average}, together with $\partial_y(m^2)=2mh$, yield
\[
 \begin{aligned}
 \frac1{2n}\sum_if_{p,ii}&=A_p^-h_p^2+m_p\partial_y\ddot U_{\rm av}(p)\\
 &=\bigl(h_p^2-2a(p)m_p^2h_p\bigr)A_p^-
 -\int_{(p,1]}m_p\,\partial_yg_p^\ell(Y_p)\,
 (P^-_{p,\ell}A_\ell^-)(X_p)\,\nu(\dd\ell).
 \end{aligned}
\]
We multiply by $\bar A_p^-$ and take the expectation. The scalar and
rest factors are independent, $\E^-[A_p^-\bar A_p^-]=\Sigma^-(p,p)$,
and $\E^-[(P^-_{p,\ell}A_\ell^-)\bar A_p^-]=\Theta(p;\ell,p)
=\Sigma^-(\ell,p)$ by \eqref{eq:c7-theta-past}. We conclude that group
(a) equals
\begin{equation}\label{eq:c7-group-a}
 \bigl(\E^0h_p^2-2a(p)\E^0[m_p^2h_p]\bigr)\Sigma^-(p,p)
 -\int_{(p,1]}\E^0[m_p\,\partial_yg_p^\ell(Y_p)]\,
                                  \Sigma^-(\ell,p)\,\nu(\dd\ell).
\end{equation}

For group (b), we insert the definitions of $\mathfrak s_i$, $x_p$, and
$y_{pr}=y_{rp}$ and find
\[
 \begin{aligned}
 \frac1n\sum_i\mathfrak s_if_{p,i}
 &=2m_ph_p\Bigl[\epsilon\,\frac{\tau\cdot M_p}n
 -\int m_r\frac{M_r\cdot M_p}n\,\nu(\dd r)\Bigr]\\
 &=2m_ph_p\Bigl[\epsilon(A_p^-+x_p)
 -\int m_r(A_{p\wedge r}^-+y_{pr})\,\nu(\dd r)\Bigr].
 \end{aligned}
\]
We multiply by $\bar A_p^-$ and take the expectation. Since
$\E^0[\epsilon\mid\calG^0_p]=m_p$, we have
$\E^0[m_ph_p\epsilon]=\E^0[m_p^2h_p]$. The terms containing $A^-$
produce covariances of the rest self-overlap, and the terms containing
$x_p$ and $y_{pr}$ make up $T_p^{(1)}$ of \eqref{eq:cav-T1}. Therefore,
group (b) equals
\begin{equation}\label{eq:c7-group-b}
 2\E^0[m_p^2h_p]\Sigma^-(p,p)
 -2\int\E^0[m_ph_pm_r]\,\Sigma^-(p\wedge r,p)\,\nu(\dd r)
 +T_p^{(1)}.
\end{equation}
The coefficients of $\Sigma^-(p,p)$ in \eqref{eq:c7-group-a} and
\eqref{eq:c7-group-b} add up to
$\E^0[h_p^2+2(1-a(p))m_p^2h_p]$, the coefficient in (J1).

\paragraph{Step 4: the group (c).}
We average the last factor of \eqref{eq:productcovderivative} over $i$.
Using $\tau_i^2=\epsilon^2=1$, the definitions of $x_r$ and
$y_{r\ell}$, the identity \eqref{eq:c7-average}, and
$\nu([0,1])=1-s$, we have
\begin{align}
 \frac1n\sum_i\Bigl(\mathfrak s_i^2-\int\ddot U_i\dd\nu\Bigr)
 ={}&s-2\int\epsilon m_r(A_r^-+x_r)\,\nu(\dd r)
 +\iint m_rm_\ell(A_{r\wedge\ell}^-+y_{r\ell})\,\nu(\dd r)\nu(\dd\ell)
 \nonumber\\
 &+\int a(r)m_r^2A_r^-\,\nu(\dd r)
 +\int\nu(\dd r)\int_{(r,1]}g_r^\ell(Y_r)\,
 (P^-_{r,\ell}A_\ell^-)(X_r)\,\nu(\dd\ell).
 \label{eq:c7-density-average}
\end{align}
Here the constant $s$ is $1-\nu([0,1])$: the squares $\tau_i^2\epsilon^2$
in $\mathfrak s_i^2$ contribute $1$, and the constant in
\eqref{eq:c7-average} integrates to $\nu([0,1])$. We multiply
\eqref{eq:c7-density-average} by $\bar f_p\bar A_p^-/2$ and take the
expectation. The constant contributes $s\E[\bar f_p\bar A_p^-]/2=0$.
Each term containing $A^-$ factorizes into a scalar expectation times
$\E^-[A^-_q\bar A_p^-]=\Sigma^-(q,p)$, with $q=r$ or $q=r\wedge\ell$,
or, in the last integral, times
$\E^-[(P^-_{r,\ell}A_\ell^-)\bar A_p^-]=\Theta(r;\ell,p)$. The terms containing $x_r$ and $y_{r\ell}$ make
up $T_p^{(2)}$ of \eqref{eq:cav-T2}. Collecting these terms, we find
that group (c) equals
\begin{equation}\label{eq:c7-group-c}
 \begin{aligned}
 &-\int\E^0[\bar f_p\epsilon m_r]\,\Sigma^-(r,p)\,\nu(\dd r)
 +\frac12\iint\E^0[\bar f_pm_rm_\ell]\,
      \Sigma^-(r\wedge\ell,p)\,\nu(\dd r)\nu(\dd\ell)\\
 &+\frac12\int a(r)\E^0[\bar f_pm_r^2]\,\Sigma^-(r,p)\,\nu(\dd r)
 +\frac12\int\nu(\dd r)\int_{(r,1]}
       \E^0[\bar f_pg_r^\ell(Y_r)]\,\Theta(r;\ell,p)\,\nu(\dd\ell)
 +T_p^{(2)}.
 \end{aligned}
\end{equation}

\paragraph{Step 5: the group (d).}
We compute the three terms in turn. First, by
\eqref{eq:cav-explicit-derivatives},
$\sum_i|H_p^-e_i|^2=\Tr((H_p^-)^2)$, and \eqref{eq:c7-average}, we have
\[
 \frac1{2n}\sum_iA_{p,ii}^-
 =\frac{m_p^2}{n^2}\Tr\bigl((H_p^-)^2\bigr)
 +\frac1nM_p\cdot\nabla_x\ddot U_{\rm av}(p).
\]
We differentiate \eqref{eq:c7-average} in $x$, using
$\nabla_xA^-(p,x)=2H^-(p,x)M(p,x)/n$, where
$H^-(p,x)=\nabla_x^2U^-(p,x)$, and obtain
\[
 M_p\cdot\nabla_x\ddot U_{\rm av}(p)
 =-\frac{2a(p)m_p^2}n\,M_p^TH_p^-M_p
 -\int_{(p,1]}g_p^\ell(Y_p)\,M_p\cdot\nabla(P^-_{p,\ell}A_\ell^-)(X_p)
 \,\nu(\dd\ell).
\]
We multiply by $\bar f_p$ and take the expectation. Using
$\E^0[\bar f_pm_p^2]=\Var^0(f_p)=v_p^0$ and
$\E^0[\bar f_pg_p^\ell(Y_p)]=\Cov^0(f_p,g_p^\ell(Y_p))$, we find
\[
 \frac1{2n}\sum_i\E\bigl[\bar f_pA_{p,ii}^-\bigr]
 =v_p^0a_p^H-2a(p)v_p^0b_p^H
 -\int_{(p,1]}\Cov^0\bigl(f_p,g_p^\ell(Y_p)\bigr)z_{p,\ell}\,\nu(\dd\ell).
\]
Second, \eqref{eq:cav-explicit-derivatives} gives
$n^{-1}\sum_if_{p,i}A_{p,i}^-=4n^{-2}m_p^2h_p\,M_p^TH_p^-M_p$, and
therefore
\[
 \frac1{2n}\sum_i\E\bigl[2f_{p,i}A_{p,i}^-\bigr]
 =4\E^0[m_p^2h_p]\,b_p^H.
\]
Third, by the definition of $\mathfrak s_i$, we have
\[
 \frac1n\sum_i\mathfrak s_iA_{p,i}^-
 =\frac{2m_p}{n^2}\Bigl[\epsilon\,\tau^TH_p^-M_p
 -\int m_r\,M_r^TH_p^-M_p\,\nu(\dd r)\Bigr].
\]
We multiply by $\bar f_p$ and take the expectation. Conditioning
$\epsilon$ on $\calG^0_p$ and $\tau$ on $\calG^-_p$, we have
$\E^0[\bar f_pm_p\epsilon]=\E^0[\bar f_pm_p^2]=v_p^0$ and
$\E^-[\tau^TH_p^-M_p]=\E^-[M_p^TH_p^-M_p]$. We conclude that
\[
 \frac1{2n}\sum_i\E\bigl[2\mathfrak s_i\bar f_pA_{p,i}^-\bigr]
 =2v_p^0b_p^H-2\int\E^0[\bar f_pm_pm_r]\,b_{r,p}^H\,\nu(\dd r).
\]
Adding the three terms, we find that group (d) equals
\begin{equation}\label{eq:c7-group-d}
 \begin{aligned}
 &v_p^0a_p^H+\bigl[4\E^0(m_p^2h_p)+2(1-a(p))v_p^0\bigr]b_p^H\\
 &-\int_{(p,1]}\Cov^0\bigl(f_p,g_p^\ell(Y_p)\bigr)z_{p,\ell}\,\nu(\dd\ell)
 -2\int\E^0[\bar f_pm_pm_r]\,b_{r,p}^H\,\nu(\dd r),
 \end{aligned}
\end{equation}
which is $T_p^{(3)}$ of \eqref{eq:cav-T3}.

\paragraph{Step 6: comparison with the scalar response operator.}
By \eqref{eq:productcovderivative}, $\mathcal C_{pp}'(0)/\alpha$ is the
sum of \eqref{eq:c7-group-a}, \eqref{eq:c7-group-b},
\eqref{eq:c7-group-c}, and \eqref{eq:c7-group-d}. We take
$\varphi=\Sigma^-(\cdot,p)$ in (J1)--(J7) and compare;
\cref{tab:c7-matching} lists the origin of every term.

\begin{table}[ht]
\centering\small
\begin{tabular}{@{}lll@{}}
\toprule
Term & Origin in \eqref{eq:productcovderivative} & Group\\
\midrule
(J1) & coefficients of $\Sigma^-(p,p)$ in \eqref{eq:c7-group-a} and \eqref{eq:c7-group-b} & (a), (b)\\
(J2) & future part of $\partial_y\ddot U_i(p)$ in $f_{p,ii}$ & (a)\\
(J3) & integral part of $\mathfrak s_i$, times $f_{p,i}$ & (b)\\
(J4) & cross term of $\mathfrak s_i^2$ & (c)\\
(J5) & present-level part of $-\int\ddot U_i\dd\nu$ & (c)\\
(J6) & square of the integral part of $\mathfrak s_i$ & (c)\\
(J7) & future part of $-\int\ddot U_i\dd\nu$, with $\Sigma^-(\ell,p)$ for $\Theta(r;\ell,p)$ & (c)\\
last line of \eqref{eq:diagonallinear} & future part of $-\int\ddot U_i\dd\nu$, the difference $\Theta-\Sigma^-$ & (c)\\
$T_p^{(1)}$ & fluctuations $x_p$, $y_{pr}$ in the score term & (b)\\
$T_p^{(2)}$ & fluctuations $x_r$, $y_{r\ell}$ in $\mathfrak s_i^2$ & (c)\\
$T_p^{(3)}$ & all derivatives of $A_p^-$ & (d)\\
\bottomrule
\end{tabular}
\caption{The terms of \eqref{eq:diagonallinear}, with
$\varphi=\Sigma^-(\cdot,p)$ in (J1)--(J7), and the groups of
\eqref{eq:productcovderivative} from which they arise.}
\label{tab:c7-matching}
\end{table}

In detail, the integral in \eqref{eq:c7-group-a} is (J2), and the
integral in \eqref{eq:c7-group-b} is (J3), with
$\varphi(p\wedge r)=\Sigma^-(p\wedge r,p)$. The first three terms of
\eqref{eq:c7-group-c} are (J4), (J6), and (J5). In the fourth
term of \eqref{eq:c7-group-c}, we write
$\Theta(r;\ell,p)=\Sigma^-(\ell,p)+[\Theta(r;\ell,p)-\Sigma^-(\ell,p)]$
and use $\E^0[\bar f_pg_r^\ell(Y_r)]=\Cov^0(f_p,g_r^\ell(Y_r))$. The
first part is (J7), and the second part is the last line of
\eqref{eq:diagonallinear}. The remaining terms are $T_p^{(1)}$,
$T_p^{(2)}$, and $T_p^{(3)}$. The terms (J1)--(J7) add up to
$d^{-1}(J_{a,d}\Sigma^-(\cdot,p))(p)
=d^{-1}(J^{\scal}_{a,d}\Sigma^-)(p,p)$, and \eqref{eq:diagonallinear}
follows.

Atoms of $\nu$ need no separate treatment, because the computation
uses only the exact density \eqref{eq:treedensity} and the exact
parameter responses \eqref{eq:nodeTheta}. At an atom, the formulas use the
right-continuous value $a(p)$ together with the strict future
intervals $(p,1]$ and $(r,1]$, as \eqref{eq:scalarJexplicit} does, and
the double integrals over $[0,1]^2$ include their diagonal atoms. This
completes the proof.
\end{proof}

\subsection{Proof of the integrated linear response}

\begin{proof}[Proof of \cref{lem:linear}]
The proof has three steps. First, we show that the last line of
\eqref{eq:diagonallinear}, integrated over $p$, equals
$-\mathcal P/2$, and that $\mathcal P\ge0$. Second, we prove the
assertions about $\Sigma^-$ and the trace. Third, we integrate
\eqref{eq:diagonallinear}.

\paragraph{Step 1: the conditional future correction.}
If $r\ge p$, the bracket in \eqref{eq:diagonallinear} vanishes by
\eqref{eq:c7-theta-past}. As the inner integral runs over $\ell>r$,
only levels with $p,\ell>r$ contribute. For such levels, we have
$g_r^p(Y_r)=\E^0[f_p\mid\calG^0_r]$, which implies that
\[
 \Cov^0\bigl(f_p,g_r^\ell(Y_r)\bigr)
 =\Cov^0\bigl(g_r^p(Y_r),g_r^\ell(Y_r)\bigr)=\Lambda_r(p,\ell).
\]
By the law of total covariance given $\calG^-_r$, we have
\[
 \begin{aligned}
 \Sigma^-(\ell,p)
 &=\E^-\Cov^-\bigl(A_\ell^-,A_p^-\bigm|\calG^-_r\bigr)
 +\Cov^-\bigl(\E^-[A_\ell^-\mid\calG^-_r],\E^-[A_p^-\mid\calG^-_r]\bigr)\\
 &=\Sigma^-_r(p,\ell)+\Theta(r;\ell,p),
 \end{aligned}
\]
where in the second equality we used that the covariance of the
$\calG^-_r$-measurable variable $\E^-[A_\ell^-\mid\calG^-_r]$ with
$A_p^-$ equals its covariance with $\E^-[A_p^-\mid\calG^-_r]$. We
conclude that $\Theta(r;\ell,p)-\Sigma^-(\ell,p)=-\Sigma^-_r(p,\ell)$. Integrating
the last line of \eqref{eq:diagonallinear} against $\nu(\dd p)$, and
exchanging the order of integration, we obtain
\begin{equation}\label{eq:c7-correction-integrated}
 \frac12\int\nu(\dd p)\int\nu(\dd r)\int_{(r,1]}
  \Cov^0\bigl(f_p,g_r^\ell(Y_r)\bigr)
    \bigl[\Theta(r;\ell,p)-\Sigma^-(\ell,p)\bigr]\nu(\dd\ell)
 =-\frac12\mathcal P.
\end{equation}

To show that $\mathcal P\ge0$, we fix $r\in[0,1]$ and define, under the
product of the scalar and rest laws,
\[
 \Xi_r=\int_{(r,1]}
   \bigl(g_r^p(Y_r)-\E^0g_r^p(Y_r)\bigr)
   \bigl(A_p^--\E^-[A_p^-\mid\calG_r^-]\bigr)\,\nu(\dd p).
\]
We expand the square of $\Xi_r$ and use the independence of the two
factors. Since
$\E^-[(A_p^--\E^-[A_p^-\mid\calG^-_r])(A_\ell^--\E^-[A_\ell^-\mid\calG^-_r])]
=\Sigma^-_r(p,\ell)$, we have
\[
 \E\,\Xi_r^2=\iint_{p,\ell>r}\Lambda_r(p,\ell)\,\Sigma_r^-(p,\ell)
 \,\nu(\dd p)\nu(\dd\ell),
 \qquad\text{and hence}\qquad
 \mathcal P=\int\E\,\Xi_r^2\,\nu(\dd r)\ge0.
\]
Individual entries of $\Lambda_r$ and $\Sigma^-_r$ need not have a
sign; only the integrated form is a square.

\paragraph{Step 2: the covariance operator and the trace.}
Let $B=\bar A^-$ be the random function $p\mapsto\bar A_p^-$. The
function $p\mapsto A_p^-$ is continuous with values in $[0,1]$, and
$\nu([0,1])\le1$. Then $B$ is a random element of $L^2(\nu)$ with
$\|B\|_{L^2(\nu)}\le1$. Its covariance operator
$\Sigma^-=\E^-[B\otimes B]$, where $(B\otimes B)w=\langle B,w\rangle B$,
is a Bochner integral in the space of trace-class operators, since
$\E^-\|B\otimes B\|_1=\E^-\|B\|^2<\infty$. Hence $\Sigma^-$ is positive
and trace class, its integral kernel is $\Sigma^-(p,q)$, and its trace
is $\E^-\|B\|^2_{L^2(\nu)}=\int\Var^-(A_p^-)\nu(\dd p)=V^-$. By
\eqref{eq:c7-root-overlap}, $B(0)=0$; together with $\nu(\{1\})=0$,
this places $B$ in $\mathcal H_\nu$ for every realization. Since
$\mathcal H_\nu$ is closed, the vector
$\Sigma^-w=\E^-[\langle B,w\rangle B]$ lies in $\mathcal H_\nu$ for
every $w\in L^2(\nu)$, that is, the range of $\Sigma^-$ is contained in
$\mathcal H_\nu$.

By \eqref{eq:sc-J-bounded}, $J^{\scal}_{a,d}$ is bounded on
$\mathcal H_\nu$. The rank-one operator $J^{\scal}_{a,d}(B\otimes B)$
has trace $\langle B,J^{\scal}_{a,d}B\rangle$. As the trace is
continuous in the trace norm and $\Sigma^-=\E^-[B\otimes B]$, the trace
of $J^{\scal}_{a,d}\Sigma^-$ on $\mathcal H_\nu$ is
$\E^-\langle B,J^{\scal}_{a,d}B\rangle$. Since $B\in\mathcal H_\nu$
and $J^{\scal}_{a,d}$ is the compression of $J_{a,d}$, we have
$\langle B,J^{\scal}_{a,d}B\rangle=\langle B,J_{a,d}B\rangle
=\int B(p)(J_{a,d}B)(p)\,\nu(\dd p)$. We expand $(J_{a,d}B)(p)$ as the
sum (J1)--(J7) with $\varphi=B$. Each term is a kernel bounded by a
constant depending only on $\beta$, integrated against $\nu$, and
evaluated on $B$ at one of the levels $p$, $r$, $\ell$, $p\wedge r$,
or $r\wedge\ell$. As $|B|\le1$ and $\nu$ is finite, Fubini's theorem
lets us take $\E^-$ inside all integrals, and
$\E^-[B(p)B(q)]=\Sigma^-(q,p)$. We obtain
\[
 \E^-\langle B,J^{\scal}_{a,d}B\rangle
 =\int\bigl(J_{a,d}\Sigma^-(\cdot,p)\bigr)(p)\,\nu(\dd p)
 =\Tr_\nu(J^{\scal}_{a,d}\Sigma^-).
\]

\paragraph{Step 3: integration.}
Let $\alpha>0$. All terms of \eqref{eq:diagonallinear} are bounded
measurable functions of $p$. We integrate \eqref{eq:diagonallinear}
against $\nu(\dd p)$, use \eqref{eq:c7-correction-integrated} and the
definition of $\Tr_\nu$, and multiply by $\alpha$, which gives
\eqref{eq:lineartrace}. If $\alpha=0$, then $\mathcal C_{pp}'(0)=0$ by
\cref{lem:c7-endpoint-pointwise}, and both sides of
\eqref{eq:lineartrace} vanish.
\end{proof}

\section{The integrated response errors}\label{sec:c7-response}

If we bounded the error terms $T_p$ of \cref{lem:linear} one level
at a time, we would meet covariances such as $\E^-[x_r\bar A_p^-]$,
whose Cauchy--Schwarz bound involves $\Var^-(A_p^-)$. Its integral
$V^-$ is what the cavity argument is meant to control, so such a bound
is of no use. Instead, we first integrate in $p$ against the
deterministic scalar kernels. The error terms then become covariances
of centered variables of the rest system with $\nu$-averages of $A^-$,
to which the variance bound \eqref{eq:cl-BL} applies. The resulting
error is linear in $A$ and $J$.

As in \cref{subsec:c7-linear}, $\E^-$ and $\E^0$ refer to $u=0$,
$a_p^-=a_p^-(0)$, $j_p^-=j_p^-(0)$, and $\Gamma=\Gamma_{a,d}$ is as in
\cref{lem:cl-scalar}. We define
\[
 \mathfrak a=\int\Gamma(p)\,a_p^-\,\nu(\dd p),\qquad
 \mathfrak j=\int\Gamma(p)\,j_p^-\,\nu(\dd p).
\]
The quantities $A$ and $J$ below are those of \cref{sec:fnd} for the
full system with $N$ spins and the same profile $\mu$.

\begin{lemma}[Signed response estimate]\label{lem:cl-response}
Let $\mu=s\delta_0+\nu$ be a deterministic profile with
$\nu(\{1\})=0$. Then
\begin{equation}\label{eq:cl-response}
 \Bigl|\int T_p\,\nu(\dd p)\Bigr|\le\mathfrak a+24\,\mathfrak j
\end{equation}
and, with $K_{\beta,3}$ from \cref{cor:comparison},
\begin{equation}\label{eq:cl-response-final}
 \alpha\Bigl|\int T_p\,\nu(\dd p)\Bigr|
 \le192\beta^2K_{\beta,3}\bigl[A+J+N^{-2}\bigr].
\end{equation}
\end{lemma}

\begin{proof}
The proof has eight steps. In Step~1, we record two elementary
inequalities. In Steps~2--6, we split $\int T_p\,\nu(\dd p)$ into eight
pieces and bound each of them: the remaining term of $T^{(1)}$, the two
terms of $T^{(2)}$, the first line of \eqref{eq:cav-T3}, the explicit
and score parts of its future-gradient term, and its last integral over
$r\ge p$ and over $r<p$. In Step~7, we add the bounds (\cref{tab:c7-response}) and obtain
\eqref{eq:cl-response}. In Step~8, we compare $\mathfrak a+\mathfrak j$
with the full system and deduce \eqref{eq:cl-response-final}. Bounded
observables and integrable Gaussian envelopes justify the exchanges of
expectations and finite level integrals below; in Step~1, integrability
follows from $\E^-\xi^2<\infty$.

\paragraph{Step 1: two elementary inequalities.}
Let $\xi$ be a random variable of the rest system with $\E^-\xi=0$ and
$\E^-\xi^2<\infty$, let $r\in[0,1]$, and let $\psi$ be a deterministic
Borel function on $[0,1]$ such that $\psi(q)^2\le\Gamma(r)\Gamma(q)$ for all
$q$. Note that $\E^-[\xi Y]=\E^-[\xi(Y-\E^-Y)]$ for every
square-integrable $Y$, because $\xi$ is centered. Using the
Cauchy--Schwarz inequality and \eqref{eq:cl-BL} at $u=0$, we find
\begin{equation}\label{eq:cl-CS}
 \Bigl|\E^-\Bigl[\xi\int\psi A^-\dd\nu\Bigr]\Bigr|
 \le(\E^-\xi^2)^{1/2}\Var^-\Bigl(\int\psi A^-\dd\nu\Bigr)^{1/2}
 \le2\bigl(\Gamma(r)\,\mathfrak j\,\E^-\xi^2\bigr)^{1/2},
\end{equation}
where in the last inequality we used
$4\int\psi^2j^-\dd\nu\le4\Gamma(r)\int\Gamma j^-\dd\nu=4\Gamma(r)\mathfrak j$.
Next, for every Borel function $\chi\ge0$, we have
\begin{equation}\label{eq:cl-min}
 \int\sqrt{\Gamma j^-}\dd\nu\le\sqrt{\mathfrak j},\qquad
 \iint\chi(r\wedge\ell)\,\nu(\dd r)\nu(\dd\ell)\le2\int\chi\dd\nu ,
\end{equation}
where the first bound follows from the Cauchy--Schwarz inequality and
$\nu([0,1])\le1$, and the second from
$\chi(r\wedge\ell)\le\chi(r)+\chi(\ell)$.

\paragraph{Step 2: the term $T^{(1)}$.}
The first term of \eqref{eq:cav-T1} vanishes, since
$\E^-[x_p\mid\calG_p^-]=0$ by \eqref{eq:cl-xy} and $\bar A_p^-$ is
$\calG_p^-$-measurable. In the second term, if $r\ge p$, then
$y_{pr}=M_p\cdot(M_r-M_p)/n$ has conditional mean zero given
$\calG^-_p$, and the term vanishes. If $r<p$, then
$\E^-[x_r\mid\calG_p^-]=M_p\cdot M_r/n-A_r^-=y_{pr}$, so that
$\E^-[y_{pr}\bar A_p^-]=\E^-[x_r\bar A_p^-]=\E^-[x_rA_p^-]$, since
$x_r$ is centered. Exchanging the integrals over $p$ and $r$, we obtain
\[
 \int T_p^{(1)}\,\nu(\dd p)
 =-2\int\E^-\Bigl[x_r\int_{(r,1]}\E^0[m_ph_pm_r]\,A_p^-\,\nu(\dd p)\Bigr]
   \nu(\dd r).
\]
For fixed $r$, the kernel $\psi(p)=\ind_{\{p>r\}}\E^0[m_ph_pm_r]$
satisfies $\psi^2\le\Gamma(r)\Gamma$ by \eqref{eq:cl-k1}, and
$\E^-x_r^2=j_r^-$ by \eqref{eq:cl-xy}. By \eqref{eq:cl-CS}, the inner
expectation is at most $2(\Gamma(r)\mathfrak j\,j_r^-)^{1/2}$ in
absolute value. Integrating over $r$ and using \eqref{eq:cl-min}, we
conclude that
\[
 \Bigl|\int T_p^{(1)}\,\nu(\dd p)\Bigr|
 \le4\sqrt{\mathfrak j}\int\sqrt{\Gamma j^-}\dd\nu\le4\mathfrak j.
\]

\paragraph{Step 3: the term $T^{(2)}$.}
Since $x_r$ is centered, the first term of \eqref{eq:cav-T2},
integrated in $p$, equals
\[
 -\int\E^-\Bigl[x_r\int\E^0[\bar f_p\epsilon m_r]\,A_p^-\,\nu(\dd p)\Bigr]
 \nu(\dd r).
\]
The kernel $\psi(p)=\E^0[\bar f_p\epsilon m_r]$ satisfies
$\psi^2\le\Gamma(r)\Gamma$ by \eqref{eq:cl-k1}. Using \eqref{eq:cl-CS} and \eqref{eq:cl-min} as in Step~2, but
without the factor $2$, we bound
this term by $2\sqrt{\mathfrak j}\int\sqrt{\Gamma j^-}\dd\nu
\le2\mathfrak j$. In the second term of \eqref{eq:cav-T2}, the variable
$y_{r\ell}$ is centered with $\E^-y_{r\ell}^2\le j^-_{r\wedge\ell}$ by
\eqref{eq:cl-xy}, and the kernel $\psi(p)=\E^0[\bar f_pm_rm_\ell]$
satisfies $\psi^2\le\Gamma(r\wedge\ell)\Gamma$ by \eqref{eq:cl-k2}.
Using \eqref{eq:cl-CS} at the level $r\wedge\ell$ and both parts of
\eqref{eq:cl-min}, with $\chi=\sqrt{\Gamma j^-}$, we bound this term
by
\[
 \frac12\iint2\sqrt{\mathfrak j}\,\chi(r\wedge\ell)\,
 \nu(\dd r)\nu(\dd\ell)
 \le2\sqrt{\mathfrak j}\int\chi\dd\nu\le2\mathfrak j.
\]

\paragraph{Step 4: the first line of $T^{(3)}$.}
By \cref{lem:fnd-posterior}, as recalled in \cref{subsec:c7-row}, we
have $0\preceq H_p^-\preceq C_p^-$. For positive semidefinite matrices
$0\preceq H\preceq C$, we have
$\Tr(H^2)\le\Tr(HC)\le\Tr(C^2)$, because
$\Tr(H(C-H))=\Tr(H^{1/2}(C-H)H^{1/2})\ge0$ and
$\Tr((C-H)C)=\Tr(C^{1/2}(C-H)C^{1/2})\ge0$. Further,
$M_p^TH_p^-M_p\le M_p^TC_p^-M_p$. We conclude that
\[
 0\le a_p^H\le a_p^-,\qquad 0\le b_p^H\le j_p^-.
\]
In addition, \eqref{eq:cl-k3} implies that $0\le v_p^0\le\Gamma(p)$ and
$0\le\E^0[m_p^2h_p]\le\Gamma(p)$, and we have $0\le1-a(p)\le1$. The first line
of \eqref{eq:cav-T3} is therefore nonnegative and at most
$\Gamma(p)a_p^-+(4+2)\Gamma(p)j_p^-$. Integrating against $\nu$, we
find that it contributes at most $\mathfrak a+6\mathfrak j$.

\paragraph{Step 5: the future gradient in $T^{(3)}$.}
We first compute $z_{p,\ell}$ for $\ell\ge p$. Given $\calG_p^-$, the
future of the rest system consists of the increments
$(X_r-X_p)_{r\ge p}$ and the terminal spin $\tau$. By the edge
likelihood \eqref{eq:cav-edge-likelihood} on $[p,1]$ and the terminal
Gibbs probability, as in \eqref{eq:treedensity}, their conditional law
has the density
\[
 \exp\Bigl\{\mathcal E^-(\tau)+X_1\cdot\tau-a(p)U^-(p,X_p)
 -\int_{(p,1]}U^-(r,X_r)\,\nu(\dd r)\Bigr\}
\]
with respect to driftless Gaussian increments and counting measure on
$\tau$. If we vary $X_p=x$ with the increments fixed, every $X_r$ with
$r\ge p$ moves by the same vector. The gradient in $x$ of the
logarithm of this density is therefore the score
$\tau-a(p)M_p-\int_{(p,1]}M_r\,\nu(\dd r)$, while the gradient in $x$
of $A_\ell^-=A^-(\ell,X_\ell)$ is $2H_\ell^-M_\ell/n$. We may
differentiate $(P^-_{p,\ell}A_\ell^-)(x)$ under the integral, because at fixed $N$ the density and its $x$-gradient have an
integrable Gaussian envelope. This gives
\[
 \nabla(P^-_{p,\ell}A_\ell^-)(X_p)
 =\E^-\Bigl[\frac2nH_\ell^-M_\ell
 +A_\ell^-\Bigl(\tau-a(p)M_p-\int_{(p,1]}M_r\,\nu(\dd r)\Bigr)
 \Bigm|\calG^-_p\Bigr].
\]
Taking the inner product with $M_p/n$ and using $a(p)+\nu((p,1])=1$,
we obtain
\begin{equation}\label{eq:c7-future-gradient}
 \begin{gathered}
 z_{p,\ell}=\frac2{n^2}\E^-[M_p^TH_\ell^-M_\ell]+\E^-[\zeta_pA_\ell^-],
 \\
 \zeta_p=\frac{M_p\cdot(\tau-M_p)}n
   -\int_{(p,1]}\frac{M_p\cdot(M_r-M_p)}n\,\nu(\dd r).
 \end{gathered}
\end{equation}

We bound the explicit term. For $\ell\ge p$, \eqref{eq:c7-total-cov}
implies that $\E^-[C_\ell^-\mid\calG_p^-]\preceq C_p^-$. Recall also
that $H_\ell^-\preceq C_\ell^-$. Using the Cauchy--Schwarz inequality
for the positive semidefinite form of $H_\ell^-$ and then for $\E^-$,
we find
\[
 \frac2{n^2}\bigl|\E^-[M_p^TH_\ell^-M_\ell]\bigr|
 \le\frac2{n^2}\bigl(\E^-[M_p^TC_\ell^-M_p]\,
 \E^-[M_\ell^TC_\ell^-M_\ell]\bigr)^{1/2}
 \le2\sqrt{j_p^-j_\ell^-},
\]
where in the last inequality we conditioned on $\calG^-_p$ in the first
factor. Next, both parts of $\zeta_p$ have conditional mean zero given
$\calG_p^-$, since $\E^-[\tau\mid\calG^-_p]=M_p$ and
$\E^-[M_r\mid\calG^-_p]=M_p$ for $r\ge p$. The first part has
conditional second moment $M_p^TC_p^-M_p/n^2$. By
\eqref{eq:c7-total-cov}, each $M_p\cdot(M_r-M_p)/n$ has conditional
second moment at most $M_p^TC_p^-M_p/n^2$. As $\nu((p,1])\le1$, the
conditional Minkowski inequality bounds the conditional $L^2$ norm of
the second part by $(M_p^TC_p^-M_p)^{1/2}/n$. We conclude that
\[
 \E^-\zeta_p=0,\qquad\E^-\zeta_p^2\le4j_p^-.
\]

Let
$\mathcal K(p,\ell)=\ind_{\{\ell>p\}}\Cov^0\bigl(f_p,g_p^\ell(Y_p)\bigr)$.
By \eqref{eq:cl-k4}, we have $\mathcal K(p,\ell)^2\le\Gamma(p)\Gamma(\ell)$.
The future-gradient term of \eqref{eq:cav-T3}, integrated in $p$, is
$-\iint\mathcal K(p,\ell)z_{p,\ell}\,\nu(\dd p)\nu(\dd\ell)$. Inserting
\eqref{eq:c7-future-gradient}, we split it into an explicit part and a
score part. Using \eqref{eq:cl-min}, we find that the explicit part is
at most
\[
 2\iint\sqrt{\Gamma(p)j_p^-}\,\sqrt{\Gamma(\ell)j_\ell^-}\,
 \nu(\dd p)\nu(\dd\ell)
 =2\Bigl(\int\sqrt{\Gamma j^-}\dd\nu\Bigr)^2\le2\mathfrak j.
\]
The score part is
$-\int\E^-[\zeta_p\int\mathcal K(p,\ell)A_\ell^-\nu(\dd\ell)]\nu(\dd p)$.
Using \eqref{eq:cl-CS} with $\xi=\zeta_p$, $\psi=\mathcal K(p,\cdot)$,
and the level $p$, and then \eqref{eq:cl-min}, we bound it by
$\int2(\Gamma(p)\mathfrak j\cdot4j_p^-)^{1/2}\nu(\dd p)
\le4\sqrt{\mathfrak j}\int\sqrt{\Gamma j^-}\dd\nu\le4\mathfrak j$.

\paragraph{Step 6: the last integral in $T^{(3)}$.}
If $r\ge p$, conditioning on $\calG^-_p$ gives
$\E^-[M_r^TH_p^-M_p]=\E^-[M_p^TH_p^-M_p]$, so that
$b^H_{r,p}=b^H_p\in[0,j_p^-]$. By \eqref{eq:cl-k2}, the kernel
satisfies $|\E^0[\bar f_pm_pm_r]|\le\sqrt{\Gamma(p)\Gamma(p\wedge r)}
=\Gamma(p)$. This part contributes at most
$2\int\Gamma(p)j_p^-\nu([p,1])\,\nu(\dd p)\le2\mathfrak j$. If $r<p$,
the Cauchy--Schwarz inequalities of Step~5 imply that
\[
 |b^H_{r,p}|
 \le n^{-2}\bigl(\E^-[M_r^TC_p^-M_r]\,\E^-[M_p^TC_p^-M_p]\bigr)^{1/2}
 \le\sqrt{j_r^-j_p^-},
\]
where we used $\E^-[C_p^-\mid\calG^-_r]\preceq C_r^-$ in the first
factor. The kernel is at most $\sqrt{\Gamma(r)\Gamma(p)}$ by
\eqref{eq:cl-k2}, and by \eqref{eq:cl-min}, this part also contributes
at most $2(\int\sqrt{\Gamma j^-}\dd\nu)^2\le2\mathfrak j$.

\paragraph{Step 7: the sum.}
\Cref{tab:c7-response} collects the bounds of Steps~2--6, with their
running total.

\begin{table}[ht]
\centering\small
\begin{tabular}{@{}llll@{}}
\toprule
Piece of $\int T_p\,\nu(\dd p)$ & Step & Bound & Running total\\
\midrule
$T^{(1)}$, second term & 2 & $4\mathfrak j$ & $4\mathfrak j$\\
$T^{(2)}$, first term & 3 & $2\mathfrak j$ & $6\mathfrak j$\\
$T^{(2)}$, second term & 3 & $2\mathfrak j$ & $8\mathfrak j$\\
$T^{(3)}$, first line & 4 & $\mathfrak a+6\mathfrak j$ & $\mathfrak a+14\mathfrak j$\\
$T^{(3)}$, future gradient, explicit part & 5 & $2\mathfrak j$ & $\mathfrak a+16\mathfrak j$\\
$T^{(3)}$, future gradient, score part & 5 & $4\mathfrak j$ & $\mathfrak a+20\mathfrak j$\\
$T^{(3)}$, last integral, $r\ge p$ & 6 & $2\mathfrak j$ & $\mathfrak a+22\mathfrak j$\\
$T^{(3)}$, last integral, $r<p$ & 6 & $2\mathfrak j$ & $\mathfrak a+24\mathfrak j$\\
\bottomrule
\end{tabular}
\caption{Bounds for the pieces of the integrated error terms in the
proof of \cref{lem:cl-response}.}
\label{tab:c7-response}
\end{table}

Summing the bounds, we obtain
$|\int T_p\,\nu(\dd p)|\le\mathfrak a+(4+2+2+6+2+4+2+2)\mathfrak j
=\mathfrak a+24\mathfrak j$, which is \eqref{eq:cl-response}.

\paragraph{Step 8: transfer to the full system.}
Fix $q\in[0,1]$, and consider three descendants splitting at $q$, with
rest spins $\tau^1,\tau^2,\tau^3$. Given $\calG_q$, they are
conditionally independent with conditional mean $M_q$ and conditional
covariance $C_q^-$. The difference $D=\tau^1-\tau^2$ has conditional
mean zero and conditional second-moment matrix $2C_q^-$, and $\tau^3$
is conditionally independent of it, with second-moment matrix
$C_q^-+M_qM_q^T$. Then
$\E[(D\cdot\tau^3)^2\mid\calG_q]=\Tr\bigl(2C_q^-(C_q^-+M_qM_q^T)\bigr)
=2\Tr((C_q^-)^2)+2M_q^TC_q^-M_q$, and
\[
 a_q^-(u)+j_q^-(u)
 =\frac12\E_{3,q,u}\Bigl[\Bigl(\frac{(\tau^1-\tau^2)\cdot\tau^3}n\Bigr)^2
 \Bigr]\qquad(0\le u\le1).
\]
The observable on the right is nonnegative and bounded, and it does not
depend on $u$. By \cref{cor:comparison} with $k=3$, its mean at $u=0$
is at most $K_{\beta,3}$ times its mean at $u=1$.

At $u=1$, we use the notation of \cref{sec:fnd} for the full system:
$m_q\in\R^N$ is the full conditional mean, $C_q$ the full conditional
covariance, and $a_q^{\rm th}$ and $j_q$ are as defined there. Then
$C_q^-$ is the principal $n\times n$ block of $C_q$, and
$(M_q,0)=m_q-m_N(q)e_N$, where $e_N$ is the last coordinate vector, so
that $M_q^TC_q^-M_q=(M_q,0)^TC_q(M_q,0)$. The
sum of the squared entries of a principal block is at most that of the
whole matrix, which implies that $\Tr((C_q^-)^2)\le\Tr(C_q^2)$. Since $C_q\succeq0$, we
have
\[
 (M_q,0)^TC_q(M_q,0)\le2m_q^TC_qm_q+2m_N(q)^2(C_q)_{NN},
\]
and $m_N(q)^2(C_q)_{NN}\le1$. Dividing by $n^2$ and taking expectations,
we find
\[
 a_q^-(1)\le\Bigl(\frac Nn\Bigr)^2a_q^{\rm th},\qquad
 j_q^-(1)\le2\Bigl(\frac Nn\Bigr)^2j_q+\frac2{n^2}.
\]
We multiply the comparison by $\Gamma(q)\in[0,1]$, integrate against
$\nu$, and use these bounds, $\nu([0,1])\le1$, and $n\ge N/2$. This
yields
\[
 \mathfrak a+\mathfrak j
 \le K_{\beta,3}\Bigl[\Bigl(\frac Nn\Bigr)^2A+2\Bigl(\frac Nn\Bigr)^2J
 +\frac2{n^2}\Bigr]
 \le8K_{\beta,3}\bigl[A+J+N^{-2}\bigr].
\]
Since $\alpha\le\beta^2$, \eqref{eq:cl-response} implies that
$\alpha|\int T_p\,\nu(\dd p)|\le24\beta^2(\mathfrak a+\mathfrak j)$.
Together with the previous bound, this proves
\eqref{eq:cl-response-final}, with $192=24\cdot8$.
\end{proof}

\chapter{Closing the cavity inequality}\label{sec:c7-closure}

In this section, we complete the cavity argument of
\cref{sec:c7,sec:c7-endpoint}. The main result is
\cref{prop:c7-closure}: if the scalar response operator at a minimizing
profile satisfies $J^{\scal}_{a,d}\preceq(1+\eps_N)I$ with
$0\le\eps_N\le\rho/4$, then the self-overlap variance satisfies
$(3\rho/4)V\le C_{\rm cav}[N^{-1}+A+J+H_3]$. In \cref{cor:c7-defect},
we absorb the cubic moment $H_3$ and bound the defect $\calD_N(\rho)$ by
$C_\beta\log^2(2+N)/(N\lambda\rho^2)$ for $1/N\le\rho\le1$, wherever
the defect is defined and some minimizer satisfies the stability
assumption. In \cref{sec:b7}, we use this corollary to bound the contribution of
the stable interval in the proof of \cref{thm:main}, and we verify the
stability assumption there with \cref{prop:lz-stability}.

We rely on the following earlier results. \Cref{sec:c7}
provides the laws of finitely many descendants with
\cref{lem:streaming}, the comparison of nonnegative observables in
\cref{cor:comparison}, the three-defect bound of
\cref{lem:cav-threeproducts}, and the Taylor remainder of
\cref{lem:remainder}. \Cref{sec:c7-endpoint} provides the integrated
linear response of \cref{lem:linear} and the signed response estimate
of \cref{lem:cl-response}, and \cref{lem:Jdom} states that
$J^{\scal}_{a,d}$ is bounded and self-adjoint. We also use the contact
identities
\eqref{eq:in-contact}, the conditional decomposition in the proof of
\cref{lem:fnd-decomp}, the trace bound \eqref{eq:in-Lvar}, the bound on
$J$ in \cref{prop:th-j}, the cubic moment bound of \cref{cor:wi-cub},
and the defect identity \eqref{eq:in-envelope}.

\paragraph{Standing assumptions and notation.} We keep the setting of
\cref{sec:c7}: $N\ge2$, $n=N-1$, $0<\rho\le1$, $t=1-\rho$,
$c=\beta^2\rho$, and $0<s=\lambda/\beta\le1/2$. The parameters $\alpha$, $d$, and
$b(u)$ are given by \eqref{eq:cavityparameters}. The profile
$\mu=s\delta_0+\nu$ is deterministic, with cumulative function $a$, and
each statement says whether $\mu$ must minimize $\mathcal F_{t,s}$ or
satisfy $\nu(\{1\})=0$. Constants $C_\beta$ depend only on $\beta$; they
do not depend on $N$, $\rho$, $s$, the profile, or the positions and
masses of its atoms.

We write $\E_1$, $\Var_1$, and $\Cov_1$ for expectation, variance, and
covariance under $\Pp_u$ at $u=1$. Recall from \cref{subsec:c7-row}
that $\Pp_1$ is the tilted law of \cref{subsec:fnd-law} for the full
system, so that $\E_1=\E$. At $u=0$, we write $\E^-$ and $\Var^-$ for
the rest system, as in \cref{sec:c7-endpoint}. We also recall the rest
self-overlap $A^-_p=|M_p|^2/n$, the conditional magnetization
$m_N(p)$ of the last spin, the descendant laws $\E_{k,q,u}$, the overlap
defects $\Delta_{\mathsf{ab}}$ and $D_{\mathsf{ab}}$ of
\eqref{eq:defectdefinition}, and the integrated third moment
$\mathcal M_3(u)$ of \eqref{eq:cav-M3}. In \cref{subsec:c7-three}, we
defined the covariance
$\mathcal C_{pp}(u)=\Cov_u(m_N(p)^2,A^-_p)$ and its Taylor remainder
$\mathcal R_{\rm cav}$. In \cref{subsec:c7-linear}, we defined the
covariance $\Sigma^-$ of the rest self-overlap, the centered rest
self-overlap $\bar A^-_p$, the error terms $T_p$ of
\eqref{eq:cav-T1}--\eqref{eq:cav-T3}, and the quantity $\mathcal P$ of
\eqref{eq:Ppositive}. \Cref{lem:linear} uses the rest variance
$V^-=\int\Var^-(A^-_p)\,\nu(\dd p)$ and shows that $\mathcal P\ge0$.

\paragraph{Exchangeability.} The tilted law of \cref{subsec:fnd-law} is
invariant under a simultaneous permutation of the $N$ coordinates of
the couplings, the auxiliary field, and the terminal spin. Indeed, let
$\pi$ be a permutation of $\{1,\dots,N\}$, and let $W^\pi$ and $x^\pi$
denote the matrix and the vector with permuted coordinates. The
reference law of $W$ is invariant under $W\mapsto W^\pi$, and since the
terminal condition in \eqref{eq:fnd-vector-pde} and the Laplacian are
invariant, uniqueness gives $U(q,x^\pi;W^\pi)=U(q,x;W)$. The tilt, the
drift in \eqref{eq:fnd-diffusion}, and the Gibbs law of the terminal
spin transform accordingly. In particular, the vector $(m_1(p),\dots,m_N(p))$ of the
conditional magnetizations $m_i(p)=\E[\sigma_i\mid\calG_p]$ is
exchangeable. As $S_p=N^{-1}\sum_im_i(p)^2$ is a symmetric function of
this vector, the pairs $(S_p,m_i(p)^2)$, $1\le i\le N$, all have the
same law. For every deterministic profile and every $p\in[0,1]$, this
implies that
\begin{equation}\label{eq:c7-exchange}
 \E_1m_N(p)^2=\E S_p,\qquad
 \Cov_1\bigl(S_p,m_N(p)^2\bigr)
 =\frac1N\sum_{i=1}^N\Cov\bigl(S_p,m_i(p)^2\bigr)
 =\Var(S_p).
\end{equation}
At $u=1$, the vector $M_p$ consists of the first $n$ coordinates of
$m_p$, and we have
\begin{equation}\label{eq:c7-site-split}
 NS_p=m_N(p)^2+nA^-_p .
\end{equation}
Exchangeability of the sites is used in the same way in the cavity
method of \cite[Section~1.6]{TalagrandVolI2011}.

\section{The mechanism of the contraction}\label{sec:c7-mechanism}

Before giving the proofs, we explain how the inequality closes. In
\cref{lem:cav-trace} below, we prove the identity
\[
 V=\frac{V^{\rm site}}N+\frac nN\Bigl[
 \frac\alpha d\Tr_\nu(J^{\scal}_{a,d}\Sigma^-)
 -\frac\alpha2\mathcal P+\alpha\int T_p\,\nu(\dd p)
 +\mathcal R_{\rm cav}\Bigr],
\]
where $0\le V^{\rm site}\le1/4$. The rest variance $V^-=\Tr\Sigma^-$
enters this identity through the trace pairing
$\Tr_\nu(J^{\scal}_{a,d}\Sigma^-)$. The other terms in the
bracket are either nonpositive, as $-\alpha\mathcal P/2$ is, or errors
controlled by $A$, $J$, $H_3$, and $N^{-1}$
(\cref{lem:cl-response,lem:remainder,lem:c7-thermal}). Under the
stability assumption $J^{\scal}_{a,d}\preceq(1+\eps_N)I$, the first term in
the bracket
is at most $(\alpha/d)(1+\eps_N)V^-$.

The contraction comes from the factor $\alpha/d$. We consider two
descendants $\mathsf a,\mathsf b$ that split at the level
$q_{\mathsf{ab}}$, and fix their rest spins
$\tau^{\mathsf a},\tau^{\mathsf b}$. Under the reference Gaussian law,
the fields acting on $\epsilon_{\mathsf a}$ and $\epsilon_{\mathsf b}$
are the incident row plus the last coordinate of the auxiliary field.
The row contributes covariance $(\alpha u/n)\,\tau^{\mathsf a}\cdot
\tau^{\mathsf b}=\alpha uR^-_{\mathsf{ab}}$, and the last field
contributes $b(u)q_{\mathsf{ab}}$. The total covariance is
\[
 b(u)q_{\mathsf{ab}}+\alpha uR^-_{\mathsf{ab}}
 =d\,q_{\mathsf{ab}}+\alpha u\,(R^-_{\mathsf{ab}}-q_{\mathsf{ab}}).
\]
Only the part $\alpha u(R^-_{\mathsf{ab}}-q_{\mathsf{ab}})$ depends on
the fluctuations of the rest overlap. Its coefficient $\alpha u$ is at
most $\alpha$, while the coefficient of $q_{\mathsf{ab}}$ is $d$. By
\eqref{eq:cavityparameters}, the ratio of the two satisfies
\begin{equation}\label{eq:c7-alpha-d}
 \frac\alpha d=\frac{\beta^2tn/N}{\beta^2(N-t)/N}
 =t\,\frac{N-1}{N-t}\le t=1-\rho .
\end{equation}
The remaining part of the field has deterministic covariance; at
$u=1$, the last field has variance coefficient $b(1)=c=\beta^2\rho$.
Inserting this bound on the first term into the identity above, we
find $V\le t(1+\eps_N)V^-+{}$errors.

To close the inequality, we must compare $V^-$ with $V$, and the
comparison must have coefficient exactly one. Indeed, suppose that we
only knew $V^-\le KV+{}$errors for a constant $K>1$, which is the kind
of bound that a comparison of the two endpoints through
\cref{cor:comparison} gives, with $K=K_{\beta,k}$. We would then have
$V\le tK(1+\eps_N)V+{}$errors, which bounds $V$ only if
$tK(1+\eps_N)<1$. This fails whenever $\rho\le1-1/K$. With coefficient
one, \cref{lem:cl-star} implies that $V^-\le V+{}$errors, and hence
$V\le(1-\rho)(1+\eps_N)V+{}$errors. As
$(1-\rho)(1+\rho/4)\le1-3\rho/4$, a margin of order $\rho$ remains once
$\eps_N\le\rho/4$.

The comparison with coefficient one has two parts. At the independent
endpoint $u=0$, the rest variance is at most the second moment
$\E^-(A^-_p-p)^2$, which is the mean of a product of two overlap defects
on four descendants that split at $p$. The $u$-derivative of this mean
consists of terms with three overlap defects, and
\cref{lem:cav-threeproducts} bounds them by the third moment
$\mathcal M_3$. At the full endpoint $u=1$, exchangeability and the
contact identity $\E S_p=p$ show that the value of this mean is
$\Var(S_p)$, up to an error $1/(4n^2)$. Finally, \cref{lem:c7-thermal}
bounds $\mathcal M_3$ by $A$, $J$, $H_3$, and $N^{-3}$, and
\cref{cor:wi-cub} bounds $H_3$ by a small multiple of $V$ plus a
thermal term.

The rest of this section carries out this plan. In
\cref{sec:c7-coefficient-one}, we prove \cref{lem:cl-star} and the
trace identity of \cref{lem:cav-trace}. In \cref{sec:c7-third-moment},
we bound $\mathcal M_3$ in \cref{lem:c7-thermal}. In
\cref{sec:c7-closing}, we prove \cref{prop:c7-closure} and
\cref{cor:c7-defect}.

\section{Comparison with coefficient one and the trace identity}
\label{sec:c7-coefficient-one}

\begin{lemma}[Four-descendant variance comparison]\label{lem:cl-star}
Let $\mu=s\delta_0+\nu$ minimize $\mathcal F_{t,s}$. Then
\begin{equation}\label{eq:cl-star}
 V^-\le V+C_\beta\sup_{0\le u\le1}\mathcal M_3(u)+\frac1{4n^2}.
\end{equation}
\end{lemma}

\begin{proof}
For each level $p$, we represent $\E_u(A^-_p-p)^2$ as the mean of a
product of two overlap defects on four descendants. The proof has two
steps. In the first step, we compare the values of this mean at $u=0$
and $u=1$ using \cref{lem:streaming,lem:cav-threeproducts}. In the
second step, we evaluate it at $u=1$ using exchangeability and the
contact identity. Only the second step uses that $\mu$ is a minimizer.

\paragraph{Step 1: The interpolation.} Fix $p\in[0,1]$ and consider
four descendants that split at $p$. Each of the pairs $(1,2)$ and
$(3,4)$ splits at $p$, and we have $\Delta_{12}=R^-_{12}-p$ and
$\Delta_{34}=R^-_{34}-p$. We set
\[
 F_p=\Delta_{12}\Delta_{34}=(R^-_{12}-p)(R^-_{34}-p).
\]
Given $\calG_p$, the four spins are conditionally independent, each
with the conditional law of $\sigma$ given $\calG_p$ under $\Pp_u$. In
particular, $\E_{4,p,u}[R^-_{12}\mid\calG_p]=|M_p|^2/n=A^-_p$, and the
same holds for $R^-_{34}$. As the two pairs are conditionally
independent given $\calG_p$, we have
$\E_{4,p,u}[F_p\mid\calG_p]=(A^-_p-p)^2$. We define
\[
 \mathcal V_p(u)=\E_{4,p,u}F_p=\E_u(A^-_p-p)^2,\qquad
 \mathcal V(u)=\int\mathcal V_p(u)\,\nu(\dd p).
\]

Since $|F_p|\le4$, \cref{lem:streaming} implies that
$\mathcal V_p\in C^2([0,1])$, with $\mathcal V_p'(u)$ given by
\eqref{eq:streaming}. The right side of \eqref{eq:streaming} is
$\alpha/2$ times a signed combination of terms of two kinds. A term of
the first kind is the expectation, under a tree that extends the four
descendants, of $F_p$ times one further overlap defect and a spin factor
$\epsilon_{\mathsf x}\epsilon_{\mathsf y}$ of modulus one. The term of
the second kind is $s\,\mathcal V_p(u)\int\E_{2,r,u}D_{12}\,\nu(\dd r)$,
the product of $\mathcal V_p(u)$ with an expectation over an independent
pair of descendants. Each term contains three overlap defects, two of
which come from $F_p$. The absolute values of the coefficients, together
with $\nu(\dd p)$, form a mark measure of the kind considered in
\cref{subsec:c7-three}. Using \cref{lem:cav-threeproducts}, with four
initial leaves and one differentiation, and $\alpha\le\beta^2$, we find
\begin{equation}\label{eq:cl-star-derivative}
 \int|\mathcal V_p'(u)|\,\nu(\dd p)\le C_\beta\mathcal M_3(u)
 \qquad(0\le u\le1).
\end{equation}

For each $p$, we have
$\mathcal V_p(0)=\mathcal V_p(1)-\int_0^1\mathcal V_p'(u)\dd u$. At
fixed $N$, the derivative $\mathcal V_p'(u)$ is bounded uniformly in $p$
and $u$, because every term in \eqref{eq:streaming} is. We integrate
the identity against $\nu(\dd p)$, bound
$|\int_0^1\mathcal V_p'(u)\dd u|$ by $\int_0^1|\mathcal V_p'(u)|\dd u$,
and exchange the integrals in $p$ and $u$ by Tonelli's theorem. Using
\eqref{eq:cl-star-derivative}, we conclude that
\begin{equation}\label{eq:cl-star-step1}
 \mathcal V(0)\le\mathcal V(1)+C_\beta\sup_{0\le u\le1}\mathcal M_3(u).
\end{equation}
At $u=0$, the variance of $A^-_p$ under the rest law is at most its
second moment about $p$, which gives $V^-\le\mathcal V(0)$. Note that
this bound does not require $\E^-A^-_p=p$. We do not know whether this
identity holds, because the profile $\mu$ need not minimize the functional of
the rest system.

\paragraph{Step 2: The full endpoint.} We write
$f^{\rm site}_p=m_N(p)^2$. By \eqref{eq:c7-site-split}, we have
$nA^-_p=NS_p-f^{\rm site}_p$. Let $p\in\supp\nu$. By \eqref{eq:c7-exchange}
and the contact identity $\E S_p=p$ of \eqref{eq:in-contact}, we have
$\E_1f^{\rm site}_p=p$, and hence $\E_1A^-_p=(Np-p)/n=p$. Using
\eqref{eq:c7-exchange} again, we find
\[
\begin{aligned}
 \Var_1(A^-_p)
 &=\frac{N^2\Var(S_p)-2N\Cov_1(S_p,f^{\rm site}_p)
         +\Var_1(f^{\rm site}_p)}{n^2}\\
 &=\frac{N^2-2N}{n^2}\Var(S_p)+\frac{\Var_1(f^{\rm site}_p)}{n^2}
 \le\Var(S_p)+\frac1{4n^2}.
\end{aligned}
\]
In the last step, we used $N^2-2N=n^2-1$ and
$\Var_1(f^{\rm site}_p)\le1/4$, which holds because
$0\le f^{\rm site}_p\le1$. As $\E_1A^-_p=p$, the left side equals
$\mathcal V_p(1)$. Integrating against $\nu$, which is carried by
$\supp\nu$ and has mass at most one, we find
$\mathcal V(1)\le V+1/(4n^2)$. Together with \eqref{eq:cl-star-step1}
and $V^-\le\mathcal V(0)$, this proves \eqref{eq:cl-star}.
\end{proof}

The next lemma expresses $V$ itself through the cavity interpolation.
It uses exchangeability but not the minimizing property.

\begin{lemma}[Trace identity for the full variance]\label{lem:cav-trace}
Let the profile be deterministic with $\nu(\{1\})=0$, and let
\[
 V^{\rm site}=\int\Var_1\bigl(m_N(p)^2\bigr)\,\nu(\dd p).
\]
Then $0\le V^{\rm site}\le1/4$ and
\begin{equation}\label{eq:finaltrace}
 V=\frac{V^{\rm site}}N+\frac nN\Bigl[
 \frac\alpha d\Tr_\nu(J^{\scal}_{a,d}\Sigma^-)
 -\frac\alpha2\mathcal P+\alpha\int T_p\,\nu(\dd p)
 +\mathcal R_{\rm cav}\Bigr].
\end{equation}
\end{lemma}

\begin{proof}
Each variance in the definition of $V^{\rm site}$ lies in $[0,1/4]$,
because $0\le m_N(p)^2\le1$, and $\nu([0,1])\le1$; this proves the bounds
on $V^{\rm site}$.

Fix $p\in[0,1]$. By \eqref{eq:c7-exchange} and \eqref{eq:c7-site-split},
we have
\begin{equation}\label{eq:c7-trace-site}
 \Var(S_p)=\Cov_1\bigl(S_p,m_N(p)^2\bigr)
 =\frac1N\Var_1\bigl(m_N(p)^2\bigr)+\frac nN\,\mathcal C_{pp}(1).
\end{equation}
At $u=0$, the law $\Pp_0$ is the product of the law of the rest system
and the scalar law. The variable $m_N(p)$ depends only on the scalar
component and $A^-_p$ only on the rest component; hence
$\mathcal C_{pp}(0)=0$. Recall from \cref{subsec:c7-three} that
$\mathcal C_{pp}\in C^2([0,1])$. By Taylor's formula with integral
remainder, we have
\[
 \mathcal C_{pp}(1)=\mathcal C_{pp}'(0)
 +\int_0^1(1-u)\,\mathcal C_{pp}''(u)\dd u .
\]
At fixed $N$, all terms in this identity are bounded uniformly in $p$,
and we may integrate it against $\nu(\dd p)$. The integral of the
remainder is $\mathcal R_{\rm cav}$ by definition. As $\nu(\{1\})=0$,
\cref{lem:linear} applies, and \eqref{eq:lineartrace} gives the integral
of $\mathcal C_{pp}'(0)$. Integrating \eqref{eq:c7-trace-site} against
$\nu(\dd p)$, recalling that $V=\int\Var(S_p)\,\nu(\dd p)$, and inserting
these two integrals, we obtain \eqref{eq:finaltrace}.
\end{proof}

\section{The third moment of the rest overlap}\label{sec:c7-third-moment}

The remaining input to \cref{lem:remainder} and \cref{lem:cl-star} is
the integrated third moment $\mathcal M_3$ of the rest overlap defect.
We bound it by the thermal quantities $A$ and $J$ and the cubic moment
$H_3$ of the full system.

\begin{lemma}[Third moment of the rest overlap]\label{lem:c7-thermal}
For every deterministic profile,
\[
 \sup_{0\le u\le1}\mathcal M_3(u)
 \le K_{\beta,2}\bigl[64(A+J)+16H_3+32N^{-3}\bigr],
\]
where $K_{\beta,2}=e^{21\beta^2}$ is the constant of
\cref{cor:comparison} for two leaves, and $A$, $J$, and $H_3$ are
computed for the full system with the same profile.
\end{lemma}

\begin{proof}
The proof has three steps: an exact formula for the conditional second
moment of $R_q-S_q$ in the full system, the passage from the full
overlap to the rest overlap at $u=1$, and the comparison along the
interpolation.

\paragraph{Step 1: The full overlap.} Fix $q\in[0,1]$, and let
$\sigma^1,\sigma^2$ be conditionally independent given $\calG_q$, each
with the conditional law of $\sigma$, so that
$R_q=\sigma^1\cdot\sigma^2/N$. With $\delta^{\mathsf a}=\sigma^{\mathsf
a}-m_q$, we have
\[
 R_q-S_q=\frac{\delta^1\cdot\delta^2+m_q\cdot\delta^1+m_q\cdot\delta^2}N.
\]
Given $\calG_q$, the vectors $\delta^1$ and $\delta^2$ are independent
and centered, with covariance $C_q$. The three summands are pairwise
orthogonal in conditional $L^2$; for instance,
$\E[(\delta^1\cdot\delta^2)(m_q\cdot\delta^1)\mid\calG_q]=0$, because
$\delta^2$ is centered and independent of $\delta^1$. We find
\[
 \E\bigl[(R_q-S_q)^2\bigm|\calG_q\bigr]
 =\frac{\Tr C_q^2+2m_q^TC_qm_q}{N^2}.
\]
This is the conditional orthogonality in the proof of
\cref{lem:fnd-decomp} (see also \cite[Section~2.5, before (2.17)]{AL});
it does not use the contact identity and holds for every deterministic
profile. Taking expectations, we have
$\E(R_q-S_q)^2=a^{\rm th}_q+2j_q$. As $|R_q|\le1$ and $0\le S_q\le1$, we
have $|R_q-S_q|\le2$ and hence $|R_q-S_q|^3\le2(R_q-S_q)^2$. Integrating
against $\nu$, we conclude that
\begin{equation}\label{eq:c7-thermal-full}
 \int\E|R_q-S_q|^3\,\nu(\dd q)\le2(A+2J).
\end{equation}

\paragraph{Step 2: The rest overlap at $u=1$.} Since $\Pp_1$ is the full
tilted law, two descendants that split at $q$ have, under
$\E_{2,q,1}$, the joint law of the pair $(\sigma^1,\sigma^2)$ of
Step~1. Their full and rest overlaps satisfy
$NR_q=nR^-_{12}+\epsilon_1\epsilon_2$, which implies that
\[
 |R^-_{12}-R_q|=\frac{|R^-_{12}-\epsilon_1\epsilon_2|}N\le\frac2N .
\]
By the convexity of $x\mapsto|x|^3$, we have
$|x+y|^3\le4(|x|^3+|y|^3)$. Applying this first to
$(R^-_{12}-R_q)+(R_q-q)$ and then to $(R_q-S_q)+(S_q-q)$, we find
\[
 |R^-_{12}-q|^3\le32N^{-3}+16|R_q-S_q|^3+16|S_q-q|^3 .
\]
We take expectations and integrate against $\nu$, whose mass is at most
one. Using \eqref{eq:c7-thermal-full} and
$H_3=\int\E|S_q-q|^3\,\nu(\dd q)$, we obtain
\[
 \mathcal M_3(1)\le32N^{-3}+32A+64J+16H_3
 \le64(A+J)+16H_3+32N^{-3}.
\]

\paragraph{Step 3: Comparison.} For each $q$, the function
$|R^-_{12}-q|^3$ is nonnegative, bounded by $8$, and depends only on the
spins of two leaves that split at $q$. By \cref{cor:comparison} with
$k=2$, we have
$\E_{2,q,u}|R^-_{12}-q|^3\le K_{\beta,2}\E_{2,q,1}|R^-_{12}-q|^3$ for
every $u\in[0,1]$. Integrating in $q$ against $\nu$, we obtain
$\mathcal M_3(u)\le K_{\beta,2}\mathcal M_3(1)$, and the lemma follows
from Step~2.
\end{proof}

\section{Closing the variance inequality}\label{sec:c7-closing}

We now combine the trace identity, the error bounds, and the comparison
with coefficient one. The stability assumption below concerns the
scalar response operator of \cref{sec:sc} with the cavity coefficient
$d$ of \eqref{eq:cavityparameters}, acting on the scalar transport space
$\mathcal H_\nu$ of the minimizing profile.

\begin{proposition}[Self-overlap contraction]\label{prop:c7-closure}
Let $N\ge2$, $0<\rho\le1$, and $0<s\le1/2$, and let $\mu=s\delta_0+\nu$
minimize $\mathcal F_{t,s}$, with cumulative function $a$. Suppose that
the scalar response with coefficient $d=\beta^2(1-t/N)$ satisfies
\begin{equation}\label{eq:c7-stab}
 J^{\scal}_{a,d}\preceq(1+\eps_N)I\ \text{ on }\mathcal H_\nu,
 \qquad 0\le\eps_N\le\rho/4 .
\end{equation}
Then, with a constant $C_{\rm cav}$ depending only on $\beta$,
\begin{equation}\label{eq:c7-C1}
 \frac{3\rho}{4}\,V\le C_{\rm cav}\bigl[N^{-1}+A+J+H_3\bigr].
\end{equation}
\end{proposition}

\begin{proof}
The proof has three steps. We first collect the bounds on the error
terms in \eqref{eq:finaltrace} and on $V^-$. We then bound the trace
pairing using \eqref{eq:c7-stab}. Finally, we insert the comparison
\eqref{eq:cl-star} into \eqref{eq:finaltrace} and absorb the resulting
multiple of $V$.

\paragraph{Step 1: The errors.} As $\mu$ is a minimizer,
\eqref{eq:in-contact} implies that $\supp\nu\subset[0,1)$. Hence
$\nu(\{1\})=0$, and \cref{lem:linear,lem:cl-response,lem:cav-trace}
apply. Let
\[
 \mathfrak e=H_3+A+J+N^{-3}.
\]
We claim that
\begin{equation}\label{eq:c7-errors}
 \alpha\Bigl|\int T_p\,\nu(\dd p)\Bigr|\le C_\beta\bigl[A+J+N^{-2}\bigr],
 \qquad
 |\mathcal R_{\rm cav}|\le C_\beta\mathfrak e,
 \qquad
 V^-\le V+C_\beta\mathfrak e+\frac1{4n^2}.
\end{equation}
The first bound is \eqref{eq:cl-response-final}, with
$C_\beta=192\beta^2K_{\beta,3}$. For the second, by \cref{lem:remainder}, we have
$|\mathcal R_{\rm cav}|\le C_\beta\sup_w\mathcal M_3(w)$, and
\cref{lem:c7-thermal} bounds this supremum by $64K_{\beta,2}\mathfrak e$.
For the third, we insert the same bound on $\mathcal M_3$ into
\eqref{eq:cl-star}.

\paragraph{Step 2: The trace pairing.} Let $\bar A^-$ be the random
function $p\mapsto\bar A^-_p$. By \cref{lem:linear}, the trace pairing is
$\Tr_\nu(J^{\scal}_{a,d}\Sigma^-)=\E^-\langle\bar A^-,J^{\scal}_{a,d}\bar
A^-\rangle_{L^2(\nu)}$. The function $\bar A^-$ lies in
$\mathcal H_\nu$ for every realization, since $\bar A^-_0=0$ by
\eqref{eq:c7-root-overlap} and $\nu(\{1\})=0$. By Fubini's theorem, we
also have
$\E^-\|\bar A^-\|^2_{L^2(\nu)}=\int\Var^-(A^-_p)\,\nu(\dd p)=V^-$. By \cref{lem:Jdom}, $J^{\scal}_{a,d}$ is a
bounded self-adjoint operator on $\mathcal H_\nu$, and
\eqref{eq:c7-stab} states that
$\langle v,J^{\scal}_{a,d}v\rangle_{L^2(\nu)}\le(1+\eps_N)\|v\|^2_{L^2(\nu)}$
for every $v\in\mathcal H_\nu$. Taking $v=\bar A^-$ and averaging over
the rest system, we obtain
\[
 \Tr_\nu(J^{\scal}_{a,d}\Sigma^-)\le(1+\eps_N)V^-.
\]
Equivalently, we may test \eqref{eq:c7-stab} on the eigenvectors of
$\Sigma^-$ with nonzero eigenvalues, which lie in $\mathcal H_\nu$. We multiply the last display
by $\alpha/d$, which lies in $[0,t]$ by \eqref{eq:c7-alpha-d}. Using
$(1+\eps_N)V^-\ge0$, we find
\begin{equation}\label{eq:c7-trace-bound}
 \frac\alpha d\Tr_\nu(J^{\scal}_{a,d}\Sigma^-)\le t(1+\eps_N)V^-,
\end{equation}
whatever the sign of the trace.

\paragraph{Step 3: Absorption.} The term $-\alpha\mathcal P/2$ in
\eqref{eq:finaltrace} is nonpositive, because $\mathcal P\ge0$ by
\cref{lem:linear}. By \eqref{eq:c7-errors} and
\eqref{eq:c7-trace-bound}, the bracket in \eqref{eq:finaltrace} is at
most
\[
 t(1+\eps_N)V^-+C_\beta\bigl[A+J+N^{-2}+\mathfrak e\bigr],
\]
and this upper bound is nonnegative. As $0<n/N\le1$, the product of the
bracket with $n/N$ is also at most this upper bound. Using $V^{\rm site}\le1/4$ in \eqref{eq:finaltrace},
we have
\[
 V\le\frac1{4N}+t(1+\eps_N)V^-+C_\beta\bigl[A+J+N^{-2}+\mathfrak e\bigr].
\]
The stability assumption \eqref{eq:c7-stab} implies that
\[
 t(1+\eps_N)\le(1-\rho)\Bigl(1+\frac\rho4\Bigr)
 =1-\frac{3\rho}4-\frac{\rho^2}4\le1-\frac{3\rho}4\le1 .
\]
Inserting the bound on $V^-$ from \eqref{eq:c7-errors}, using
$t(1+\eps_N)V\le(1-3\rho/4)V$, and bounding $t(1+\eps_N)$ by one in the
error terms, we obtain
\[
 V\le\Bigl(1-\frac{3\rho}4\Bigr)V+\frac1{4N}+\frac1{4n^2}
 +C_\beta\bigl[A+J+N^{-2}+\mathfrak e\bigr].
\]
We move $(1-3\rho/4)V$ to the left side and note that each of
$(4n^2)^{-1}$, $N^{-2}$, and $N^{-3}$ is at most $N^{-1}$; for the first,
we used $n\ge N/2$. This proves \eqref{eq:c7-C1}, with a constant $C_{\rm cav}$
depending only on $\beta$.
\end{proof}

By \cref{cor:wi-cub}, $H_3$ is at most an arbitrarily small multiple of
$V$ plus a thermal term. Absorbing this multiple in \eqref{eq:c7-C1}, we
obtain the bound on the defect that we use in \cref{sec:b7}.

\begin{corollary}[Defect on the stable interval]\label{cor:c7-defect}
Let $N\ge2$, $1/N\le\rho\le1$, and $0<s=\lambda/\beta\le1/2$.
\begin{enumerate}[label=\textup{(\alph*)}]
\item Every minimizer $\mu=s\delta_0+\nu$ of $\mathcal F_{t,s}$ at these
parameters that satisfies \eqref{eq:c7-stab} for some $\eps_N$, with $a$
its cumulative function, obeys
\begin{equation}\label{eq:c7-defect}
 \beta(A+2J+V)\le C_\beta\frac{\log^2(2+N)}{N\lambda\rho^2},
\end{equation}
where $A$, $J$, and $V$ are computed at $\mu$, and $C_\beta$ depends only
on $\beta$.
\item If the defect $\calD_N(\rho)$ is defined at $\rho$ by
\eqref{eq:in-envelope} and some minimizer at $\rho$ satisfies
\eqref{eq:c7-stab}, then $\calD_N(\rho)$ satisfies the bound
\eqref{eq:c7-defect}.
\item If $1/N\le\rho_0\le1$ and, for every $\rho\in[\rho_0,1]$, every
minimizer at $\rho$ satisfies \eqref{eq:c7-stab}, then
\begin{equation}\label{eq:c7-integrated}
 \int_{\rho_0}^1\calD_N(\rho)\dd\rho
 \le C_\beta\frac{\log^2(2+N)}{N\lambda\rho_0}.
\end{equation}
\end{enumerate}
\end{corollary}

\begin{proof}
We first prove (a). Fix $C_{\rm cav}$ from \eqref{eq:c7-C1}, and let
$\kappa_\beta=\min\{1/4,(4C_{\rm cav})^{-1}\}$. As $\beta>1$ and
$\rho\ge1/N$, the parameter $M=Nc=\beta^2N\rho$ satisfies
$1<\beta^2\le M\le\beta^2N$, and $\kappa_\beta\rho\in(0,1/4]$. Hence the
hypotheses of \cref{cor:wi-cub} hold with $z=\kappa_\beta\rho$, and
\eqref{eq:wi-cub} gives
\[
 H_3\le\kappa_\beta\rho V
   +\frac{425\log(2M)}{Ms}\log\frac2{\kappa_\beta\rho}.
\]
We insert this bound into \eqref{eq:c7-C1}. The resulting term
$C_{\rm cav}\kappa_\beta\rho V$ is at most $\rho V/4$, and moving it to
the left side leaves $(\rho/2)V$ there. We find
\[
 \rho V\le2C_{\rm cav}\Bigl[N^{-1}+A+J
   +\frac{425\log(2M)}{Ms}\log\frac2{\kappa_\beta\rho}\Bigr].
\]
By \eqref{eq:in-Lvar} and \eqref{eq:th-J}, we have $A\le(Ms)^{-1}$ and
$J\le341(1+\log_+M)^2/(Ms)$. As $1\le M\le\beta^2N$ and $\rho\ge1/N$, we
also have
\[
 \log(2M)\le\log(2\beta^2)+\log N,\qquad
 1+\log_+M\le1+\log\beta^2+\log N,
\]
and $\log(2/(\kappa_\beta\rho))\le\log(2/\kappa_\beta)+\log N$. Each of
these three quantities is at most $C_\beta\log(2+N)$, because
$\log(2+N)\ge\log4>1$. Recalling that
$Ms=\beta^2N\rho\cdot\lambda/\beta=\beta N\rho\lambda$, we find
\[
 \beta V\le\frac{C_\beta}{N\rho}
   +\frac{C_\beta\log^2(2+N)}{N\lambda\rho^2},\qquad
 \beta(A+2J)\le\frac{C_\beta\log^2(2+N)}{N\lambda\rho}.
\]
Finally, $\rho\le1$, $\lambda=\beta s\le\beta/2$, and $\log(2+N)\ge1$
give
\[
 \frac1{N\rho}=\frac{\lambda\rho}{N\lambda\rho^2}
 \le\frac\beta2\cdot\frac{\log^2(2+N)}{N\lambda\rho^2},
 \qquad
 \frac{\log^2(2+N)}{N\lambda\rho}\le\frac{\log^2(2+N)}{N\lambda\rho^2},
\]
which proves \eqref{eq:c7-defect}.

For (b), recall from \eqref{eq:in-envelope} that, at every $\rho$ where
$\calD_N(\rho)$ is defined, $\calD_N(\rho)=\beta(A+2J+V)$ for every
minimizer at $\rho$. We evaluate the right side at a minimizer that
satisfies \eqref{eq:c7-stab}, and (a) gives the bound for
$\calD_N(\rho)$.

For (c), recall from \cref{sec:fnd} that $\calD_N$ is a nonnegative
measurable function defined at almost every level, and that a minimizer
exists at every $\rho$. Under the hypothesis of (c), part (b) applies at
almost every $\rho\in[\rho_0,1]$, and gives
$\calD_N(\rho)\le C_\beta\log^2(2+N)/(N\lambda\rho^2)$ there. This bound
holds at almost every level whichever minimizer is used, so no
measurable choice of minimizers is needed to integrate it. As
$\int_{\rho_0}^1\rho^{-2}\dd\rho=\rho_0^{-1}-1\le\rho_0^{-1}$,
integrating yields \eqref{eq:c7-integrated}.
\end{proof}


\part{Proof of the main theorem}\label{part:proof}
\chapter{Proof of the main theorem}\label{sec:b7}

In this section, we prove \cref{thm:main}. By \eqref{eq:in-envelope},
four times the fractional-moment error $e_{N,\beta}(\lambda)$ is the
integral of the defect $\calD_N(\rho)$ over $\rho\in(0,1]$. We split
this integral at a cutoff $\rho_0$. Above the cutoff, we verify the
stability assumption of the cavity argument with
\cref{prop:lz-stability} and bound the defect with
\cref{cor:c7-defect}. Below the cutoff, we bound the defect in \cref{prop:sm-small}, using the
quantile-shell estimate of \cref{prop:wi-vfs} and the root localization
of \cref{cor:lz-omega}. The cutoff is $\rho_0=\min\{1,K\eta^{1/7}\}$,
where $K$ depends only on $\beta$ and
$\eta=e_{N,\beta}(\lambda)+L_s\lambda^5$ is the trial budget. This budget
contains the unknown error, and the two estimates combine
into an implicit inequality for $e_{N,\beta}(\lambda)$. In
\cref{subsec:b7-general}, we choose the tilt $\lambda=N^{-7/50}$,
normalize this inequality by $N^{-7/10}$, and close it. Then
\cref{prop:fnd-mean-bias} converts the bound on
$e_{N,\beta}(N^{-7/50})$ into the bound on the mean bias. In
\cref{rem:b7-exponents}, we explain how the exponent $7/10$ arises.

We use results from \cref{sec:fnd} (the defect identity
\eqref{eq:in-envelope}, the bound \eqref{eq:fnd-moving}, the
trial budget of \cref{lem:fnd-trial}, the trace bound \eqref{eq:in-Lvar},
the decomposition of \cref{lem:fnd-decomp}, and the reduction of the
mean bias in \cref{prop:fnd-mean-bias}), from \cref{sec:wi,sec:th}
(\cref{prop:wi-vfs,prop:th-j}), from \cref{sec:lz-root}
(\cref{cor:lz-omega}), from \cref{sec:sc} (\cref{prop:lz-stability}), and
from \cref{sec:c7-closure} (\cref{cor:c7-defect}). Before the proof of
\cref{thm:main}, we list these inputs with the parameters at which they
are used.

\paragraph{Standing assumptions and notation.} Throughout this section,
$\beta>1$ is fixed, and $C_\beta$ denotes a constant depending only on
$\beta$ whose value may change from line to line. The constants
$\lambda_0\le\min\{1,\beta/4\}$ and $L_s\ge1$ are those of
\cref{sec:fnd}; in particular, \eqref{eq:fnd-scalar-trial-cost} holds
for $0<\lambda\le2\lambda_0$. Recall that $c=\beta^2\rho$, $M=Nc$,
$s=\lambda/\beta$, and $\mu=s\delta_0+\nu$, and that
$\log_+x=\log(\max\{1,x\})$ for $x\ge0$. As in \cref{sec:fnd}, a
\emph{trial budget at $\rho$} is a number $\eta\ge0$ that bounds the left
side of \eqref{eq:fnd-trial}. Since every minimizer $\gamma_*$ of
$F_\rho$ attains $\min_\gamma F_\rho(\gamma)$, such an $\eta$ is a trial
budget for every minimizer at $\rho$ in the sense of
\eqref{eq:loc-trialbudget}.

\section{The interval near \texorpdfstring{$\rho=0$}{rho = 0}}
\label{subsec:b7-small}

We estimate the three terms of $\calD_N(\rho)=\beta(A+2J+V)$ separately.
By the trace bound \eqref{eq:in-Lvar} and
\eqref{eq:th-J}, we have at every minimizer
\begin{equation}\label{eq:sm-AJ}
 A\le\frac1{Ms},\qquad
 J\le\frac{341(1+\log_+M)^2}{Ms}.
\end{equation}
To estimate the self-overlap variance $V$, we combine
\cref{prop:wi-vfs} on dyadic quantile shells with the root localization
of \cref{cor:lz-omega}. This follows Steps~2--5 of the proof of
\cite[Proposition~7.7]{AL}, with \cref{prop:wi-vfs} in place of
\cite[Lemma~3.3]{AL} and \cref{cor:lz-omega} in place of
\cite[Lemma~7.6]{AL}.

\begin{proposition}[The interval near $\rho=0$]\label{prop:sm-small}
Let $N\ge2$, $0<\lambda\le\lambda_0$, and $0<\rho\le1$, and let $\eta>0$
be a trial budget at $\rho$, that is, a bound for the left side of
\eqref{eq:fnd-trial}. Let
$\omega(\rho,\eta)=\min\{1,(\eta/\rho)^{1/6}\}$ be as in
\eqref{eq:lz-omega}. Then every minimizing profile $\mu$ at these
parameters satisfies
\begin{align}
 \beta(A+2J+V)
 &\le C_\beta\left[
 \frac{\log^2(2+N)}{N\rho\lambda}
 +(N\rho)^{-2/3}
 +(N\rho)^{-2/3}\lambda^{-1/3}\omega(\rho,\eta)^{2/3}
 \right],\label{eq:sm-pointwise}\\
 \beta(A+2J+V)&\le C_\beta(N\rho)^{-1/2},
 \label{eq:sm-moving}
\end{align}
where $A$, $J$, and $V$ are computed at $\mu$. In particular, the same
bounds hold for $\calD_N(\rho)$ whenever $\calD_N(\rho)$ is defined. If
$0<\rho_0\le1$ and $\eta$ is a trial budget at every $\rho\in(0,\rho_0]$,
then
\begin{equation}\label{eq:sm-integrated}
 \int_0^{\rho_0}\calD_N(\rho)\dd\rho
 \le C_\beta\left[
 N^{-1}+\frac{\log^3(2+N)}{N\lambda}
 +N^{-2/3}\rho_0^{1/3}
 +N^{-2/3}\lambda^{-1/3}\eta^{1/9}\rho_0^{2/9}
 \right].
\end{equation}
The constants depend only on $\beta$; in particular, they do not depend
on $N$, $\rho$, $\lambda$, $\eta$, the minimizer, or the positions and
masses of the atoms of $\nu$.
\end{proposition}

\begin{proof}
The proof has five steps. We partition the quantile variable into dyadic
shells and bound the diameter of the image of each shell by the root
localization. We then apply \cref{prop:wi-vfs} on each shell and sum
over the shells to bound $V$. Finally, we deduce the pointwise bounds
and integrate them in $\rho$.

\paragraph{Step 1: Shells and localization.} Fix a minimizing profile
$\mu=s\delta_0+\nu$, and let $Q$ be the quantile function of $\nu$ from
\eqref{eq:fnd-quantile}. Recall that $Q$ pushes Lebesgue measure on
$(0,1-s)$ forward to $\nu$ and that $a(Q(v))\ge s+v$ for
$0<v<1-s$. In particular,
\begin{equation}\label{eq:sm-V-quantile}
 V=\int\Var(S_q)\,\nu(\dd q)=\int_0^{1-s}\Var(S_{Q(v)})\dd v .
\end{equation}
For the dyadic levels $u=2^js<1$, $j\ge0$, of which there are finitely
many, we consider the shells
\[
 I_u=\{v\in(0,1-s):u\le s+v<2u\},\qquad h_u=|I_u|.
\]
(Here $u$ denotes a level; it is unrelated to the parameter of the
cavity interpolation in \cref{sec:c7}.) For each $v\in(0,1-s)$, we have
$s<s+v<1$, and there exists exactly one $j\ge0$ such that
$2^js\le s+v<2^{j+1}s$; for this $j$, we have $u=2^js<1$. Hence the
shells partition $(0,1-s)$. Each $I_u$ is an interval of positive length
$h_u\le u$, and $a(Q(v))\ge s+v\ge u$ for $v\in I_u$. In other words,
$I_u$ has threshold $u$ in the sense of \cref{sec:wi-setup}. A shell may contain only
part of the quantile interval of an atom of $\nu$, which
\cref{prop:wi-vfs} allows.

We write $\omega=\omega(\rho,\eta)$, and we let
$\ell(I_u)=\operatorname*{ess\,sup}_{v,w\in I_u}|Q(v)-Q(w)|$ be the
diameter of the image of $I_u$, as in \cref{sec:wi-interval}. Since
$\mu$ is a minimizer, $0<\lambda\le\lambda_0$, and $\eta$ is a trial
budget for $\gamma_*=\beta\mu([0,\cdot\,])$, \cref{cor:lz-omega} implies
that $0\le Q(v)\le C_Q(s+v+\omega)$ for $0<v<1-s$, with $C_Q$ depending
only on $\beta$. For $v\in I_u$, we have $s+v<2u$. As $Q\ge0$, this
gives
\[
 \ell(I_u)\le C_Q(2u+\omega)\le2C_Q(u+\omega).
\]

\paragraph{Step 2: The variance on one shell.} We use \cref{prop:wi-vfs}
on $I_u$, with threshold $u$, length $h_u$, diameter $\ell(I_u)$, and
$\epsilon_u=(Mu)^{-1}$. The first term in \eqref{eq:wi-vfs} is at most
$\epsilon_u(1+\log_+(h_u/\epsilon_u))\le(Mu)^{-1}(1+\log_+(Mu^2))$,
because $h_u/\epsilon_u=Muh_u\le Mu^2$. For the second term, we use
$h_u\le u$ and $(u+\omega)^{2/3}\le u^{2/3}+\omega^{2/3}$ and find
\[
\begin{aligned}
 (\epsilon_u\ell(I_u))^{2/3}h_u^{1/3}
 &\le(2C_Q)^{2/3}(Mu)^{-2/3}(u+\omega)^{2/3}u^{1/3}\\
 &\le(2C_Q)^{2/3}M^{-2/3}
       \bigl(u^{1/3}+u^{-1/3}\omega^{2/3}\bigr).
\end{aligned}
\]
We conclude that
\begin{equation}\label{eq:sm-shell}
 \int_{I_u}\Var(S_{Q(v)})\dd v
 \le C_\beta\left[
 \frac{1+\log_+(Mu^2)}{Mu}
 +M^{-2/3}\bigl(u^{1/3}+u^{-1/3}\omega^{2/3}\bigr)
 \right].
\end{equation}
The term $u^{-1/3}\omega^{2/3}$ accounts for the range of levels that
the localization permits beyond the shell's own mass.

\paragraph{Step 3: Summation over the shells.} For $u=2^js$, we have
$\log_+(Mu^2)\le\log_+(Ms^2)+2j\log2$. Using
$\sum_{j\ge0}2^{-j}=\sum_{j\ge0}j2^{-j}=2$, we see that the first terms
of \eqref{eq:sm-shell} sum to at most
\[
 \frac{2(1+\log_+(Ms^2))+4\log2}{Ms}
 \le\frac{5(1+\log_+(Ms^2))}{Ms},
\]
where we used $4\log2<3$. Let $c_3=(1-2^{-1/3})^{-1}$. Summing downward
from the largest level, which is less than one, we have
$\sum_{u<1}u^{1/3}\le c_3$, and summing upward from $s$, we have
$\sum_{u<1}u^{-1/3}\le c_3s^{-1/3}$. As the shells partition
$(0,1-s)$, summing \eqref{eq:sm-shell} and using
\eqref{eq:sm-V-quantile} yields
\begin{equation}\label{eq:sm-V}
 V\le C_\beta\left[
 \frac{1+\log_+(Ms^2)}{Ms}
 +M^{-2/3}\bigl(1+s^{-1/3}\omega^{2/3}\bigr)
 \right].
\end{equation}

\paragraph{Step 4: The pointwise bounds.} We insert \eqref{eq:sm-AJ} and
\eqref{eq:sm-V} into $\beta(A+2J+V)$ and use
\[
 Ms=\beta N\rho\lambda,\qquad
 \beta M^{-2/3}=\beta^{-1/3}(N\rho)^{-2/3},\qquad
 \beta M^{-2/3}s^{-1/3}=(N\rho)^{-2/3}\lambda^{-1/3}.
\]
As $M\le\beta^2N$ and $Ms^2=N\rho\lambda^2\le N$, both $1+\log_+M$ and
$1+\log_+(Ms^2)$ are at most $C_\beta\log(2+N)$. Using also
$\log(2+N)\ge1$, we find that $\beta A$, $2\beta J$, and $\beta$ times the
first term of \eqref{eq:sm-V} are at most
$C_\beta\log^2(2+N)/(N\rho\lambda)$, while the other two terms of
\eqref{eq:sm-V} give the last two terms of \eqref{eq:sm-pointwise}. This
proves \eqref{eq:sm-pointwise}. No lower bound on $M$ is needed here.

For \eqref{eq:sm-moving}, recall the first bound in
\eqref{eq:fnd-moving}: at every minimizer,
$\int\E(R_q-q)^2\,\nu(\dd q)\le C\min\{1,M^{-1/2}\}$ with a numerical
constant $C$. The measure $\nu$ is carried by the contact levels, where
\eqref{eq:in-decomp} gives $\E(R_q-q)^2=a^{\rm th}_q+2j_q+\Var(S_q)$.
Integrating against $\nu$, we find
$\beta(A+2J+V)=\beta\int\E(R_q-q)^2\,\nu(\dd q)\le C\beta M^{-1/2}
=C(N\rho)^{-1/2}$, which is \eqref{eq:sm-moving}.

At every $\rho$ at which $\calD_N(\rho)$ is defined, a minimizer
exists, and $\calD_N(\rho)=\beta(A+2J+V)$ for every minimizer at $\rho$
by \eqref{eq:in-envelope}. This implies that the two bounds also hold
for $\calD_N(\rho)$.

\paragraph{Step 5: Integration.} Suppose that $\eta$ is a trial budget
at every $\rho\in(0,\rho_0]$. By Step~4, the bounds \eqref{eq:sm-pointwise}
and \eqref{eq:sm-moving} hold for $\calD_N(\rho)$ at almost every
$\rho\in(0,\rho_0]$. As they hold for every minimizer, no measurable
choice of minimizers is needed to integrate them. We first use
\eqref{eq:sm-moving}, whose singularity at $\rho=0$ is integrable:
\[
 \int_0^{\min\{\rho_0,1/N\}}\calD_N(\rho)\dd\rho
 \le C_\beta N^{-1/2}\int_0^{1/N}\rho^{-1/2}\dd\rho
 =2C_\beta N^{-1}.
\]
If $\rho_0\le1/N$, this proves \eqref{eq:sm-integrated}. If
$\rho_0>1/N$, we integrate \eqref{eq:sm-pointwise} over $[1/N,\rho_0]$.
The first term contributes at most
\[
 \frac{\log^2(2+N)}{N\lambda}\int_{1/N}^{\rho_0}\frac{\dd\rho}\rho
 =\frac{\log^2(2+N)\log(N\rho_0)}{N\lambda}
 \le\frac{\log^3(2+N)}{N\lambda},
\]
and the second at most
$N^{-2/3}\int_0^{\rho_0}\rho^{-2/3}\dd\rho=3N^{-2/3}\rho_0^{1/3}$. For the
third term, we use
$\omega(\rho,\eta)^{2/3}=\min\{1,(\eta/\rho)^{1/9}\}\le(\eta/\rho)^{1/9}$,
which holds for all $\rho>0$. With this bound, the third term
contributes at most
\[
 N^{-2/3}\lambda^{-1/3}\eta^{1/9}
       \int_0^{\rho_0}\rho^{-2/3-1/9}\dd\rho
 =\frac92N^{-2/3}\lambda^{-1/3}\eta^{1/9}\rho_0^{2/9}.
\]
Adding the contributions, we obtain \eqref{eq:sm-integrated}. The
constants in Steps~1--5 depend only on $\beta$, through $C_Q$, through
the absolute constants of \cref{prop:wi-vfs}, \cref{prop:th-j}, and
\eqref{eq:fnd-moving}, and through powers of $\beta$.
\end{proof}

\section{Closure at the critical tilt}
\label{subsec:b7-general}\label{subsec:b7-main}

We now fix the tilt. We take $\lambda=N^{-7/50}$, so that the scalar
cost $\lambda^5$ in the trial budget equals $N^{-7/10}$. The cutoff
$\rho_0$ is defined through the trial budget, which contains the
unknown error $e_{N,\beta}(\lambda)$, and we do not assume any
preliminary bound on this error. After division by $N^{-7/10}$, the
terms that depend on the budget become sublinear powers of the single
variable $x=\eta/N^{-7/10}$, and the resulting inequality bounds $x$.
We follow the closure in the proof of \cite[Proposition~7.7]{AL}:
there, the trial budget \cite[(7.44)]{AL} also contains the unknown
error, and the resulting inequality \cite[(7.52)]{AL} is closed in
Step~6 because the error enters only through sublinear powers. Here the
cutoff produces the leading term $N^{-2/3}\eta^{1/21}$.

\paragraph{The inputs.} The proof of \cref{thm:main} uses the following
results, at the parameters listed. In each case, the hypotheses are
verified in the step indicated.

\begin{center}
\small
\begin{tabular}{@{}>{\raggedright\arraybackslash}p{0.2\textwidth}
 >{\raggedright\arraybackslash}p{0.4\textwidth}
 >{\raggedright\arraybackslash}p{0.32\textwidth}@{}}
\toprule
Result & Conclusion used & Parameters in the proof\\
\midrule
\Cref{lem:fnd-trial} &
$\eta=e+L_s\lambda^5$ is a trial budget at every $\rho\in(0,1]$ &
$N\ge N_1$, $\lambda=N^{-7/50}\le\lambda_0$ (Step~1)\\
\eqref{eq:in-envelope} &
$4e=\int_0^1\calD_N(\rho)\dd\rho$, and $\calD_N(\rho)=\beta(A+2J+V)$ at
every minimizer, for almost every $\rho$ &
$s=\lambda/\beta\le1/4$ (Steps~1 and~4)\\
\Cref{prop:lz-stability} &
$J^{\scal}_{a,d}\preceq(1+\eps_N)I$ with
$\eps_N\le C_{\rm st}[\rho^{-3/4}\eta^{1/4}+\lambda^2+N^{-1}]$, at every
minimizer &
$N\ge2$, $\rho\in[\rho_0,1]$, trial budget $\eta$ (Step~3)\\
\Cref{cor:c7-defect}(c) &
$\int_{\rho_0}^1\calD_N\le C_\beta\log^2(2+N)/(N\lambda\rho_0)$ &
$1/N\le\rho_0\le1$, $s\le1/2$, stability on $[\rho_0,1]$ (Step~3)\\
\Cref{prop:sm-small} &
the integrated bound \eqref{eq:sm-integrated} &
$0<\rho_0\le1$, trial budget $\eta>0$ on $(0,\rho_0]$ (Step~4)\\
\Cref{prop:fnd-mean-bias} &
$(P_\beta-p_{N,\beta})/\beta\le32L_s\lambda^5+2e$ &
$N\ge N_1$, $0<\lambda\le\lambda_0$ (Step~6)\\
\bottomrule
\end{tabular}
\end{center}

\begin{proof}[Proof of \cref{thm:main}]
The proof has six steps. We fix the parameters and the common trial
budget, define the cutoff, bound the integral of the defect above and
below the cutoff, close the resulting inequality after normalization,
and convert the bound on the fractional-moment error into the bound on
the mean bias.

\paragraph{Step 1: Parameters and the trial budget.} Let
$N_1=\max\{2,\lceil\lambda_0^{-50/7}\rceil\}$, which depends only on
$\beta$. For every integer $N\ge N_1$, we have $N^{-7/50}\le\lambda_0$,
and we set
\begin{equation}\label{eq:b7-parameters}
 \lambda=N^{-7/50},\qquad r=N^{-7/10}=\lambda^5,\qquad
 e=e_{N,\beta}(\lambda),\qquad \eta=e+L_sr.
\end{equation}
Then $0<\lambda\le\lambda_0\le\min\{1,\beta/4\}$ and
$s=\lambda/\beta\le1/4$. By \cref{lem:fnd-trial}, $\eta$ is a trial
budget at every $\rho\in(0,1]$. The number $e$ is finite. Indeed, the
constrained value $E_\beta(\lambda)$ is finite, and
$F^c_{N,\beta}/\beta$ is a Lipschitz function of finitely many Gaussian
variables (\cref{sec:fnd}), which has finite exponential moments of
every order; these moments enter \eqref{eq:fnd-error-definition}. As
$e\ge0$ by \eqref{eq:in-envelope} and $L_s\ge1$, we have
\begin{equation}\label{eq:b7-eta-lower}
 \infty>\eta\ge L_s\lambda^5\ge\lambda^5>0.
\end{equation}

\paragraph{Step 2: The cutoff.} Let $C_{\rm st}\ge1$ be the constant of
\cref{prop:lz-stability}. We define
\begin{equation}\label{eq:b7-cutoff}
 K=12C_{\rm st},\qquad \rho_0=\min\{1,K\eta^{1/7}\}.
\end{equation}
In both cases $\rho_0<1$ and $\rho_0=1$, we have
\[
 \rho_0\le K\eta^{1/7},\qquad
 \rho_0\ge\min\{1,K\lambda^{5/7}\}\ge\lambda^{5/7}=N^{-1/10}\ge N^{-1},
\]
where we used \eqref{eq:b7-eta-lower}, $K\ge1$, and $\lambda\le1$.

\paragraph{Step 3: The stable interval.} Suppose first that $\rho_0<1$,
and let
$\rho\in[\rho_0,1]$. Then $\rho\ge K\eta^{1/7}\ge K\lambda^{5/7}$. Since
$N^{-1}\le\lambda\le1$, this implies
\[
 \lambda^2\le\lambda^{5/7}\le\frac\rho K,\qquad
 N^{-1}\le\lambda^{5/7}\le\frac\rho K,\qquad
 \rho^{-3/4}\eta^{1/4}\le\frac\rho{K^{7/4}}\le\frac\rho K,
\]
where the last bound follows from $\rho^{7/4}\ge K^{7/4}\eta^{1/4}$ and
$K\ge1$. Let $\mu$ be any minimizer at $\rho$, with cumulative function
$a$. The hypotheses of \cref{prop:lz-stability} hold: $N\ge2$,
$0<\rho\le1$, $0<\lambda\le\lambda_0$, and $\eta$ is a trial budget for
$\mu$. Applying the proposition with $d=\beta^2(1-t/N)$, the cavity
coefficient of \eqref{eq:cavityparameters}, we have
\[
 J^{\scal}_{a,d}\preceq(1+\eps_N)I\ \text{ on }\mathcal H_\nu,\qquad
 0\le\eps_N\le C_{\rm st}\bigl[\rho^{-3/4}\eta^{1/4}+\lambda^2+N^{-1}\bigr]
 \le\frac{3C_{\rm st}\rho}K=\frac\rho4,
\]
where $\eps_N\ge0$ is the number defined in \cref{prop:lz-stability}.
This is the stability condition \eqref{eq:c7-stab}, at every minimizer
and every $\rho\in[\rho_0,1]$. As $N\ge2$, $1/N\le\rho_0\le1$, and
$s\le1/2$, \cref{cor:c7-defect}(c) applies. Using
$\rho_0=K\eta^{1/7}\ge\lambda^{5/7}$ and
$\lambda^{12/7}=N^{-6/25}$, we obtain
\begin{equation}\label{eq:b7-closure}
 \int_{\rho_0}^1\calD_N(\rho)\dd\rho
 \le C_\beta\frac{\log^2(2+N)}{N\lambda\rho_0}
 \le C_\beta\frac{\log^2(2+N)}{N\lambda^{12/7}}
 =C_\beta N^{-19/25}\log^2(2+N).
\end{equation}
When $\rho_0=1$, the integral is zero, so \eqref{eq:b7-closure} holds in
both cases.

\paragraph{Step 4: The interval near zero.} The hypotheses of
\cref{prop:sm-small} hold on $(0,\rho_0]$: we have $N\ge2$,
$0<\lambda\le\lambda_0$, and $0<\rho_0\le1$. By Step~1 and
\eqref{eq:b7-eta-lower}, $\eta>0$ is a trial budget at every
$\rho\in(0,\rho_0]$. By Step~2, we have $\rho_0\le K\eta^{1/7}$ in both
cases $\rho_0<1$ and $\rho_0=1$, which implies that
\[
 \rho_0^{1/3}\le K^{1/3}\eta^{1/21},\qquad
 \eta^{1/9}\rho_0^{2/9}\le K^{2/9}\eta^{1/9+2/63}=K^{2/9}\eta^{1/7}.
\]
We insert these bounds into \eqref{eq:sm-integrated}, add
\eqref{eq:b7-closure}, and use $4e=\int_0^1\calD_N(\rho)\dd\rho$ from
\eqref{eq:in-envelope}. This gives
\[
\begin{aligned}
 e\le C_\beta\Bigl[
 &N^{-1}+N^{-43/50}\log^3(2+N)+N^{-19/25}\log^2(2+N)\\
 &+N^{-2/3}\eta^{1/21}+N^{-2/3}\lambda^{-1/3}\eta^{1/7}\Bigr],
\end{aligned}
\]
where we used $N^{-1}\lambda^{-1}=N^{-43/50}$. As $1>43/50>19/25$ and
$\log(2+N)\ge1$, the first three terms are at most
$3N^{-19/25}\log^3(2+N)$. We conclude that
\begin{equation}\label{eq:b7-normalized-input}
 e\le C_\beta\left[
 N^{-19/25}\log^3(2+N)+N^{-2/3}\eta^{1/21}
                     +N^{-2/3}\lambda^{-1/3}\eta^{1/7}\right].
\end{equation}

\paragraph{Step 5: Normalization and closure.} Let
\[
 x=\frac\eta r=\frac er+L_s ,
\]
which is finite by \eqref{eq:b7-eta-lower} and satisfies $x\ge L_s\ge1$.
Then $\eta=rx$, and the terms of the bound for $e$ in Step~4 have the
following sizes after division by $r=N^{-7/10}$.
\begin{center}
\small
\begin{tabular}{@{}lll@{}}
\toprule
Term & Origin & Size after division by $r$\\
\midrule
$N^{-1}$ & $\rho\le1/N$, by \eqref{eq:sm-moving} & $N^{-3/10}$\\
$N^{-43/50}\log^3(2+N)$ & first term of \eqref{eq:sm-pointwise} &
$N^{-4/25}\log^3(2+N)$\\
$N^{-19/25}\log^2(2+N)$ & stable interval, \eqref{eq:b7-closure} &
$N^{-3/50}\log^2(2+N)$\\
$N^{-2/3}\eta^{1/21}$ & second term of \eqref{eq:sm-pointwise} & $x^{1/21}$\\
$N^{-2/3}\lambda^{-1/3}\eta^{1/7}$ & third term of \eqref{eq:sm-pointwise} & $N^{-1/50}x^{1/7}$\\
\bottomrule
\end{tabular}
\end{center}
The entries in the last column are exact. For the two terms that depend
on the budget, we have
\[
 N^{-2/3}\eta^{1/21}=N^{-2/3}r^{1/21}x^{1/21}
 =N^{-2/3-1/30}x^{1/21}=rx^{1/21}
\]
and
\[
 N^{-2/3}\lambda^{-1/3}\eta^{1/7}
 =N^{-2/3+7/150-1/10}x^{1/7}=N^{-18/25}x^{1/7}=rN^{-1/50}x^{1/7},
\]
since $2/3+1/30=7/10$ and $18/25=7/10+1/50$. For the first term of
\eqref{eq:b7-normalized-input}, we have $N^{-19/25}/r=N^{-3/50}$.
Dividing \eqref{eq:b7-normalized-input} by $r$ and adding $L_s$ to both
sides, we find
\begin{equation}\label{eq:b7-normalized}
 x\le L_s+C_\beta\left[
 N^{-3/50}\log^3(2+N)+x^{1/21}+N^{-1/50}x^{1/7}\right]
 \le C_\beta x^{1/7}.
\end{equation}
In the last inequality, we used that
$\sup_{N\ge2}N^{-3/50}\log^3(2+N)$ is finite, because $\log(2+N)$ grows
more slowly than any power of $N$. We also used $N^{-1/50}\le1$ and
the bounds $1\le x^{1/7}$ and $x^{1/21}\le x^{1/7}$, which hold because
$x\ge1$.
Dividing \eqref{eq:b7-normalized} by $x^{1/7}>0$, we find
$x^{6/7}\le C_\beta$, that is, $x\le C_\beta^{7/6}$. Then
\[
 e\le\eta=rx\le C_\beta N^{-7/10}.
\]
This argument uses no preliminary bound on $e$ and no smallness of
$\eta$; it uses only that $x$ is finite.

\paragraph{Step 6: The mean bias.} By \cref{prop:fnd-mean-bias}, which
applies since $0<\lambda\le\lambda_0$, and by Step~5, we have
\[
 0\le\frac{P_\beta-p_{N,\beta}}\beta
 \le32L_s\lambda^5+2e
 \le(32L_s+2C_\beta)N^{-7/10}.
\]
Multiplying by $\beta$, we obtain
$P_\beta-p_{N,\beta}\le C_\beta N^{-7/10}$ for every $N\ge N_1$. For the
finitely many integers $2\le N<N_1$, the mean bias $P_\beta-p_{N,\beta}$
is finite. The number
\[
 C^{\rm small}_\beta=\max_{2\le N<N_1}N^{7/10}(P_\beta-p_{N,\beta}),
\]
with the maximum interpreted as zero when $N_1=2$, is therefore finite
and depends only on $\beta$. Replacing $C_\beta$ by $\max\{C_\beta,C^{\rm small}_\beta\}$, we
obtain the bound of \cref{thm:main} for every $N\ge2$. This completes the
proof.
\end{proof}

\begin{remark}\label{rem:b7-exponents}
The leading term in \eqref{eq:b7-normalized-input} is
$N^{-2/3}\eta^{1/21}$. The exponent $1/21$ is the product of the exponent
$1/3$ of $\rho_0$ in the integrated shell estimate
\eqref{eq:sm-integrated} and the exponent $1/7$ of $\eta$ in the cutoff
\eqref{eq:b7-cutoff}, which comes from the stability threshold
$\rho^{-3/4}\eta^{1/4}\lesssim\rho$. If $\eta$ is of order $N^{-a}$,
the leading term is of the same order when $a=2/3+a/21$, that is,
$20a/21=2/3$ and $a=7/10$. The scalar cost $\lambda^5$ in the trial
budget is then of the same order when $\lambda=N^{-7/50}$. For this
choice, the localization correction has the extra factor $N^{-1/50}$
after normalization, and the stable interval and the two elementary
terms in \eqref{eq:sm-integrated} contribute at most
$N^{-3/50}\log^3(2+N)$, as the table in Step~5 shows. The bound
\eqref{eq:b7-closure} on the stable interval has only a squared
logarithm; the third logarithm comes from integrating the first term of
\eqref{eq:sm-pointwise} over $[1/N,\rho_0]$. These contributions have
strict polynomial slack, and the leading term carries no logarithm.
After the closure, $L_sN^{-7/10}\le\eta\le C_\beta N^{-7/10}$, so that
for all sufficiently large $N$ the cutoff satisfies $\rho_0<1$ and is
comparable to $N^{-1/10}$.
\end{remark}


\backmatter

\setlength{\bibsep}{3pt plus 1pt minus 1pt}
\bibliographystyle{plainnat}
\bibliography{refs}
\end{document}
